\documentclass[11pt]{article}

\usepackage[margin=1in]{geometry}
\usepackage{fontspec}
\usepackage{microtype}
\usepackage{setspace}
\usepackage{amsmath,amssymb,amsfonts,amsthm,mathtools}
\usepackage{bm}
\usepackage{booktabs}
\usepackage{longtable}
\usepackage{array}
\usepackage{tabularx}
\usepackage{mathrsfs}
\usepackage{stmaryrd}
\usepackage{tikz-cd}
\usepackage{hyperref}
\usepackage{cleveref}
\usepackage{titlesec}
\usepackage{parskip}
\usepackage{xcolor}

\setmainfont{TeX Gyre Pagella}
\setsansfont{TeX Gyre Heros}
\setmonofont{DejaVu Sans Mono}
\setstretch{1.08}

\definecolor{ink}{HTML}{18212A}
\definecolor{accent}{HTML}{315C55}
\definecolor{muted}{HTML}{66717C}

\hypersetup{
  colorlinks=true,
  linkcolor=accent,
  citecolor=accent,
  urlcolor=accent,
  pdftitle={Representational Simplicity},
  pdfauthor={Flyxion}
}

\titleformat{\section}
  {\normalfont\Large\bfseries}
  {\thesection}
  {0.8em}
  {}

\titleformat{\subsection}
  {\normalfont\large\bfseries}
  {\thesubsection}
  {0.8em}
  {}

\titleformat{\subsubsection}
  {\normalfont\normalsize\bfseries}
  {\thesubsubsection}
  {0.8em}
  {}

\newtheorem{definition}{Definition}[section]
\newtheorem{theorem}[definition]{Theorem}
\newtheorem{proposition}[definition]{Proposition}
\newtheorem{lemma}[definition]{Lemma}
\newtheorem{corollary}[definition]{Corollary}
\newtheorem{conjecture}[definition]{Conjecture}
\newtheorem{criterion}[definition]{Criterion}
\newtheorem{principle}[definition]{Principle}
\newtheorem{assumption}[definition]{Assumption}
\newtheorem{construction}[definition]{Construction}
\newtheorem{observation}[definition]{Observation}

\theoremstyle{definition}
\newtheorem{example}[definition]{Example}
\newtheorem{counterexample}[definition]{Counterexample}
\newtheorem{remark}[definition]{Remark}
\newtheorem{protocol}[definition]{Protocol}

\DeclareMathOperator{\Adm}{Adm}
\DeclareMathOperator{\Ans}{Ans}
\DeclareMathOperator{\Req}{Req}
\DeclareMathOperator{\Res}{Res}
\DeclareMathOperator{\Esc}{Esc}
\DeclareMathOperator{\Audit}{Audit}
\DeclareMathOperator{\Supp}{Supp}
\DeclareMathOperator{\Vol}{Vol}
\DeclareMathOperator{\KL}{D_{\mathrm{KL}}}
\DeclareMathOperator{\colim}{colim}
\DeclareMathOperator{\argmin}{arg\,min}
\DeclareMathOperator{\argmax}{arg\,max}
\DeclareMathOperator{\Var}{Var}
\DeclareMathOperator{\Cov}{Cov}
\DeclareMathOperator{\Tr}{Tr}
\DeclareMathOperator{\rank}{rank}
\DeclareMathOperator{\im}{im}
\DeclareMathOperator{\Hom}{Hom}
\DeclareMathOperator{\Ob}{Ob}
\DeclareMathOperator{\Ext}{Ext}
\DeclareMathOperator{\Prob}{Pr}
\DeclareMathOperator{\IG}{IG}
\DeclareMathOperator{\Prog}{Prog}
\DeclareMathOperator{\Burden}{Burden}
\DeclareMathOperator{\closure}{cl}
\DeclareMathOperator{\dist}{dist}
\DeclareMathOperator{\Risk}{Risk}
\DeclareMathOperator{\Reach}{Reach}
\DeclareMathOperator{\Corr}{Corr}
\DeclareMathOperator{\Appeal}{Appeal}
\DeclareMathOperator{\Comp}{Comp}

\newcommand{\R}{\mathbb{R}}
\newcommand{\N}{\mathbb{N}}
\newcommand{\E}{\mathbb{E}}
\newcommand{\F}{\mathcal{F}}
\newcommand{\G}{\mathcal{G}}
\newcommand{\K}{\mathcal{K}}
\newcommand{\C}{\mathcal{C}}
\newcommand{\A}{\mathcal{A}}
\newcommand{\B}{\mathcal{B}}
\newcommand{\M}{\mathcal{M}}
\newcommand{\Q}{\mathcal{Q}}
\newcommand{\D}{\mathcal{D}}
\newcommand{\Hh}{\mathcal{H}}
\newcommand{\Z}{\mathcal{Z}}
\newcommand{\ThetaSpace}{\boldsymbol{\Theta}}
\newcommand{\thetaState}{\boldsymbol{\theta}}
\newcommand{\FieldPhi}{\Phi}
\newcommand{\FieldV}{\mathbf{v}}
\newcommand{\FieldS}{S}
\newcommand{\given}{\,\middle|\,}
\newcommand{\ind}{\mathbf{1}}
\newcommand{\eps}{\varepsilon}
\newcommand{\AC}{\mathfrak{A}}
\newcommand{\CLIO}{\mathsf{CLIO}}
\newcommand{\HYDRA}{\mathsf{HYDRA}}
\newcommand{\RSVP}{\mathsf{RSVP}}
\newcommand{\PERSCEN}{\mathsf{PERSCEN}}
\newcommand{\RAT}{\mathsf{RAT}}
\newcommand{\CoM}{\mathsf{CoM}}


\title{
  \textbf{Representational Simplicity}\\[0.4em]
  \large Sufficient Projection, Obstruction-Sensitive Routing,
  and Repair Under Finite Attention
}


\author{
  Flyxion\\
  \small Independent Researcher
}

\date{August 2026}

\begin{document}

\maketitle

\begin{abstract}
This treatise develops a mathematical architecture for many-to-one intellectual communication under finite recipient attention. The central problem is not merely to rank incoming messages, retrieve related documents, or suppress repetition. It is to transform a heterogeneous submission into a representation that preserves every distinction required by authorized downstream decisions, to assemble that representation with participant-specific, archive-specific, and recipient-specific context, to detect when the resulting local structures fail to form a coherent global state, and to allocate scarce human attention without making unfamiliarity itself a reason for suppression.

The proposed architecture integrates three previously distinct frameworks. The Attention Compiler supplies the institutional and conversational problem: many participants submit questions, objections, conjectures, evidence, corrections, and partially formed ideas toward a recipient whose attention budget is finite. CLIO supplies a theory of admissibility-relative projection: a submission may be compressed only through a quotient sufficient for the full class of routing, explanation, testing, escalation, audit, and appeal operations the institution authorizes. HYDRA supplies a theory of structured assembly: the projected submission, participant-state estimate, prerequisite graph, archive neighborhood, provenance record, recipient model, and ambient relevance fields must glue into a coherent routing state before any disposition is treated as warranted.

The resulting system is formalized as a context-indexed composite
\[
\AC(x_i\mid r,t)
=
D\!\left[
\Hh
\left(
\colim_{\C_{ir}}
\A_{ir}
\left(
q_{c_{ir}}(x_i),
K_r,
\theta_i,
G,
M_{i,0:j}
\right);
\FieldPhi_t,\FieldV_t,\FieldS_t
\right)
\right],
\]
where \(q_{c_{ir}}\) is a CLIO projection indexed by participant, recipient, archive, and time; \(\colim_{\C_{ir}}\A_{ir}\) is a HYDRA gluing of locally defined conversational structures; \(\Hh\) is a transition operator governed by relevance, memory, and RSVP field semantics; and \(D\) returns an authorized disposition such as answer, clarification, prerequisite instruction, adversarial test, escalation, random audit, or appeal review.

Two principal failure classes are derived. Projection collapse occurs when the quotient identifies submissions that require different authorized dispositions. Gluing obstruction occurs when locally coherent archive, participant, prerequisite, and recipient states cannot be assembled under the declared compatibility maps. These failures are shown to require different forms of repair. Projection collapse may require restoration of a discarded distinction or expansion of the quotient space. Gluing obstruction requires revision of compatibility maps, contextual indexing, or the architecture of assembly. Neither failure is corrected reliably by greater smoothing, stronger similarity matching, or repeated application of the same routing rule.

The treatise derives context-relative \(F\)-sufficiency, fiberwise collapse observables, disposition divergence, provenance preservation conditions, sequential epistemic calibration, prerequisite-path optimization, parameter-mediated coupling, obstruction-sensitive routing, recipient-relative escalation value, exploration reserves, false-suppression estimators, and queueing conditions for stable attention allocation. It concludes that intellectual triage is warranted only when the representation of a submission is sufficient for every authorized disposition and the structures used to interpret that representation admit coherent gluing under auditable compatibility conditions.
\end{abstract}

\section{The Problem of Finite Attention and Transformative Access}

\subsection{Many-to-one communication as a constrained transformation problem}

Let \(\mathcal{I}\) be a potentially unbounded set of participants and let \(r\) be a recipient whose available attention is finite. Participant \(i\in\mathcal{I}\) produces a sequence of submissions
\[
x_{i,1},x_{i,2},\ldots
\]
at arrival times
\[
t_{i,1},t_{i,2},\ldots.
\]
Each submission may contain several logically distinct components. It may contain a request for information, a disagreement with a recorded claim, a conjecture, evidence, a report of an observation, a correction, an analogy, a demand for clarification, a personal request, or an underdeveloped structure whose eventual status cannot yet be determined.

The recipient has an attention budget
\[
A_r([t,t+\Delta t])
\]
over every finite interval \([t,t+\Delta t]\). The raw arrival stream need not respect this budget. If submissions arrive with average rate \(\lambda_r\), and if direct examination of a raw submission requires random attention cost \(C_{\mathrm{raw}}\), then unmediated communication is stable only when
\[
\lambda_r\,\E[C_{\mathrm{raw}}] < A_r,
\]
where \(A_r\) is interpreted as available attention time per unit time. In an open communication environment this inequality will generally fail. The number of potential participants can grow independently of the recipient's cognitive capacity, while the cost of reconstructing the background of each submission can remain substantial even when the submission itself is short.

A chronological inbox does not solve this problem. It imposes an ordering but performs no transformation. It treats the first repeated question received during an interval as prior to a later counterexample merely because of arrival time. It does not distinguish the cost of answering a familiar question from the value of examining an unfamiliar claim. It does not preserve previously performed explanatory labor except insofar as the recipient manually recalls and reuses it. Most importantly, it does not determine whether two superficially similar submissions may safely receive the same response.

A similarity-ranked inbox also fails to solve the full problem. Similarity can identify neighboring texts, embeddings, topics, or archive nodes, but similarity does not establish substitutability. Two submissions can be close in lexical, semantic, or embedding distance while differing in a single premise that changes the correct disposition. Conversely, two submissions can be textually remote while expressing the same normalized claim. A routing architecture therefore requires a stronger relation than resemblance.

\begin{definition}[Authorized disposition class]
For a recipient \(r\), let
\[
\D_r
=
\{
d_{\mathrm{ans}},
d_{\mathrm{clar}},
d_{\mathrm{teach}},
d_{\mathrm{test}},
d_{\mathrm{esc}},
d_{\mathrm{audit}},
d_{\mathrm{appeal}},
d_{\mathrm{defer}}
\}
\]
be the set of dispositions the system is authorized to produce. These represent, respectively, answering from the archive, requesting clarification, supplying prerequisite instruction, initiating adversarial or diagnostic testing, escalating to the recipient, selecting for audit, initiating appeal review, and deferring under explicitly reported uncertainty.
\end{definition}

The disposition set is not merely a list of interface actions. It defines the functions relative to which a compressed representation of a submission must be sufficient. A quotient that preserves enough information to retrieve an answer may still discard information needed to distinguish misunderstanding from disagreement. A quotient sufficient for prerequisite assignment may still be insufficient for escalation because it has removed provenance, evidential status, or recipient-relative consequence.

Let
\[
\F_{\mathrm{AC},r}
\]
denote the class of all authorized downstream functions used by the Attention Compiler for recipient \(r\). This class may contain functions for claim extraction, disposition assignment, prerequisite selection, novelty estimation, coherence estimation, recipient relevance, consequence estimation, expected joint value, audit selection, and appeal reconstruction.

The governing condition is not that a submission be compressed accurately in some general semantic sense. The condition is that compression preserve every distinction used by every function in \(\F_{\mathrm{AC},r}\).

\begin{definition}[Transformative access]
A participant has transformative access to recipient \(r\) when the participant may enter a governed process that preserves the original submission, reconstructs its claims charitably, identifies archive-relative and recipient-relative structure, repairs remediable prerequisite deficiencies, exposes the strongest surviving residual, and permits escalation whenever continued progress is expected to require the reciprocal attention of \(r\).
\end{definition}

Transformative access differs from immediate personal access. Immediate access sends the raw submission directly to the recipient. Transformative access gives the submission entry into a process whose output may or may not require the recipient, but whose decision is constrained by sufficiency, provenance, audit, and appeal.

This distinction avoids a false opposition between universal communication and finite attention. Universal immediate access to one finite recipient is impossible under unbounded arrival. Universal entry into a transformation process is possible in principle, provided that the process remains stable and does not obtain stability by suppressing precisely the submissions it cannot already recognize.

\subsection{The submission as a structured epistemic object}

A raw submission is not treated as a single string. Let
\[
x_i
=
(\tau_i,\sigma_i,\eta_i,\gamma_i,\zeta_i),
\]
where \(\tau_i\) is the observed textual or multimodal expression, \(\sigma_i\) is the submitter-declared speech act when one is supplied, \(\eta_i\) is provenance available at submission time, \(\gamma_i\) is interaction history, and \(\zeta_i\) contains explicitly declared constraints or intended recipients.

The system does not observe the participant's complete intended meaning. Let
\[
\vartheta_i
\]
denote a latent intended epistemic object. The observed submission is generated through a communication channel
\[
x_i \sim P(\,\cdot\mid \vartheta_i,\ell_i,e_i),
\]
where \(\ell_i\) represents the participant's available expressive language and \(e_i\) represents environmental and interface conditions. The same intended object can produce different surface expressions under different vocabularies, levels of mathematical preparation, languages, disabilities, media, or degrees of urgency.

A system that equates polish with coherence implicitly assumes that the communication channel is uniform across participants. That assumption is generally false. The purpose of charitable reconstruction is to infer a candidate normalized object without treating conventional terminology as a prerequisite for having had the relevant idea.

Let
\[
Q(x_i)
=
\{q_{i,1},\ldots,q_{i,n_i}\}
\]
be a candidate set of normalized claims extracted from \(x_i\). Let
\[
U(x_i)
=
\{u_{i,1},\ldots,u_{i,m_i}\}
\]
be the unresolved questions inferred from the submission, and let
\[
E(x_i)
=
\{e_{i,1},\ldots,e_{i,k_i}\}
\]
be its evidential items. These sets remain provisional. Their provenance must record whether an element was explicitly stated, reconstructed by the system, inferred from context, or supplied by an archive comparison.

\begin{definition}[Provenance typing]
Every derived conversational object \(o\) receives a provenance type
\[
\operatorname{prov}(o)
\in
\{
\mathsf{declared},
\mathsf{reconstructed},
\mathsf{inferred},
\mathsf{retrieved},
\mathsf{tested},
\mathsf{recipient\text{-}authored},
\mathsf{participant\text{-}revised}
\}.
\]
A transformation is provenance preserving when the type and derivational ancestry of every output object can be recovered from the output record.
\end{definition}

Let
\[
\chi_i
\]
be the provenance ledger for submission \(x_i\). If the system maps \(x_i\) into a compressed state \(m_i\), provenance preservation requires a reconstruction map
\[
R_{\chi}:m_i\longrightarrow \mathcal{P}(\chi_i)
\]
that identifies the source and transformation history of every claim used in a disposition.

This does not require reconstruction of every character of the original submission from \(m_i\). The original object may be stored separately and referenced immutably. It requires that no system inference be silently attributed to the participant or recipient, and that no archive statement be presented as though it were a newly derived conclusion.

\subsection{Context-indexed interpretation}

No submission has one context-independent routing value. Let the routing context for participant \(i\), recipient \(r\), and time \(t\) be
\[
c_{ir}(t)
=
\bigl(
K_r(t),
\theta_i(t),
G_r(t),
\Gamma_i(t),
B_r(t),
\Xi(t)
\bigr).
\]
Here \(K_r(t)\) is the recipient's archive and unresolved-question structure, \(\theta_i(t)\) is the provisional participant-state estimate, \(G_r(t)\) is the prerequisite graph, \(\Gamma_i(t)\) is the interaction and provenance history, \(B_r(t)\) is the recipient's current attention and responsibility boundary, and \(\Xi(t)\) is the broader institutional state relevant to routing.

The same proposition can be answered for one participant, taught to another, escalated for a third, and audited for a fourth. These dispositions need not contradict one another because the authorized target function is context indexed:
\[
f_r:X\times\mathcal{C}\longrightarrow\D_r.
\]

A context-free quotient
\[
q:X\longrightarrow M
\]
is generally too coarse for this purpose. The appropriate object is a family of projections
\[
q_c:X\longrightarrow M_c
\]
indexed by contexts \(c\in\mathcal{C}\), or equivalently a single dependent map
\[
q:
\coprod_{c\in\mathcal C}X_c
\longrightarrow
\coprod_{c\in\mathcal C}M_c.
\]

\begin{definition}[Context-relative disposition equivalence]
For a fixed context \(c\), define
\[
x\sim_c y
\]
when every authorized disposition function agrees on \(x\) and \(y\):
\[
x\sim_c y
\iff
\forall f\in\F_{\mathrm{AC},r},
\quad
f(x,c)=f(y,c).
\]
\end{definition}

The relation \(\sim_c\) is an equivalence relation whenever every \(f\in\F_{\mathrm{AC},r}\) is a well-defined function on \(X\times\{c\}\).

\begin{proposition}
For every fixed context \(c\), the relation \(\sim_c\) is reflexive, symmetric, and transitive.
\end{proposition}

\begin{proof}
For every \(x\in X\) and every \(f\in\F_{\mathrm{AC},r}\),
\[
f(x,c)=f(x,c),
\]
so \(x\sim_c x\), establishing reflexivity.

If \(x\sim_c y\), then
\[
f(x,c)=f(y,c)
\]
for every authorized function \(f\). Equality is symmetric, hence
\[
f(y,c)=f(x,c)
\]
for every \(f\), and therefore \(y\sim_c x\).

If \(x\sim_c y\) and \(y\sim_c z\), then for every authorized \(f\),
\[
f(x,c)=f(y,c)
\]
and
\[
f(y,c)=f(z,c).
\]
By transitivity of equality,
\[
f(x,c)=f(z,c),
\]
so \(x\sim_c z\). Therefore \(\sim_c\) is an equivalence relation.
\end{proof}

The canonical context-relative quotient is
\[
M_c
=
X/{\sim_c},
\qquad
\pi_c(x)
=
[x]_c.
\]

This quotient is exact by construction relative to the declared function class. It is not computable merely because it is definable. Determining whether two submissions receive the same value under every authorized downstream function may itself require the full processing that the quotient was intended to simplify. CLIO must therefore distinguish the canonical quotient from an estimator of that quotient.

\subsection{CLIO as the projection layer}

Let
\[
\CLIO_c:X\longrightarrow \widehat{M}_c
\]
be an implementable approximation to the canonical quotient \(\pi_c\). Its output is a structured packet
\[
m_i
=
\CLIO_{c_{ir}}(x_i)
=
(a_i,p_i,r_i,u_i,\chi_i,\nu_i).
\]
Here \(a_i\) is the portion disposed of by recorded material, \(p_i\) is a candidate prerequisite route, \(r_i\) is the unresolved residual, \(u_i\) is classification uncertainty, \(\chi_i\) is the provenance ledger, and \(\nu_i\) contains novelty coordinates or witnesses not eliminated during normalization.

The map \(\CLIO_c\) is not licensed to collapse two submissions merely because they are semantically similar. Its target is disposition sufficiency.

\begin{definition}[Attention-compiler sufficiency]
A projection \(q_c:X\to M_c\) is sufficient for the Attention Compiler in context \(c\) when
\[
q_c(x)=q_c(y)
\implies
f(x,c)=f(y,c)
\]
for every
\[
f\in\F_{\mathrm{AC},r}.
\]
\end{definition}

Equivalently, every authorized function must factor through \(q_c\). That is, for each \(f\in\F_{\mathrm{AC},r}\), there must exist a function
\[
\bar f_c:M_c\longrightarrow Y_f
\]
such that
\[
f(\,\cdot\,,c)
=
\bar f_c\circ q_c.
\]

\begin{theorem}[Factorization criterion for conversational sufficiency]
Let \(q_c:X\to M_c\) be surjective. The following conditions are equivalent:
\[
q_c(x)=q_c(y)
\implies
f(x,c)=f(y,c)
\]
for every \(x,y\in X\), and the existence of a unique map
\[
\bar f_c:M_c\to Y_f
\]
satisfying
\[
f(\,\cdot\,,c)=\bar f_c\circ q_c.
\]
\end{theorem}

\begin{proof}
Assume first that
\[
q_c(x)=q_c(y)
\implies
f(x,c)=f(y,c).
\]
For \(m\in M_c\), choose any \(x\in X\) such that \(q_c(x)=m\), which is possible because \(q_c\) is surjective, and define
\[
\bar f_c(m)=f(x,c).
\]
To show that this is well defined, suppose \(y\in X\) is another point with \(q_c(y)=m\). Then
\[
q_c(x)=q_c(y),
\]
so the assumed condition gives
\[
f(x,c)=f(y,c).
\]
Thus the value of \(\bar f_c(m)\) does not depend on the chosen representative. For every \(x\in X\),
\[
(\bar f_c\circ q_c)(x)
=
\bar f_c(q_c(x))
=
f(x,c).
\]
Hence the factorization exists.

For uniqueness, suppose
\[
g_c:M_c\to Y_f
\]
also satisfies
\[
f(\,\cdot\,,c)=g_c\circ q_c.
\]
For any \(m\in M_c\), choose \(x\in X\) with \(q_c(x)=m\). Then
\[
g_c(m)
=
g_c(q_c(x))
=
f(x,c)
=
\bar f_c(q_c(x))
=
\bar f_c(m).
\]
Thus \(g_c=\bar f_c\).

Conversely, assume that
\[
f(\,\cdot\,,c)=\bar f_c\circ q_c.
\]
If \(q_c(x)=q_c(y)\), then
\[
f(x,c)
=
\bar f_c(q_c(x))
=
\bar f_c(q_c(y))
=
f(y,c).
\]
Therefore the quotient is sufficient for \(f\).
\end{proof}

The theorem converts an intuitive demand into an architectural requirement. Every authorized operation must consume the projected packet rather than silently recovering discarded features from the raw submission. If one operation bypasses the quotient because it needs information the quotient omitted, then either the operation lies outside the declared function class or the quotient was not sufficient for the architecture actually implemented.

\subsection{HYDRA as the assembly layer}

CLIO alone cannot determine a final disposition. The projected packet must be interpreted relative to archive nodes, prerequisite relations, participant competencies, recipient responsibilities, provenance constraints, and current institutional state.

Let
\[
\C_{ir}
\]
be an indexing category whose objects represent local conversational contexts relevant to submission \(x_i\) and recipient \(r\). Typical objects include an archive neighborhood, a prerequisite neighborhood, a participant-state neighborhood, a provenance neighborhood, a recipient-question neighborhood, and an institutional-governance neighborhood.

Let
\[
\A_{ir}:\C_{ir}\longrightarrow\mathscr{E}
\]
be a diagram in a category \(\mathscr{E}\) of structured epistemic states. Each object \(c\in\Ob(\C_{ir})\) is assigned a local state
\[
\A_{ir}(c),
\]
and each morphism
\[
\alpha:c\to c'
\]
is assigned a compatibility map
\[
\A_{ir}(\alpha):
\A_{ir}(c)\to\A_{ir}(c').
\]

HYDRA attempts to assemble the local states into a global conversational state
\[
Z_{ir}
=
\colim_{\C_{ir}}\A_{ir}.
\]

The use of a colimit means that the global state is not a simple concatenation. Local states are identified according to declared compatibility maps. An archive claim reconstructed as equivalent to a participant claim may be glued to it. A prerequisite node demonstrated by the participant may be identified with a competency-state witness. A provenance conflict may prevent two locally coherent claims from being treated as the same global object.

\begin{definition}[Conversational gluing]
Conversational gluing is the construction of a global state \(Z_{ir}\) together with compatible structure maps
\[
\iota_c:\A_{ir}(c)\longrightarrow Z_{ir}
\]
such that for every morphism \(\alpha:c\to c'\),
\[
\iota_{c'}
\circ
\A_{ir}(\alpha)
=
\iota_c,
\]
and such that \(Z_{ir}\) is universal among objects satisfying these compatibility equations.
\end{definition}

If such a colimit exists in \(\mathscr E\), it is unique up to canonical isomorphism. Its existence in an abstract category does not prove that the intended semantic compatibility maps are correct. A badly designed category can force incompatible claims to become identified, while a weak diagram can omit the morphism that would reveal a contradiction. The mathematical existence of a colimit must therefore be separated from the epistemic admissibility of the diagram used to construct it.

\subsection{The composite architecture}

The integrated Attention Compiler is the composite
\[
\AC(x_i\mid r,t)
=
D_r
\circ
\Hh_t
\circ
\HYDRA_{c_{ir}}
\circ
\CLIO_{c_{ir}}
(x_i),
\]
where
\[
\CLIO_{c_{ir}}:
X\to \widehat M_{c_{ir}},
\]
\[
\HYDRA_{c_{ir}}:
\widehat M_{c_{ir}}
\to Z_{ir},
\]
\[
\Hh_t:
Z_{ir}
\to Z_{ir}',
\]
and
\[
D_r:
Z_{ir}'
\to\D_r.
\]

Expanded, the composite is
\[
\AC(x_i\mid r,t)
=
D_r\!\left[
\Hh_t
\left(
\colim_{\C_{ir}}
\A_{ir}
\left(
\CLIO_{c_{ir}}(x_i),
K_r(t),
\theta_i(t),
G_r(t),
M_{i,0:j}
\right);
\FieldPhi_t,\FieldV_t,\FieldS_t
\right)
\right].
\]

The architecture contains two logically independent compression steps. CLIO compresses the raw submission into a disposition-sufficient packet. HYDRA compresses a distributed family of local contextual states into a global state through compatibility-constrained identification. The first operation risks collapsing distinctions within a fiber. The second risks imposing an invalid gluing or encountering an obstruction to the proposed assembly.

\begin{principle}[Dual accountability]
A disposition produced by the integrated system is warranted only if the submission projection is sufficient for every authorized function used in producing that disposition and the contextual assembly is admissible under the declared compatibility structure.
\end{principle}

This principle yields two proof obligations:
\[
\mathsf{PO}_{\mathrm{proj}}:
\quad
\forall x,y,
\quad
q_c(x)=q_c(y)
\implies
\forall f\in\F_{\mathrm{AC},r},
\quad
f(x,c)=f(y,c),
\]
and
\[
\mathsf{PO}_{\mathrm{glue}}:
\quad
\text{the local states used by }D_r
\text{ admit a coherent and provenance-compatible assembly}.
\]

Neither obligation implies the other.

\begin{proposition}[Independence of projection and gluing correctness]
Projection sufficiency does not imply admissible gluing, and admissible gluing does not imply projection sufficiency.
\end{proposition}

\begin{proof}
For the first direction, suppose \(q_c\) is the identity map on \(X\). It preserves every distinction and is therefore sufficient for every function on \(X\). Let two local archive sections assign contradictory statuses to the same proposition, while the compatibility diagram identifies them without preserving the distinction between assertion and refutation. The projection is sufficient, but the gluing is epistemically inadmissible.

For the second direction, let the contextual diagram contain a single projected object and only its identity morphism. Its colimit exists trivially and is coherent relative to the supplied diagram. Let \(q_c\) collapse two submissions \(x\) and \(y\) such that the correct disposition of \(x\) is archival answer while the correct disposition of \(y\) is escalation. The gluing of the projected object is unobstructed, but the projection is insufficient. Therefore neither condition implies the other.
\end{proof}

The proposition explains why one undifferentiated confidence score cannot govern the architecture. Projection confidence estimates whether the reduced submission preserved disposition-relevant distinctions. Gluing confidence estimates whether contextual structures are mutually compatible. A system may be highly confident about one and uncertain about the other.

\subsection{Attention as an external capacity constraint}

The integrated architecture is not complete until it respects the recipient's finite capacity. Let
\[
E_i
\]
be the event that submission \(i\) is escalated. Let
\[
C_i^{(r)}
\]
be the recipient attention cost of examining the transformed escalation packet. Stability requires
\[
\lambda_r
\Prob(E_i)
\E[C_i^{(r)}\mid E_i]
<
A_r.
\]

This condition is necessary but not sufficient. A system can satisfy it by escalating nothing. Let
\[
V_i
\]
be the event that submission \(i\) contains a residual whose human examination would have positive recipient-relative value, and let
\[
D_i
\]
be the event that it is discarded or delayed beyond a declared horizon. The principal epistemic safety quantity is
\[
\Prob(D_i\mid V_i).
\]

Since \(V_i\) is not directly observable for discarded submissions, the system must estimate this quantity through an exploration and audit policy. Let
\[
A_r=A_r^{(p)}+A_r^{(x)},
\]
where \(A_r^{(p)}\) is allocated to predicted high-value escalations and \(A_r^{(x)}\) is allocated to random, uncertainty-directed, diversity-directed, or appeal-directed exploration.

The architecture must therefore satisfy a joint condition of congestion stability and bounded estimated suppression:
\[
\lambda_r
\Prob(E_i)
\E[C_i^{(r)}\mid E_i]
<
A_r,
\]
\[
\widehat{\Prob}(D_i\mid V_i)
\leq
\eps_{\mathrm{supp}},
\]
and
\[
A_r^{(x)}
\geq
\xi A_r
\]
for a declared exploration fraction \(\xi\in(0,1)\).

These conditions interact with CLIO and HYDRA. Projection collapse increases false suppression by making valuable and nonvaluable residuals indistinguishable. Gluing obstruction increases deferral or misclassification by preventing locally coherent evidence from entering a stable global decision. Exploration provides observations needed to distinguish these failure modes from ordinary prediction error.

\subsection{The main thesis}

The integrated framework can now be stated in its strongest compact form.

\begin{theorem}[Conditional validity of epistemic triage]
Let \(x_i\) be a submission addressed to recipient \(r\) in context \(c_{ir}\). Suppose that the CLIO projection \(q_{c_{ir}}\) is sufficient for the full authorized function class \(\F_{\mathrm{AC},r}\), that the HYDRA diagram \(\A_{ir}\) admits an epistemically admissible and provenance-compatible global assembly, that the transition operator \(\Hh_t\) preserves the admissibility constraints encoded in that assembly, and that the disposition operator \(D_r\) acts only on information available through the warranted composite. Then the disposition
\[
\AC(x_i\mid r,t)
\]
is invariant under replacement of \(x_i\) by any submission in the same context-relative sufficiency class and under replacement of the local diagram by any canonically equivalent admissible presentation.
\end{theorem}

\begin{proof}
Let \(x_i\) and \(y_i\) satisfy
\[
q_{c_{ir}}(x_i)
=
q_{c_{ir}}(y_i).
\]
Because \(q_{c_{ir}}\) is sufficient for every authorized downstream function, every operation used after projection factors through \(q_{c_{ir}}\). Hence the projected packets supplied to the contextual assembly are identical in every disposition-relevant coordinate.

Let \(\A_{ir}\) and \(\A_{ir}'\) be canonically equivalent admissible diagrams presenting the same contextual data and compatibility relations. Their colimits are canonically isomorphic:
\[
\colim_{\C_{ir}}\A_{ir}
\cong
\colim_{\C_{ir}'}\A_{ir}'.
\]
Because the global assembly is assumed provenance compatible, the canonical isomorphism preserves the provenance-typed objects used by the transition and disposition operators.

The transition operator preserves the admissibility constraints encoded in the assembly, and the disposition operator depends only on the warranted composite. Therefore the two computations produce the same disposition:
\[
\AC(x_i\mid r,t)
=
\AC(y_i\mid r,t).
\]
The result is likewise invariant under canonically equivalent admissible presentations of the local diagram.
\end{proof}

The theorem is deliberately conditional. It does not assert that a deployed system automatically satisfies its premises. It identifies what would have to be shown before the result of intellectual triage could be treated as structurally warranted rather than merely predicted.

The rest of this treatise develops those premises separately. Section 2 constructs the context-relative sufficiency theory underlying CLIO. Section 3 derives collapse observables and false-suppression bounds. Section 4 develops participant-state estimation and prerequisite-sensitive repair. Section 5 formalizes HYDRA assembly and distinguishes successful gluing from forced identification. Section 6 develops parameter-mediated operational coupling through PERSCEN, RAT, CoM, and RSVP field semantics. Section 7 derives obstruction-sensitive routing. Section 8 distinguishes repair from representational expansion. Section 9 develops recipient-relative escalation and stable attention allocation. Section 10 gives audit and appeal estimators for unobservable suppression. Section 11 studies distributional novelty and collective recurrence. Section 12 states governance conditions and impossibility results. Section 13 assembles the full architecture and proves the principal conditional guarantees.

\section{Context-Relative Sufficiency and the Geometry of CLIO}

\subsection{Representational simplicity as relative invariance}

The phrase \emph{representational simplicity} does not refer merely to a small number of coordinates, a short description, a low-dimensional embedding, or a compressed storage format. A representation is simple only relative to a declared family of operations. Removing a coordinate is simplifying when every authorized operation is invariant under its removal. The same removal is destructive when at least one authorized operation distinguishes states that the representation identifies.

Let \(X\) be a measurable state space of submissions and let \(\mathcal C\) be a space of conversational contexts. A context
\[
c\in\mathcal C
\]
contains every external condition under which representational sufficiency is evaluated. In the Attention Compiler, a context may be written
\[
c
=
(K_r,\theta_i,G_r,\Gamma_i,B_r,\Xi,t),
\]
where \(K_r\) is the recipient-relative archive, \(\theta_i\) is the provisional participant-state estimate, \(G_r\) is the prerequisite graph, \(\Gamma_i\) is the interaction and provenance history, \(B_r\) is the recipient's responsibility and attention boundary, \(\Xi\) is the institutional state, and \(t\) is time.

Let
\[
\F_c
=
\{f_{\alpha,c}:X\to Y_\alpha\}_{\alpha\in I_c}
\]
be the family of downstream functions authorized in context \(c\). The codomain \(Y_\alpha\) may differ across functions. One function may return a disposition, another a prerequisite path, another a novelty vector, another a provenance judgment, and another a probability distribution over audit outcomes.

A representation
\[
q_c:X\to M_c
\]
is simple relative to \(\F_c\) when it preserves exactly the distinctions required by the functions in \(\F_c\), subject to whatever approximation tolerance and governance constraints have been declared.

\begin{definition}[Exact contextual sufficiency]
A map \(q_c:X\to M_c\) is exactly sufficient for \(\F_c\) when
\[
q_c(x)=q_c(y)
\implies
f_{\alpha,c}(x)=f_{\alpha,c}(y)
\]
for every \(x,y\in X\) and every \(\alpha\in I_c\).
\end{definition}

\begin{definition}[Contextual indistinguishability]
For a fixed context \(c\), define
\[
x\equiv_c y
\iff
\forall \alpha\in I_c,
\quad
f_{\alpha,c}(x)=f_{\alpha,c}(y).
\]
The equivalence class of \(x\) is denoted
\[
[x]_c
=
\{y\in X:y\equiv_c x\}.
\]
\end{definition}

The relation \(\equiv_c\) expresses operational indistinguishability under the complete authorized function family. It is finer than equality of final dispositions whenever intermediate functions require distinctions not visible in the final output. Two submissions may both be escalated while requiring different provenance packets, different prerequisite histories, or different uncertainty reports. If those differences are used by an authorized operation, then the submissions do not belong to the same \(\equiv_c\)-class even though the final disposition agrees.

Let
\[
M_c^{*}
=
X/{\equiv_c}
\]
be the canonical contextual quotient, and let
\[
\pi_c^{*}:X\to M_c^{*},
\qquad
\pi_c^{*}(x)=[x]_c.
\]

\begin{theorem}[Canonical sufficiency quotient]
For every fixed context \(c\), the quotient map
\[
\pi_c^{*}:X\to M_c^{*}
\]
is exactly sufficient for \(\F_c\). Every exactly sufficient map
\[
q_c:X\to M_c
\]
factors through a refinement of \(\pi_c^{*}\). More precisely, there exists a unique map
\[
h_c:q_c(X)\to M_c^{*}
\]
such that
\[
\pi_c^{*}
=
h_c\circ q_c.
\]
\end{theorem}

\begin{proof}
Suppose
\[
\pi_c^{*}(x)=\pi_c^{*}(y).
\]
Then
\[
[x]_c=[y]_c,
\]
so \(x\equiv_c y\). By definition,
\[
f_{\alpha,c}(x)=f_{\alpha,c}(y)
\]
for every \(\alpha\in I_c\). Therefore \(\pi_c^{*}\) is exactly sufficient.

Now let \(q_c:X\to M_c\) be exactly sufficient. Define
\[
h_c(q_c(x))
=
[x]_c.
\]
To prove that \(h_c\) is well defined, suppose
\[
q_c(x)=q_c(y).
\]
Exact sufficiency of \(q_c\) implies
\[
f_{\alpha,c}(x)=f_{\alpha,c}(y)
\]
for every \(\alpha\), hence \(x\equiv_c y\), and consequently
\[
[x]_c=[y]_c.
\]
Thus the value of \(h_c\) does not depend on the chosen representative.

For every \(x\in X\),
\[
(h_c\circ q_c)(x)
=
h_c(q_c(x))
=
[x]_c
=
\pi_c^{*}(x).
\]
Hence
\[
\pi_c^{*}=h_c\circ q_c.
\]

For uniqueness, suppose \(g_c:q_c(X)\to M_c^{*}\) also satisfies
\[
\pi_c^{*}=g_c\circ q_c.
\]
For every \(m\in q_c(X)\), choose \(x\in X\) with \(q_c(x)=m\). Then
\[
g_c(m)
=
g_c(q_c(x))
=
\pi_c^{*}(x)
=
h_c(q_c(x))
=
h_c(m).
\]
Therefore \(g_c=h_c\).
\end{proof}

The factorization direction is important. An exactly sufficient representation \(q_c\) may retain distinctions beyond those required by \(\F_c\). The map \(h_c\) removes these unnecessary distinctions and reaches the canonical quotient. Thus \(q_c\) refines \(\pi_c^{*}\), while \(\pi_c^{*}\) is the coarsest exactly sufficient partition of \(X\).

\begin{definition}[Exact representational simplicity]
An exactly sufficient representation \(q_c:X\to M_c\) is exactly simple relative to \(\F_c\) when
\[
q_c(x)=q_c(y)
\iff
x\equiv_c y.
\]
\end{definition}

\begin{corollary}[Uniqueness up to relabeling]
If \(q_c:X\to M_c\) and \(q_c':X\to M_c'\) are both surjective and exactly simple relative to \(\F_c\), then there exists a unique bijection
\[
\varphi_c:M_c\to M_c'
\]
such that
\[
q_c'
=
\varphi_c\circ q_c.
\]
\end{corollary}

\begin{proof}
Because \(q_c\) and \(q_c'\) identify exactly the same pairs of points, define
\[
\varphi_c(q_c(x))
=
q_c'(x).
\]
If \(q_c(x)=q_c(y)\), then \(x\equiv_c y\), so \(q_c'(x)=q_c'(y)\). Hence \(\varphi_c\) is well defined.

Surjectivity of \(q_c'\) gives surjectivity of \(\varphi_c\). If
\[
\varphi_c(q_c(x))
=
\varphi_c(q_c(y)),
\]
then
\[
q_c'(x)=q_c'(y),
\]
which implies \(x\equiv_c y\), and therefore \(q_c(x)=q_c(y)\). Thus \(\varphi_c\) is injective. The identity
\[
q_c'=\varphi_c\circ q_c
\]
holds by construction. Uniqueness follows from surjectivity of \(q_c\).
\end{proof}

This result provides the precise sense in which an exact simple representation is canonical. Its coordinates, labels, serialization, and internal arrangement may vary, but its partition of the source space cannot vary without either retaining unnecessary distinctions or collapsing necessary ones.

\subsection{The lattice of contextual partitions}

Let \(\operatorname{Part}(X)\) denote the set of partitions of \(X\). For partitions \(P,Q\in\operatorname{Part}(X)\), write
\[
P\preceq Q
\]
when \(P\) refines \(Q\). Thus every block of \(P\) is contained in a block of \(Q\). A finer partition retains more distinctions, while a coarser partition identifies more states.

Every representation \(q:X\to M\) induces a partition
\[
P_q
=
\{q^{-1}(m):m\in q(X)\}.
\]
Every authorized function \(f_{\alpha,c}\) likewise induces
\[
P_{\alpha,c}
=
\{f_{\alpha,c}^{-1}(y):y\in f_{\alpha,c}(X)\}.
\]

The canonical sufficiency partition is the common refinement
\[
P_c^{*}
=
\bigwedge_{\alpha\in I_c}P_{\alpha,c},
\]
where two states occupy the same block exactly when every authorized function agrees on them.

\begin{proposition}[Partition criterion]
A representation \(q_c\) is exactly sufficient for \(\F_c\) if and only if
\[
P_{q_c}\preceq P_c^{*}.
\]
It is exactly simple if and only if
\[
P_{q_c}=P_c^{*}.
\]
\end{proposition}

\begin{proof}
Suppose \(P_{q_c}\preceq P_c^{*}\). If \(q_c(x)=q_c(y)\), then \(x\) and \(y\) belong to the same block of \(P_{q_c}\). Since this block lies within a block of \(P_c^{*}\), every authorized function agrees on \(x\) and \(y\). Thus \(q_c\) is sufficient.

Conversely, suppose \(q_c\) is sufficient. Every fiber \(q_c^{-1}(m)\) contains only states on which every authorized function agrees. Therefore every fiber lies within a block of \(P_c^{*}\), giving
\[
P_{q_c}\preceq P_c^{*}.
\]

Exact simplicity requires both sufficiency and the absence of distinctions beyond those required by \(\F_c\). This is equivalent to equality of the induced partitions.
\end{proof}

The partition lattice reveals two opposed errors. If
\[
P_{q_c}\not\preceq P_c^{*},
\]
then at least one fiber crosses a required decision boundary. This is projection collapse. If
\[
P_{q_c}\prec P_c^{*},
\]
then \(q_c\) is sufficient but unnecessarily fine. This is representational surplus.

\begin{definition}[Projection collapse]
A pair \((x,y)\) is collapsed by \(q_c\) when
\[
q_c(x)=q_c(y)
\]
but
\[
x\not\equiv_c y.
\]
\end{definition}

\begin{definition}[Representational surplus]
A pair \((x,y)\) exhibits representational surplus under \(q_c\) when
\[
x\equiv_c y
\]
but
\[
q_c(x)\neq q_c(y).
\]
\end{definition}

Collapse threatens correctness. Surplus threatens efficiency, privacy, portability, and generalization. A system that retains participant identity, institutional status, writing style, or irrelevant biographical attributes after every authorized function has been declared invariant to them is not representationally simple even if its decisions happen to be correct.

\begin{principle}[No gratuitous distinction]
If a distinction cannot alter any authorized disposition, explanation, audit, appeal, or provenance obligation in context \(c\), then retaining that distinction inside the operational representation requires an independent justification.
\end{principle}

This principle is not equivalent to deleting the original submission. The original may remain in an immutable provenance store. It states that irrelevant distinctions should not remain active coordinates in the representation used to make decisions.

\subsection{Context refinement and the instability of equivalence}

Contexts change. An answer recorded at time \(t\) may be superseded at time \(t'\). A recipient may acquire a new responsibility. A participant may demonstrate a competency previously uncertain. A new authorized function may be added during governance review.

Let
\[
c\preceq_{\mathcal C}c'
\]
mean that \(c'\) is a functional refinement of \(c\). Formally, assume
\[
\F_c\subseteq\F_{c'}.
\]
The refined context imposes every distinction required previously and possibly additional distinctions.

\begin{proposition}[Monotonic refinement of contextual equivalence]
If
\[
\F_c\subseteq\F_{c'},
\]
then
\[
x\equiv_{c'}y
\implies
x\equiv_c y.
\]
Equivalently,
\[
P_{c'}^{*}\preceq P_c^{*}.
\]
\end{proposition}

\begin{proof}
If \(x\equiv_{c'}y\), then every function in \(\F_{c'}\) agrees on \(x\) and \(y\). Since every function in \(\F_c\) belongs to \(\F_{c'}\), every function in \(\F_c\) also agrees. Therefore \(x\equiv_c y\). The partition relation follows directly.
\end{proof}

Adding authorized operations can split equivalence classes but cannot merge classes unless some former requirements are withdrawn. A representation sufficient under \(c\) need not remain sufficient under \(c'\).

\begin{corollary}[Obsolescence under function-class expansion]
Let \(q_c\) be exactly simple for \(\F_c\). If there exist \(x,y\in X\) such that
\[
x\equiv_c y
\]
but
\[
x\not\equiv_{c'}y,
\]
then \(q_c\) is not sufficient for \(\F_{c'}\).
\end{corollary}

\begin{proof}
Exact simplicity under \(c\) gives
\[
x\equiv_c y
\implies
q_c(x)=q_c(y).
\]
Since \(x\not\equiv_{c'}y\), some function in \(\F_{c'}\) distinguishes them. Thus \(q_c\) collapses a distinction required under \(c'\).
\end{proof}

This is the formal origin of representational debt. The representation was not necessarily defective when constructed. It may have been exactly simple for the declared function class at that time. The debt appears when the system begins using the representation for a larger function class without expanding the quotient.

\begin{definition}[Representational debt]
Given contexts \(c\preceq_{\mathcal C}c'\), the representational debt of \(q_c\) relative to \(c'\) is the family of distinctions
\[
\mathfrak D(q_c;c')
=
\left\{
(x,y):
q_c(x)=q_c(y),
\;
x\not\equiv_{c'}y
\right\}.
\]
\end{definition}

A system incurs operational debt when it deploys \(q_c\) under \(c'\) without testing whether
\[
\mathfrak D(q_c;c')=\varnothing.
\]

\subsection{Recipient relativity}

Let \(r\) and \(s\) be distinct recipients. Their archives, responsibilities, expertise, unresolved questions, and attention boundaries may differ. Consequently,
\[
\F_{\mathrm{AC},r}
\neq
\F_{\mathrm{AC},s}.
\]
A quotient sufficient for one recipient may be insufficient for the other.

Let
\[
c_{ir}
\]
and
\[
c_{is}
\]
be the corresponding contexts for the same participant \(i\) and submission \(x_i\).

\begin{example}[Recipient-relative novelty]
Suppose a submission contains proposition \(q\). Recipient \(r\) has already proved and archived \(q\), while recipient \(s\) has not encountered it and is actively working on a problem for which \(q\) supplies a missing lemma. For \(r\), the relevant residual may be empty:
\[
\Res(x_i\mid K_r)=\varnothing.
\]
For \(s\),
\[
\Res(x_i\mid K_s)=\{q\}.
\]
The same raw submission therefore belongs to different context-relative equivalence classes:
\[
[x_i]_{c_{ir}}
\neq
[x_i]_{c_{is}}.
\]
\end{example}

Recipient relativity prevents the system from treating intellectual value as an intrinsic scalar attached to a message. Let
\[
V_{ir}
\]
denote the value realized if participant \(i\)'s transformed residual is examined by recipient \(r\). The system can estimate only
\[
\widehat V_{ir}
=
\E[V_{ir}\mid I_i,K_r,\theta_i,t].
\]
No recipient-independent value \(\widehat V_i\) can substitute for \(\widehat V_{ir}\) unless an additional invariance argument is supplied.

The same applies to sufficiency. A global quotient
\[
q:X\to M
\]
sufficient for every possible recipient must refine the meet of every recipient-specific partition:
\[
P_q
\preceq
\bigwedge_{r\in\mathcal R}
P_r^{*}.
\]
As the recipient class grows, this meet may approach the identity partition, eliminating meaningful compression.

\begin{theorem}[Universal-recipient degeneration]
Let \(\mathcal R\) be a recipient class. Suppose that for every pair of distinct submissions \(x\neq y\), there exists some recipient \(r\in\mathcal R\) and some authorized function
\[
f\in\F_{\mathrm{AC},r}
\]
such that
\[
f(x,c_{ir})\neq f(y,c_{yr}).
\]
Then every representation sufficient for all recipients in \(\mathcal R\) is injective.
\end{theorem}

\begin{proof}
Let \(q:X\to M\) be sufficient for every recipient. Suppose
\[
q(x)=q(y).
\]
If \(x\neq y\), the hypothesis provides a recipient \(r\) and an authorized function \(f\) distinguishing them. Sufficiency for \(r\) would require
\[
q(x)=q(y)
\implies
f(x,c_{ir})=f(y,c_{yr}),
\]
which contradicts the choice of \(r\) and \(f\). Therefore \(x=y\), so \(q\) is injective.
\end{proof}

The theorem shows that context-free universal sufficiency can destroy the possibility of compression. CLIO must therefore be indexed by the authorized recipient and task context, or it must retain enough latent structure to refine its output when a new context is supplied.

\subsection{Conditional and adaptive projections}

A context-indexed family
\[
\{q_c:X\to M_c\}_{c\in\mathcal C}
\]
can be represented as one conditional map
\[
q:X\times\mathcal C\to\mathcal M,
\]
where
\[
q(x,c)=q_c(x).
\]
The codomain \(\mathcal M\) may be a disjoint union
\[
\mathcal M
=
\coprod_{c\in\mathcal C}M_c
\]
or a common ambient space containing context-specific subspaces.

A conditional projection raises a new sufficiency condition. It must preserve distinctions not only among submissions within one context but also among contexts whenever contextual variation changes an authorized result.

\begin{definition}[Joint context sufficiency]
A conditional representation
\[
q:X\times\mathcal C\to\mathcal M
\]
is jointly sufficient for a function family \(\F\) on \(X\times\mathcal C\) when
\[
q(x,c)=q(y,c')
\implies
f(x,c)=f(y,c')
\]
for every \(f\in\F\).
\end{definition}

The system need not require joint sufficiency if context identifiers remain available outside the quotient. It may instead use a pair
\[
(q_c(x),c)
\]
as the operational representation. This makes explicit that the same projected coordinate in two contexts does not imply interchangeable treatment.

\begin{proposition}[Indexed sufficiency]
Suppose each
\[
q_c:X\to M_c
\]
is sufficient for \(\F_c\). Then
\[
\widetilde q(x,c)
=
(c,q_c(x))
\]
is sufficient for the union of context-indexed functions
\[
\F
=
\bigcup_{c\in\mathcal C}
\{(x,c)\mapsto f_c(x):f_c\in\F_c\}.
\]
\end{proposition}

\begin{proof}
Suppose
\[
\widetilde q(x,c)
=
\widetilde q(y,c').
\]
Equality of ordered pairs gives \(c=c'\) and
\[
q_c(x)=q_c(y).
\]
For every \(f_c\in\F_c\), sufficiency of \(q_c\) implies
\[
f_c(x)=f_c(y).
\]
Thus every context-indexed authorized function agrees.
\end{proof}

This construction is simple but potentially expensive. Carrying the full context identifier can retain more information than necessary. A further quotient over contexts is possible when two contexts induce identical authorized partitions and identical downstream behavior.

\begin{definition}[Operational context equivalence]
Contexts \(c,c'\in\mathcal C\) are operationally equivalent, written
\[
c\approx_{\mathcal C}c',
\]
when there exists a codomain-compatible identification of their function families such that
\[
f_{\alpha,c}(x)
=
f_{\alpha,c'}(x)
\]
for every submission \(x\) and every corresponding authorized operation \(\alpha\).
\end{definition}

Contexts may then be represented by their equivalence classes
\[
[c]_{\mathcal C},
\]
provided provenance, governance, and temporal obligations do not independently require the original context identity.

\subsection{Approximate sufficiency}

Exact sufficiency is often unattainable in learned systems. The source space may be continuous, the target functions may be uncertain, the archive may be incomplete, and the correct disposition may itself be contestable. Approximate sufficiency must therefore be defined relative to losses and distributions rather than asserted through vague semantic preservation.

Let
\[
P_c
\]
be a probability measure on \(X\) in context \(c\). For each authorized function \(f_{\alpha,c}\), let
\[
L_\alpha:
Y_\alpha\times Y_\alpha\to[0,\infty]
\]
be a declared loss. Given a representation \(q_c:X\to M_c\), define its optimal reconstruction risk for \(f_{\alpha,c}\) as
\[
\mathcal R_{\alpha}(q_c\mid c)
=
\inf_{g_\alpha:M_c\to Y_\alpha}
\E_{X\sim P_c}
\left[
L_\alpha
\left(
f_{\alpha,c}(X),
g_\alpha(q_c(X))
\right)
\right].
\]

\begin{definition}[Approximate contextual sufficiency]
The representation \(q_c\) is \(\bm{\eps}\)-sufficient for
\[
\F_c
=
\{f_{\alpha,c}\}_{\alpha\in I_c}
\]
when
\[
\mathcal R_\alpha(q_c\mid c)
\leq
\eps_\alpha
\]
for every \(\alpha\in I_c\), where
\[
\bm{\eps}
=
(\eps_\alpha)_{\alpha\in I_c}
\]
is a declared tolerance profile.
\end{definition}

A single scalar tolerance can conceal asymmetrical risks. Confusing two equivalent prerequisite explanations may be low cost. Confusing archival answer with escalation may be much more serious. The tolerance should therefore remain operation specific.

For weighted operation priorities \(w_\alpha(c)\geq0\), one may define aggregate risk
\[
\mathcal R_{\F}(q_c\mid c)
=
\sum_{\alpha\in I_c}
w_\alpha(c)
\mathcal R_\alpha(q_c\mid c),
\]
when the index set is finite or the sum converges. Aggregate risk does not replace per-operation constraints when some failures are categorically prohibited.

\begin{definition}[Protected operation]
An authorized function \(f_{\alpha,c}\) is protected when its tolerance is fixed at
\[
\eps_\alpha=0
\]
or when its loss is constrained by a separate hard bound that cannot be traded against improvements in other operations.
\end{definition}

Provenance attribution, appeal availability, and certain safety classifications are natural protected operations. A reduction in average routing cost cannot compensate for fabricated attribution.

\subsection{Fiberwise variation and the collapse functional}

Approximate risk can hide where insufficiency occurs. CLIO therefore requires a fiberwise diagnostic.

For \(m\in M_c\), let
\[
\mu_{c,m}
\]
be a regular conditional distribution over the fiber
\[
q_c^{-1}(m).
\]
For an authorized function \(f_{\alpha,c}\), let
\[
P_{\alpha,c}(\,\cdot\mid x)
\]
be the distribution of its warranted output given the full submission \(x\). This distribution may be degenerate when the warranted output is deterministic.

Define the fiberwise disposition divergence
\[
\Lambda_{\alpha,c}(m)
=
\E_{\substack{
X_1,X_2\sim\mu_{c,m}
}}
\left[
\KL
\left(
P_{\alpha,c}(\,\cdot\mid X_1)
\middle\|
P_{\alpha,c}(\,\cdot\mid X_2)
\right)
\right].
\]

A symmetrized form avoids the asymmetry of \(\KL\):
\[
\Lambda_{\alpha,c}^{\mathrm{sym}}(m)
=
\frac12
\E_{X_1,X_2\sim\mu_{c,m}}
\left[
\KL(P_1\|P_2)
+
\KL(P_2\|P_1)
\right],
\]
where
\[
P_j
=
P_{\alpha,c}(\,\cdot\mid X_j).
\]

Alternatively, Jensen-Shannon divergence gives a bounded quantity:
\[
\Lambda_{\alpha,c}^{\mathrm{JS}}(m)
=
\E_{X_1,X_2\sim\mu_{c,m}}
\left[
\operatorname{JS}(P_1,P_2)
\right].
\]

\begin{proposition}[Fiber homogeneity criterion]
Assume the warranted output of \(f_{\alpha,c}\) is deterministic. Then
\[
\Lambda_{\alpha,c}^{\mathrm{JS}}(m)=0
\]
if and only if \(f_{\alpha,c}\) is constant almost surely on the fiber \(q_c^{-1}(m)\).
\end{proposition}

\begin{proof}
For deterministic output, each conditional distribution is a point mass:
\[
P_{\alpha,c}(\,\cdot\mid x)
=
\delta_{f_{\alpha,c}(x)}.
\]
The Jensen-Shannon divergence between two point masses is zero exactly when the point masses are identical, which occurs exactly when their support values agree. Thus
\[
\operatorname{JS}
\left(
\delta_{f_{\alpha,c}(X_1)},
\delta_{f_{\alpha,c}(X_2)}
\right)
=
0
\]
exactly when
\[
f_{\alpha,c}(X_1)
=
f_{\alpha,c}(X_2).
\]
The expectation is zero if and only if this equality holds almost surely for pairs drawn from the fiber.
\end{proof}

Define the total protected collapse risk as
\[
\Lambda_c^{\mathrm{prot}}
=
\sum_{\alpha\in I_c^{\mathrm{prot}}}
w_\alpha
\E_{M\sim q_{c\#}P_c}
\left[
\Lambda_{\alpha,c}^{\mathrm{JS}}(M)
\right],
\]
where \(q_{c\#}P_c\) is the pushforward of \(P_c\) through \(q_c\).

Exact sufficiency for every protected deterministic operation implies
\[
\Lambda_c^{\mathrm{prot}}=0.
\]
The converse holds up to null sets when all relevant fiberwise conditional distributions exist and all protected functions receive positive weight.

\subsection{Representational simplicity as constrained optimization}

Let \(\mathfrak Q_c\) be a hypothesis class of candidate context-relative projections. Let
\[
\Comp(q_c)
\]
measure operational representational complexity. Depending on the architecture, this may involve dimension, description length, mutual information, computational cost, memory cost, or the number of active distinctions.

The CLIO objective should not minimize complexity without constraints. Unconstrained minimization yields the constant representation
\[
q_c(x)=m_0,
\]
which is maximally simple in storage and maximally destructive whenever any authorized function is nonconstant.

The appropriate optimization is
\[
q_c^{*}
\in
\argmin_{q\in\mathfrak Q_c}
\Comp(q)
\]
subject to
\[
\mathcal R_\alpha(q\mid c)
\leq
\eps_\alpha
\qquad
\forall\alpha\in I_c,
\]
\[
\Lambda_{\alpha,c}(m)
\leq
\delta_{\alpha,c}(m)
\qquad
\text{for protected fibers},
\]
and
\[
\operatorname{ProvLoss}(q)=0.
\]

A Lagrangian relaxation is
\[
\mathcal L_{\CLIO}(q;c)
=
\Comp(q)
+
\sum_{\alpha\in I_c}
\lambda_\alpha
\mathcal R_\alpha(q\mid c)
+
\sum_{\alpha\in I_c^{\mathrm{prot}}}
\gamma_\alpha
\E_M[\Lambda_{\alpha,c}(M)]
+
\beta\operatorname{ProvLoss}(q).
\]

The multipliers are institutional commitments, not natural constants. Choosing them determines which forms of error the architecture is willing to exchange for simplicity.

\begin{proposition}[Constant-quotient criterion]
The constant representation is sufficient for \(\F_c\) if and only if every authorized function in \(\F_c\) is constant on \(X\).
\end{proposition}

\begin{proof}
If the constant representation is sufficient, then for all \(x,y\in X\),
\[
q_c(x)=q_c(y).
\]
Sufficiency implies
\[
f_{\alpha,c}(x)=f_{\alpha,c}(y)
\]
for every authorized function, so every function is constant.

Conversely, if every authorized function is constant, then any two submissions receive identical output under every function. Therefore identifying all submissions in one fiber loses no authorized distinction.
\end{proof}

This proposition formalizes why compression alone cannot be the objective. Simplicity is meaningful only after the authorized distinction structure has been fixed.

\subsection{Information complexity and its limitations}

When \(X\) and \(M_c\) are random variables with a specified joint distribution induced by \(P_c\) and \(q_c\), one candidate complexity measure is
\[
\Comp_{\mathrm{MI}}(q_c)
=
I(X;q_c(X)).
\]
For finite discrete spaces and deterministic \(q_c\),
\[
I(X;q_c(X))
=
H(q_c(X)).
\]
Minimizing mutual information subject to sufficiency constraints seeks the least informative operational representation that still supports the authorized function class.

Let
\[
F_c(X)
=
\left(
f_{\alpha,c}(X)
\right)_{\alpha\in I_c}
\]
denote the joint authorized-output variable when this product is well defined.

\begin{theorem}[Information lower bound]
If \(q_c\) is exactly sufficient for the joint deterministic authorized-output variable \(F_c(X)\), then
\[
I(X;q_c(X))
\geq
I(X;F_c(X)).
\]
If \(q_c\) is exactly simple and differs from \(F_c(X)\) only by a bijective relabeling of its attainable joint outputs, then equality holds.
\end{theorem}

\begin{proof}
Exact sufficiency implies that there exists a deterministic map \(g_c\) satisfying
\[
F_c(X)=g_c(q_c(X)).
\]
Thus
\[
X\longrightarrow q_c(X)\longrightarrow F_c(X)
\]
forms a Markov chain. By the data-processing inequality,
\[
I(X;F_c(X))
\leq
I(X;q_c(X)).
\]

If \(q_c(X)\) and \(F_c(X)\) are related by a measurable bijection on their supports, each is a deterministic function of the other. Applying data processing in both directions gives equality.
\end{proof}

Mutual information requires a specified probability law. It does not define representational simplicity independently of how often states occur. A rare distinction can contribute little to average information while remaining essential to a protected operation. The fiberwise and per-operation constraints therefore cannot be replaced by mutual-information minimization.

For continuous deterministic projections, mutual information may be infinite or sensitive to measure-theoretic conventions. In such settings one may use rate-distortion quantities, discretized operational partitions, minimum description length, or computational complexity, but every substitute must preserve the distinction between average compressibility and protected correctness.

\subsection{Randomized projections}

A learned representation may be stochastic:
\[
Q_c(dm\mid x),
\]
a Markov kernel from \(X\) to \(M_c\). The downstream operation then receives
\[
M\sim Q_c(\,\cdot\mid X).
\]

Exact pointwise sufficiency is too strong when different representations may be sampled for the same input. A natural statistical condition is conditional independence:
\[
F_c(X)\perp X\mid M.
\]
Equivalently, once \(M\) is known, the original submission supplies no additional information about the authorized output variable.

\begin{definition}[Stochastic contextual sufficiency]
A kernel \(Q_c(dm\mid x)\) is stochastically sufficient for an authorized output \(Y_c=F_c(X)\) when
\[
P(Y_c\in B\mid X,M)
=
P(Y_c\in B\mid M)
\]
almost surely for every measurable set \(B\).
\end{definition}

When \(Y_c\) is a deterministic function of \(X\), stochastic sufficiency requires that every representation value \(m\) receiving positive probability from two inputs \(x\) and \(y\) be compatible with the same authorized output, except on a null set.

\begin{proposition}[Support separation]
Let \(Y_c=f_c(X)\) be deterministic. If \(Q_c\) is stochastically sufficient for \(Y_c\), then for almost every pair \(x,y\) with
\[
f_c(x)\neq f_c(y),
\]
the representation kernels have disjoint support up to null sets:
\[
Q_c(\,\cdot\mid x)
\perp
Q_c(\,\cdot\mid y).
\]
\end{proposition}

\begin{proof}
Suppose there exists a measurable representation region \(B\subseteq M_c\) with positive probability under both kernels. Then observing \(M\in B\) would remain compatible with two different deterministic values of \(Y_c\). Knowledge of \(X\) would resolve which value is correct, contradicting
\[
Y_c\perp X\mid M.
\]
Therefore inputs with distinct authorized outputs cannot assign positive probability to a common representation region, except on null sets.
\end{proof}

Randomization does not permit the system to blur protected decision boundaries without cost. It can support privacy, exploration, or robust estimation only while maintaining separation of outputs that must remain distinguishable.

\subsection{Provenance as a nonquotientable coordinate}

Some information is required not because it changes the substantive answer but because it changes what the system may truthfully claim about that answer. Let
\[
\operatorname{prov}_c(x)
\]
be the provenance structure of submission \(x\). Two submissions may express the same normalized proposition while one reports direct observation and the other repeats an archive inference. If the system must distinguish those sources, provenance belongs to the authorized function class.

Let
\[
f_{\mathrm{prov},c}(x)
=
\operatorname{prov}_c(x).
\]
Any sufficient quotient must satisfy
\[
q_c(x)=q_c(y)
\implies
\operatorname{prov}_c(x)
=
\operatorname{prov}_c(y),
\]
unless provenance is retained through a separate channel available to every operation that requires it.

This yields two architectures. In an internal-provenance architecture,
\[
q_c(x)
=
(\bar q_c(x),\operatorname{prov}_c(x)).
\]
In an external-ledger architecture,
\[
q_c(x)
=
\bar q_c(x),
\]
while an immutable reference
\[
\ell(x)
\]
permits authorized functions to retrieve provenance from a separate ledger.

\begin{proposition}[External-ledger sufficiency]
Suppose the operational representation is
\[
\widetilde q_c(x)
=
(\bar q_c(x),\ell(x)),
\]
and every provenance-sensitive function can recover
\[
\operatorname{prov}_c(x)
\]
from \(\ell(x)\). Then \(\bar q_c\) may collapse provenance distinctions without making \(\widetilde q_c\) insufficient, provided no authorized operation receives \(\bar q_c(x)\) without access to \(\ell(x)\).
\end{proposition}

\begin{proof}
The relevant quotient is \(\widetilde q_c\), not \(\bar q_c\) alone. If
\[
\widetilde q_c(x)
=
\widetilde q_c(y),
\]
then
\[
\bar q_c(x)=\bar q_c(y)
\]
and
\[
\ell(x)=\ell(y).
\]
Recoverability from the ledger implies equality of the provenance values required downstream. If \(\bar q_c\) preserves every remaining authorized distinction, then the pair is sufficient. If any operation receives only \(\bar q_c(x)\), provenance-sensitive factorization fails for that operation.
\end{proof}

This proposition identifies a general architectural rule: information may be removed from one representation only if every operation that needs it has reliable access to another representation in which it remains preserved.

\subsection{CLIO estimation}

The canonical quotient \(\pi_c^{*}\) is defined by the true authorized functions. A deployed system observes only training data, archive judgments, audit results, appeals, and delayed recipient evaluations. CLIO must estimate a sufficient partition from incomplete evidence.

Let
\[
\mathcal S_n
=
\{(x_j,c_j,\mathbf y_j)\}_{j=1}^{n}
\]
be an observed corpus, where
\[
\mathbf y_j
=
(y_{\alpha,j})_{\alpha\in I_{c_j}^{\mathrm{obs}}}
\]
contains whichever authorized outputs were observed for example \(j\). Let
\[
\Pi
\]
be a hypothesis class of conditional projections
\[
q_\psi:X\times\mathcal C\to\mathcal M,
\qquad
\psi\in\Psi.
\]

Define the empirical objective
\[
\widehat{\mathcal L}_n(\psi)
=
\widehat{\Comp}_n(q_\psi)
+
\sum_{\alpha}
\lambda_\alpha
\widehat{\mathcal R}_{\alpha,n}(q_\psi)
+
\sum_{\alpha\in I^{\mathrm{prot}}}
\gamma_\alpha
\widehat{\Lambda}_{\alpha,n}(q_\psi)
+
\beta
\widehat{\operatorname{ProvLoss}}_n(q_\psi).
\]

Let
\[
\mathcal L(\psi)
=
\E[\widehat{\mathcal L}_n(\psi)]
\]
denote the corresponding population objective under the sampling process.

A consistency claim requires more than uniform convergence.

\begin{assumption}[Metric compactness]
The parameter space \((\Psi,d_\Psi)\) is compact, or the relevant sublevel sets of \(\mathcal L\) are compact.
\end{assumption}

\begin{assumption}[Uniform convergence]
\[
\sup_{\psi\in\Psi}
\left|
\widehat{\mathcal L}_n(\psi)
-
\mathcal L(\psi)
\right|
\xrightarrow{p}
0.
\]
\end{assumption}

\begin{assumption}[Identification]
The population objective has a unique minimizer
\[
\psi^{*}
\]
up to a declared equivalence relation of representation relabelings.
\end{assumption}

\begin{assumption}[Separation]
For every neighborhood \(U\) of the equivalence class \([\psi^{*}]\), there exists \(\delta_U>0\) such that
\[
\inf_{\psi\notin U}
\mathcal L(\psi)
\geq
\mathcal L(\psi^{*})+\delta_U.
\]
\end{assumption}

\begin{theorem}[Consistency of CLIO estimation]
Under metric compactness, uniform convergence, identification, and separation, every sequence of empirical minimizers
\[
\widehat\psi_n
\in
\argmin_{\psi\in\Psi}
\widehat{\mathcal L}_n(\psi)
\]
converges in probability to the equivalence class \([\psi^{*}]\).
\end{theorem}

\begin{proof}
Let \(U\) be any neighborhood of \([\psi^{*}]\), and let \(\delta_U>0\) be supplied by separation. Define the event
\[
E_n
=
\left\{
\sup_{\psi\in\Psi}
\left|
\widehat{\mathcal L}_n(\psi)-\mathcal L(\psi)
\right|
<
\frac{\delta_U}{3}
\right\}.
\]
Uniform convergence implies
\[
\Prob(E_n)\to1.
\]

On \(E_n\), for every \(\psi\notin U\),
\[
\widehat{\mathcal L}_n(\psi)
>
\mathcal L(\psi)-\frac{\delta_U}{3}
\geq
\mathcal L(\psi^{*})
+
\frac{2\delta_U}{3}.
\]
Also,
\[
\widehat{\mathcal L}_n(\psi^{*})
<
\mathcal L(\psi^{*})
+
\frac{\delta_U}{3}.
\]
Therefore
\[
\widehat{\mathcal L}_n(\psi)
>
\widehat{\mathcal L}_n(\psi^{*})
+
\frac{\delta_U}{3}
\]
for every \(\psi\notin U\). No empirical minimizer can lie outside \(U\) on \(E_n\). Hence
\[
\Prob(\widehat\psi_n\notin U)
\leq
1-\Prob(E_n)
\to0.
\]
Since \(U\) was arbitrary, \(\widehat\psi_n\) converges in probability to \([\psi^{*}]\).
\end{proof}

The theorem establishes consistency for the minimizer of the declared population objective. It does not establish that the population objective identifies the epistemically correct quotient. If the loss omits an authorized distinction, if audit labels are biased, if the observed context distribution excludes a relevant region, or if the hypothesis class cannot express the required partition, the estimator may converge perfectly to an insufficient representation.

\begin{definition}[Target misspecification]
CLIO is target misspecified in context \(c\) when the population minimizer \(q_{\psi^{*}}\) is not sufficient for the true authorized function class \(\F_c\).
\end{definition}

\begin{corollary}[Consistency does not imply sufficiency]
There exist CLIO estimators satisfying all conditions of the consistency theorem whose limiting projection is insufficient for the true authorized function class.
\end{corollary}

\begin{proof}
Choose an observed loss that evaluates only a strict subset
\[
\F_c^{\mathrm{obs}}
\subsetneq
\F_c.
\]
Let \(x,y\) be indistinguishable under every observed function but distinguishable under some unobserved protected function. A hypothesis class containing a projection that identifies \(x\) and \(y\) may have a unique, separated population minimizer under the observed objective. The estimator converges to that minimizer while collapsing the distinction required by the unobserved function.
\end{proof}

\subsection{The impossibility of self-certifying projection}

A projection cannot generally certify its own sufficiency using only the information it retains. Once two source states have been identified, evidence of their difference may no longer exist inside the quotient.

\begin{theorem}[No internal detection of arbitrary collapse]
Let \(q:X\to M\) be noninjective. Then there exists a function
\[
f:X\to\{0,1\}
\]
for which \(q\) is insufficient, but no diagnostic depending only on \(q(x)\) can distinguish the collapsed pair witnessing the insufficiency.
\end{theorem}

\begin{proof}
Since \(q\) is noninjective, choose distinct \(x,y\in X\) with
\[
q(x)=q(y).
\]
Define
\[
f(x)=0,
\qquad
f(y)=1,
\]
and assign arbitrary values elsewhere. Then \(q\) is insufficient for \(f\).

Let
\[
d:M\to Z
\]
be any internal diagnostic. Since
\[
q(x)=q(y),
\]
we have
\[
d(q(x))=d(q(y)).
\]
Therefore \(d\) cannot distinguish the pair whose distinction \(f\) requires.
\end{proof}

This impossibility result explains why the Attention Compiler requires immutable raw submissions, external audits, random review, appeals, and delayed outcome comparison. The projected representation cannot be the sole source of evidence about whether its own fibers were homogeneous.

\subsection{Representational simplicity and reversibility}

Exact simplicity does not require lossless reconstruction of the original submission. It requires reversibility only with respect to authorized distinctions. Let
\[
\pi_c^{*}:X\to M_c^{*}
\]
be the canonical quotient. A section
\[
s_c:M_c^{*}\to X
\]
selects one representative from each equivalence class when such a selection exists. Generally,
\[
s_c\circ\pi_c^{*}\neq\operatorname{id}_X,
\]
because the original point cannot be recovered from its equivalence class. Nevertheless,
\[
\pi_c^{*}\circ s_c
=
\operatorname{id}_{M_c^{*}}.
\]

For every authorized function \(f_{\alpha,c}\), there exists
\[
\bar f_{\alpha,c}:M_c^{*}\to Y_\alpha
\]
such that
\[
f_{\alpha,c}
=
\bar f_{\alpha,c}\circ\pi_c^{*}.
\]
Therefore
\[
f_{\alpha,c}\circ s_c\circ\pi_c^{*}
=
\bar f_{\alpha,c}\circ
\pi_c^{*}\circ
s_c\circ
\pi_c^{*}
=
\bar f_{\alpha,c}\circ\pi_c^{*}
=
f_{\alpha,c}.
\]

The selected representative may differ from the original submission while remaining identical under every authorized operation. This is operational reversibility.

\begin{definition}[Operational reconstruction]
A reconstruction map
\[
r_c:M_c\to X
\]
is operationally exact for \(\F_c\) when
\[
f_{\alpha,c}(r_c(q_c(x)))
=
f_{\alpha,c}(x)
\]
for every \(x\in X\) and every authorized function \(f_{\alpha,c}\).
\end{definition}

Operational reconstruction is weaker than literal reconstruction and stronger than superficial semantic resemblance.

\subsection{The CLIO obligation}

The entire CLIO layer can now be condensed into one obligation.

\begin{principle}[CLIO obligation]
A projection used by the Attention Compiler may identify two submissions only when every authorized downstream operation, including provenance, audit, and appeal operations, is invariant under that identification in the active participant-recipient context.
\end{principle}

The principle has an exact form:
\[
q_c(x)=q_c(y)
\implies
x\equiv_c y.
\]
It has an approximate form:
\[
\mathcal R_\alpha(q_c\mid c)\leq\eps_\alpha
\]
for every authorized operation. It has a fiberwise diagnostic:
\[
\Lambda_{\alpha,c}(m)
\]
for every operational fiber. It has a temporal vulnerability:
\[
\F_c\subsetneq\F_{c'}
\]
can render a formerly sufficient quotient obsolete. It has an estimation vulnerability: consistency for an observed objective does not imply sufficiency for the true function class. It has an epistemic limitation: no noninjective quotient can detect every possible failure from within its own retained coordinates.

Representational simplicity is therefore not the smallest representation a system can construct. It is the smallest representation whose lost distinctions have been justified relative to a declared and auditable future of use.

\section{Projection Collapse, Disposition Divergence, and False Suppression}

\subsection{Collapse as a failure of authorized invariance}

Section 2 defined contextual sufficiency through invariance of every authorized downstream operation over the fibers of a projection. The present section studies the complementary object: the structure, magnitude, and institutional consequences of a failure of that invariance.

Fix a participant \(i\), recipient \(r\), and context
\[
c=c_{ir}(t).
\]
Let
\[
q_c:X\to M_c
\]
be the CLIO projection currently used by the Attention Compiler. For every \(m\in M_c\), define the fiber
\[
X_{c,m}
=
q_c^{-1}(m).
\]
The projection treats every element of \(X_{c,m}\) as representationally interchangeable. Sufficiency requires that this imposed interchangeability agree with the equivalence relation induced by the authorized function class.

Let
\[
\F_c
=
\{f_{\alpha,c}:X\to Y_\alpha\}_{\alpha\in I_c}.
\]
For each operation \(\alpha\), define the disagreement relation
\[
\Delta_{\alpha,c}
=
\left\{
(x,y)\in X\times X:
f_{\alpha,c}(x)\neq f_{\alpha,c}(y)
\right\}.
\]
Define the projection kernel relation
\[
\ker(q_c)
=
\left\{
(x,y)\in X\times X:
q_c(x)=q_c(y)
\right\}.
\]

\begin{definition}[Operation-specific collapse relation]
The collapse relation of \(q_c\) relative to operation \(\alpha\) is
\[
\mathfrak C_{\alpha,c}(q_c)
=
\ker(q_c)\cap\Delta_{\alpha,c}.
\]
\end{definition}

\begin{definition}[Total authorized collapse relation]
The total authorized collapse relation is
\[
\mathfrak C_c(q_c)
=
\bigcup_{\alpha\in I_c}
\mathfrak C_{\alpha,c}(q_c).
\]
\end{definition}

Thus,
\[
(x,y)\in\mathfrak C_c(q_c)
\]
exactly when the representation identifies \(x\) and \(y\) while at least one authorized operation distinguishes them.

\begin{proposition}[Collapse criterion for insufficiency]
The projection \(q_c\) is exactly sufficient for \(\F_c\) if and only if
\[
\mathfrak C_c(q_c)=\varnothing.
\]
\end{proposition}

\begin{proof}
Suppose \(q_c\) is exactly sufficient. If
\[
q_c(x)=q_c(y),
\]
then
\[
f_{\alpha,c}(x)=f_{\alpha,c}(y)
\]
for every \(\alpha\). Therefore no pair in \(\ker(q_c)\) belongs to any disagreement relation \(\Delta_{\alpha,c}\), and
\[
\mathfrak C_c(q_c)=\varnothing.
\]

Conversely, suppose
\[
\mathfrak C_c(q_c)=\varnothing.
\]
Let \(x,y\in X\) satisfy
\[
q_c(x)=q_c(y).
\]
If some authorized function distinguished them, then
\[
(x,y)\in
\ker(q_c)\cap\Delta_{\alpha,c}
\subseteq
\mathfrak C_c(q_c),
\]
contradicting emptiness. Hence every authorized function agrees on \(x\) and \(y\), which is exact sufficiency.
\end{proof}

The binary existence of a collapsed pair is enough to refute exact sufficiency. It does not indicate how frequently the collapse occurs, which operations are affected, how consequential the disagreement is, or whether the collapse systematically suppresses one class of submissions. These require quantitative observables.

\subsection{Warranted dispositions and observed dispositions}

The system's observed disposition must be distinguished from the disposition warranted by complete information.

Let
\[
D_c^{*}:X\to\mathcal D_r
\]
be the warranted disposition function in context \(c\). This is an idealized target. It represents the disposition that would be selected under complete access to the authorized evidence, correct application of the declared rules, and no projection loss.

Let
\[
\widehat D_c:M_c\to\mathcal D_r
\]
be the implemented disposition function operating on the projected representation. The deployed system produces
\[
D_c^{q}(x)
=
\widehat D_c(q_c(x)).
\]

A projection-induced disposition error occurs when
\[
D_c^{q}(x)\neq D_c^{*}(x).
\]
This error can arise because the quotient is insufficient, because \(\widehat D_c\) is a poor decoder over an otherwise sufficient quotient, or because both failures occur simultaneously.

\begin{definition}[Bayes disposition risk]
Let \(P_c\) be the context-relative submission distribution, and let
\[
L_D:\mathcal D_r\times\mathcal D_r\to[0,\infty]
\]
be a disposition loss. The optimal risk achievable through \(q_c\) is
\[
\mathcal R_D^{*}(q_c\mid c)
=
\inf_{\widehat D_c}
\E_{X\sim P_c}
\left[
L_D
\left(
D_c^{*}(X),
\widehat D_c(q_c(X))
\right)
\right].
\]
\end{definition}

This quantity isolates representational limitation. It gives the lowest possible disposition risk among all decoders that see only \(q_c(X)\). If it is positive, no improvement to the downstream classifier can eliminate the error without changing the representation.

Assume temporarily that \(\mathcal D_r\) is finite and that
\[
L_D(d,d')
=
\ind[d\neq d']
\]
is zero-one loss. For \(m\in M_c\), define
\[
p_c(d\mid m)
=
\Prob(D_c^{*}(X)=d\mid q_c(X)=m).
\]

\begin{theorem}[Fiberwise Bayes disposition error]
Under zero-one loss, the optimal projection-constrained disposition risk is
\[
\mathcal R_D^{*}(q_c\mid c)
=
\E_{M}
\left[
1-
\max_{d\in\mathcal D_r}
p_c(d\mid M)
\right],
\]
where
\[
M=q_c(X).
\]
\end{theorem}

\begin{proof}
Condition on \(M=m\). Any decoder must choose one disposition
\[
\widehat D_c(m)=d.
\]
Its conditional error probability is
\[
\Prob(D_c^{*}(X)\neq d\mid M=m)
=
1-p_c(d\mid m).
\]
This is minimized by choosing any modal disposition
\[
d_m^{*}
\in
\argmax_{d\in\mathcal D_r}
p_c(d\mid m).
\]
The minimum conditional error is therefore
\[
1-
\max_{d\in\mathcal D_r}
p_c(d\mid m).
\]
Taking expectation over \(M\) proves the result.
\end{proof}

\begin{corollary}[Zero Bayes risk and disposition sufficiency]
Under the preceding assumptions,
\[
\mathcal R_D^{*}(q_c\mid c)=0
\]
if and only if \(D_c^{*}\) is almost surely constant on almost every fiber of \(q_c\).
\end{corollary}

\begin{proof}
The integrand
\[
1-\max_d p_c(d\mid M)
\]
is nonnegative. Its expectation is zero exactly when it is zero almost surely. This occurs exactly when, for almost every \(m\), one disposition has conditional probability one on the fiber.
\end{proof}

The theorem reveals an irreducible error floor. If one fiber contains warranted answers, escalations, and appeals, a decoder must assign one output or randomize among outputs. It cannot recover which disposition applies to a particular source state because that distinction has already been removed.

\subsection{Disposition entropy}

The heterogeneity of warranted dispositions within a fiber can be measured by conditional entropy:
\[
H(D_c^{*}\mid M)
=
-\E_M
\sum_{d\in\mathcal D_r}
p_c(d\mid M)
\log p_c(d\mid M).
\]

\begin{definition}[Disposition-collapse entropy]
The disposition-collapse entropy of \(q_c\) is
\[
\mathscr H_D(q_c\mid c)
=
H(D_c^{*}(X)\mid q_c(X)).
\]
\end{definition}

If the warranted disposition is deterministic from the full submission, then
\[
H(D_c^{*}\mid X)=0.
\]
Consequently,
\[
I(D_c^{*};X\mid M)
=
H(D_c^{*}\mid M).
\]
The collapse entropy is therefore exactly the authorized disposition information present in the original submission but absent from the projection.

\begin{proposition}[Information identity for deterministic dispositions]
If \(D_c^{*}\) is a deterministic function of \(X\), then
\[
\mathscr H_D(q_c\mid c)
=
I(D_c^{*};X\mid q_c(X)).
\]
\end{proposition}

\begin{proof}
By definition,
\[
I(D_c^{*};X\mid M)
=
H(D_c^{*}\mid M)
-
H(D_c^{*}\mid X,M).
\]
Since \(D_c^{*}\) is determined by \(X\),
\[
H(D_c^{*}\mid X,M)=0.
\]
Substituting \(M=q_c(X)\) yields the result.
\end{proof}

This identity gives a precise meaning to the claim that the quotient discarded disposition-relevant information. It does not merely discard information about \(X\). It discards information about the warranted action.

\begin{corollary}[Data-processing bound]
For deterministic \(D_c^{*}\),
\[
I(D_c^{*};q_c(X))
\leq
I(D_c^{*};X)
=
H(D_c^{*}).
\]
Moreover,
\[
H(D_c^{*})
=
I(D_c^{*};q_c(X))
+
\mathscr H_D(q_c\mid c).
\]
\end{corollary}

The total warranted-disposition information decomposes into the portion preserved by the representation and the portion remaining unresolved within its fibers.

\subsection{Generalized disposition divergence}

Conditional entropy assumes a common finite disposition space and gives no direct account of asymmetric consequences. Confusing answer with clarification may be less damaging than confusing escalation with discard. Introduce a cost matrix
\[
C_c:\mathcal D_r\times\mathcal D_r\to[0,\infty],
\]
where
\[
C_c(d^{*},d)
\]
is the cost of producing \(d\) when \(d^{*}\) is warranted.

For fiber \(m\), define its optimal conditional cost:
\[
\Lambda_{D,c}(m)
=
\min_{d\in\mathcal D_r}
\sum_{d^{*}\in\mathcal D_r}
C_c(d^{*},d)
p_c(d^{*}\mid m).
\]

\begin{definition}[Cost-sensitive collapse functional]
The global cost-sensitive collapse of \(q_c\) is
\[
\Lambda_{D,c}(q_c)
=
\E_M[\Lambda_{D,c}(M)].
\]
\end{definition}

This is exactly the Bayes risk under the declared cost matrix:
\[
\Lambda_{D,c}(q_c)
=
\mathcal R_D^{*}(q_c\mid c).
\]

The cost matrix can encode protected asymmetries. Let
\[
d_{\mathrm{esc}}
\]
denote escalation and
\[
d_{\mathrm{discard}}
\]
denote discard or indefinite delay. One may impose
\[
C_c(d_{\mathrm{esc}},d_{\mathrm{discard}})
\gg
C_c(d_{\mathrm{clar}},d_{\mathrm{teach}}).
\]
The projection is then penalized much more strongly for fibers that mix valuable escalation candidates with apparently routine submissions.

\subsection{False suppression as a one-sided collapse event}

Let
\[
V_c(x)\in\{0,1\}
\]
indicate whether \(x\) contains a genuinely valuable residual relative to recipient \(r\) and context \(c\). Let
\[
S_c^{*}(x)\in\{0,1\}
\]
indicate whether the warranted outcome preserves access to eventual human consideration. Let
\[
\widehat S_c(q_c(x))\in\{0,1\}
\]
be the implemented preservation decision.

Define the false-suppression event
\[
\mathsf{FS}
=
\left\{
V_c(X)=1,
\quad
\widehat S_c(q_c(X))=0
\right\}.
\]
The conditional false-suppression rate is
\[
\operatorname{FSR}(q_c,\widehat S_c\mid c)
=
\Prob
\left(
\widehat S_c(q_c(X))=0
\middle|
V_c(X)=1
\right).
\]

Even the optimal preservation decoder may have positive false suppression if valuable and nonvaluable submissions share a fiber.

For \(m\in M_c\), define
\[
v_c(m)
=
\Prob(V_c(X)=1\mid q_c(X)=m).
\]

Suppose preserving a submission incurs review cost \(c_R>0\), while suppressing a valuable submission incurs epistemic cost \(c_S>0\). The expected conditional cost of preserving fiber \(m\) is
\[
\mathcal C_{\mathrm{preserve}}(m)
=
c_R.
\]
The expected conditional cost of suppressing it is
\[
\mathcal C_{\mathrm{suppress}}(m)
=
c_S v_c(m).
\]
The Bayes-optimal policy preserves when
\[
c_R
\leq
c_Sv_c(m),
\]
or equivalently,
\[
v_c(m)
\geq
\frac{c_R}{c_S}.
\]

\begin{proposition}[Projection-constrained suppression threshold]
Under the preceding binary cost model, the optimal projection-constrained policy is
\[
\widehat S_c^{*}(m)
=
\ind
\left[
v_c(m)\geq\frac{c_R}{c_S}
\right].
\]
\end{proposition}

\begin{proof}
For each fiber, choose the action with smaller conditional expected cost. Preservation costs \(c_R\), while suppression costs \(c_Sv_c(m)\). Preservation is optimal exactly under the stated inequality.
\end{proof}

The result exposes the moral significance of fiber composition. A valuable submission can be suppressed even under a Bayes-optimal policy if it lies in a fiber with sufficiently low valuable-submission prevalence. The error does not require malicious design or a poor classifier. It can follow from rational cost minimization over an insufficient representation.

\begin{definition}[Minority-value collapse]
A fiber \(m\) exhibits minority-value collapse when
\[
0<v_c(m)<\frac{c_R}{c_S}.
\]
\end{definition}

Every valuable submission in such a fiber is suppressed under the cost-optimal projection-constrained policy, even though the fiber contains genuine value.

\begin{theorem}[Irreducible false suppression under mixed fibers]
Suppose a set \(B\subseteq M_c\) has positive probability and satisfies
\[
0<v_c(m)<\frac{c_R}{c_S}
\]
for every \(m\in B\). Then the Bayes-optimal projection-constrained policy has false-suppression rate bounded below by
\[
\operatorname{FSR}
\geq
\frac{
\int_B v_c(m)\,dP_M(m)
}{
\Prob(V_c(X)=1)
}.
\]
\end{theorem}

\begin{proof}
For every \(m\in B\), the optimal policy suppresses the fiber. The joint probability of a valuable submission and representation in \(B\) is
\[
\Prob(V_c=1,M\in B)
=
\int_B
\Prob(V_c=1\mid M=m)\,dP_M(m)
=
\int_Bv_c(m)\,dP_M(m).
\]
Every such submission is falsely suppressed. Dividing by
\[
\Prob(V_c=1)
\]
gives the stated conditional lower bound.
\end{proof}

No improvement to \(\widehat S_c\) can eliminate this lower bound without increasing review cost or refining the projection. The system must split the mixed fibers, change the cost threshold, or allocate exploration to recover suppressed minority structure.

\subsection{Escalation and suppression as recipient-relative events}

Value is indexed by recipient. Let
\[
V_{ir}(x,c)
\]
be the value of direct or eventual consideration of submission \(x\) by recipient \(r\). Let
\[
\Delta_{\mathrm{joint}}(x,i,r)
\]
denote the gain expected specifically from interaction between participant \(i\) and recipient \(r\).

A submission may be valuable but not escalation-worthy for a particular recipient. Conversely, a familiar claim may warrant escalation because the participant possesses complementary evidence or a construction that interacts unusually well with the recipient's unresolved work.

Define the warranted escalation indicator
\[
E_{ir}^{*}(x)
=
\ind
\left[
\E\left[
\Delta_{\mathrm{joint}}(x,i,r)
\middle|
I_x,K_r,\theta_i,t
\right]
>
C_{ir}
\right],
\]
where \(C_{ir}\) is the anticipated recipient attention cost.

A quotient \(q_{c_{ir}}\) is escalation sufficient when
\[
q_{c_{ir}}(x)=q_{c_{ir}}(y)
\implies
E_{ir}^{*}(x)=E_{ir}^{*}(y).
\]

\begin{observation}
Semantic normalization can be sufficient for claim retrieval while being insufficient for escalation. Two submissions may normalize to the same proposition while differing in evidence, timing, authorship, complementary expertise, or relation to the recipient's unresolved problems.
\end{observation}

This is why the residual packet must retain more than normalized propositional content. It must retain every coordinate used to estimate recipient-relative complementarity.

\subsection{Pairwise collapse observables}

The previous quantities condition on representation fibers. A pairwise formulation can compare the geometry of the source and projected spaces directly.

Let
\[
d_X:X\times X\to[0,\infty)
\]
measure source dissimilarity, and let
\[
d_M:M_c\times M_c\to[0,\infty)
\]
measure representational dissimilarity. Let
\[
d_{\F,c}(x,y)
\]
be an authorized operational distance. One construction is
\[
d_{\F,c}(x,y)
=
\sum_{\alpha\in I_c}
w_\alpha
d_\alpha
\left(
f_{\alpha,c}(x),
f_{\alpha,c}(y)
\right),
\]
where \(d_\alpha\) is a codomain-appropriate distance.

Define the compression discrepancy
\[
\Delta_q(x,y)
=
d_{\F,c}(x,y)
-
d_M(q_c(x),q_c(y)).
\]
Large positive values indicate that operationally important differences have been compressed more strongly than the representation geometry records.

For a hard quotient with
\[
d_M(m,m)=0,
\]
any collapsed pair satisfies
\[
\Delta_q(x,y)
=
d_{\F,c}(x,y).
\]

\begin{definition}[Pairwise operational collapse energy]
The pairwise collapse energy is
\[
\mathcal E_{\mathrm{coll}}(q_c)
=
\E_{X,Y\sim P_c}
\left[
d_{\F,c}(X,Y)
\ind[q_c(X)=q_c(Y)]
\right].
\]
\end{definition}

\begin{proposition}
If every authorized codomain distance satisfies
\[
d_\alpha(a,b)=0
\iff
a=b,
\]
and every weight is strictly positive, then
\[
\mathcal E_{\mathrm{coll}}(q_c)=0
\]
if and only if \(q_c\) is sufficient almost surely for \(\F_c\).
\end{proposition}

\begin{proof}
The integrand is nonnegative. If \(q_c\) is sufficient, then every collapsed pair agrees under every authorized function, so
\[
d_{\F,c}(X,Y)=0
\]
whenever \(q_c(X)=q_c(Y)\), and the energy is zero.

Conversely, if the energy is zero, then the integrand vanishes almost surely. For almost every pair with equal representation,
\[
d_{\F,c}(X,Y)=0.
\]
Strict positivity of the weights and definiteness of the codomain distances imply equality under every authorized function. Thus sufficiency holds almost surely.
\end{proof}

The energy can be decomposed by operation:
\[
\mathcal E_{\mathrm{coll}}(q_c)
=
\sum_{\alpha\in I_c}
w_\alpha
\mathcal E_{\alpha,c}(q_c),
\]
where
\[
\mathcal E_{\alpha,c}(q_c)
=
\E
\left[
d_\alpha
\left(
f_{\alpha,c}(X),
f_{\alpha,c}(Y)
\right)
\ind[q_c(X)=q_c(Y)]
\right].
\]

This decomposition identifies which authorized operations are damaged by the quotient.

\subsection{Directional collapse and protected minorities}

Average pairwise energy may remain small when collapse disproportionately affects a rare participant class. Let
\[
A\in\mathcal G
\]
be a participant-group variable, where \(\mathcal G\) may encode any audited class relevant to access, communication form, discipline, language, or institutional position.

Define group-conditional collapse energy
\[
\mathcal E_{\mathrm{coll}}(q_c\mid A=g)
=
\E
\left[
d_{\F,c}(X,Y)
\ind[q_c(X)=q_c(Y)]
\middle|
A_X=g
\right].
\]

Define group-conditional false suppression
\[
\operatorname{FSR}_g
=
\Prob
\left(
\widehat S_c(q_c(X))=0
\middle|
V_c(X)=1,
A=g
\right).
\]

An aggregate constraint
\[
\operatorname{FSR}\leq\eps
\]
does not imply
\[
\operatorname{FSR}_g\leq\eps
\]
for every \(g\). A common class can dominate the aggregate while a rare class suffers nearly total suppression.

\begin{definition}[Suppression disparity]
For groups \(g,h\in\mathcal G\), define
\[
\Delta_{\mathrm{FSR}}(g,h)
=
\left|
\operatorname{FSR}_g
-
\operatorname{FSR}_h
\right|.
\]
\end{definition}

A governance constraint may require
\[
\sup_{g,h\in\mathcal G}
\Delta_{\mathrm{FSR}}(g,h)
\leq
\delta_{\mathrm{group}},
\]
subject to the limitation that group definitions themselves must not become unjustified active coordinates in ordinary routing.

The correct architecture can keep protected-group information in an audit channel while excluding it from routine disposition functions except where explicitly authorized.

\subsection{Unknown value and partial identification}

The event \(V_c=1\) is not directly observed for submissions that never receive sufficient examination. Let
\[
R_c\in\{0,1\}
\]
indicate whether a submission's latent value is eventually reviewed. The observed value label is
\[
V_c^{\mathrm{obs}}
=
R_cV_c.
\]
When \(R_c=0\), the value remains unknown, not zero.

A naive estimator using only reviewed cases estimates
\[
\Prob(D\mid V,R=1),
\]
not
\[
\Prob(D\mid V).
\]
These coincide only under strong selection assumptions.

Let
\[
\pi_R(x,c)
=
\Prob(R_c=1\mid X=x,c)
\]
be the review propensity.

\begin{assumption}[Positivity of review]
There exists \(\eta>0\) such that
\[
\pi_R(x,c)\geq\eta
\]
for every submission class whose false-suppression rate is to be estimated.
\end{assumption}

\begin{assumption}[Conditional review ignorability]
Conditional on recorded audit covariates \(W\),
\[
V_c\perp R_c\mid W.
\]
\end{assumption}

Under these assumptions, inverse-propensity weighting can estimate population quantities.

Let
\[
Z
=
\ind[\widehat S_c(q_c(X))=0]
\]
be the suppression indicator. Then
\[
\E
\left[
\frac{R_cV_cZ}{\pi_R(W)}
\right]
=
\E[V_cZ],
\]
and
\[
\E
\left[
\frac{R_cV_c}{\pi_R(W)}
\right]
=
\E[V_c].
\]

\begin{proposition}[Inverse-propensity false-suppression identity]
Under positivity and conditional review ignorability,
\[
\operatorname{FSR}
=
\frac{
\E\left[
R_cV_cZ/\pi_R(W)
\right]
}{
\E\left[
R_cV_c/\pi_R(W)
\right]
}.
\]
\end{proposition}

\begin{proof}
By iterated expectation,
\[
\E
\left[
\frac{R_cV_cZ}{\pi_R(W)}
\right]
=
\E
\left[
\E
\left[
\frac{R_cV_cZ}{\pi_R(W)}
\middle|
W,V_c,Z
\right]
\right].
\]
Conditional ignorability gives
\[
\E[R_c\mid W,V_c,Z]
=
\E[R_c\mid W]
=
\pi_R(W).
\]
Therefore
\[
\E
\left[
\frac{R_cV_cZ}{\pi_R(W)}
\middle|
W,V_c,Z
\right]
=
V_cZ.
\]
Thus the numerator equals
\[
\E[V_cZ].
\]
The same argument gives
\[
\E
\left[
\frac{R_cV_c}{\pi_R(W)}
\right]
=
\E[V_c].
\]
Their ratio is
\[
\frac{\Prob(V_c=1,Z=1)}{\Prob(V_c=1)}
=
\Prob(Z=1\mid V_c=1).
\]
\end{proof}

Conditional ignorability is often implausible because review is directed precisely by variables correlated with latent value. Random audit selection is therefore structurally important. It creates a component of \(R_c\) whose propensity is known by design.

\subsection{Random audit as fiber sampling}

Let
\[
A_c^{(x)}
\]
be an exploration budget. Suppose a fraction of this budget samples projected fibers according to a known distribution
\[
\omega_c(m)>0.
\]
Within a sampled fiber \(m\), the audit selects raw submissions according to
\[
\omega_c(x\mid m)>0.
\]

The resulting inclusion probability is
\[
\pi_{\mathrm{audit}}(x)
=
\omega_c(q_c(x))
\omega_c(x\mid q_c(x)).
\]

The audit observes the raw submission, reconstructs the warranted disposition under enhanced review, and compares it with the projected disposition. This produces direct samples of within-fiber heterogeneity.

For operation \(\alpha\), an unbiased Horvitz-Thompson estimator of pairwise collapse energy can be formed from sampled pairs:
\[
\widehat{\mathcal E}_{\alpha,c}
=
\frac{1}{N}
\sum_{j=1}^{N}
\frac{
d_\alpha
\left(
f_{\alpha,c}(X_j),
f_{\alpha,c}(Y_j)
\right)
\ind[q_c(X_j)=q_c(Y_j)]
}{
\pi_{\mathrm{pair}}(X_j,Y_j)
},
\]
with normalization adjusted to the sampling design.

The audit should oversample large, uncertain, historically heterogeneous, and institutionally consequential fibers while retaining a nonzero probability of selecting every eligible fiber. Pure uncertainty sampling is insufficient because the projection may be confidently wrong within a collapsed fiber.

\begin{principle}[Fiber audit principle]
The unit of representational audit is not only the individual decision. It is the fiber within which the representation claims interchangeability.
\end{principle}

\subsection{Appeals as endogenous counterexamples}

An appeal occurs when a participant challenges the projected disposition, reconstructed claim, prerequisite assignment, or assertion of prior resolution. Let
\[
A_i=1
\]
indicate an appeal. Successful appeals provide counterexamples to the current quotient or decoder.

Appeals are not random samples. Their probability depends on participant persistence, access, confidence, available time, understanding of the process, and belief that appeal may succeed. Consequently, an appeal-success rate cannot directly estimate population collapse.

Nevertheless, an individual successful appeal can refute exact sufficiency.

\begin{proposition}[Appeal witness]
Suppose submissions \(x\) and \(y\) satisfy
\[
q_c(x)=q_c(y),
\]
the system assigns the same projected disposition to both, and an appeal establishes that their warranted dispositions differ. Then \(q_c\) is not sufficient for the disposition function.
\end{proposition}

\begin{proof}
The established difference in warranted dispositions gives
\[
D_c^{*}(x)\neq D_c^{*}(y),
\]
while equality of projected representations gives
\[
q_c(x)=q_c(y).
\]
This violates the definition of sufficiency.
\end{proof}

Thus random audit supports rate estimation, while appeal supplies adversarial witnesses. Both are required.

\subsection{Collapse localization}

Once insufficiency is detected, repair requires identifying the missing distinction. Let
\[
\varphi:X\to\mathcal U
\]
be a candidate feature omitted from \(q_c\). Define an augmented representation
\[
q_c^{+\varphi}(x)
=
(q_c(x),\varphi(x)).
\]

The usefulness of \(\varphi\) for disposition recovery can be measured by conditional mutual information:
\[
I(D_c^{*};\varphi(X)\mid q_c(X)).
\]
If this quantity is zero, \(\varphi\) supplies no additional disposition information beyond the current quotient. If it is positive, the feature explains some within-fiber heterogeneity.

\begin{proposition}[Refinement monotonicity]
For any augmentation \(\varphi\),
\[
H(D_c^{*}\mid q_c(X),\varphi(X))
\leq
H(D_c^{*}\mid q_c(X)).
\]
\end{proposition}

\begin{proof}
Conditioning cannot increase entropy.
\end{proof}

The reduction
\[
H(D_c^{*}\mid q_c(X))
-
H(D_c^{*}\mid q_c(X),\varphi(X))
=
I(D_c^{*};\varphi(X)\mid q_c(X))
\]
quantifies how much collapse the feature resolves.

A minimal repair seeks a feature family
\[
\Phi^{*}
=
\{\varphi_1,\ldots,\varphi_k\}
\]
minimizing representational cost subject to a collapse constraint:
\[
\Phi^{*}
\in
\argmin_{\Phi}
\Comp(\Phi)
\]
subject to
\[
H(D_c^{*}\mid q_c(X),\Phi(X))
\leq
\eps_D
\]
and corresponding constraints for every protected operation.

This is representational expansion. It differs from retraining the same decoder over the same quotient. If
\[
\mathcal R_D^{*}(q_c\mid c)>0,
\]
decoder-only retraining cannot remove the irreducible component.

\subsection{The decomposition of observed error}

Let
\[
\widehat D_c
\]
be the deployed decoder. Its total expected loss can be decomposed into irreducible projection risk and decoder excess risk:
\[
\mathcal R_D(\widehat D_c\circ q_c)
=
\mathcal R_D^{*}(q_c)
+
\mathcal E_D(\widehat D_c\mid q_c),
\]
where
\[
\mathcal E_D(\widehat D_c\mid q_c)
=
\mathcal R_D(\widehat D_c\circ q_c)
-
\mathcal R_D^{*}(q_c)
\geq0.
\]

The first term requires representation repair. The second can be reduced by improving estimation or optimization while retaining the same representation.

\begin{principle}[Error localization principle]
No routing repair should be prescribed before determining whether the observed error arises from the decoder, the quotient, or the contextual gluing stage.
\end{principle}

Repeatedly optimizing \(\widehat D_c\) against an insufficient \(q_c\) can make the system more confidently approximate the best decision available from inadequate information. It cannot restore the missing distinction.

\subsection{Collapse propagation through the integrated architecture}

Projection collapse occurs before HYDRA gluing. Its effects can therefore propagate through every later stage.

Let
\[
x,y\in X
\]
satisfy
\[
q_c(x)=q_c(y)
\]
but
\[
D_c^{*}(x)\neq D_c^{*}(y).
\]
Assume the contextual diagram supplied to HYDRA depends on the raw submission only through \(q_c(x)\) and shared context \(c\). Then the two HYDRA diagrams are identical:
\[
\A_{x,c}
=
\A_{y,c}.
\]
Their colimits are therefore canonically identical, and every deterministic downstream transition produces the same result.

\begin{theorem}[Irrecoverability after complete quotient mediation]
Suppose every post-CLIO operation depends on \(x\) only through \(q_c(x)\) and fixed context \(c\). If
\[
q_c(x)=q_c(y),
\]
then the integrated architecture produces
\[
\AC(x\mid r,t)=\AC(y\mid r,t).
\]
Consequently, if the warranted dispositions differ, no downstream HYDRA, RSVP, or disposition operation can correct the error without access to an additional distinction.
\end{theorem}

\begin{proof}
Equality of projected representations and equality of fixed context imply equality of every local object and compatibility map constructed from them. Therefore the HYDRA diagrams coincide. Their global assemblies coincide. Applying the same transition operator and disposition map to equal global states yields equal outputs.
\end{proof}

The theorem establishes a strict ordering of repair. A downstream gluing mechanism cannot reconstruct a distinction removed upstream unless another channel retains it. HYDRA may reveal incompatibility among what remains, but it cannot infer an arbitrary distinction absent from every local section.

\subsection{The collapse obligation}

The central result of this section can be stated as a requirement on every deployed projection.

\begin{principle}[Collapse obligation]
For every operational fiber \(m\), the Attention Compiler must maintain evidence that the fiber is sufficiently homogeneous with respect to every protected disposition, provenance, audit, and appeal function, or it must assign uncertainty and exploration probability adequate to test that claim.
\end{principle}

Exact homogeneity is characterized by
\[
\mathfrak C_c(q_c)=\varnothing.
\]
Average retained disposition information is characterized by
\[
H(D_c^{*}\mid q_c(X)).
\]
Cost-sensitive irreducible error is characterized by
\[
\Lambda_{D,c}(q_c).
\]
One-sided epistemic danger is characterized by
\[
\Prob
\left(
\widehat S_c(q_c(X))=0
\middle|
V_c(X)=1
\right).
\]
Group-conditioned suppression is characterized by
\[
\operatorname{FSR}_g.
\]
The information supplied by a candidate repair feature is characterized by
\[
I(D_c^{*};\varphi(X)\mid q_c(X)).
\]

These quantities answer different questions and must not be collapsed into one confidence score. A representation may have low average disposition entropy while catastrophically suppressing a rare valuable class. It may preserve final dispositions while destroying provenance. It may perform well on reviewed cases while hiding unknown value in unaudited fibers. It may admit a highly accurate decoder while retaining an irreducible error caused by projection itself.

Representational simplicity is therefore constrained not only by what the system usually needs, but by the distinctions whose loss would make the system unable to recognize its own most consequential mistakes.

\section{Participant-State Estimation and Prerequisite-Sensitive Repair}

\subsection{Routing under latent epistemic state}

The same surface submission can require different treatment depending on what its author understands, intends, rejects, or is attempting to accomplish. A compressed question may be evidence of expertise, confusion, terminological mismatch, strategic ambiguity, or familiarity with a different formal tradition. Correct routing therefore depends on variables not directly observable in the submitted expression.

Let
\[
\theta_i(t)\in\ThetaSpace
\]
denote the participant state relevant to participant \(i\) at conversational time \(t\). This state is not a model of the participant's complete psychology. It contains only variables needed by authorized conversational operations.

Write
\[
\theta_i(t)
=
\left(
\kappa_i(t),
b_i(t),
\delta_i(t),
\iota_i(t),
\upsilon_i(t),
\omega_i(t)
\right).
\]
Here \(\kappa_i(t)\) is demonstrated competency, \(b_i(t)\) is proposition-relative commitment or rejection, \(\delta_i(t)\) is the participant's operative definitional state, \(\iota_i(t)\) is inferred conversational intention, \(\upsilon_i(t)\) is presentation preference, and \(\omega_i(t)\) is residual uncertainty about the preceding coordinates.

These coordinates are separated because no one of them licenses inference about the others. A participant can understand a proposition while rejecting it. A participant can prefer computational explanations while possessing strong proof competency. A participant can use unfamiliar vocabulary while applying the relevant concept correctly. A participant can produce technically polished language while failing elementary transfer tests.

Let
\[
Y_{i,1:t}
=
(y_{i,1},\ldots,y_{i,t})
\]
be the observable interaction history. The system maintains a posterior
\[
P_t^i(d\theta)
=
\Prob
\left(
\theta_i(t)\in d\theta
\middle|
Y_{i,1:t}
\right).
\]

\begin{definition}[Operational participant model]
An operational participant model is a posterior distribution over the least latent state required to choose authorized conversational interventions. It is not authorized to support unrelated classification of the participant.
\end{definition}

The restriction to operational necessity is part of representational simplicity. If a personal attribute cannot affect any authorized routing, explanation, audit, or appeal function, it should not remain an active coordinate of \(\theta_i\).

\subsection{State dynamics}

Participant state changes during the conversation. A prerequisite explanation may increase competency. A failed application may lower confidence in an earlier mastery estimate. A clarification may alter the inferred operative definition without changing competency. A participant may revise a proposition after adversarial testing.

Let
\[
a_t\in\mathcal A_{\mathrm{conv}}
\]
be the intervention selected at time \(t\). Let
\[
o_{t+1}
\]
be the resulting observation. A controlled state-space model has transition kernel
\[
P
\left(
\theta_{t+1}
\middle|
\theta_t,a_t
\right)
\]
and observation kernel
\[
P
\left(
o_{t+1}
\middle|
\theta_{t+1},a_t
\right).
\]

The prediction step is
\[
P_{t+1|t}^i(d\theta')
=
\int_{\ThetaSpace}
P(d\theta'\mid\theta,a_t)
P_t^i(d\theta).
\]
The update step is
\[
P_{t+1}^i(d\theta')
=
\frac{
P(o_{t+1}\mid\theta',a_t)
P_{t+1|t}^i(d\theta')
}{
\int_{\ThetaSpace}
P(o_{t+1}\mid\vartheta,a_t)
P_{t+1|t}^i(d\vartheta)
}.
\]

This is a controlled Bayesian filter. Its value is not that every conversational variable is literally Bayesian. Its value is that uncertainty, intervention, learning, and revision are represented explicitly rather than replaced by a permanent label.

\begin{principle}[Locality of participant inference]
An observation may update only those participant-state coordinates for which the observation model declares evidential relevance.
\end{principle}

A preference gesture can update \(\upsilon_i\), but it cannot directly update proof competency \(\kappa_i\). A demonstrated counterexample can update competency and commitment, but it does not by itself reveal the participant's preferred explanatory style.

\subsection{Competency as a graph-indexed random field}

Let
\[
G=(V,E)
\]
be a prerequisite graph. Each vertex \(v\in V\) represents a competency, concept, method, distinction, or operation. An edge
\[
u\to v
\]
means that demonstrated use of \(u\) is normally required for reliable use of \(v\).

For participant \(i\), define
\[
K_i(v,t)\in\{0,1\}
\]
as the latent indicator that the participant can reliably use competency \(v\) at time \(t\). The system maintains
\[
d_i(v,t)
=
\Prob(K_i(v,t)=1\mid Y_{i,1:t}).
\]

The vector
\[
d_i(t)
=
(d_i(v,t))_{v\in V}
\]
is a probabilistic field over the prerequisite graph. It is not necessarily monotone along edges. A programmer may demonstrate operational use of a concept while lacking its usual formal prerequisite. An autodidact may have gaps inconsistent with a conventional curriculum. Graph edges therefore encode defeasible dependency expectations, not mandatory biographical histories.

Let \(q\) be a diagnostic task targeted at competency \(v\). Let
\[
Y_q
\]
be the response. If
\[
p_1(y)
=
\Prob(Y_q=y\mid K_i(v)=1)
\]
and
\[
p_0(y)
=
\Prob(Y_q=y\mid K_i(v)=0),
\]
then Bayes' rule gives
\[
d_i'(v)
=
\frac{
p_1(y)d_i(v)
}{
p_1(y)d_i(v)+p_0(y)(1-d_i(v))
}.
\]

The posterior odds form is
\[
\frac{d_i'(v)}{1-d_i'(v)}
=
\frac{p_1(y)}{p_0(y)}
\frac{d_i(v)}{1-d_i(v)}.
\]
The likelihood ratio
\[
\frac{p_1(y)}{p_0(y)}
\]
measures the diagnostic force of the observation.

\begin{definition}[Transfer-sensitive diagnostic]
A task \(q\) is transfer sensitive for competency \(v\) when successful performance requires applying \(v\) in a context not identical to the context in which the participant encountered its explanation.
\end{definition}

Recognition tasks often have weak likelihood ratios because memorized vocabulary can produce correct responses without usable competency. Transfer tasks can provide stronger evidence when the response model distinguishes genuine application from surface recall.

\subsection{Comprehension and assent}

Let
\[
B_i(p,t)\in\{-1,0,+1\}
\]
represent the participant's commitment to proposition \(p\), where \(-1\) denotes rejection, \(0\) suspension or uncertainty, and \(+1\) acceptance. Let
\[
U_i(p,t)\in\{0,1\}
\]
represent whether the participant understands \(p\) sufficiently to state its content, consequences, and relevant alternatives.

The system maintains separate posteriors
\[
u_i(p,t)
=
\Prob(U_i(p,t)=1\mid Y_{i,1:t})
\]
and
\[
b_i(s\mid p,t)
=
\Prob(B_i(p,t)=s\mid Y_{i,1:t}).
\]

\begin{principle}[Comprehension-assent separation]
No routing rule may infer failed comprehension solely from disagreement, nor infer assent solely from demonstrated comprehension.
\end{principle}

A participant passes a prerequisite associated with proposition \(p\) when they demonstrate the competency needed to use \(p\), not when they accept \(p\) as true.

Let
\[
\operatorname{Pass}_i(p,t)
=
\ind[u_i(p,t)\geq\tau_U].
\]
This gate contains no condition on
\[
b_i(+1\mid p,t).
\]

\begin{proposition}[Ideological neutrality of competency gating]
Suppose prerequisite passage depends only on \(u_i(p,t)\) and task-relevant competency variables. Then two participants with equal comprehension and competency posteriors receive equal prerequisite status even when their commitment distributions differ.
\end{proposition}

\begin{proof}
Let participants \(i\) and \(j\) satisfy
\[
u_i(p,t)=u_j(p,t)
\]
and have equal task-relevant competency posteriors. If the passage function contains no commitment coordinate, its inputs are identical for \(i\) and \(j\). Therefore its outputs are identical.
\end{proof}

The proposition is elementary, but its violation is a central gatekeeping failure. If rejection increases the inferred probability of incompetence without independent performance evidence, the system transforms contested premises into compulsory ideology.

\subsection{Operative definitions}

Participants may use the same term under different definitions or different terms for the same structure. Let
\[
\mathcal L
\]
be a vocabulary and let
\[
\delta_i(t):\mathcal L\to\mathcal S
\]
map terms into semantic or operational structures.

The archive maintains a recipient-relative definitional map
\[
\delta_r^K(t):\mathcal L\to\mathcal S.
\]
A terminological mismatch occurs when
\[
\delta_i(t)(w)\neq\delta_r^K(t)(w)
\]
for some term \(w\), while the participant's underlying claim may remain coherent under their own map.

Let a normalized claim have the form
\[
q_i
=
\varphi
\left(
\delta_i(w_1),
\ldots,
\delta_i(w_k)
\right).
\]
Substituting archive definitions without checking equivalence can change the claim:
\[
\varphi
\left(
\delta_r^K(w_1),
\ldots,
\delta_r^K(w_k)
\right)
\neq
q_i.
\]

\begin{definition}[Definition-preserving reconstruction]
A reconstruction \(R_c(x_i)\) is definition preserving when every term substitution is accompanied by a witnessed map between the participant's operative structure and the archive structure, and the normalized claim commutes under that map.
\end{definition}

Suppose
\[
h_j:
\delta_i(w_j)\to\delta_r^K(w_j)
\]
are candidate translation maps. Definition preservation requires an appropriate commutative relation
\[
\varphi_r\circ(h_1,\ldots,h_k)
=
h_{\mathrm{out}}\circ\varphi_i.
\]
When no such relation is established, the system should report terminological uncertainty rather than silently replacing the participant's claim.

\subsection{Presentation preference as a separate channel}

Let
\[
\mathcal S_v
\]
be the set of admissible presentation styles for competency node \(v\). Examples include proof-oriented, computational, geometric, historical, experimental, and worked-example presentations.

Define
\[
\upsilon_i(s\mid v,t)
\]
as the estimated utility of style \(s\in\mathcal S_v\) for participant \(i\). A swipe-style signal
\[
z_t\in\{-1,+1\}
\]
may update the selected style weight through
\[
\upsilon_i'(s\mid v)
=
(1-\eta)\upsilon_i(s\mid v)
+
\eta z_t.
\]

This update is intentionally shallow. It changes selection among already admissible explanations. It does not change the missing-prerequisite set.

Let
\[
M_i(t)
\subseteq V
\]
be the current missing-prerequisite set. The separation requirement is
\[
\frac{\partial M_i(t)}{\partial \upsilon_i}=0
\]
in the structural sense that presentation preference is not an input to the function computing competency deficiency.

\begin{proposition}[Preference-invariance of competency status]
If the competency update kernel is conditionally independent of preference signals given demonstrated task performance, then changing only the participant's presentation-preference history leaves the competency posterior unchanged.
\end{proposition}

\begin{proof}
Let \(Z_{1:t}\) be preference observations and \(Y_{1:t}\) demonstrated-performance observations. Assume
\[
K_i(v,t)\perp Z_{1:t}\mid Y_{1:t}.
\]
Then
\[
\Prob(K_i(v,t)=1\mid Y_{1:t},Z_{1:t})
=
\Prob(K_i(v,t)=1\mid Y_{1:t}).
\]
Therefore changes to \(Z_{1:t}\) alone do not alter the competency posterior.
\end{proof}

\subsection{Required concepts and missing prerequisites}

Let
\[
R_c(x_i)\subseteq V
\]
be the competency nodes required to evaluate submission \(x_i\) in context \(c\). This set depends on the claims being tested and the operations the system intends to perform. It is not the set of everything related to the submission.

Let
\[
C_i(t;\tau)
=
\{v\in V:d_i(v,t)\geq\tau_v\}
\]
be the set of demonstrated competencies above node-specific confidence thresholds.

Let
\[
\closure_G(C_i)
\]
be the defeasible competency closure inferred through the graph. Because edges need not represent logical implication, closure should be governed by confidence propagation rather than unconditional reachability.

For an edge \(u\to v\), let
\[
\omega_{uv}
=
\Prob(K_i(u)=1\mid K_i(v)=1)
\]
or another empirically justified dependency weight. A path
\[
p=(v_0,\ldots,v_k)
\]
has reliability
\[
\Omega(p)
=
\prod_{j=0}^{k-1}\omega_{v_jv_{j+1}}.
\]
An inferred competency enters the closure only when the best supporting path exceeds a declared threshold.

Define
\[
\closure_G^\tau(C_i)
=
\left\{
v\in V:
\sup_{p:C_i\leadsto v}
\Omega(p)\geq\tau_v^{\mathrm{infer}}
\right\}.
\]

The missing-prerequisite set is
\[
M_i(x_i,c,t)
=
R_c(x_i)
\setminus
\closure_G^\tau(C_i(t;\tau)).
\]

A large \(M_i\) does not imply that the participant must traverse every node. Some nodes may be bypassed through direct demonstration, alternative formalizations, or equivalent competencies.

\subsection{Alternative prerequisite paths}

Let
\[
\mathcal P_i(x_i,c)
\]
be the set of candidate curricular subgraphs whose addition to the participant's demonstrated competency state covers the required set. A route
\[
P\subseteq V
\]
is feasible when
\[
R_c(x_i)
\subseteq
\closure_G^\tau(C_i\cup P).
\]

Define route cost
\[
J_i(P)
=
\operatorname{Burden}_i(P)
+
\lambda_R\operatorname{Redundancy}_i(P)
+
\lambda_D\operatorname{DropoutRisk}_i(P)
+
\lambda_T\operatorname{Time}_i(P)
+
\lambda_U\operatorname{Uncertainty}_i(P).
\]

The route-selection problem is
\[
P_i^{*}
\in
\argmin_{P\in\mathcal P_i(x_i,c)}
J_i(P)
\]
subject to
\[
R_c(x_i)
\subseteq
\closure_G^\tau(C_i\cup P).
\]

The route is not necessarily a simple path. It may be a directed acyclic subgraph, a collection of mutually reinforcing nodes, or a compressed strongly connected module.

If \(G\) contains a strongly connected component
\[
S\subseteq V,
\]
the competencies in \(S\) may be mutually reinforcing. Let
\[
G/{\sim_{\mathrm{SCC}}}
\]
be the condensation graph formed by contracting every strongly connected component. This graph is acyclic. Route planning can occur over the condensation graph, with local instructional policies operating inside each module.

\begin{proposition}[Existence of a minimum route in the finite case]
Suppose \(G\) is finite, the feasible route set is nonempty, and \(J_i(P)\) is real valued. Then a minimum-cost feasible route exists.
\end{proposition}

\begin{proof}
A finite graph has finitely many vertex subsets. Therefore the family of feasible routes is finite. A real-valued function on a nonempty finite set attains its minimum.
\end{proof}

Existence does not imply computational ease. Depending on the route constraints and cost structure, the optimization can contain set cover, Steiner subgraph, or resource-constrained path problems.

\subsection{Prerequisite minimality}

A prerequisite route can become an instrument of exclusion if it contains nodes not required for the present evaluation.

\begin{definition}[Locally necessary prerequisite]
A node \(v\in P\) is locally necessary for route \(P\) when
\[
R_c(x_i)
\not\subseteq
\closure_G^\tau(C_i\cup(P\setminus\{v\})).
\]
\end{definition}

\begin{definition}[Irredundant route]
A feasible route \(P\) is irredundant when every \(v\in P\) is locally necessary.
\end{definition}

\begin{proposition}[Redundancy elimination]
Every finite feasible route contains an irredundant feasible subroute.
\end{proposition}

\begin{proof}
Begin with a finite feasible route \(P\). If it is irredundant, the result follows. Otherwise, there exists \(v\in P\) such that \(P\setminus\{v\}\) remains feasible. Remove \(v\). Repeat. Since the route is finite and each step strictly decreases its cardinality, the process terminates at an irredundant feasible subroute.
\end{proof}

Irredundancy is weaker than minimum burden. Two irredundant routes may differ greatly in length, difficulty, cultural accessibility, or compatibility with the participant's existing competencies.

\subsection{Bypass through demonstration}

A participant may bypass a prerequisite node by demonstrating the operation it represents.

Let
\[
\mathcal Q(v)
\]
be a family of diagnostic tasks for competency \(v\). Let
\[
\operatorname{Pass}(y,q)
\in\{0,1\}
\]
indicate whether response \(y\) demonstrates the required transfer under task \(q\).

For posterior threshold \(\tau_v\), the node is bypassed when
\[
d_i'(v,t)\geq\tau_v.
\]

\begin{protocol}[Demonstration bypass]
Given a candidate missing node \(v\), select a diagnostic task \(q\in\mathcal Q(v)\) that also advances the participant's original inquiry. Observe \(y\), update \(d_i(v,t)\), and remove \(v\) from the route if the posterior crosses \(\tau_v\).
\end{protocol}

The diagnostic should be chosen for information gain and conversational progress, not maximum difficulty.

\subsection{Prompt selection as sequential experimental design}

Let
\[
\mathcal Q_t
\]
be the set of admissible next prompts. For \(q\in\mathcal Q_t\), let \(Y_q\) be the random response. The expected information gain about participant state is
\[
\IG_t(q)
=
H(\theta_i(t)\mid Y_{1:t})
-
\E_{Y_q}
\left[
H(\theta_i(t)\mid Y_{1:t},Y_q)
\right].
\]

Let
\[
\Prog_t(q)
\]
be expected progress on the submitted inquiry, and let
\[
\Burden_t(q)
\]
be expected participant cost. Let
\[
\Risk_t(q)
\]
measure the probability that the prompt induces avoidable frustration, repeats demonstrated material, or distorts the participant's claim.

The next prompt is selected by
\[
q_{t+1}^{*}
\in
\argmax_{q\in\mathcal Q_t}
\left[
\IG_t(q)
+
\alpha\Prog_t(q)
-
\beta\Burden_t(q)
-
\gamma\Risk_t(q)
\right].
\]

\begin{proposition}[Nonnegativity of expected information gain]
For every admissible prompt \(q\),
\[
\IG_t(q)\geq0.
\]
\end{proposition}

\begin{proof}
Expected information gain equals conditional mutual information:
\[
\IG_t(q)
=
I(\theta_i(t);Y_q\mid Y_{1:t}).
\]
Conditional mutual information is nonnegative.
\end{proof}

High information gain alone is not sufficient. A humiliating entrance examination may reveal much about the participant while contributing nothing to the inquiry. The progress and burden terms prevent diagnostic efficiency from becoming the sole objective.

\subsection{Dual-purpose prompts}

A prompt has dual purpose when it both advances the inquiry and distinguishes among competing participant-state hypotheses.

Suppose the system entertains hypotheses
\[
H_1:
\text{the participant confuses observer uncertainty with process stochasticity},
\]
and
\[
H_2:
\text{the participant understands the distinction but rejects the model}.
\]
Asking for a definition may not distinguish them. Asking the participant to identify where uncertainty resides in their own proposed mechanism can both advance the proposal and produce different likelihoods under \(H_1\) and \(H_2\).

Let
\[
\ell_j(y\mid q)
=
\Prob(Y_q=y\mid H_j).
\]
A prompt is discriminating when these likelihoods differ substantially. One measure is expected log Bayes factor:
\[
\mathcal B(q)
=
\E_{Y_q\sim H_1}
\left[
\log
\frac{\ell_1(Y_q\mid q)}
{\ell_2(Y_q\mid q)}
\right]
=
\KL
\left(
P(Y_q\mid H_1)
\middle\|
P(Y_q\mid H_2)
\right).
\]

A dual-purpose prompt seeks high \(\mathcal B(q)\) while maintaining positive \(\Prog_t(q)\).

\subsection{Identifiability of participant state}

Not every latent distinction can be recovered from conversation.

\begin{definition}[Observational equivalence of participant states]
States \(\theta,\theta'\in\ThetaSpace\) are observationally equivalent under intervention policy \(\Pi\) when
\[
P_\theta(Y_{1:\infty}\mid\Pi)
=
P_{\theta'}(Y_{1:\infty}\mid\Pi).
\]
\end{definition}

If two states are observationally equivalent, no amount of interaction under \(\Pi\) can distinguish them.

\begin{definition}[Policy-relative identifiability]
A participant-state coordinate is identifiable under policy \(\Pi\) when no two states differing in that coordinate are observationally equivalent under \(\Pi\).
\end{definition}

\begin{theorem}[Impossibility under observational equivalence]
If \(\theta\) and \(\theta'\) are observationally equivalent under every authorized prompt policy, then no estimator based only on authorized interaction history can consistently distinguish them.
\end{theorem}

\begin{proof}
Every authorized interaction history has the same probability law under \(\theta\) and \(\theta'\). Therefore any estimator has the same output distribution under both states. It cannot converge in probability to distinct correct values under the two states.
\end{proof}

The system must retain uncertainty rather than invent distinctions unsupported by observable behavior.

\subsection{Calibration error}

For competency \(v\), predictions \(d_i(v,t)\) should be calibrated. Among cases assigned probability \(p\), approximately fraction \(p\) should demonstrate the competency under an appropriate future test.

Let
\[
K_i(v,t)\in\{0,1\}
\]
be the later validated competency indicator. Define calibration function
\[
\operatorname{Cal}_v(p)
=
\E[K_i(v,t)\mid d_i(v,t)=p]-p.
\]

An integrated calibration error is
\[
\operatorname{ICE}_v
=
\int_0^1
|\operatorname{Cal}_v(p)|
\,d\mu_v(p),
\]
where \(\mu_v\) is the distribution of predicted probabilities.

Systematic overconfidence produces unnecessary bypass and possible misunderstanding. Systematic underconfidence produces excessive prerequisites and exclusion.

Group-conditional calibration can be defined as
\[
\operatorname{Cal}_{v,g}(p)
=
\E[K_i(v,t)\mid d_i(v,t)=p,A_i=g]-p.
\]
Aggregate calibration does not guarantee calibration within participant classes or communication styles.

\subsection{Prerequisite assignment as a decision under asymmetric loss}

Let the system choose between
\[
a_0=\text{bypass}
\]
and
\[
a_1=\text{assign prerequisite}.
\]
Let
\[
K\in\{0,1\}
\]
indicate actual competency.

Assigning a prerequisite when \(K=1\) imposes redundant burden \(c_{\mathrm{red}}\). Bypassing when \(K=0\) incurs misunderstanding risk \(c_{\mathrm{miss}}\).

The expected loss of bypass is
\[
L(a_0)
=
c_{\mathrm{miss}}(1-d_i(v,t)).
\]
The expected loss of assignment is
\[
L(a_1)
=
c_{\mathrm{red}}d_i(v,t).
\]
Bypass is optimal when
\[
c_{\mathrm{miss}}(1-d_i)
\leq
c_{\mathrm{red}}d_i,
\]
which is equivalent to
\[
d_i
\geq
\frac{c_{\mathrm{miss}}}
{c_{\mathrm{miss}}+c_{\mathrm{red}}}.
\]

\begin{proposition}[Cost-sensitive competency threshold]
Under the preceding loss model, the optimal bypass threshold is
\[
\tau_v^{*}
=
\frac{c_{\mathrm{miss}}}
{c_{\mathrm{miss}}+c_{\mathrm{red}}}.
\]
\end{proposition}

The threshold is not a universal constant. When misunderstanding has severe consequences, it rises. When redundant instruction is highly exclusionary or costly, it falls. Governance must make these asymmetries visible.

\subsection{The danger of self-confirming prerequisite graphs}

The prerequisite graph is learned partly from observed participant trajectories. Routing decisions influence which trajectories become observable.

Suppose the system believes \(u\) is prerequisite for \(v\). It routes nearly every participant through \(u\) before allowing attempts at \(v\). The resulting data contain few observations of participants attempting \(v\) without \(u\). The system then interprets the absence of successful bypass cases as evidence that \(u\) is necessary.

Let
\[
A_{uv}=1
\]
indicate that the system assigns \(u\) before \(v\). Let
\[
Y_v
\]
be performance on \(v\). The causal quantity of interest is
\[
\E[Y_v\mid\operatorname{do}(A_{uv}=0)]
-
\E[Y_v\mid\operatorname{do}(A_{uv}=1)].
\]
Observed association
\[
\E[Y_v\mid A_{uv}=0]
-
\E[Y_v\mid A_{uv}=1]
\]
does not identify this effect when assignment depends on participant state.

\begin{principle}[Prerequisite exploration]
Every defeasible prerequisite edge must receive some controlled bypass probability, subject to safety constraints, so that its necessity can be empirically tested rather than inferred from a policy that already presupposes it.
\end{principle}

Let
\[
\epsilon_{uv}>0
\]
be the minimum exploration probability for safe bypass of edge \(u\to v\). Without such exploration, the graph can become an institutional fossilization of its own earlier routing assumptions.

\subsection{Repair as state transformation}

Prerequisite instruction is a repair operation only when the representational and architectural state is adequate and the participant lacks a competency needed to continue.

Let
\[
\mathcal R_v
\]
be an instructional repair operator for competency \(v\). It acts on participant state:
\[
\theta_i'
=
\mathcal R_v(\theta_i,e_v),
\]
where \(e_v\) is the instructional intervention.

The repair is successful when it increases the probability of demonstrated use:
\[
d_i'(v)>d_i(v),
\]
while preserving unrelated competencies and the participant's claim unless revision is voluntarily adopted.

\begin{definition}[Non-destructive prerequisite repair]
A repair operator \(\mathcal R_v\) is non-destructive relative to a protected state set \(W\) when
\[
\operatorname{proj}_W
\mathcal R_v(\theta_i,e_v)
=
\operatorname{proj}_W\theta_i
\]
except where the participant explicitly revises a protected coordinate.
\end{definition}

The protected set can include original claim provenance, recorded disagreement, identity of the intended recipient, and distinctions not implicated in competency \(v\).

Prerequisite instruction should not silently rewrite rejection as acceptance, uncertainty as assent, or unconventional terminology as archive terminology.

\subsection{Repair, rerouting, and expansion}

Three responses to difficulty must be separated.

Repair changes participant competency while preserving the representational space and prerequisite topology:
\[
\theta_i
\mapsto
\theta_i'.
\]

Rerouting changes the chosen path through the existing graph:
\[
P_i
\mapsto
P_i'.
\]

Expansion changes the graph or representation because the existing structure cannot express the required distinction:
\[
G
\hookrightarrow
G^{+},
\qquad
M
\hookrightarrow
M^{+}.
\]

A repeated failure after demonstrated comprehension is evidence against prescribing more of the same repair.

Let
\[
e_t
\]
be the residual error after intervention \(t\). Suppose the repair operator predicts contraction:
\[
\E[\|e_{t+1}\|\mid e_t]
\leq
\rho\|e_t\|,
\qquad
0\leq\rho<1.
\]
Persistent low-period behavior
\[
e_{t+k}\approx e_t
\]
for small \(k\) violates the expected contraction pattern.

\begin{criterion}[Expansion trigger]
If the participant has demonstrated the relevant competencies, the claim remains stable under definition-preserving reconstruction, and repeated repair produces persistent structured residuals rather than contraction, then the system must test whether the prerequisite graph or CLIO quotient requires expansion.
\end{criterion}

This criterion does not prove that the participant's claim is correct. It proves that continued prerequisite assignment is no longer justified solely by the existing repair model.

\subsection{A stopping rule for prerequisite interrogation}

Without a stopping rule, the system can continue asking diagnostic questions indefinitely.

Let
\[
V_t
\]
be the expected value of improved routing after one additional prompt. Let
\[
C_t
\]
be its expected participant burden. Let
\[
R_t
\]
be the risk of distortion or disengagement.

Continue diagnostic interaction only when
\[
V_t>C_t+R_t.
\]

A more explicit value-of-information rule is
\[
\operatorname{VOI}(q)
=
\E
\left[
\min_a
\E[L(a,\theta)\mid Y_{1:t}]
\right]
-
\E_{Y_q}
\left[
\min_a
\E[L(a,\theta)\mid Y_{1:t},Y_q]
\right].
\]
Prompt \(q\) is warranted only when
\[
\operatorname{VOI}(q)
>
\Burden(q)+\Risk(q).
\]

\begin{definition}[Diagnostic exhaustion]
A conversation reaches diagnostic exhaustion when no authorized prompt has positive net value under the declared value-of-information criterion.
\end{definition}

At diagnostic exhaustion, the system must answer, defer with uncertainty, escalate, audit, or initiate an expansion review. It may not conceal uncertainty behind an endless prerequisite route.

\subsection{Participant correction of the state model}

The participant must be able to challenge \(\theta_i(t)\). Let
\[
\widetilde\theta_i(t)
\]
be the state summary exposed to the participant, restricted to information safe and useful to display.

A correction
\[
c_i^{\mathrm{corr}}
\]
becomes evidence in the posterior update:
\[
P(\theta_i\mid Y_{1:t},c_i^{\mathrm{corr}})
\propto
P(c_i^{\mathrm{corr}}\mid\theta_i)
P(\theta_i\mid Y_{1:t}).
\]

The correction need not be accepted uncritically. Self-report can be uncertain. It must nevertheless enter as evidence rather than be treated as adversarial merely because it conflicts with the current model.

\begin{principle}[Contestability of latent-state inference]
Every participant-state coordinate that materially affects access must be either directly contestable or indirectly testable through an authorized demonstration procedure.
\end{principle}

A hidden, untestable variable should not determine whether a submission approaches scarce attention.

\subsection{The prerequisite safety condition}

Let
\[
P_i^{*}
\]
be the selected prerequisite route. Define
\[
\operatorname{Need}(v\mid x_i,c)
\]
as the event that competency \(v\) is genuinely required for an authorized evaluation of \(x_i\).

A false prerequisite occurs when
\[
v\in P_i^{*}
\]
but
\[
\operatorname{Need}(v\mid x_i,c)=0.
\]
A missed prerequisite occurs when
\[
v\notin
\closure_G^\tau(C_i\cup P_i^{*})
\]
while
\[
\operatorname{Need}(v\mid x_i,c)=1.
\]

Define
\[
\operatorname{FPR}_{\mathrm{prereq}}
=
\Prob
\left(
v\in P_i^{*}
\middle|
\operatorname{Need}(v\mid x_i,c)=0
\right)
\]
and
\[
\operatorname{FNR}_{\mathrm{prereq}}
=
\Prob
\left(
v\notin
\closure_G^\tau(C_i\cup P_i^{*})
\middle|
\operatorname{Need}(v\mid x_i,c)=1
\right).
\]

The first measures unnecessary gatekeeping. The second measures underpreparation. They are asymmetric and should be reported separately.

\subsection{The participant-state obligation}

The Attention Compiler may use a participant model only under a constrained inference obligation.

\begin{principle}[Participant-state obligation]
Every participant-state inference affecting routing must remain local to the current conversational purpose, probabilistic where evidence is incomplete, revisable under later performance, separated from assent and preference, exposed sufficiently for correction, and prevented from becoming a global rank of the participant.
\end{principle}

The mathematical consequences are immediate. Competency is represented by
\[
d_i(v,t)
=
\Prob(K_i(v,t)=1\mid Y_{1:t}),
\]
not by a permanent label. Commitment is represented separately by
\[
b_i(s\mid p,t).
\]
Preference updates \(\upsilon_i\) but does not enter the missing-prerequisite computation. Required concepts form
\[
R_c(x_i),
\]
while missing prerequisites form
\[
M_i
=
R_c(x_i)
\setminus
\closure_G^\tau(C_i).
\]
The selected route minimizes declared burden under a coverage constraint. Diagnostic prompts maximize information and progress while subtracting burden and risk. Persistent structured failure after demonstrated comprehension triggers a test for expansion rather than automatic repetition of repair.

The purpose of participant-state estimation is not to decide what kind of person has submitted an idea. It is to estimate the smallest intervention needed to bring the idea into a form where the archive, the participant, and the intended recipient can encounter one another without avoidable loss.

\section{HYDRA Assembly and the Construction of Global Conversational State}

\subsection{From projected packets to distributed contextual structure}

CLIO produces a context-relative representation of a submission. That representation is not yet a global conversational state. A routing decision may depend simultaneously on the projected claim, the participant's demonstrated competencies, the archive's treatment of neighboring propositions, the recipient's unresolved questions, the provenance of evidence, the history of prior transformations, and the institutional rules governing escalation.

These structures are not naturally reducible to one flat vector without first declaring how their coordinates correspond. An archive node and a participant claim may use different vocabularies. A prerequisite node may be linked to several alternative competency witnesses. A provenance record may distinguish two propositions that semantic normalization treats as equivalent. A recipient-relative novelty judgment may depend on an unresolved question absent from the general archive.

HYDRA is introduced as the assembly layer that combines these local structures while preserving their declared compatibilities and exposing their incompatibilities.

Fix a submission \(x_i\), recipient \(r\), and context
\[
c=c_{ir}(t).
\]
Let
\[
m_i=q_c(x_i)
\]
be the CLIO packet. Let
\[
\mathscr E
\]
be a category of typed epistemic states. An object of \(\mathscr E\) contains a state space together with claim types, provenance annotations, admissibility constraints, and any local transition structure required by later stages. A morphism in \(\mathscr E\) preserves the declared protected structure.

\begin{definition}[Typed epistemic state]
A typed epistemic state is a tuple
\[
E
=
\left(
|E|,
\tau_E,
\operatorname{prov}_E,
\Adm_E,
\mathcal T_E
\right),
\]
where \(|E|\) is the underlying carrier, \(\tau_E\) assigns semantic or operational types, \(\operatorname{prov}_E\) assigns provenance data, \(\Adm_E\) specifies admissible continuations, and \(\mathcal T_E\) is a local transition structure.
\end{definition}

\begin{definition}[Structure-preserving epistemic morphism]
A morphism
\[
f:E\to E'
\]
is structure preserving when it respects every protected component:
\[
\tau_{E'}\circ f
=
\tau_E,
\]
\[
\operatorname{prov}_{E'}\circ f
=
\operatorname{prov}_E
\]
up to an explicitly recorded provenance transformation,
\[
f(\Adm_E(e))
\subseteq
\Adm_{E'}(f(e)),
\]
and
\[
f\circ\mathcal T_E
=
\mathcal T_{E'}\circ f
\]
whenever transition preservation is required.
\end{definition}

The equalities may be weakened to coherent natural transformations in architectures where types or provenance change under authorized translation. Such weakening must be explicit. A system inference cannot become participant-authored merely because two claims are semantically close.

\subsection{The contextual indexing category}

Let
\[
\C_c
\]
be the contextual indexing category. Its objects name the local structures used in the present routing problem. Denote representative objects by
\[
c_{\mathrm{sub}},
\quad
c_{\mathrm{arc}},
\quad
c_{\mathrm{pre}},
\quad
c_{\mathrm{part}},
\quad
c_{\mathrm{rec}},
\quad
c_{\mathrm{prov}},
\quad
c_{\mathrm{gov}}.
\]
These correspond to the projected submission, archive neighborhood, prerequisite neighborhood, participant-state neighborhood, recipient-state neighborhood, provenance ledger, and governance state.

This notation does not require \(\C_c\) to contain exactly seven objects. Each neighborhood may be decomposed into many local objects. A large archive may contribute one object per answer node or one object per locally coherent cluster. A prerequisite graph may contribute a subcategory rather than a single node.

Morphisms in \(\C_c\) represent declared relations among local contexts. A morphism
\[
\alpha:
c_{\mathrm{sub}}\to c_{\mathrm{arc}}
\]
may state that a reconstructed participant claim maps to an archive proposition. A morphism
\[
\beta:
c_{\mathrm{pre}}\to c_{\mathrm{part}}
\]
may connect a required competency to a demonstrated performance witness. A morphism
\[
\gamma:
c_{\mathrm{sub}}\to c_{\mathrm{prov}}
\]
may connect a claim to the record of how it entered the conversation.

The category contains compositional information. If a submission claim maps to an archive node and the archive node maps to a recipient question, their composite
\[
c_{\mathrm{sub}}
\xrightarrow{\alpha}
c_{\mathrm{arc}}
\xrightarrow{\delta}
c_{\mathrm{rec}}
\]
represents a derived route from the submission into the recipient's unresolved-question structure.

\begin{definition}[Context diagram]
A context diagram is a functor
\[
\A_c:\C_c\to\mathscr E.
\]
It assigns a typed epistemic state to every local context and a structure-preserving morphism to every declared contextual relation.
\end{definition}

Functoriality requires
\[
\A_c(\operatorname{id}_u)
=
\operatorname{id}_{\A_c(u)}
\]
and
\[
\A_c(\beta\circ\alpha)
=
\A_c(\beta)\circ\A_c(\alpha).
\]
This prevents the result of translating through several intermediate contexts from silently differing from the declared direct translation.

\subsection{Cocones and global conversational states}

A global state receiving information from every local state is represented by a cocone.

\begin{definition}[Contextual cocone]
A cocone from \(\A_c\) to an epistemic state \(Z\) is a family of morphisms
\[
\iota_u:\A_c(u)\to Z
\]
for every object \(u\in\Ob(\C_c)\), satisfying
\[
\iota_v\circ\A_c(\alpha)
=
\iota_u
\]
for every morphism
\[
\alpha:u\to v.
\]
\end{definition}

The compatibility equation says that an element entering the global state directly from \(u\) must agree with the same element first translated into \(v\) and then inserted into the global state.

A cocone can be overly large. It can contain distinctions not generated by any local state. HYDRA seeks a universal cocone.

\begin{definition}[HYDRA assembly]
A HYDRA assembly of \(\A_c\) is a colimit
\[
Z_c
=
\colim_{\C_c}\A_c
\]
with canonical structure maps
\[
\iota_u:\A_c(u)\to Z_c.
\]
For every cocone
\[
j_u:\A_c(u)\to W,
\]
there exists a unique morphism
\[
h:Z_c\to W
\]
such that
\[
j_u=h\circ\iota_u
\]
for every \(u\).
\end{definition}

The colimit is the least global object generated by the local states subject to the declared identifications. It introduces no additional distinctions beyond those required to receive the diagram.

\begin{proposition}[Uniqueness of HYDRA assembly]
If
\[
(Z,\iota_u)
\]
and
\[
(Z',\iota_u')
\]
are both colimits of \(\A_c\), then there exists a unique isomorphism
\[
\varphi:Z\to Z'
\]
such that
\[
\varphi\circ\iota_u
=
\iota_u'
\]
for every \(u\).
\end{proposition}

\begin{proof}
By the universal property of \(Z\), the cocone \((Z',\iota_u')\) induces a unique morphism
\[
\varphi:Z\to Z'
\]
satisfying
\[
\varphi\circ\iota_u=\iota_u'.
\]
By the universal property of \(Z'\), the cocone \((Z,\iota_u)\) induces a unique morphism
\[
\psi:Z'\to Z
\]
satisfying
\[
\psi\circ\iota_u'=\iota_u.
\]
Then
\[
(\psi\circ\varphi)\circ\iota_u
=
\psi\circ\iota_u'
=
\iota_u.
\]
Both \(\psi\circ\varphi\) and \(\operatorname{id}_Z\) mediate the cocone from \(\A_c\) to \(Z\). Uniqueness gives
\[
\psi\circ\varphi
=
\operatorname{id}_Z.
\]
Similarly,
\[
\varphi\circ\psi
=
\operatorname{id}_{Z'}.
\]
Therefore \(\varphi\) is an isomorphism, and uniqueness follows from the universal property.
\end{proof}

The proposition gives invariance under representational presentation. Two implementations may serialize the global state differently while representing the same assembly up to canonical isomorphism.

\subsection{Explicit construction in a concrete category}

When \(\mathscr E=\mathbf{Set}\), the colimit can be constructed explicitly. Form the disjoint union
\[
U_c
=
\coprod_{u\in\Ob(\C_c)}
\A_c(u).
\]
An element of \(U_c\) is a tagged pair
\[
(u,a),
\qquad
a\in\A_c(u).
\]

Generate an equivalence relation \(\sim_{\A_c}\) by declaring
\[
(u,a)
\sim_{\A_c}
\left(
v,\A_c(\alpha)(a)
\right)
\]
for every morphism
\[
\alpha:u\to v.
\]
Take the reflexive, symmetric, and transitive closure. Then
\[
Z_c
=
U_c/{\sim_{\A_c}}
\]
is a colimit in \(\mathbf{Set}\).

\begin{theorem}[Set-valued colimit construction]
The quotient
\[
U_c/{\sim_{\A_c}}
\]
with canonical maps
\[
\iota_u(a)=[(u,a)]
\]
satisfies the universal property of the colimit.
\end{theorem}

\begin{proof}
For a morphism \(\alpha:u\to v\),
\[
\iota_v(\A_c(\alpha)(a))
=
[(v,\A_c(\alpha)(a))]
=
[(u,a)]
=
\iota_u(a).
\]
Thus the canonical maps form a cocone.

Let
\[
j_u:\A_c(u)\to W
\]
be any other cocone. Define
\[
h:Z_c\to W
\]
by
\[
h([(u,a)])=j_u(a).
\]
To show that \(h\) is well defined, it is enough to check the generating identifications. If
\[
(u,a)
\sim_{\A_c}
(v,\A_c(\alpha)(a)),
\]
then cocone compatibility gives
\[
j_v(\A_c(\alpha)(a))
=
j_u(a).
\]
Therefore equivalent representatives receive the same image. The map satisfies
\[
h(\iota_u(a))
=
h([(u,a)])
=
j_u(a).
\]
Uniqueness follows because every equivalence class is represented by some \((u,a)\), so any mediating map must take
\[
[(u,a)]
\]
to \(j_u(a)\).
\end{proof}

This construction shows that HYDRA assembly is itself a quotient. CLIO quotients the source submission according to authorized operational equivalence. HYDRA quotients a disjoint union of local contextual states according to declared compatibility morphisms. The two quotients have different proof obligations.

CLIO must justify why two source states may be identified. HYDRA must justify why two locally situated objects denote one global object.

\subsection{Typed colimits and protected distinctions}

The set-valued construction can identify objects that should remain distinct when type and provenance are ignored. Suppose the participant states proposition \(p\), while the archive contains a refutation of \(p\). A semantic normalizer may map both to the same underlying proposition token. If the global quotient identifies assertion, refutation, and mention, the resulting state becomes incoherent.

Let the type set contain
\[
\mathcal T
=
\{
\mathsf{asserted},
\mathsf{denied},
\mathsf{queried},
\mathsf{quoted},
\mathsf{inferred},
\mathsf{superseded}
\}.
\]
A typed element is
\[
(a,\tau(a)).
\]
An epistemic morphism may identify two underlying proposition contents only when their type relation is preserved explicitly.

\begin{definition}[Type-safe identification]
An identification
\[
(u,a)\sim(v,b)
\]
is type safe when either
\[
\tau_u(a)=\tau_v(b)
\]
or the diagram contains an explicit authorized type-conversion witness
\[
\omega_{\tau}:
\tau_u(a)\Rightarrow\tau_v(b).
\]
\end{definition}

A conversion from \(\mathsf{queried}\) to \(\mathsf{answered}\) may be authorized when an answer node and supporting derivation are supplied. A conversion from \(\mathsf{participant\text{-}asserted}\) to \(\mathsf{recipient\text{-}authored}\) is not authorized merely because the recipient's archive contains a related statement.

\begin{definition}[Provenance-safe identification]
An identification is provenance safe when the global equivalence class retains the distinct derivational paths by which its representatives entered the assembly.
\end{definition}

Instead of reducing provenance to one label, the global state can attach a provenance subgraph
\[
\operatorname{Prov}_{Z_c}([a])
=
\left\{
(u,a_u,\gamma_u):
\iota_u(a_u)=[a]
\right\}.
\]
This preserves multiple origins of globally related content.

\begin{proposition}[Provenance recoverability]
Suppose every canonical map
\[
\iota_u:\A_c(u)\to Z_c
\]
is accompanied by an immutable tagged inclusion record. Then every global element \(z\in Z_c\) determines the set of local representatives and paths contributing to it.
\end{proposition}

\begin{proof}
By construction, each global element is an equivalence class of tagged local representatives. The immutable inclusion record retains the tag \(u\), local element \(a\), and relation path used to enter the class. Collecting all records whose target is \(z\) recovers the contributor set.
\end{proof}

The proposition does not imply that every local state can be reconstructed from the global state alone. It implies that the assembly preserves the ancestry of every global identification.

\subsection{Pushouts as elementary conversational gluing}

A common local pattern involves two states sharing an overlap:
\[
A
\xleftarrow{f}
O
\xrightarrow{g}
B.
\]
Here \(O\) contains the structure recognized as common to \(A\) and \(B\). The pushout
\[
A\sqcup_O B
\]
glues \(A\) and \(B\) along \(O\).

For example, let \(A\) be a participant claim graph, \(B\) an archive answer graph, and \(O\) the normalized proposition both are asserted to concern. The pushout combines the graphs while identifying the two images of \(O\).

In \(\mathbf{Set}\),
\[
A\sqcup_O B
=
(A\sqcup B)/{\sim},
\]
where
\[
f(o)\sim g(o)
\]
for every \(o\in O\).

\begin{proposition}[Pushout mediation]
Let
\[
j_A:A\to W,
\qquad
j_B:B\to W
\]
satisfy
\[
j_A\circ f
=
j_B\circ g.
\]
Then there exists a unique
\[
h:A\sqcup_O B\to W
\]
such that
\[
h\circ\iota_A=j_A,
\qquad
h\circ\iota_B=j_B.
\]
\end{proposition}

This elementary gluing pattern appears whenever two local structures are declared to describe the same object on an overlap. Its validity depends entirely on whether \(O\), \(f\), and \(g\) encode the correct shared structure.

\subsection{False overlaps}

A false overlap occurs when two local states are glued through an object that removes a distinction relevant to their compatibility.

Suppose participant claim \(A\) states
\[
p\land h,
\]
while archive claim \(B\) states
\[
p\land\neg h.
\]
If the overlap object contains only \(p\), then both local states map consistently from \(p\), and their pushout exists. The pushout does not automatically mark the conflict between \(h\) and \(\neg h\).

The categorical construction has performed exactly what the diagram requested. The diagram was insufficient.

\begin{definition}[Semantically adequate overlap]
An overlap object \(O\) between local states \(A\) and \(B\) is semantically adequate for an authorized function family \(\F\) when every distinction shared by \(A\) and \(B\) that can change an authorized global operation is represented in \(O\).
\end{definition}

This is a sufficiency condition on overlaps. HYDRA inherits the same representational obligation as CLIO, but it applies to interfaces between local states rather than to the original submission.

Let
\[
q_{AB}:A\times B\to O
\]
represent the overlap extraction. It is adequate only if
\[
q_{AB}(a,b)=q_{AB}(a',b')
\]
implies agreement under every compatibility-sensitive operation.

\begin{counterexample}[Existing colimit with invalid semantic assembly]
Let \(A=\{p\land h\}\), \(B=\{p\land\neg h\}\), and \(O=\{p\}\). Let both overlap maps forget the \(h\)-coordinate. The pushout exists in \(\mathbf{Set}\), but a global operation testing consistency of \(h\) receives an assembly containing incompatible local extensions. Therefore existence of the pushout does not imply semantic coherence.
\end{counterexample}

\subsection{Formal existence and epistemic admissibility}

The preceding counterexample motivates a distinction between formal and epistemic gluing.

\begin{definition}[Formally assembled state]
A formally assembled state is a categorical colimit of the supplied context diagram.
\end{definition}

\begin{definition}[Epistemically admissible assembly]
A formally assembled state is epistemically admissible when the diagram's objects, overlaps, compatibility maps, type transformations, and provenance relations are sufficient for every authorized global operation.
\end{definition}

Formal existence is relative to the category. Epistemic admissibility is relative to the authorized function class and the truth conditions governing the represented claims.

Let
\[
\F_c^{\mathrm{global}}
\]
be the family of authorized operations on assembled states. A diagram presentation
\[
\A_c
\]
is assembly sufficient when every global operation can be computed invariantly from its colimit without requiring distinctions omitted from the local objects or overlaps.

\begin{definition}[Assembly sufficiency]
A context diagram \(\A_c\) is sufficient for \(\F_c^{\mathrm{global}}\) when any two source-level contextual families that induce isomorphic presented diagrams receive identical outputs under every function in \(\F_c^{\mathrm{global}}\).
\end{definition}

This condition prevents diagram construction from hiding relevant contextual differences before gluing begins.

\subsection{Forced identification}

A compatibility map can be too strong. Suppose two archive nodes are merely related, but the diagram maps both from one overlap element and thereby forces them into one equivalence class.

\begin{definition}[Forced identification]
A forced identification occurs when the colimit identifies local elements \(a\in\A_c(u)\) and \(b\in\A_c(v)\) through the equivalence relation generated by the diagram, while some protected global operation distinguishes their intended referents.
\end{definition}

Forced identification is HYDRA's analogue of CLIO projection collapse.

Let
\[
\iota_u(a)=\iota_v(b)
\]
in \(Z_c\), and let
\[
F:Z_c^{\mathrm{intended}}\to Y
\]
be a protected operation on the intended global structure. If
\[
F(a)\neq F(b),
\]
then the diagram has collapsed a global distinction.

\begin{definition}[Assembly-collapse relation]
The assembly-collapse relation is
\[
\mathfrak C_{\mathrm{asm}}(\A_c)
=
\left\{
((u,a),(v,b)):
\iota_u(a)=\iota_v(b),
\;
(u,a)\not\equiv_{\F_c^{\mathrm{global}}}(v,b)
\right\}.
\]
\end{definition}

\begin{proposition}[Assembly sufficiency criterion]
The HYDRA assembly is sufficient for protected global operations if and only if
\[
\mathfrak C_{\mathrm{asm}}(\A_c)=\varnothing.
\]
\end{proposition}

The proof is identical in structure to the CLIO collapse criterion. Equal global representation must imply equality under every protected global function.

\subsection{Under-identification and fragmented global state}

A diagram can also be too weak. Two local elements may denote the same operational object but remain disconnected because the relevant compatibility morphism is absent.

\begin{definition}[Under-identification]
An under-identification occurs when local elements should be globally identified under every authorized operation but remain distinct in the assembled state.
\end{definition}

Under-identification produces representational surplus. It may cause repeated answers, duplicated prerequisite nodes, inflated novelty, or failure to recognize that several participants have independently reached the same residual.

Let
\[
(u,a)\equiv_{\F_c^{\mathrm{global}}}(v,b)
\]
but
\[
\iota_u(a)\neq\iota_v(b).
\]
Then the global state retains an unnecessary distinction.

Forced identification threatens correctness. Under-identification threatens efficiency and collective recurrence detection. A representationally simple assembly avoids both.

\subsection{The exact global quotient}

Let
\[
U_c
=
\coprod_{u\in\Ob(\C_c)}
\A_c(u)
\]
be the disjoint union of all local states. Let
\[
\equiv_c^{\mathrm{global}}
\]
be the equivalence relation induced by every authorized global operation:
\[
(u,a)\equiv_c^{\mathrm{global}}(v,b)
\iff
\forall F\in\F_c^{\mathrm{global}},
\quad
F(u,a)=F(v,b).
\]

The canonical globally simple state is
\[
Z_c^{*}
=
U_c/{\equiv_c^{\mathrm{global}}}.
\]

The implemented HYDRA diagram induces its own equivalence relation
\[
\sim_{\A_c}.
\]

\begin{theorem}[Exact assembly criterion]
The implemented assembly is exactly simple for the global operation class if and only if
\[
\sim_{\A_c}
=
\equiv_c^{\mathrm{global}}.
\]
It is sufficient when
\[
\sim_{\A_c}
\subseteq
\equiv_c^{\mathrm{global}},
\]
and it is free of representational surplus when
\[
\equiv_c^{\mathrm{global}}
\subseteq
\sim_{\A_c}.
\]
\end{theorem}

\begin{proof}
The relation \(\sim_{\A_c}\) records which local elements the implemented assembly identifies. The relation \(\equiv_c^{\mathrm{global}}\) records which elements may be identified without changing any authorized global operation.

If
\[
\sim_{\A_c}
\subseteq
\equiv_c^{\mathrm{global}},
\]
every implemented identification is authorized, so the assembly is sufficient.

If
\[
\equiv_c^{\mathrm{global}}
\subseteq
\sim_{\A_c},
\]
every operationally unnecessary distinction is removed, so there is no representational surplus.

Both properties hold exactly when the relations are equal.
\end{proof}

This theorem is the HYDRA analogue of the canonical CLIO quotient theorem. CLIO seeks the coarsest source representation sufficient for authorized operations. HYDRA seeks the coarsest global assembly that identifies exactly the locally situated objects interchangeable under authorized global operations.

\subsection{Compatibility as a typed relation}

Not every relationship should be encoded as equality. Two claims may support, contradict, generalize, refine, instantiate, translate, or condition one another.

Let
\[
\mathcal R_{\mathrm{epi}}
=
\{
\mathsf{equal},
\mathsf{supports},
\mathsf{contradicts},
\mathsf{refines},
\mathsf{generalizes},
\mathsf{instantiates},
\mathsf{translates},
\mathsf{depends},
\mathsf{supersedes}
\}
\]
be a set of epistemic relation types.

A relation-typed morphism is
\[
\alpha:u\xrightarrow{\rho}v,
\qquad
\rho\in\mathcal R_{\mathrm{epi}}.
\]
Only selected relation types induce equality in the colimit quotient. A contradiction relation must remain a relation between distinct global objects. Treating every edge as an identification would erase the very structure the graph was intended to preserve.

\begin{definition}[Identification-admitting relation]
A relation type \(\rho\) is identification admitting when its semantics warrant treating source and target as two presentations of one global object.
\end{definition}

Usually
\[
\mathsf{equal}
\]
admits identification. A verified structure-preserving translation may admit identification relative to a restricted function class. Support, contradiction, dependence, and supersession do not ordinarily admit equality.

The assembly category can encode nonidentifying relations through enriched hom-objects, labeled spans, profunctors, or a graph-valued epistemic state. The choice depends on the implementation, but the distinction between relation and equality is mandatory.

\subsection{Spans, correspondences, and uncertain matching}

Semantic correspondence is often uncertain. A participant claim may probably, but not certainly, match an archive node.

A deterministic morphism
\[
f:A\to B
\]
forces every source element to one target. A span
\[
A\xleftarrow{p}R\xrightarrow{q}B
\]
represents a relation without requiring a total function. Elements of \(R\) witness candidate correspondences.

Let
\[
w:R\to[0,1]
\]
assign confidence to each correspondence witness. A thresholded assembly can include only relations with
\[
w(r)\geq\tau_{\mathrm{glue}}.
\]
A probabilistic assembly can retain several candidate diagrams
\[
\A_c^{(1)},\ldots,\A_c^{(k)}
\]
with posterior weights
\[
\Prob(\A_c^{(j)}\mid x_i,c).
\]

The system then maintains a distribution over global states:
\[
\Prob(Z_c\mid x_i,c)
=
\sum_j
\Prob
\left(
Z_c=\colim\A_c^{(j)}
\middle|
\A_c^{(j)}
\right)
\Prob(\A_c^{(j)}\mid x_i,c).
\]

\begin{definition}[Assembly uncertainty]
Assembly uncertainty is uncertainty over local correspondence, compatibility maps, or global states induced by unresolved matching among contextual structures.
\end{definition}

This uncertainty must not be folded into participant incompetence. A failure to match a claim to the archive can arise because the archive lacks the concept, because the participant uses an unfamiliar formulation, or because the matching model is inadequate.

\subsection{Naturality of contextual translation}

Suppose two representation systems encode the same local contextual family. Let
\[
\A_c:\C_c\to\mathscr E
\]
and
\[
\B_c:\C_c\to\mathscr E'
\]
be diagrams, and let
\[
F:\mathscr E\to\mathscr E'
\]
be a translation functor.

A natural transformation
\[
\eta:F\circ\A_c\Rightarrow\B_c
\]
assigns maps
\[
\eta_u:F(\A_c(u))\to\B_c(u)
\]
such that for every
\[
\alpha:u\to v,
\]
the square commutes:
\[
\B_c(\alpha)\circ\eta_u
=
\eta_v\circ F(\A_c(\alpha)).
\]

Naturality means that translating a local state and then following a contextual relation agrees with first following the relation and then translating. Without naturality, routing can depend on arbitrary ordering of translation steps.

\begin{principle}[Translation coherence]
Every representational translation used across archive, participant, prerequisite, and recipient states must commute with the contextual relations whose meaning it claims to preserve.
\end{principle}

This principle is especially important when a language model reconstructs participant terminology into archive terminology. A translation that preserves isolated claims but changes their dependency relations is not contextually faithful.

\subsection{Incremental assembly}

Conversations evolve. Let
\[
\C_0\hookrightarrow\C_1\hookrightarrow\cdots
\]
be an increasing sequence of context categories as new claims, answers, competencies, and provenance records enter the interaction.

Let
\[
Z_t
=
\colim_{\C_t}\A_t.
\]
An inclusion
\[
J_t:\C_t\hookrightarrow\C_{t+1}
\]
and a compatible extension of diagrams induce a canonical comparison map
\[
\zeta_t:Z_t\to Z_{t+1}.
\]

\begin{definition}[Monotone assembly extension]
An incremental assembly is monotone when every comparison map
\[
\zeta_t
\]
is injective on protected global distinctions.
\end{definition}

Monotonicity does not require every old conclusion to remain true. It requires the record of the old state and its provenance not to disappear when revised. A superseded claim can remain present with status \(\mathsf{superseded}\).

\begin{proposition}[Protected-history preservation]
If every local extension map is injective on protected provenance and the new compatibility relations do not identify distinct protected histories, then the induced comparison map
\[
\zeta_t:Z_t\to Z_{t+1}
\]
is injective on protected history classes.
\end{proposition}

\begin{proof}
Suppose two protected classes in \(Z_t\) become equal in \(Z_{t+1}\). Then the equivalence relation generated at stage \(t+1\) contains a chain identifying their local representatives. If all local extension maps are injective and no newly added compatibility relation identifies distinct protected histories, no such chain exists. Therefore distinct protected classes remain distinct.
\end{proof}

\subsection{Revision without erasure}

Let \(p_t\) be a claim accepted at time \(t\), and let later evidence support \(\neg p\). A destructive update replaces \(p_t\) with \(\neg p\). A provenance-preserving update creates a new state containing
\[
(p,\mathsf{accepted},t),
\qquad
(\neg p,\mathsf{accepted},t'),
\qquad
(p,\mathsf{superseded},t').
\]

The global archive can then answer both present-tense and historical questions. The current accepted claim is \(\neg p\), while the system also records that \(p\) was previously accepted.

This is necessary for appeal. A participant challenging a disposition must be able to reconstruct which archive state and compatibility maps produced it at the time.

Let
\[
Z_c(t)
\]
be the assembled state used for a disposition. The audit record should preserve a content-addressed reference
\[
\operatorname{hash}(Z_c(t),\A_c(t),q_c(t)).
\]
Later revisions produce a new state without rewriting the earlier decision basis.

\subsection{Assembly and operational coupling}

A colimit assembles structures through declared identifications. It does not by itself establish that the later transition dynamics are operationally coupled.

Suppose
\[
Z_c
=
\colim\A_c.
\]
A transition
\[
\Hh:Z_c\to Z_c'
\]
may decompose into independent component operations even though its input was globally assembled. Conversely, components may be operationally coupled even when stored separately.

The assembly question is:
\[
\text{Do the local states form a coherent global object?}
\]
The coupling question is:
\[
\text{Does the transition depend nonseparably on their joint configuration?}
\]

These questions are related but not equivalent.

\begin{proposition}[Assembly does not imply operational coupling]
There exists a nontrivial HYDRA assembly whose transition operator is additively separable over the images of its local components.
\end{proposition}

\begin{proof}
Let
\[
Z_c=A\oplus B
\]
be a coproduct assembly with no identifications beyond tagged inclusion. Define
\[
\Hh(a,b)
=
h_A(a)+h_B(b).
\]
The assembly is nontrivial because it contains both local states, but the transition is additively separable.
\end{proof}

\begin{proposition}[Operational coupling does not imply categorical identification]
There exists a transition over separately represented states whose output contains nonseparable interactions.
\end{proposition}

\begin{proof}
Let \(A=B=\R\), retain them as separate coordinates, and define
\[
\Hh(a,b)=ab.
\]
The mixed partial is
\[
\frac{\partial^2\Hh}{\partial a\,\partial b}=1,
\]
so the transition contains an interaction, even though no categorical identification of \(a\) and \(b\) has occurred.
\end{proof}

HYDRA assembly must therefore be evaluated separately from the operational coupling developed in Section 6.

\subsection{The assembly obligation}

The HYDRA layer is governed by a distinct representational obligation.

\begin{principle}[HYDRA assembly obligation]
A global conversational state may identify local objects only through compatibility relations sufficient for every protected global operation, must preserve nonidentifying epistemic relations as relations, must retain recoverable provenance for every identification, and must expose uncertainty or obstruction whenever the declared local structures do not support a coherent assembly.
\end{principle}

The exact target is
\[
Z_c^{*}
=
\left(
\coprod_{u\in\Ob(\C_c)}
\A_c(u)
\right)
/{\equiv_c^{\mathrm{global}}}.
\]
The implemented diagram generates
\[
Z_c
=
\left(
\coprod_{u\in\Ob(\C_c)}
\A_c(u)
\right)
/{\sim_{\A_c}}.
\]
Sufficiency requires
\[
\sim_{\A_c}
\subseteq
\equiv_c^{\mathrm{global}}.
\]
Absence of representational surplus requires
\[
\equiv_c^{\mathrm{global}}
\subseteq
\sim_{\A_c}.
\]
Exact representational simplicity requires
\[
\sim_{\A_c}
=
\equiv_c^{\mathrm{global}}.
\]

Formal existence of a colimit proves only that the supplied diagram can be assembled in the chosen category. It does not prove that the diagram supplied the correct overlaps, that the compatibility maps were semantically adequate, that provenance was preserved, or that contradictory local extensions were recognized. These are epistemic proof obligations imposed on diagram construction.

CLIO governs which distinctions may be removed from the submission. HYDRA governs which locally situated structures may be treated as parts or presentations of one conversational state. The first controls compression before context is assembled. The second controls identification during assembly. Together they determine whether the state entering the transition layer is representationally simple or merely representationally convenient.

\section{Operational Coupling in PERSCEN, RAT, CoM, and RSVP}

\subsection{Assembly and transition as distinct mathematical problems}

Section 5 treated HYDRA as an assembly architecture. Local participant, archive, prerequisite, recipient, and provenance states were organized into a diagram and glued into a global state
\[
Z_c
=
\colim_{\C_c}\A_c.
\]
The existence of this global state does not determine how its components participate in subsequent dynamics.

Let
\[
Z_t
\]
be the assembled conversational state at time \(t\). HYDRA applies a transition
\[
Z_{t+1}
=
\Hh
\left(
Z_t;
\FieldPhi_t,
\FieldV_t,
\FieldS_t
\right).
\]
The ambient fields
\[
(\FieldPhi_t,\FieldV_t,\FieldS_t)
\]
represent the RSVP scalar, vector, and entropy coordinates governing the transition environment.

A transition is operationally coupled when the contribution of one component depends essentially on the state of another component under a declared family of admissible separable models. Coupling is not established merely because several inputs appear in one formula, occupy one object, or are evaluated during one time step.

\begin{definition}[Declared separability family]
Let
\[
\mathcal X
=
X_1\times\cdots\times X_k
\]
be a product input space and let
\[
\mathcal Y
\]
be an output space. A declared separability family
\[
\F_{\mathrm{sep}}
\subseteq
\{f:\mathcal X\to\mathcal Y\}
\]
is a restricted class of functions representing the forms of composition the architecture agrees to treat as operationally separable.
\end{definition}

The family must be declared before coupling is tested. If every function is allowed, the test is vacuous.

\begin{definition}[Approximate operational separability]
A transition
\[
\Hh:\mathcal X\to\mathcal Y
\]
is \(\eps_{\mathrm{sep}}\)-separable relative to loss \(L\), distribution \(P\), and family \(\F_{\mathrm{sep}}\) when
\[
\inf_{\widetilde\Hh\in\F_{\mathrm{sep}}}
\E_{X\sim P}
\left[
L
\left(
\Hh(X),
\widetilde\Hh(X)
\right)
\right]
\leq
\eps_{\mathrm{sep}}.
\]
\end{definition}

\begin{definition}[Operational coupling]
The transition is operationally coupled relative to
\[
(\F_{\mathrm{sep}},L,P,\eps_{\mathrm{sep}})
\]
when
\[
\inf_{\widetilde\Hh\in\F_{\mathrm{sep}}}
\E
\left[
L
\left(
\Hh(X),
\widetilde\Hh(X)
\right)
\right]
>
\eps_{\mathrm{sep}}.
\]
\end{definition}

Operational coupling is therefore relative, not absolute. A function can fail additive separability while satisfying multiplicative separability. A transition can be coupled under one coordinate system and separable after an authorized invertible change of variables. Every claim of coupling must specify what forms of decomposition were tested and why those forms are operationally relevant.

\subsection{The vacuity of unrestricted factorization}

Suppose one proposes that
\[
\Hh(P,R,M)
=
h
\left(
h_P(P),
h_R(R),
h_M(M)
\right)
\]
defines separability for some functions
\[
h,h_P,h_R,h_M.
\]

\begin{proposition}[Vacuity of unrestricted factorization]
Every function
\[
\Hh:P\times R\times M\to Y
\]
admits a factorization of the preceding form.
\end{proposition}

\begin{proof}
Choose
\[
h_P=\operatorname{id}_P,
\qquad
h_R=\operatorname{id}_R,
\qquad
h_M=\operatorname{id}_M,
\]
and define
\[
h(p,r,m)=\Hh(p,r,m).
\]
Then
\[
h
\left(
h_P(P),
h_R(R),
h_M(M)
\right)
=
\Hh(P,R,M).
\]
\end{proof}

The proposition shows that factorization has content only when the outer map, component maps, dimensionalities, ranks, or interaction order are restricted.

A basic family is additive:
\[
\F_{\mathrm{add}}
=
\left\{
\widetilde\Hh:
\widetilde\Hh(P,R,M)
=
h_0+h_P(P)+h_R(R)+h_M(M)
\right\}.
\]

A pairwise family allows second-order interactions:
\[
\F_{\mathrm{pair}}
=
\left\{
\widetilde\Hh:
\widetilde\Hh
=
h_0
+
h_P(P)
+
h_R(R)
+
h_M(M)
+
h_{PR}(P,R)
+
h_{PM}(P,M)
+
h_{RM}(R,M)
\right\}.
\]

A low-rank family can be written
\[
\widetilde\Hh(P,R,M)
=
\sum_{\ell=1}^{k}
a_\ell(P)b_\ell(R)c_\ell(M),
\]
with declared rank bound \(k\).

A compositionally separable sequential family can require each stage to depend only on the output of its immediate predecessor and a declared local parameter:
\[
z_{j+1}
=
T_j(z_j;\eta_j),
\qquad
\eta_j\perp z_{<j-1}
\mid z_j.
\]

These families test different notions of independence.

\subsection{Functional ANOVA decomposition}

Assume
\[
P=(P,R,M)
\]
follows a product reference measure
\[
\mu
=
\mu_P\otimes\mu_R\otimes\mu_M,
\]
and suppose
\[
\Hh\in L^2(\mu).
\]
Then \(\Hh\) admits an orthogonal functional ANOVA decomposition:
\[
\Hh
=
\Hh_{\varnothing}
+
\Hh_P
+
\Hh_R
+
\Hh_M
+
\Hh_{PR}
+
\Hh_{PM}
+
\Hh_{RM}
+
\Hh_{PRM}.
\]

The components satisfy zero-mean conditions over each of their arguments. For example,
\[
\int
\Hh_{PR}(p,r)\,d\mu_P(p)
=
0,
\]
and
\[
\int
\Hh_{PR}(p,r)\,d\mu_R(r)
=
0.
\]

The additive projection is
\[
\Pi_{\mathrm{add}}\Hh
=
\Hh_{\varnothing}
+
\Hh_P
+
\Hh_R
+
\Hh_M.
\]
The residual interaction is
\[
\Hh_{\mathrm{int}}
=
\Hh
-
\Pi_{\mathrm{add}}\Hh.
\]

\begin{definition}[Interaction energy]
The additive interaction energy is
\[
\mathcal I_{\mathrm{add}}(\Hh)
=
\left\|
\Hh-\Pi_{\mathrm{add}}\Hh
\right\|_{L^2(\mu)}^2.
\]
\end{definition}

Orthogonality gives
\[
\mathcal I_{\mathrm{add}}(\Hh)
=
\|\Hh_{PR}\|_2^2
+
\|\Hh_{PM}\|_2^2
+
\|\Hh_{RM}\|_2^2
+
\|\Hh_{PRM}\|_2^2.
\]

\begin{proposition}[Additive separability criterion]
Under the functional ANOVA assumptions,
\[
\mathcal I_{\mathrm{add}}(\Hh)=0
\]
if and only if \(\Hh\) is additively separable almost everywhere.
\end{proposition}

\begin{proof}
If \(\Hh\) is additively separable, it belongs to the additive subspace, so its orthogonal projection onto that subspace equals itself and the residual norm is zero.

Conversely, if the residual norm is zero, then
\[
\Hh=\Pi_{\mathrm{add}}\Hh
\]
almost everywhere. The right side is an additive function of \(P\), \(R\), and \(M\).
\end{proof}

The reference measure matters. A transition may appear separable on a narrow observed distribution because the system never visits states where interactions become active.

\subsection{Local mixed-partial diagnostics}

Assume
\[
\Hh:\R^{d_P}\times\R^{d_R}\times\R^{d_M}\to\R^{d_Y}
\]
is twice continuously differentiable.

For scalar output, additive separability implies vanishing mixed partials:
\[
\frac{\partial^2\Hh}
{\partial P_j\partial R_k}
=
0,
\]
\[
\frac{\partial^2\Hh}
{\partial P_j\partial M_\ell}
=
0,
\]
and
\[
\frac{\partial^2\Hh}
{\partial R_k\partial M_\ell}
=
0.
\]

\begin{proposition}[Mixed-partial necessity]
If
\[
\Hh(P,R,M)
=
h_0+h_P(P)+h_R(R)+h_M(M),
\]
then every cross-component mixed partial vanishes wherever the derivatives exist.
\end{proposition}

\begin{proof}
Differentiating \(h_P(P)\) with respect to an \(R\) or \(M\) coordinate gives zero. The same holds for the other component functions. Therefore every mixed derivative involving two distinct component groups is zero.
\end{proof}

Under suitable domain conditions, vanishing mixed partials is also sufficient.

\begin{theorem}[Mixed-partial sufficiency on a rectangular domain]
Let
\[
\Omega
=
\Omega_P\times\Omega_R\times\Omega_M
\]
be an open connected rectangular domain. Let
\[
\Hh:\Omega\to\R
\]
be twice continuously differentiable. If every cross-component mixed partial vanishes on \(\Omega\), then
\[
\Hh(P,R,M)
=
h_0+h_P(P)+h_R(R)+h_M(M)
\]
for suitable functions \(h_P,h_R,h_M\).
\end{theorem}

\begin{proof}
Fix a reference point
\[
(P_0,R_0,M_0)\in\Omega.
\]
Define
\[
h_0
=
\Hh(P_0,R_0,M_0),
\]
\[
h_P(P)
=
\Hh(P,R_0,M_0)-h_0,
\]
\[
h_R(R)
=
\Hh(P_0,R,M_0)-h_0,
\]
and
\[
h_M(M)
=
\Hh(P_0,R_0,M)-h_0.
\]
Set
\[
G(P,R,M)
=
\Hh(P,R,M)
-
h_0-h_P(P)-h_R(R)-h_M(M).
\]

The vanishing mixed partials imply that
\[
\nabla_P\Hh(P,R,M)
\]
is independent of \(R\) and \(M\). Hence
\[
\nabla_PG(P,R,M)=0.
\]
Similarly,
\[
\nabla_RG(P,R,M)=0
\]
and
\[
\nabla_MG(P,R,M)=0.
\]
Thus \(G\) is constant on the connected domain. At the reference point,
\[
G(P_0,R_0,M_0)=0.
\]
Therefore \(G=0\) throughout \(\Omega\), proving additive separability.
\end{proof}

This test is coordinate sensitive. An invertible transformation of inputs or output can alter mixed partials. It establishes separability only in the declared operational coordinates.

\subsection{PERSCEN as personalized graph dynamics}

Let
\[
G_i(t)
=
(V_i(t),E_i(t),A_i(t),F_i(t))
\]
be the participant-specific graph maintained by PERSCEN. Here \(V_i\) is the node set, \(E_i\) the edge set, \(A_i\) an adjacency or relation tensor, and \(F_i\) the node-feature field.

For participant \(i\), node \(a\), mode \(m\), and feature vectors
\[
e_{a,1},\ldots,e_{a,N_f},
\]
a recorded graph-learning update has the form
\[
[A_a^{(1)}]_{m,:}
=
\operatorname{MLP}_m
\left(
[
e_{a,1},
\ldots,
e_{a,N_f},
\operatorname{onehot}(m)
]
\right).
\]

The graph state evolves through
\[
G_i(t+1)
=
\mathcal G
\left(
G_i(t),
m_i(t),
\theta_i(t),
K_r(t)
\right).
\]

PERSCEN is participant specific because the same archive node can occupy different positions in different competency and relevance geometries. This does not authorize arbitrary personalization of truth conditions. The personalized graph determines routing salience and instructional adjacency, not whether a proposition is true.

Let
\[
P_t
=
\Psi_P(G_i(t),m_i(t))
\]
be the perception state supplied downstream.

Operational coupling can enter PERSCEN when its graph transformation depends on relevance or memory:
\[
G_i(t+1)
=
\mathcal G
\left(
G_i(t);
R_t,M_t
\right).
\]
Even if later combination is additive, the perception state already contains cross-layer dependence.

\begin{definition}[Upstream parameter coupling]
A stage output \(P_t\) is upstream coupled to \(R_t\) and \(M_t\) when
\[
P_t
=
\Psi_P
\left(
G_i(t;R_t,M_t),
m_i(t)
\right)
\]
and no member of the declared separability family can remove this dependence within tolerance.
\end{definition}

\subsection{RAT as a cue-conditioned relevance field}

Let \(x\in\R^d\) be a local representational coordinate. RAT defines a cue-density field
\[
\rho_c(x)
=
\exp
\left[
-\frac12
(x-\mu_c)^T
\Sigma_c^{-1}
(x-\mu_c)
\right],
\]
where
\[
\mu_c\in\R^d
\]
is a cue center and
\[
\Sigma_c
\]
is a positive-definite covariance or relevance-shape matrix.

The relevance-directed flow is
\[
\dot x
=
\nabla\rho_c(x).
\]
Differentiating gives
\[
\nabla\rho_c(x)
=
-\rho_c(x)
\Sigma_c^{-1}
(x-\mu_c).
\]
Thus
\[
\dot x
=
-\rho_c(x)
\Sigma_c^{-1}
(x-\mu_c).
\]

The flow points toward \(\mu_c\) in the metric induced by \(\Sigma_c^{-1}\), with magnitude modulated by the density itself.

Let
\[
V_c(x)
=
-\rho_c(x).
\]
Then
\[
\dot x
=
-\nabla V_c(x),
\]
so the dynamics are gradient descent on \(V_c\), equivalently gradient ascent on \(\rho_c\).

\begin{proposition}[Monotonic cue-density ascent]
Along any nonstationary RAT trajectory,
\[
\frac{d}{dt}\rho_c(x(t))
=
\|\nabla\rho_c(x(t))\|^2
>
0.
\]
\end{proposition}

\begin{proof}
By the chain rule,
\[
\frac{d}{dt}\rho_c(x(t))
=
\nabla\rho_c(x(t))^T\dot x(t).
\]
Since
\[
\dot x(t)=\nabla\rho_c(x(t)),
\]
the result is
\[
\frac{d}{dt}\rho_c(x(t))
=
\|\nabla\rho_c(x(t))\|^2.
\]
This is positive whenever the gradient is nonzero.
\end{proof}

With fixed \(\mu_c\) and \(\Sigma_c\), RAT can enter a larger transition additively:
\[
\dot x
=
F_P(x,P_t)
+
F_M(x,M_t)
+
\nabla\rho_c(x).
\]
At this displayed level, the relevance contribution is additively separable.

The important question is how \(\mu_c\) and \(\Sigma_c\) are generated.

\subsection{Parameter-mediated coupling}

Suppose
\[
\mu_c
=
\mu(P_t,M_t),
\qquad
\Sigma_c
=
\Sigma(P_t,M_t).
\]
Then the relevance term is
\[
R(x,P_t,M_t)
=
-\rho
\left(
x;
\mu(P_t,M_t),
\Sigma(P_t,M_t)
\right)
\Sigma(P_t,M_t)^{-1}
\left(
x-\mu(P_t,M_t)
\right).
\]

The top-level transition may still be written as a sum:
\[
\dot x
=
F_P(x,P_t)
+
F_M(x,M_t)
+
R(x,P_t,M_t).
\]
This notation does not make the transition additively separable in \(P_t\) and \(M_t\), because \(R\) contains their joint dependence.

\begin{theorem}[Hidden coupling through shared parameters]
Let
\[
\Hh(P,M)
=
h_P(P)+h_M(M)+r(\vartheta(P,M)),
\]
where
\[
\vartheta:P\times M\to\Theta
\]
is twice differentiable and \(r:\Theta\to\R\) is twice differentiable. Then
\[
\frac{\partial^2\Hh}{\partial P\,\partial M}
=
J_P\vartheta^T
H_r
J_M\vartheta
+
\nabla r^T
\frac{\partial^2\vartheta}{\partial P\,\partial M},
\]
with the appropriate tensor contractions. Consequently, the displayed additive sum can possess nonzero cross-component interaction entirely through parameter generation.
\end{theorem}

\begin{proof}
The component terms \(h_P(P)\) and \(h_M(M)\) contribute no mixed derivative. For the parameterized term, the chain rule gives
\[
\frac{\partial}{\partial P}
r(\vartheta(P,M))
=
\nabla r(\vartheta)^T
J_P\vartheta.
\]
Differentiating with respect to \(M\) yields one term from differentiating \(\nabla r\):
\[
J_P\vartheta^T
H_r
J_M\vartheta,
\]
and one term from differentiating \(J_P\vartheta\):
\[
\nabla r^T
\frac{\partial^2\vartheta}
{\partial P\,\partial M}.
\]
Their sum is the stated expression.
\end{proof}

\begin{corollary}[False negative under top-level additive audit]
An audit that tests only whether the final transition is written as a sum of named layer contributions can classify a parameter-coupled system as separable even when its mixed derivative is nonzero.
\end{corollary}

The correct audit boundary includes the parameter-construction graph. Every parameter used by a nominally local term must be traced to its generating variables.

For RAT, this means testing
\[
\mu(P,M),
\qquad
\Sigma(P,M),
\]
not only the final appearance of
\[
\nabla\rho_c(x).
\]

\subsection{Differentiating the RAT field through its center}

For clarity, assume \(\Sigma\) is fixed and
\[
\mu=\mu(P,M).
\]
Let
\[
A=\Sigma^{-1},
\qquad
d=x-\mu(P,M),
\]
so
\[
\rho=\exp\left(-\frac12d^TAd\right),
\]
and
\[
R=-\rho Ad.
\]

The derivative with respect to \(\mu\) is
\[
\frac{\partial\rho}{\partial\mu}
=
\rho Ad.
\]
For the vector field,
\[
\frac{\partial R}{\partial\mu}
=
-\frac{\partial\rho}{\partial\mu}(Ad)^T
+
\rho A.
\]
Substituting gives
\[
\frac{\partial R}{\partial\mu}
=
\rho
\left[
A-(Ad)(Ad)^T
\right].
\]

By the chain rule,
\[
\frac{\partial R}{\partial P}
=
\frac{\partial R}{\partial\mu}
\frac{\partial\mu}{\partial P},
\]
and
\[
\frac{\partial R}{\partial M}
=
\frac{\partial R}{\partial\mu}
\frac{\partial\mu}{\partial M}.
\]

If both derivatives of \(\mu\) are nonzero, the relevance field responds to perception and memory even though neither appears explicitly in the Gaussian expression after \(\mu\) is instantiated.

The mixed derivative contains terms involving
\[
\frac{\partial^2R}{\partial\mu^2},
\qquad
\frac{\partial\mu}{\partial P},
\qquad
\frac{\partial\mu}{\partial M},
\qquad
\frac{\partial^2\mu}{\partial P\,\partial M}.
\]
Thus shared cue construction creates a direct mechanism for operational coupling.

\subsection{CoM as a causal memory recurrence}

Let
\[
M_t
\]
be the Chain of Memory state. A recorded recurrence is
\[
M_{t+1}
=
\phi(M_t,u_t,c_t),
\]
where \(u_t\) is the current update and \(c_t\) is contextual input.

If
\[
c_t
=
c(P_t,R_t,Z_t),
\]
then the memory recurrence is coupled to perception, relevance, and global conversational state.

The local influence of memory state \(M_t\) on output \(y\) can be measured by
\[
I(M_t\to y)
=
\frac{\partial y}{\partial M_t}.
\]
For a recurrent trajectory,
\[
M_{t+k}
=
\phi_{t+k-1}\circ\cdots\circ\phi_t(M_t),
\]
the influence propagates through
\[
\frac{\partial M_{t+k}}{\partial M_t}
=
\prod_{j=t}^{t+k-1}
\frac{\partial\phi_j}{\partial M_j},
\]
with products ordered according to composition.

For a final output
\[
y_{t+k}
=
g(M_{t+k},Z_{t+k}),
\]
the memory contribution is
\[
\frac{\partial y_{t+k}}{\partial M_t}
=
\frac{\partial g}{\partial M_{t+k}}
\frac{\partial M_{t+k}}{\partial M_t}
+
\frac{\partial g}{\partial Z_{t+k}}
\frac{\partial Z_{t+k}}{\partial M_t}.
\]

The second term records indirect influence through the assembled conversational state.

\subsection{Memory as stabilized residue}

In a passive store, retrieval leaves state unchanged:
\[
\operatorname{read}(M_t)=M_t.
\]
In a residue-field memory, traversal reinforces the traversed path:
\[
M_t^{+}
=
\mathcal R_{\mathrm{read}}(M_t;\gamma_t),
\]
where \(\gamma_t\) is the retrieval trajectory.

If
\[
\mathcal R_{\mathrm{read}}(M_t;\gamma_t)\neq M_t,
\]
then observation of memory is an intervention on memory.

A simple reinforcement model is
\[
M_t^{+}
=
M_t+\eta K_{\gamma_t},
\]
where \(K_{\gamma_t}\) is a trace kernel associated with the accessed path.

Repeated retrieval gives
\[
M_{t+n}
=
M_t
+
\eta
\sum_{j=0}^{n-1}
K_{\gamma_{t+j}}
\]
before decay or normalization.

This creates path dependence. The probability of future retrieval may depend on accumulated trace strength:
\[
\Prob(\gamma_{t+1}=\gamma\mid M_t)
\propto
\exp
\left(
\beta
\langle
M_t,K_\gamma
\rangle
\right).
\]
Retrieval then reinforces the probability of further retrieval, producing a feedback loop.

\begin{proposition}[Self-reinforcing retrieval]
Suppose one trace \(\gamma\) is selected with softmax probability proportional to
\[
\exp(\beta m_\gamma),
\]
and retrieval updates
\[
m_\gamma^{+}=m_\gamma+\eta
\]
while leaving competing trace strengths fixed. Then the probability of selecting \(\gamma\) strictly increases after retrieval whenever
\[
\beta\eta>0.
\]
\end{proposition}

\begin{proof}
Let
\[
p_\gamma
=
\frac{e^{\beta m_\gamma}}
{e^{\beta m_\gamma}+\sum_{\delta\neq\gamma}e^{\beta m_\delta}}.
\]
After reinforcement,
\[
p_\gamma^{+}
=
\frac{e^{\beta(m_\gamma+\eta)}}
{e^{\beta(m_\gamma+\eta)}+\sum_{\delta\neq\gamma}e^{\beta m_\delta}}.
\]
Multiplication of the numerator's exponential term by
\[
e^{\beta\eta}>1
\]
strictly increases the normalized probability.
\end{proof}

This mechanism can make an archive self-confirming. Frequently retrieved answers become easier to retrieve again. Unfamiliar interpretations receive less traversal and therefore less reinforcement.

\subsection{The failure of naive intervention}

A naive coupling test holds \(M_t\) fixed while varying \(P_t\). In a residue-field memory, evaluating the system at \(M_t\) may itself update \(M_t\). Repeated evaluations do not share the same memory condition.

Let the observed transition be
\[
Z_{t+1}
=
\Hh(P_t,R_t,M_t),
\]
but let evaluation implement
\[
M_t^{+}
=
\mathcal R_{\mathrm{read}}(M_t;P_t,R_t),
\]
followed by
\[
Z_{t+1}
=
\Hh(P_t,R_t,M_t^{+}).
\]

Then the measured effect of changing \(P_t\) includes both the direct effect and the induced memory update.

\[
\frac{dZ_{t+1}}{dP_t}
=
\frac{\partial\Hh}{\partial P_t}
+
\frac{\partial\Hh}{\partial M_t^{+}}
\frac{\partial\mathcal R_{\mathrm{read}}}{\partial P_t}.
\]

\begin{proposition}[Intervention contamination]
If
\[
\frac{\partial\mathcal R_{\mathrm{read}}}{\partial P_t}\neq0
\]
and
\[
\frac{\partial\Hh}{\partial M_t^{+}}\neq0,
\]
then varying \(P_t\) while nominally holding the stored memory input fixed does not identify the direct effect
\[
\frac{\partial\Hh}{\partial P_t}.
\]
\end{proposition}

The intervention protocol must distinguish direct transition dependence from read-induced memory modification.

\subsection{Counterfactual memory substitution}

One possible protocol constructs an external memory snapshot
\[
\widetilde M
\]
that does not update under read. Define a frozen evaluator
\[
\Hh^{\mathrm{freeze}}(P,R,\widetilde M).
\]
The direct controlled contrast is
\[
\Delta_P^{\mathrm{freeze}}
=
\Hh^{\mathrm{freeze}}(P_1,R,\widetilde M)
-
\Hh^{\mathrm{freeze}}(P_0,R,\widetilde M).
\]

The natural-system contrast is
\[
\Delta_P^{\mathrm{nat}}
=
\Hh
\left(
P_1,R,
\mathcal R_{\mathrm{read}}(M;P_1,R)
\right)
-
\Hh
\left(
P_0,R,
\mathcal R_{\mathrm{read}}(M;P_0,R)
\right).
\]

Their difference,
\[
\Delta_P^{\mathrm{mem}}
=
\Delta_P^{\mathrm{nat}}
-
\Delta_P^{\mathrm{freeze}},
\]
estimates the contribution mediated through read-induced memory change, subject to fidelity of the frozen evaluator.

The protocol introduces its own question: whether freezing memory destroys the system whose coupling is being studied. If memory semantics require read reinforcement, then
\[
\Hh^{\mathrm{freeze}}
\]
is an altered architecture rather than a neutral measurement device.

\begin{definition}[Intervention fidelity]
A counterfactual intervention is faithful when it changes the targeted component while preserving every structural law not logically dependent on that component's natural variation.
\end{definition}

Frozen memory may fail this criterion. A more faithful method may clone the full pre-read state for every intervention branch, allowing each branch one natural read:
\[
M_t^{(0)}=M_t,
\qquad
M_t^{(1)}=M_t,
\]
then
\[
M_{t,+}^{(j)}
=
\mathcal R_{\mathrm{read}}
\left(
M_t^{(j)};P_j,R
\right).
\]
The comparison is made across independently cloned trajectories rather than repeated reads of one mutable state.

\subsection{Sequential composition of named HYDRA layers}

A recorded HYDRA pipeline has the form
\[
\Hh
=
\operatorname{GLU}_{\RSVP}
\circ
\mathcal M
\circ
\mathcal T
\circ
\mathcal F_a
\circ
\mathcal G_a
\circ
\rho_c.
\]

Let
\[
z_0=x,
\]
\[
z_1=\rho_c(z_0),
\]
\[
z_2=\mathcal G_a(z_1),
\]
\[
z_3=\mathcal F_a(z_2),
\]
\[
z_4=\mathcal T(z_3),
\]
\[
z_5=\mathcal M(z_4),
\]
and
\[
z_6
=
\operatorname{GLU}_{\RSVP}(z_5).
\]
Then
\[
\Hh(x)=z_6.
\]

This sequential representation differs from a simultaneous three-input abstraction
\[
\Hh(P,R,M).
\]
In the sequential pipeline, later components receive transformed outputs of earlier stages. Dependence is guaranteed by composition, but operational coupling requires more than ordinary causal succession.

Let each stage have parameters
\[
z_{j+1}
=
T_j(z_j;\eta_j).
\]
If
\[
\eta_j
\]
depends on nonadjacent earlier states,
\[
\eta_j
=
\eta_j(z_0,\ldots,z_{j-1}),
\]
then the pipeline contains skip-mediated coupling.

The total Jacobian is
\[
\frac{d z_6}{d z_0}
=
\prod_{j=0}^{5}
\frac{\partial z_{j+1}}{\partial z_j}
\]
only when stage parameters are independent of earlier states except through \(z_j\). With parameter dependence, the total derivative contains additional paths.

For stage \(j\),
\[
\frac{d z_{j+1}}{d z_0}
=
\frac{\partial T_j}{\partial z_j}
\frac{d z_j}{d z_0}
+
\frac{\partial T_j}{\partial\eta_j}
\frac{d\eta_j}{d z_0}.
\]
The second term is parameter-mediated influence bypassing the displayed stage input.

\begin{definition}[Pipeline-locality condition]
A sequential pipeline is parameter local when
\[
\eta_j\perp(z_0,\ldots,z_{j-1})
\mid z_j
\]
for every stage \(j\), or deterministically when
\[
\eta_j=\eta_j(z_j)
\]
with no additional dependence on earlier states.
\end{definition}

Failure of pipeline locality is one mechanism by which HYDRA can be operationally coupled even when every stage looks locally simple.

\subsection{Adjacent-stage substitution tests}

For adjacent stages
\[
z_j\xrightarrow{T_j}z_{j+1}\xrightarrow{T_{j+1}}z_{j+2},
\]
construct a substituted intermediate state
\[
\widetilde z_{j+1}
\]
drawn from another trajectory or generated by a controlled perturbation.

The downstream sensitivity is
\[
\Delta_{j+1}
=
T_{j+1}(\widetilde z_{j+1};\eta_{j+1})
-
T_{j+1}(z_{j+1};\eta_{j+1}).
\]

To test whether the meaning of \(z_{j+1}\) depends on hidden alignment with upstream state \(z_j\), compare ordinary substitution with context-preserving substitution:
\[
\Delta_{j+1}^{\mathrm{ordinary}}
\]
and
\[
\Delta_{j+1}^{\mathrm{aligned}}.
\]

If a numerically similar intermediate state from another trajectory produces inadmissible output because it lacks shared provenance or parameter alignment, the interface between stages carries more structure than the displayed state variable records.

\begin{definition}[Interface sufficiency]
The interface representation \(z_{j+1}\) is sufficient for downstream stage \(T_{j+1}\) when any two upstream histories producing the same \(z_{j+1}\) induce the same downstream behavior under every authorized continuation.
\end{definition}

This is CLIO sufficiency applied inside HYDRA. A pipeline can appear modular while its interface states omit hidden dependencies required by later stages.

\begin{proposition}[Hidden-state criterion for nonmodularity]
Suppose two upstream histories \(h\) and \(h'\) produce the same displayed interface state:
\[
z_{j+1}(h)=z_{j+1}(h').
\]
If a later stage produces different outputs when the histories are preserved,
\[
T_{j+1}(z_{j+1};h)
\neq
T_{j+1}(z_{j+1};h'),
\]
then the displayed interface is insufficient and the stages are not modular relative to the authorized continuation.
\end{proposition}

The missing history may be encoded in shared parameters, mutable memory, provenance, or ambient fields.

\subsection{RSVP-governed transition fields}

Let the assembled state occupy a domain
\[
\Omega\subseteq\R^d.
\]
RSVP supplies a scalar field
\[
\FieldPhi:\Omega\times\R_{\geq0}\to\R,
\]
a vector field
\[
\FieldV:\Omega\times\R_{\geq0}\to\R^d,
\]
and an entropy or dispersion field
\[
\FieldS:\Omega\times\R_{\geq0}\to\R.
\]

A generic coupled field system may be written
\[
\partial_t\FieldPhi
=
F_\Phi
\left(
\FieldPhi,
\FieldV,
\FieldS,
Z
\right),
\]
\[
\partial_t\FieldV
=
F_V
\left(
\FieldPhi,
\FieldV,
\FieldS,
Z
\right),
\]
\[
\partial_t\FieldS
=
F_S
\left(
\FieldPhi,
\FieldV,
\FieldS,
Z
\right).
\]

A structural relaxation functional can be defined by
\[
\mathscr S[\FieldS]
=
\int_\Omega
\Psi(\FieldS(x,t))
\,dx.
\]
Under a diffusion-like evolution
\[
\partial_t\FieldS
=
\gamma\Delta\FieldS
\]
with no-flux boundary condition
\[
\nabla\FieldS\cdot n=0
\quad
\text{on }\partial\Omega,
\]
the quadratic energy
\[
\mathscr E_S(t)
=
\frac12
\int_\Omega
\FieldS(x,t)^2
\,dx
\]
satisfies
\[
\frac{d\mathscr E_S}{dt}
=
-\gamma
\int_\Omega
\|\nabla\FieldS(x,t)\|^2
\,dx.
\]

\begin{proposition}[Diffusive dissipation]
Under the preceding diffusion equation and no-flux boundary condition,
\[
\frac{d\mathscr E_S}{dt}
\leq0.
\]
\end{proposition}

\begin{proof}
Differentiate:
\[
\frac{d\mathscr E_S}{dt}
=
\int_\Omega
\FieldS\partial_t\FieldS
\,dx
=
\gamma
\int_\Omega
\FieldS\Delta\FieldS
\,dx.
\]
Integration by parts gives
\[
\int_\Omega
\FieldS\Delta\FieldS
\,dx
=
-\int_\Omega
\|\nabla\FieldS\|^2
\,dx
+
\int_{\partial\Omega}
\FieldS\nabla\FieldS\cdot n\,d\sigma.
\]
The boundary term vanishes under the no-flux condition, yielding
\[
\frac{d\mathscr E_S}{dt}
=
-\gamma
\int_\Omega
\|\nabla\FieldS\|^2
\,dx
\leq0.
\]
\end{proof}

This is a mathematical structural model of relaxation. It does not establish that the conversational fields are physical entropy, fluid velocity, or gravitational potential.

\subsection{Field-mediated coupling}

Even if PERSCEN, RAT, and CoM contribute separately to a local state, shared RSVP fields can mediate interaction.

Suppose
\[
\FieldPhi_t
=
\Phi(P_t,R_t,M_t),
\]
\[
\FieldV_t
=
V(P_t,R_t,M_t),
\]
and
\[
\FieldS_t
=
S(P_t,R_t,M_t).
\]
The transition is
\[
Z_{t+1}
=
\Hh_0(Z_t)
+
B_\Phi\FieldPhi_t
+
B_V\FieldV_t
+
B_S\FieldS_t.
\]

This expression is additive in the field values. It is not necessarily separable in \(P_t,R_t,M_t\) because the fields are joint functions of them.

\begin{proposition}[Mediator coupling]
Let
\[
Y
=
h_0+
B\,\eta(P,R),
\]
where \(B\neq0\). If
\[
\frac{\partial^2\eta}{\partial P\,\partial R}\neq0
\]
or if the first derivatives interact through a nonlinear \(B\), then \(Y\) is not additively separable in \(P\) and \(R\) on any neighborhood where the resulting mixed derivative is nonzero.
\end{proposition}

Field mediation can therefore produce global operational coupling without direct pairwise terms in the final equation.

\subsection{Useful coupling and harmful coupling}

Nonseparability is not automatically beneficial. A coupled transition can improve contextual coherence, or it can make errors propagate across layers.

Let
\[
J_{\mathrm{task}}(\Hh)
\]
measure task performance,
\[
J_{\mathrm{robust}}(\Hh)
\]
measure robustness,
\[
J_{\mathrm{audit}}(\Hh)
\]
measure auditability, and
\[
\mathcal I(\Hh)
\]
measure interaction energy.

Useful coupling cannot be defined by
\[
\mathcal I(\Hh)>0
\]
alone. A comparative criterion is required.

\begin{definition}[Functionally useful coupling]
Coupling is functionally useful relative to a separable baseline class when every sufficiently accurate separable approximation incurs a declared deficit:
\[
\inf_{\widetilde\Hh\in\F_{\mathrm{sep}}}
\left[
J_{\mathrm{task}}(\Hh)
-
J_{\mathrm{task}}(\widetilde\Hh)
\right]
>
\delta_{\mathrm{use}}
\]
subject to matched complexity, data, and attention constraints.
\end{definition}

A harmful coupling can increase sensitivity:
\[
\left\|
\frac{\partial Z_{t+1}}{\partial M_t}
\right\|
\gg1,
\]
causing a small memory distortion to alter relevance, graph state, and final disposition.

Define local amplification
\[
\mathcal A_t
=
\|J_{\Hh}(Z_t)\|_{\mathrm{op}}.
\]
If
\[
\mathcal A_t>1
\]
persistently, perturbations can grow. Coupling may then reduce modular containment of errors.

\subsection{The operational-coupling obligation}

\begin{principle}[Operational-coupling obligation]
HYDRA may be described as operationally coupled only relative to a declared separability family, coordinate system, loss, state distribution, and tolerance, and only after the audit includes shared parameter generation, mutable memory effects, intermediate-interface sufficiency, and field-mediated dependence.
\end{principle}

The relevant evidence can include residual interaction
\[
\inf_{\widetilde\Hh\in\F_{\mathrm{sep}}}
\E[L(\Hh,\widetilde\Hh)],
\]
mixed derivatives
\[
\frac{\partial^2\Hh}{\partial P\,\partial R},
\quad
\frac{\partial^2\Hh}{\partial P\,\partial M},
\quad
\frac{\partial^2\Hh}{\partial R\,\partial M},
\]
parameter paths
\[
P,M\to\mu_c,\Sigma_c\to\nabla\rho_c,
\]
memory-mediated paths
\[
P,R\to\mathcal R_{\mathrm{read}}(M)\to\Hh,
\]
pipeline skip paths
\[
z_k\to\eta_j\to z_{j+1},
\]
and field-mediated paths
\[
P,R,M\to(\FieldPhi,\FieldV,\FieldS)\to Z_{t+1}.
\]

No single diagnostic is sufficient in every architecture. Mixed partials are local and coordinate sensitive. separability risk depends on the declared family and observed distribution. Intervention can alter residue-field memory. Top-level equations can hide shared-parameter dependence. Categorical assembly can succeed while transition dynamics remain separable.

The correct conclusion is not that HYDRA is coupled because its components have been placed under one name. The correct conclusion is that HYDRA supplies several explicit mechanisms through which operational coupling can arise, and each mechanism generates a separate empirical and mathematical proof obligation.

\section{Gluing Obstructions and Obstruction-Sensitive Routing}

\subsection{Local coherence does not imply global coherence}

A conversational architecture can contain several locally coherent structures that fail when assembled. A participant claim may agree with one archive node under one vocabulary and with another archive node under a second vocabulary, while the two translations disagree on their common overlap. A prerequisite path may be valid locally but conflict with a competency equivalence used elsewhere. Two pieces of evidence may independently support a proposition while relying on mutually incompatible provenance assumptions.

These failures are not ordinary projection collapse. Projection collapse occurs when a representation identifies states that an authorized operation must distinguish. A gluing obstruction occurs when local states and their pairwise compatibility data fail to determine one coherent global state.

Let
\[
X_c
\]
be the contextual space associated with a submission and let
\[
\mathcal U
=
\{U_a\}_{a\in A}
\]
be a cover by local conversational regions. A region may represent an archive neighborhood, participant interpretation, prerequisite module, provenance neighborhood, recipient-relative problem context, or governance context.

Let
\[
\mathscr F
\]
be a presheaf assigning admissible epistemic states to every region:
\[
U
\longmapsto
\mathscr F(U).
\]
For an inclusion
\[
V\subseteq U,
\]
there is a restriction map
\[
\rho_{UV}:
\mathscr F(U)\to\mathscr F(V).
\]

The presheaf laws require
\[
\rho_{UU}
=
\operatorname{id}_{\mathscr F(U)}
\]
and
\[
\rho_{UW}
=
\rho_{VW}\circ\rho_{UV}
\]
whenever
\[
W\subseteq V\subseteq U.
\]

A section
\[
s_a\in\mathscr F(U_a)
\]
is a locally coherent interpretation or state over region \(U_a\).

\begin{definition}[Pairwise compatibility]
Local sections \(s_a\in\mathscr F(U_a)\) and \(s_b\in\mathscr F(U_b)\) are compatible when
\[
\rho_{U_a,U_a\cap U_b}(s_a)
=
\rho_{U_b,U_a\cap U_b}(s_b).
\]
\end{definition}

\begin{definition}[Global gluing]
A family
\[
\{s_a\in\mathscr F(U_a)\}_{a\in A}
\]
admits a global gluing when there exists
\[
s\in\mathscr F(X_c)
\]
such that
\[
\rho_{X_c,U_a}(s)=s_a
\]
for every \(a\in A\).
\end{definition}

If \(\mathscr F\) is a sheaf, every pairwise-compatible family has a unique global gluing. In an implemented conversational system, however, the assigned structure may be only a presheaf, the overlap maps may be incomplete, or the local compatibility judgments may themselves be uncertain.

\subsection{The sheaf condition}

\begin{definition}[Epistemic sheaf]
A presheaf \(\mathscr F\) is an epistemic sheaf when every cover
\[
U=\bigcup_{a\in A}U_a
\]
satisfies locality and gluing. Locality requires that two global sections agreeing on every \(U_a\) are equal. Gluing requires that every compatible family of local sections arises from a global section.
\end{definition}

Formally, locality is
\[
\left[
\forall a,
\quad
\rho_{U,U_a}(s)
=
\rho_{U,U_a}(s')
\right]
\implies
s=s'.
\]
Gluing is
\[
\left[
\forall a,b,
\quad
s_a|_{U_a\cap U_b}
=
s_b|_{U_a\cap U_b}
\right]
\implies
\exists s\in\mathscr F(U)
\quad
s|_{U_a}=s_a.
\]

\begin{proposition}[Uniqueness of global gluing]
If \(\mathscr F\) satisfies locality and a compatible family admits two global gluings \(s,s'\), then
\[
s=s'.
\]
\end{proposition}

\begin{proof}
Both global sections restrict to \(s_a\) on every region \(U_a\). Therefore
\[
s|_{U_a}
=
s'|_{U_a}
\]
for every \(a\). Locality implies \(s=s'\).
\end{proof}

The sheaf condition is an architectural ideal. It says that local conversational states contain enough overlap information to support unique global reconstruction whenever they are mutually compatible.

\subsection{Failure of pairwise compatibility}

For local sections \(s_a,s_b\), define the overlap discrepancy
\[
\delta_{ab}
=
\rho_{a,ab}(s_a)
-
\rho_{b,ab}(s_b)
\]
when \(\mathscr F(U_a\cap U_b)\) has an additive structure.

Pairwise compatibility is equivalent to
\[
\delta_{ab}=0.
\]

For metric-valued sections, define
\[
e_{ab}
=
d_{ab}
\left(
\rho_{a,ab}(s_a),
\rho_{b,ab}(s_b)
\right).
\]
Then exact compatibility is
\[
e_{ab}=0,
\]
while approximate compatibility may require
\[
e_{ab}\leq\eps_{ab}.
\]

The total pairwise discrepancy can be written
\[
E_{\mathrm{pair}}
=
\sum_{a<b}
w_{ab}e_{ab}^2.
\]

A nonzero pairwise discrepancy localizes an immediate incompatibility. It does not require higher-order analysis.

\begin{example}[Definition mismatch]
Let \(U_a\) contain the participant's use of a term \(w\), and let \(U_b\) contain the archive's use of the same term. Suppose
\[
s_a(w)=\delta_i(w),
\qquad
s_b(w)=\delta_r^K(w),
\]
with
\[
\delta_i(w)\neq\delta_r^K(w).
\]
If the overlap map treats the lexical token \(w\) as sufficient evidence of semantic identity, then
\[
e_{ab}>0.
\]
The obstruction is a failed translation on the overlap, not evidence that the participant lacks the concept.
\end{example}

\subsection{Pairwise agreement with higher-order failure}

Pairwise agreement need not ensure a globally coherent assignment when correspondences contain transformations rather than literal equality.

Suppose local states are related by transition maps
\[
g_{ab}
\]
on overlaps. Pairwise compatibility may mean
\[
s_b
=
g_{ab}s_a.
\]
On a triple overlap, consistency additionally requires
\[
g_{ca}g_{bc}g_{ab}
=
\operatorname{id}.
\]

\begin{example}[Cyclic translation inconsistency]
Let three local vocabularies use scalar coordinates. Suppose the translations are
\[
g_{12}(x)=x+1,
\]
\[
g_{23}(x)=x+1,
\]
and
\[
g_{31}(x)=x-1.
\]
Each pair admits a valid translation. Around the cycle,
\[
g_{31}\circ g_{23}\circ g_{12}(x)
=
x+1.
\]
The composition is not the identity. No single global coordinate assignment realizes all three pairwise translations simultaneously.
\end{example}

This failure resembles holonomy. Local transport around a closed path returns a different state. The residual transformation records global inconsistency invisible to any one edge.

\begin{definition}[Conversational holonomy]
For a cycle
\[
\gamma:
a_0\to a_1\to\cdots\to a_k=a_0,
\]
the conversational holonomy is
\[
\operatorname{Hol}_{\gamma}
=
g_{a_{k-1}a_k}
\circ\cdots\circ
g_{a_0a_1}.
\]
The cycle is flat when
\[
\operatorname{Hol}_{\gamma}
=
\operatorname{id}.
\]
\end{definition}

Nontrivial holonomy is evidence that local translations cannot all be treated as restrictions of one globally consistent identification.

\subsection{The Čech complex}

Assume \(\mathscr F\) is a sheaf of abelian groups or modules. The Čech cochain groups of the cover \(\mathcal U\) are
\[
\check C^p(\mathcal U,\mathscr F)
=
\prod_{a_0<\cdots<a_p}
\mathscr F
\left(
U_{a_0}\cap\cdots\cap U_{a_p}
\right).
\]

A \(0\)-cochain is a family of local sections
\[
s=(s_a)_a.
\]
Its coboundary is the \(1\)-cochain
\[
(\delta^0s)_{ab}
=
s_b|_{U_a\cap U_b}
-
s_a|_{U_a\cap U_b}.
\]
Thus
\[
\delta^0s=0
\]
exactly when the local sections agree on every pairwise overlap.

For a \(1\)-cochain \(\omega=(\omega_{ab})\), the next coboundary is
\[
(\delta^1\omega)_{abc}
=
\omega_{bc}|_{U_{abc}}
-
\omega_{ac}|_{U_{abc}}
+
\omega_{ab}|_{U_{abc}},
\]
where
\[
U_{abc}
=
U_a\cap U_b\cap U_c.
\]

The identity
\[
\delta^{p+1}\circ\delta^p=0
\]
implies that every coboundary is a cocycle.

Define
\[
Z^p
=
\ker\delta^p,
\qquad
B^p
=
\im\delta^{p-1}.
\]
The Čech cohomology group is
\[
\check H^p(\mathcal U,\mathscr F)
=
Z^p/B^p.
\]

\subsection{Obstruction classes}

Suppose pairwise translations or corrections define a \(1\)-cocycle
\[
\omega\in Z^1.
\]
If
\[
\omega\in B^1,
\]
then there exists a \(0\)-cochain \(b=(b_a)\) such that
\[
\omega=\delta^0b.
\]
The local corrections \(b_a\) absorb the pairwise discrepancies.

If
\[
[\omega]\neq0
\quad
\text{in }
\check H^1(\mathcal U,\mathscr F),
\]
no family of local corrections eliminates the discrepancy. The class
\[
[\omega]
\]
is a gluing obstruction.

\begin{definition}[First-order conversational obstruction]
A first-order conversational obstruction is a nonzero class
\[
[\omega]
\in
\check H^1(\mathcal U,\mathscr F)
\]
generated by pairwise compatibility transformations that cannot be removed by local reparameterization.
\end{definition}

\begin{theorem}[Vanishing criterion for locally correctable discrepancies]
Let
\[
\omega\in Z^1(\mathcal U,\mathscr F).
\]
There exists a family of local corrections
\[
b\in\check C^0(\mathcal U,\mathscr F)
\]
such that
\[
\omega=\delta^0b
\]
if and only if
\[
[\omega]=0
\]
in \(\check H^1(\mathcal U,\mathscr F)\).
\end{theorem}

\begin{proof}
By definition,
\[
[\omega]=0
\]
exactly when \(\omega\) belongs to the subgroup of \(1\)-coboundaries:
\[
\omega\in B^1=\im\delta^0.
\]
This holds exactly when there exists \(b\in\check C^0\) with
\[
\omega=\delta^0b.
\]
\end{proof}

The theorem separates local repair from architectural revision. If the obstruction class vanishes, local coordinate corrections can reconcile the sections. If it does not vanish, no collection of independent local adjustments suffices.

\subsection{Obstruction is not successful coupling}

A nonzero obstruction must not be misinterpreted as proof of operational coupling.

\begin{proposition}[Obstruction distinction]
The condition
\[
[\omega]\neq0
\]
proves failure of the attempted gluing under the supplied cover, coefficient structure, and compatibility maps. It does not prove that a successfully assembled transition is operationally nonseparable.
\end{proposition}

\begin{proof}
The obstruction is defined before a valid global section has been constructed. Operational coupling concerns the separability of a transition on a successfully defined global state. Since the obstructed assembly may produce no admissible global state at all, no conclusion about the separability of a valid global transition follows solely from nonvanishing of \([\omega]\).
\end{proof}

Obstruction and coupling answer different questions:
\[
[\omega]\neq0
\]
asks whether local states can be assembled under the current compatibility structure, while
\[
\inf_{\widetilde\Hh\in\F_{\mathrm{sep}}}
\E[L(\Hh,\widetilde\Hh)]
>
\eps_{\mathrm{sep}}
\]
asks whether a defined transition resists a declared separable approximation.

\subsection{Obstruction is not projection collapse}

Projection collapse concerns fibers of
\[
q_c:X\to M_c.
\]
Gluing obstruction concerns compatibility among sections over a cover.

A projection can be sufficient while the contextual assembly is obstructed. Conversely, an insufficient projection may produce an unobstructed but incorrect assembly.

\begin{proposition}[Logical independence]
There exist systems with no projection collapse and nonzero gluing obstruction, and systems with projection collapse and zero gluing obstruction.
\end{proposition}

\begin{proof}
For the first case, let \(q_c=\operatorname{id}_X\), so no distinctions are collapsed. Choose three local translation maps with nontrivial cyclic holonomy. The resulting assembly is obstructed.

For the second case, let \(q_c\) identify two submissions requiring different dispositions. Let the contextual cover contain only one local region. Every local section glues trivially, so the obstruction vanishes, while projection collapse remains.
\end{proof}

An observed routing failure must therefore be localized before repair.

\subsection{Quantitative obstruction energy}

Exact cohomology can be difficult to compute in noisy systems. A quantitative relaxation measures the distance of a cocycle from the space of coboundaries.

Let
\[
\omega\in\check C^1(\mathcal U,\mathscr F)
\]
be an observed discrepancy cochain. Define
\[
\mathcal E_{\mathrm{obs}}(\omega)
=
\inf_{b\in\check C^0}
\|
\omega-\delta^0b
\|^2_W,
\]
where \(\|\cdot\|_W\) is a weighted norm over overlaps.

\begin{definition}[Obstruction energy]
The obstruction energy is
\[
\mathcal E_{\mathrm{obs}}(\omega)
=
\operatorname{dist}^2
\left(
\omega,
\im\delta^0
\right).
\]
\end{definition}

If \(\omega\) is a cocycle, then
\[
\mathcal E_{\mathrm{obs}}(\omega)=0
\]
exactly when its cohomology class vanishes, assuming the image of \(\delta^0\) is closed in the chosen normed space.

\begin{proposition}[Least-squares local repair]
In a finite-dimensional Hilbert cochain complex, the minimizing local correction \(b^{*}\) satisfies the normal equation
\[
(\delta^0)^{*}\delta^0b^{*}
=
(\delta^0)^{*}\omega.
\]
\end{proposition}

\begin{proof}
Minimize
\[
J(b)
=
\|\omega-\delta^0b\|^2.
\]
Differentiating in direction \(h\) gives
\[
DJ(b)[h]
=
-2
\left\langle
\omega-\delta^0b,
\delta^0h
\right\rangle
=
-2
\left\langle
(\delta^0)^{*}
(\omega-\delta^0b),
h
\right\rangle.
\]
At a minimizer this vanishes for every \(h\), so
\[
(\delta^0)^{*}
(\omega-\delta^0b^{*})
=
0.
\]
Rearranging yields the normal equation.
\end{proof}

The residual
\[
\omega_{\perp}
=
\omega-\delta^0b^{*}
\]
is orthogonal to locally correctable discrepancies. It is the harmonic or obstruction-bearing component under the selected inner product.

\subsection{The Čech Laplacian}

Define the \(p\)-th cochain Laplacian
\[
\Delta_p
=
(\delta^{p-1})(\delta^{p-1})^{*}
+
(\delta^p)^{*}\delta^p.
\]
For \(p=1\),
\[
\Delta_1
=
\delta^0(\delta^0)^{*}
+
(\delta^1)^{*}\delta^1.
\]

A harmonic \(1\)-cochain satisfies
\[
\Delta_1\omega=0.
\]
In finite-dimensional Hodge theory,
\[
\ker\Delta_1
\cong
\check H^1(\mathcal U,\mathscr F).
\]

Thus the persistent component of discrepancy that is both closed and orthogonal to local corrections represents the obstruction class.

\begin{definition}[Harmonic obstruction score]
Let
\[
\Pi_{\mathrm{harm}}
\]
be the orthogonal projection onto \(\ker\Delta_1\). Define
\[
\mathcal H_{\mathrm{obs}}(\omega)
=
\|
\Pi_{\mathrm{harm}}\omega
\|^2.
\]
\end{definition}

This score distinguishes noise-like discrepancy from topologically persistent incompatibility under the chosen complex.

\subsection{Approximate sheaf consistency}

Exact equality is often unrealistic. Let each overlap carry a tolerance
\[
\eps_{ab}\geq0.
\]
A family of sections is approximately compatible when
\[
e_{ab}
\leq
\eps_{ab}
\]
for every overlap.

Define normalized discrepancy
\[
\widetilde e_{ab}
=
\frac{e_{ab}}{\eps_{ab}+\eta},
\]
where \(\eta>0\) prevents division by zero in numerical implementations.

An approximate gluing minimizes
\[
J(s)
=
\sum_a
\lambda_a
d_a
\left(
s|_{U_a},
s_a
\right)^2
+
\sum_{a<b}
w_{ab}
e_{ab}(s)^2.
\]

The minimum
\[
J^{*}
=
\inf_{s\in\mathscr F(X_c)}
J(s)
\]
measures the best global fit. A small value does not by itself guarantee epistemic admissibility. The loss may permit averaging contradictory claims into a state representing neither.

\begin{principle}[No semantic averaging of protected contradiction]
When two protected local states stand in a contradiction relation, approximate gluing must preserve the contradiction as a typed relation rather than replace both states with a numerical compromise.
\end{principle}

For scalar sensor estimates, averaging may be appropriate. For an asserted proposition and its negation, it is generally not.

\subsection{The nerve of the context cover}

The nerve
\[
N(\mathcal U)
\]
is a simplicial complex with one vertex for every region \(U_a\), one edge for every nonempty pairwise overlap, one triangle for every nonempty triple overlap, and higher simplices for higher-order intersections.

The nerve records the topology of contextual interaction. A cycle in the nerve can support nontrivial holonomy. Missing overlaps can disconnect evidence that should interact. Excessive overlap can create unnecessary compatibility constraints.

Let
\[
N_k(\mathcal U)
\]
be the set of \(k\)-simplices. A \(1\)-cochain assigns discrepancy data to edges. A \(2\)-cochain assigns higher-order consistency data to triangles.

If the nerve is a tree, every edge discrepancy in an abelian translation system can be removed by choosing local origins, because there are no cycles.

\begin{proposition}[Vanishing first obstruction on a tree]
Let \(N(\mathcal U)\) be a connected tree and let the coefficient system be constant and abelian. Then
\[
\check H^1(\mathcal U,\mathscr F)=0.
\]
\end{proposition}

\begin{proof}
Choose a root vertex. For any \(1\)-cochain \(\omega\), define a \(0\)-cochain \(b\) recursively along the unique path from the root. Set the root value to zero. For each oriented edge \(a\to b\), define
\[
b_b=b_a+\omega_{ab}.
\]
Because the graph contains no cycles, this definition is unambiguous. Then
\[
(\delta^0b)_{ab}
=
b_b-b_a
=
\omega_{ab}
\]
on every edge. Hence every \(1\)-cochain is a coboundary, so \(H^1=0\).
\end{proof}

Nontrivial first obstruction requires cycles or nonconstant coefficient structure. This gives the architecture a concrete diagnostic: incompatibility that persists only after several contextual routes close into a cycle is structurally different from an isolated pairwise mismatch.

\subsection{Obstruction localization}

A global obstruction should be localized to the smallest subcover that supports it.

Let
\[
\mathcal V\subseteq\mathcal U
\]
be a subcover or subfamily of regions. Restriction induces
\[
\operatorname{res}_{\mathcal U,\mathcal V}:
\check H^1(\mathcal U,\mathscr F)
\to
\check H^1(\mathcal V,\mathscr F).
\]

\begin{definition}[Obstruction support]
A subfamily \(\mathcal V\) supports obstruction class \([\omega]\) when
\[
\operatorname{res}_{\mathcal U,\mathcal V}[\omega]\neq0.
\]
It is minimally supporting when no proper subfamily retains a nonzero restricted class.
\end{definition}

A minimally supporting subfamily identifies the archive nodes, prerequisite modules, participant interpretations, or provenance assumptions jointly responsible for the incompatibility.

Finding a minimum obstruction support can be combinatorially difficult. Practical systems may use cycle bases, sparse harmonic representatives, or iterative deletion tests.

\subsection{Routing under obstruction}

A conventional router selects the most likely answer or path. An obstruction-sensitive router first tests whether the local structures required by that path admit coherent assembly.

Let
\[
\mathcal R_c
\]
be the set of candidate routes. Each route \(\rho\in\mathcal R_c\) induces a local cover
\[
\mathcal U_\rho
\]
and discrepancy cochain
\[
\omega_\rho.
\]

Define route utility
\[
U(\rho)
=
V(\rho)
-
C(\rho)
-
\lambda_{\mathrm{obs}}
\mathcal E_{\mathrm{obs}}(\omega_\rho)
-
\lambda_{\mathrm{coll}}
\Lambda_c(\rho),
\]
where \(V\) is expected conversational value, \(C\) is burden or attention cost, \(\mathcal E_{\mathrm{obs}}\) is obstruction energy, and \(\Lambda_c\) is projection-collapse risk.

The selected route is
\[
\rho^{*}
\in
\argmax_{\rho\in\mathcal R_c}
U(\rho)
\]
subject to protected constraints.

A route with high topical relevance may be rejected when it requires an incoherent assembly. A lower-similarity route may be preferred because it preserves provenance and admits consistent gluing.

\begin{definition}[Obstruction-sensitive disposition]
A disposition is obstruction sensitive when it conditions its confidence and action on the gluing status of every contextual structure required to justify it.
\end{definition}

The system should not declare a question answered when the answer depends on mutually incompatible archive sections.

\subsection{Disposition under unresolved obstruction}

Let
\[
O_c\in\{0,1\}
\]
indicate whether the required assembly is obstructed. A deterministic answer function
\[
D_c:Z_c\to\mathcal D_r
\]
is undefined when no admissible \(Z_c\) exists.

The architecture must therefore extend the disposition space:
\[
\mathcal D_r^{+}
=
\mathcal D_r
\cup
\{
d_{\mathrm{obstruction}},
d_{\mathrm{expansion}},
d_{\mathrm{human\text{-}reconcile}}
\}.
\]

An obstruction disposition reports that the system cannot coherently combine the relevant local states. It does not report that the participant is wrong.

\begin{principle}[Obstruction transparency]
When a routing decision fails because required local states do not glue, the reported reason must identify the assembly failure rather than translating it into participant uncertainty, low competency, or low novelty.
\end{principle}

This principle prevents an architectural defect from becoming a personal classification.

\subsection{Projection-induced false obstruction}

An apparent obstruction can be created by insufficient local representations. Suppose two local sections would agree if a discarded coordinate were restored.

Let
\[
q_a:\widetilde{\mathscr F}(U_a)\to\mathscr F(U_a)
\]
be local projections. Full sections
\[
\widetilde s_a
\]
may admit a global gluing, while projected sections
\[
s_a=q_a(\widetilde s_a)
\]
fail compatibility because the projection altered the translation maps or removed a disambiguating distinction.

Conversely, projection can hide a real obstruction by collapsing conflicting states into one local value.

\begin{definition}[Projection-stable obstruction]
An obstruction is projection stable when it persists under every admissible refinement of the local representations within the declared repair class.
\end{definition}

Before expanding the global architecture, the system should test whether the obstruction is an artifact of CLIO collapse at one of the local interfaces.

\subsection{Obstruction-induced false novelty}

Suppose a submission cannot be matched to an archive answer because two translation paths disagree. A naive novelty detector may classify the residual as new. The novelty may instead arise from failed gluing.

Let
\[
N_{\mathrm{raw}}(x)
\]
be apparent novelty after retrieval, and let
\[
O(x)
\]
measure assembly obstruction. Define obstruction-adjusted novelty only after resolving or reporting the obstruction:
\[
N_{\mathrm{adj}}(x)
=
N_{\mathrm{raw}}(x)
-
\Psi(O(x)),
\]
where \(\Psi\) does not necessarily subtract numerical value but marks the portion of novelty attribution unsupported while assembly remains unresolved.

A claim should not be declared familiar merely because one path maps it to the archive, nor novel merely because another path fails to map it. The disagreement itself is the primary result.

\subsection{Repairing a vanishing obstruction}

If
\[
[\omega]=0,
\]
there exists
\[
b\in\check C^0
\]
with
\[
\omega=\delta^0b.
\]
The local corrections \(b_a\) can be interpreted as revised definitions, coordinate changes, provenance alignments, or contextual translations.

The repaired sections are
\[
s_a'
=
s_a-b_a
\]
under additive conventions. Their pairwise discrepancy is
\[
\delta^0s'
=
\delta^0s-\delta^0b.
\]
If
\[
\omega=\delta^0s,
\]
then
\[
\delta^0s'=0.
\]

\begin{proposition}[Local correction of a coboundary discrepancy]
If the observed discrepancy is a coboundary, local correction produces pairwise-compatible sections.
\end{proposition}

This is genuine repair. The local representational spaces and cover remain adequate. Only their coordinate alignment required correction.

\subsection{Responding to a nonvanishing obstruction}

If
\[
[\omega]\neq0,
\]
no local correction within the current cochain complex eliminates the incompatibility. The architecture must change one of the structures defining the complex.

It may refine the cover:
\[
\mathcal U\rightsquigarrow\mathcal U^{+}.
\]
It may enlarge the coefficient structure:
\[
\mathscr F\hookrightarrow\mathscr F^{+}.
\]
It may add a missing intermediate node that resolves a translation:
\[
U_a\cap U_b
\hookrightarrow
U_{ab}^{+}.
\]
It may weaken an incorrect equality relation into a typed nonidentifying relation. It may revise a compatibility map. It may preserve several incompatible global branches rather than forcing one state.

These are expansion operations, not local repair.

\subsection{Branched global states}

When incompatible local states cannot be reconciled, the correct global representation may be branched.

Let
\[
Z_c^{(1)},\ldots,Z_c^{(k)}
\]
be maximal coherent assemblies. The global conversational object becomes
\[
\mathfrak Z_c
=
\{Z_c^{(1)},\ldots,Z_c^{(k)}\}
\]
with branch weights or unresolved relations.

A disposition can then be branch conditional:
\[
D_c^{(j)}
=
D(Z_c^{(j)}).
\]
If all branches produce the same disposition,
\[
D_c^{(1)}
=
\cdots
=
D_c^{(k)},
\]
the obstruction may be irrelevant to that operation even though it matters to others.

\begin{definition}[Disposition-irrelevant obstruction]
An obstruction is irrelevant to operation \(f\) when every maximal coherent branch yields the same value under \(f\).
\end{definition}

This is an \(F\)-sufficiency principle applied to obstruction. Not every unresolved incompatibility blocks every action. A system may answer a narrow factual question consistently across all branches while deferring a broader theoretical synthesis.

\begin{proposition}[Branch-invariant disposition]
If
\[
f(Z_c^{(j)})=y
\]
for every maximal coherent branch \(j\), then \(f\) is well defined on the branched state by
\[
f(\mathfrak Z_c)=y.
\]
\end{proposition}

The architecture should therefore block only operations sensitive to the unresolved branch distinction.

\subsection{Obstruction recurrence}

Let
\[
\omega_1,\omega_2,\ldots
\]
be obstruction representatives observed across submissions. Repeated support on the same contextual subcomplex suggests a structural defect rather than isolated misunderstanding.

Define support frequency
\[
F_t(\mathcal V)
=
\sum_{j\leq t}
\ind
\left[
\mathcal V
\text{ supports }
[\omega_j]
\right].
\]
A change statistic is
\[
\Delta F_t(\mathcal V)
=
F_t(\mathcal V)-F_{t-1}(\mathcal V).
\]

A high recurrence rate at one archive-prerequisite interface may indicate a missing distinction, an invalid translation, or an unstable compatibility rule.

\begin{definition}[Distributional obstruction]
A distributional obstruction is a recurrent obstruction pattern whose significance appears at the aggregate level even when no individual submission warrants escalation.
\end{definition}

This gives distributional novelty a structural interpretation. Repetition can reveal topology in the knowledge architecture.

\subsection{Obstruction and appeal}

An appeal may assert that the system reconstructed a claim incorrectly, applied an irrelevant prerequisite, or treated incompatible archive statements as one resolved answer.

Let the original route use cover \(\mathcal U\) and diagram \(\A_c\). An appeal supplies a candidate refinement
\[
\A_c\to\A_c^{\mathrm{appeal}}
\]
or a challenge to one compatibility map.

The appeal succeeds structurally when it demonstrates that the original disposition depended on a nonzero obstruction, a forced identification, or an omitted branch relevant to the disposition.

The audit record must preserve
\[
(\mathcal U,\A_c,\omega,D_c)
\]
as they existed when the decision was made. Recomputing against a revised archive does not establish whether the original decision was warranted.

\subsection{The obstruction-sensitive routing theorem}

Let \(\mathcal R_c\) be a set of candidate routes. For route \(\rho\), let
\[
q_\rho
\]
be its local projection family, let
\[
\A_\rho
\]
be its context diagram, and let
\[
\omega_\rho
\]
be its discrepancy cocycle.

\begin{theorem}[Conditional validity of an obstruction-sensitive route]
Suppose route \(\rho\) satisfies
\[
\mathfrak C_c(q_\rho)=\varnothing,
\]
\[
[\omega_\rho]=0,
\]
and its repaired global assembly is sufficient for every operation used by its disposition. Then the route's disposition is invariant under local coordinate changes represented by \(0\)-cochains and under replacement of the diagram by a canonically equivalent admissible presentation.
\end{theorem}

\begin{proof}
The absence of projection collapse ensures that every required local representation preserves the distinctions used by downstream operations.

Since
\[
[\omega_\rho]=0,
\]
there exists a \(0\)-cochain \(b\) such that
\[
\omega_\rho=\delta^0b.
\]
Applying the local correction gives compatible sections. The sheaf gluing condition then supplies a global section, unique under locality.

Changing local coordinates by another \(0\)-cochain alters the discrepancy by a coboundary and therefore leaves the cohomology class unchanged. Canonically equivalent admissible diagrams have canonically isomorphic global assemblies. Sufficiency of the repaired assembly ensures every authorized disposition function factors through this global state. Hence the resulting disposition is invariant under the stated changes.
\end{proof}

\subsection{The obstruction obligation}

\begin{principle}[Obstruction obligation]
Before an Attention Compiler treats locally coherent contextual states as one warranted global basis for action, it must establish compatibility on their overlaps, distinguish locally correctable discrepancy from nonvanishing obstruction, preserve contradictory states as typed relations when reconciliation is unwarranted, and report unresolved assembly failure as an architectural condition rather than a participant defect.
\end{principle}

The relevant hierarchy is:
\[
\text{pairwise discrepancy}
\longrightarrow
\text{cocycle condition}
\longrightarrow
\text{cohomology class}
\longrightarrow
\text{obstruction-sensitive disposition}.
\]
A pairwise discrepancy may be corrected locally. A nonclosed discrepancy indicates failure even to define coherent cycle data. A closed but nonexact discrepancy defines a genuine obstruction class. A nonzero class requires expansion, branching, or human reconciliation.

The Attention Compiler should not ask only whether a submission resembles an existing answer. It must ask whether the submission, the answer, the participant model, the prerequisite path, the provenance record, and the recipient's unresolved context can inhabit one coherent global state. When they cannot, the failure to glue is itself part of the intellectual content.

\section{Repair, Rerouting, and Representational Expansion}

\subsection{Three responses to failure}

A system that encounters disagreement, nonconvergence, classification uncertainty, or failed gluing must determine which object should change. It can alter the participant state while preserving the representation, alter the route through an existing structure, or enlarge the structure itself.

Let the full conversational state be
\[
\Xi_t
=
\left(
x_i,
q_t,
M_t,
G_t,
\theta_i(t),
\A_t,
Z_t,
\Hh_t
\right).
\]
Here \(q_t\) is the active CLIO projection, \(M_t\) its representation space, \(G_t\) the prerequisite graph, \(\theta_i(t)\) the participant-state estimate, \(\A_t\) the HYDRA diagram, \(Z_t\) the assembled state, and \(\Hh_t\) the transition operator.

A repair operation changes a state inside the existing architecture:
\[
\mathcal R:
\Xi_t\mapsto\Xi_{t+1},
\]
subject to
\[
q_{t+1}=q_t,
\qquad
M_{t+1}=M_t,
\qquad
G_{t+1}=G_t,
\]
except for ordinary parameter updates already licensed by the architecture.

A rerouting operation changes the chosen path through an existing graph or diagram:
\[
\mathcal P:
P_t\mapsto P_{t+1},
\]
while preserving the available nodes and representational distinctions.

An expansion operation changes the space of available distinctions or compatibilities:
\[
\mathcal X:
(M_t,G_t,\A_t)
\mapsto
(M_{t+1},G_{t+1},\A_{t+1}),
\]
where at least one of the inclusions
\[
M_t\hookrightarrow M_{t+1},
\]
\[
G_t\hookrightarrow G_{t+1},
\]
or
\[
\A_t\hookrightarrow\A_{t+1}
\]
is proper in an operationally relevant sense.

\begin{definition}[Repair regime]
A failure lies in a repair regime when the declared representation and compatibility architecture are sufficient for the required continuation, but the current state, parameter estimate, local alignment, or participant competency has not yet reached an admissible configuration.
\end{definition}

\begin{definition}[Expansion regime]
A failure lies in an expansion regime when no admissible change confined to the current representation, graph, or compatibility architecture can support the required continuation.
\end{definition}

\begin{definition}[Rerouting regime]
A failure lies in a rerouting regime when the current architecture contains an admissible continuation, but the selected path does not reach it efficiently or safely.
\end{definition}

These regimes are properties of a failure relative to a declared operation class. A state can lie in a repair regime for one task and an expansion regime for another.

\subsection{Residual dynamics}

Let
\[
e_t
\]
be a residual measuring the failure of the current state to satisfy the active conversational constraints. It may contain several components:
\[
e_t
=
\left(
e_t^{\mathrm{proj}},
e_t^{\mathrm{comp}},
e_t^{\mathrm{glue}},
e_t^{\mathrm{disp}},
e_t^{\mathrm{prov}}
\right).
\]
These represent projection error, participant-competency error, gluing discrepancy, disposition error, and provenance inconsistency.

A repair operator induces dynamics
\[
e_{t+1}
=
F_{\mathcal R}(e_t;\eta_t),
\]
where \(\eta_t\) contains contextual parameters and observations.

The standard repair assumption is contractivity.

\begin{assumption}[Local repair contraction]
There exists a neighborhood \(U\) of the admissible set and a constant
\[
0\leq\rho<1
\]
such that
\[
d
\left(
F_{\mathcal R}(e),
\mathcal E^{*}
\right)
\leq
\rho
d(e,\mathcal E^{*})
\]
for every \(e\in U\), where \(\mathcal E^{*}\) is the set of admissible residual states.
\end{assumption}

\begin{proposition}[Geometric convergence under contraction]
Under local repair contraction, if \(e_0\in U\) and the repair trajectory remains in \(U\), then
\[
d(e_t,\mathcal E^{*})
\leq
\rho^t
d(e_0,\mathcal E^{*}).
\]
\end{proposition}

\begin{proof}
Apply the contraction inequality recursively:
\[
d(e_1,\mathcal E^{*})
\leq
\rho d(e_0,\mathcal E^{*}),
\]
\[
d(e_2,\mathcal E^{*})
\leq
\rho d(e_1,\mathcal E^{*})
\leq
\rho^2d(e_0,\mathcal E^{*}),
\]
and by induction,
\[
d(e_t,\mathcal E^{*})
\leq
\rho^td(e_0,\mathcal E^{*}).
\]
\end{proof}

Observed failure of contraction can arise because the state lies outside the contraction region, because the repair operator is poorly tuned, because the target moves, or because no admissible target exists inside the current architecture.

\subsection{Convergence does not imply representational adequacy}

Let
\[
\pi_{\mathcal A}
\]
be a stipulated admissibility quotient. Let
\[
\widehat\pi_n
\]
be an estimator minimizing an empirical objective. Under compactness, uniform convergence, identification, and separation, one may have
\[
\widehat\pi_n
\xrightarrow{p}
\pi_{\mathcal A}.
\]

This is a consistency statement about estimation of the stipulated target. It is not a proof that the target preserves every distinction required by the true downstream function class.

Let
\[
\F_{\mathrm{train}}
\subsetneq
\F_{\mathrm{true}}.
\]
Suppose \(\pi_{\mathcal A}\) is sufficient for \(\F_{\mathrm{train}}\) but not for \(\F_{\mathrm{true}}\).

\begin{theorem}[Convergent insufficiency]
There exists a sequence of estimators converging consistently to \(\pi_{\mathcal A}\) while the limiting projection has strictly positive irreducible risk for some operation in \(\F_{\mathrm{true}}\).
\end{theorem}

\begin{proof}
Choose
\[
f^{*}\in
\F_{\mathrm{true}}\setminus
\F_{\mathrm{train}}
\]
and points \(x,y\in X\) satisfying
\[
\pi_{\mathcal A}(x)
=
\pi_{\mathcal A}(y)
\]
but
\[
f^{*}(x)\neq f^{*}(y).
\]
Consistency gives
\[
\widehat\pi_n
\xrightarrow{p}
\pi_{\mathcal A}.
\]
In the limit, \(x\) and \(y\) remain identified. Any decoder using only the limiting quotient must assign the same output to both, so it cannot reproduce \(f^{*}\) on both points. Therefore the irreducible risk for \(f^{*}\) is positive under any distribution assigning positive mass to both.
\end{proof}

The theorem explains why successful convergence can be harmful. If the target quotient is too coarse, stronger convergence removes precisely the residual variation that could have revealed misspecification.

\subsection{Oversmoothing}

Let
\[
q_t:X\to M
\]
be a sequence of progressively smoothed projections. Suppose smoothing reduces a training discrepancy
\[
L_{\mathrm{train}}(q_{t+1})
\leq
L_{\mathrm{train}}(q_t).
\]
Let
\[
\Lambda_{\mathrm{prot}}(q_t)
\]
be protected collapse risk.

\begin{definition}[Oversmoothing]
A projection update is oversmoothing when
\[
L_{\mathrm{train}}(q_{t+1})
<
L_{\mathrm{train}}(q_t)
\]
but
\[
\Lambda_{\mathrm{prot}}(q_{t+1})
>
\Lambda_{\mathrm{prot}}(q_t).
\]
\end{definition}

The update improves the optimized proxy while worsening preservation of protected distinctions.

Let
\[
S_\alpha(q)
\]
be a smoothing operator indexed by strength \(\alpha\). A regularization path is
\[
q_\alpha=S_\alpha(q_0).
\]
The derivative
\[
\frac{d}{d\alpha}
L_{\mathrm{train}}(q_\alpha)
\]
may be negative while
\[
\frac{d}{d\alpha}
\Lambda_{\mathrm{prot}}(q_\alpha)
\]
is positive.

\begin{criterion}[Oversmoothing warning]
If increased smoothing produces simultaneous reduction in training loss and increase in audited within-fiber disposition divergence, the system must suspend smoothing and test for expansion.
\end{criterion}

\subsection{Oscillation as diagnostic structure}

Suppose repeated repair produces a sequence
\[
e_0,e_1,e_2,\ldots
\]
that does not converge. Nonconvergence alone is not evidence for expansion. Random noise, unstable step size, delayed feedback, or a noncontractive implementation can also produce oscillation.

A diagnostic oscillation must be structured.

\begin{definition}[Approximate period-\(k\) recurrence]
A residual trajectory has approximate period \(k\) over window \(T\) when
\[
\frac{1}{T-k}
\sum_{t=1}^{T-k}
\|e_{t+k}-e_t\|^2
\leq
\eps_k.
\]
\end{definition}

Define the lag-\(k\) recurrence score
\[
R_k(T)
=
1-
\frac{
\sum_{t=1}^{T-k}
\|e_{t+k}-e_t\|^2
}{
\sum_{t=1}^{T-k}
\|e_t-\bar e\|^2+\eta
}.
\]
High \(R_k(T)\) indicates low-period structure.

A discrete Fourier transform can reveal spectral concentration:
\[
\widehat e(\omega_j)
=
\sum_{t=0}^{T-1}
e_t
e^{-2\pi ijt/T}.
\]
Define normalized spectral concentration
\[
C_{\mathrm{spec}}
=
\frac{
\max_{j\neq0}
\|\widehat e(\omega_j)\|^2
}{
\sum_{j\neq0}
\|\widehat e(\omega_j)\|^2
}.
\]

Persistent high recurrence and spectral concentration distinguish low-period dynamics from diffuse noise.

\subsection{Reachability-relevant oscillation}

An oscillating distinction matters only if the alternating states support different admissible futures.

Let
\[
\Reach(e)
\]
be the set of admissible continuations reachable from residual state \(e\). Suppose a period-two orbit alternates between
\[
e^{(0)}
\quad\text{and}\quad
e^{(1)}.
\]

\begin{definition}[Reachability-relevant oscillation]
An oscillation is reachability relevant when
\[
\Reach(e^{(0)})
\neq
\Reach(e^{(1)}).
\]
\end{definition}

If
\[
\Reach(e^{(0)})
=
\Reach(e^{(1)}),
\]
the oscillation may be a coordinate artifact or operationally irrelevant variation. If the reachable futures differ, collapsing the orbit to its average can destroy a consequential distinction.

Let
\[
\bar e
=
\frac12
\left(
e^{(0)}+e^{(1)}
\right).
\]
A smoothing repair may replace the orbit with \(\bar e\). This is admissible only if
\[
\Reach(\bar e)
\]
preserves the required continuations.

\begin{counterexample}[Destructive orbit averaging]
Let residual state be scalar and let admissible continuation be determined by sign:
\[
\Reach(e)
=
\begin{cases}
\{a\}, & e>0,\\
\{b\}, & e<0,\\
\varnothing, & e=0.
\end{cases}
\]
A period-two orbit alternates between \(+1\) and \(-1\). Averaging gives
\[
\bar e=0.
\]
The original orbit preserves access to \(a\) and \(b\) on alternating steps, while the smoothed state reaches neither. Smoothing eliminates both continuations.
\end{counterexample}

The example captures the danger of treating oscillation amplitude as error without examining what the alternating states encode.

\subsection{TARTAN and CLIO under persistent obstruction}

Let
\[
\mathcal T
\]
be a TARTAN tiling operator and
\[
\pi
\]
a CLIO projection. Consider the alternating update
\[
z_{t+1}
=
\pi\circ\mathcal T(z_t).
\]
Suppose the tiling introduces local distinctions required for gluing, while the projection repeatedly identifies them. The next tiling reconstructs the distinctions because local compatibility still requires them:
\[
z_t
\xrightarrow{\mathcal T}
\widetilde z_t
\xrightarrow{\pi}
z_{t+1}.
\]

A period-two pattern can arise:
\[
z^{(0)}
\xrightarrow{\mathcal T}
\widetilde z^{(0)}
\xrightarrow{\pi}
z^{(1)}
\xrightarrow{\mathcal T}
\widetilde z^{(1)}
\xrightarrow{\pi}
z^{(0)}.
\]

The oscillation may be interpreted as a disagreement between the distinctions generated by local gluing and the distinctions admitted by the global quotient.

Let
\[
P_{\mathcal T}
\]
be the partition induced by TARTAN compatibility requirements, and let
\[
P_\pi
\]
be the partition induced by the CLIO quotient. If
\[
P_\pi
\not\preceq
P_{\mathcal T},
\]
then CLIO identifies states TARTAN must distinguish.

\begin{proposition}[Partition incompatibility as a source of reconstruction cycles]
Suppose TARTAN reconstructs every distinction in partition \(P_{\mathcal T}\), while CLIO removes every distinction not present in \(P_\pi\). If
\[
P_\pi\not\preceq P_{\mathcal T},
\]
then there exists a distinction that TARTAN restores and CLIO removes under every alternating application.
\end{proposition}

\begin{proof}
The failure
\[
P_\pi\not\preceq P_{\mathcal T}
\]
means some block of \(P_\pi\) contains two states lying in different blocks of \(P_{\mathcal T}\). TARTAN separates those states because its compatibility structure requires the distinction. CLIO subsequently identifies them because they occupy one \(P_\pi\)-block. Therefore the distinction is restored and removed on every alternating cycle.
\end{proof}

This gives a structural explanation for persistent low-period behavior. The operators are not merely failing to converge numerically. They implement incompatible equivalence relations.

\subsection{Diagnostic decomposition of nonconvergence}

Let the observed nonconvergence statistic be
\[
N_T.
\]
It may arise from several latent causes:
\[
H_{\mathrm{num}},
\quad
H_{\mathrm{noise}},
\quad
H_{\mathrm{move}},
\quad
H_{\mathrm{collapse}},
\quad
H_{\mathrm{obstruct}}.
\]
These correspond to numerical instability, stochastic noise, moving target, projection collapse, and gluing obstruction.

The system maintains posterior probabilities
\[
\Prob(H_j\mid e_{0:T},\mathcal A_{0:T}).
\]

Evidence for numerical instability includes sensitivity to step size:
\[
N_T(\eta)\to0
\quad
\text{as}\quad
\eta\to0.
\]
Evidence for noise includes low spectral concentration and weak recurrence. Evidence for a moving target includes correlation between residual direction and documented context change. Evidence for projection collapse includes positive fiberwise disposition divergence. Evidence for obstruction includes persistent harmonic discrepancy or nontrivial cycle holonomy.

A repair decision should condition on these competing explanations rather than infer expansion directly from nonconvergence.

\subsection{A decision-theoretic regime classifier}

Let the action set be
\[
\mathcal A_{\mathrm{reg}}
=
\{
a_{\mathrm{repair}},
a_{\mathrm{reroute}},
a_{\mathrm{expand}},
a_{\mathrm{defer}},
a_{\mathrm{human}}
\}.
\]
Let
\[
H
\]
be the latent failure regime. Let
\[
L(a,H)
\]
be the loss of taking action \(a\) when the true regime is \(H\).

The Bayes action is
\[
a^{*}
\in
\argmin_{a\in\mathcal A_{\mathrm{reg}}}
\E[L(a,H)\mid\mathcal O_T],
\]
where \(\mathcal O_T\) contains residual history, audit evidence, obstruction measures, participant demonstrations, and route history.

Expansion generally has higher immediate cost than repair:
\[
L(a_{\mathrm{expand}},H_{\mathrm{repair}})
>
L(a_{\mathrm{repair}},H_{\mathrm{repair}}).
\]
Repeated repair under genuine expansion can have much higher cumulative cost:
\[
L(a_{\mathrm{repair}},H_{\mathrm{expand}})
\gg
L(a_{\mathrm{expand}},H_{\mathrm{expand}}).
\]

The system should therefore expand only after sufficient evidence, but it must not set the expansion threshold so high that endless repair becomes the default.

\subsection{Expected value of expansion}

Let
\[
\mathcal X_j
\]
be a candidate expansion. It may add a representation coordinate, prerequisite node, archive distinction, compatibility relation, branch, or provenance type.

Let
\[
C(\mathcal X_j)
\]
be its implementation and complexity cost. Let
\[
\Delta\Lambda_j
=
\Lambda_{\mathrm{before}}
-
\Lambda_{\mathrm{after},j}
\]
be expected reduction in protected collapse. Let
\[
\Delta O_j
\]
be expected reduction in obstruction energy. Let
\[
\Delta V_j
\]
be expected gain in conversational or intellectual value.

Define
\[
U(\mathcal X_j)
=
\alpha\Delta\Lambda_j
+
\beta\Delta O_j
+
\gamma\Delta V_j
-
C(\mathcal X_j).
\]
Expansion is warranted when
\[
\max_jU(\mathcal X_j)>0
\]
and protected constraints are satisfied.

\begin{definition}[Minimal warranted expansion]
A candidate \(\mathcal X^{*}\) is minimally warranted when it has positive utility and no proper subexpansion achieves the same protected-risk and obstruction reductions within declared tolerance.
\end{definition}

Representational simplicity requires the smallest expansion that restores the missing authorized distinction. It does not license indefinite accumulation of explanatory dimensions.

\subsection{Expansion of the CLIO quotient}

Suppose audit identifies a feature
\[
\varphi:X\to W
\]
with
\[
I(D_c^{*};\varphi(X)\mid q_c(X))>0.
\]
Define the refined quotient
\[
q_c^{+}(x)
=
(q_c(x),\varphi(x)).
\]

The induced partition satisfies
\[
P_{q_c^{+}}
\preceq
P_{q_c}.
\]
The refinement splits some old fibers but never merges them.

\begin{proposition}[Monotonic reduction of conditional disposition entropy]
For
\[
q_c^{+}(X)
=
(q_c(X),\varphi(X)),
\]
\[
H(D_c^{*}\mid q_c^{+}(X))
\leq
H(D_c^{*}\mid q_c(X)).
\]
Equality holds if and only if
\[
I(D_c^{*};\varphi(X)\mid q_c(X))=0.
\]
\end{proposition}

\begin{proof}
The chain rule gives
\[
I(D_c^{*};\varphi(X)\mid q_c(X))
=
H(D_c^{*}\mid q_c(X))
-
H(D_c^{*}\mid q_c(X),\varphi(X)).
\]
Conditional mutual information is nonnegative. Substituting
\[
q_c^{+}(X)
=
(q_c(X),\varphi(X))
\]
proves the inequality and equality condition.
\end{proof}

Expansion can reduce collapse but increase representational complexity:
\[
I(X;q_c^{+}(X))
\geq
I(X;q_c(X)).
\]
The added information is
\[
I(X;\varphi(X)\mid q_c(X)).
\]

A useful expansion maximizes recovered protected information per added representational cost:
\[
\eta_{\mathrm{exp}}
=
\frac{
I(D_c^{*};\varphi(X)\mid q_c(X))
}{
I(X;\varphi(X)\mid q_c(X))+\epsilon
}.
\]

\subsection{Expansion of the prerequisite graph}

Suppose participants repeatedly fail transition
\[
u\to v.
\]
The observed failure rate is
\[
F(u,v)
=
\frac{
N_{\mathrm{fail}}(u\to v)
}{
N_{\mathrm{attempt}}(u\to v)
}.
\]

A high failure rate may indicate poor instruction, an invalid dependency, participant heterogeneity, or a missing intermediate competency \(w\).

Graph expansion introduces
\[
u\to w\to v.
\]
Let the expanded graph be
\[
G^{+}
=
(V\cup\{w\},E^{+}).
\]

The expansion is supported when performance satisfies
\[
\Prob(\operatorname{Pass}(v)\mid\operatorname{Pass}(w))
>
\Prob(\operatorname{Pass}(v)\mid\operatorname{Pass}(u))
\]
under appropriately controlled comparisons, and when \(w\) represents a coherent transferable competency rather than a post hoc label for failure.

The alternative may be edge removal. If many participants reach \(v\) without \(u\), the dependency
\[
u\to v
\]
may be unnecessary. Expansion and simplification can occur together:
\[
G
\rightsquigarrow
G^{+}
\setminus\{u\to v\}.
\]

\subsection{Expansion of the HYDRA cover}

Suppose local regions \(U_a\) and \(U_b\) fail to glue because their overlap
\[
U_a\cap U_b
\]
is too coarse. Introduce an intermediate region \(U_{ab}^{+}\) containing the missing distinction and maps
\[
U_{ab}^{+}\to U_a,
\qquad
U_{ab}^{+}\to U_b.
\]

At the section level, the new region supplies
\[
s_{ab}^{+}\in\mathscr F^{+}(U_{ab}^{+})
\]
whose restrictions mediate the two local states.

Alternatively, if an equality relation caused forced identification, expansion may replace it with a typed relation:
\[
a=b
\quad\rightsquigarrow\quad
a\xrightarrow{\mathsf{contradicts}}b.
\]
Although the global state now contains more structure, it makes fewer false identifications.

\begin{observation}
Representational expansion can consist of weakening an identification. The expanded structure retains two objects and the relation between them where the earlier structure retained one collapsed object.
\end{observation}

\subsection{Branch expansion}

When no single global section exists but several maximal coherent assemblies do, expansion may replace one forced state with a branch family:
\[
Z
\rightsquigarrow
\mathfrak Z
=
\{Z^{(1)},\ldots,Z^{(k)}\}.
\]

Let branch probabilities be
\[
\pi_j
=
\Prob(Z^{(j)}\mid\mathcal O).
\]
For operation \(f\), the induced output distribution is
\[
\Prob(f=y\mid\mathcal O)
=
\sum_j
\pi_j
\ind[f(Z^{(j)})=y].
\]

If all branches agree, the operation can proceed. If they disagree, uncertainty is structurally warranted.

Branch expansion prevents the system from manufacturing certainty through forced gluing.

\subsection{Repair at the participant layer}

Suppose the representation and assembly are sufficient, but the participant lacks demonstrated competency \(v\). An instructional repair updates
\[
d_i(v,t)
\to
d_i(v,t+1).
\]

The repair is successful when it increases transfer performance without altering protected claim provenance or converting rejection into assent.

Let
\[
J_{\mathrm{learn}}(\mathcal R_v)
=
\E[d_i(v,t+1)-d_i(v,t)]
\]
and let
\[
J_{\mathrm{distort}}(\mathcal R_v)
\]
measure change to protected coordinates not authorized by the participant.

A valid participant repair requires
\[
J_{\mathrm{learn}}(\mathcal R_v)>0
\]
and
\[
J_{\mathrm{distort}}(\mathcal R_v)=0
\]
within tolerance.

Repeated failure of participant repair after demonstrated transfer is evidence that the failure lies elsewhere.

\subsection{Repair at the projection layer}

Projection repair adjusts a learned estimator toward an already adequate target quotient:
\[
\widehat q_t
\to
\widehat q_{t+1},
\]
while preserving the target equivalence relation.

This may reduce estimation error:
\[
d_\Pi(\widehat q_{t+1},q^{*})
<
d_\Pi(\widehat q_t,q^{*}).
\]

Projection expansion changes \(q^{*}\) itself by refining the authorized equivalence relation. Confusing these operations creates the convergence error described earlier.

\begin{definition}[Estimator repair]
Estimator repair changes the approximation to a fixed target.
\end{definition}

\begin{definition}[Target expansion]
Target expansion changes the quotient regarded as operationally adequate.
\end{definition}

No amount of estimator repair corrects target insufficiency.

\subsection{Repair at the gluing layer}

If
\[
[\omega]=0,
\]
a local \(0\)-cochain \(b\) satisfies
\[
\omega=\delta^0b.
\]
Applying \(b\) repairs local alignment without changing the cover or coefficient structure.

If
\[
[\omega]\neq0,
\]
local correction cannot succeed. Expansion must change the cover, coefficients, compatibility maps, or branch structure.

This gives an exact regime criterion in the ideal cohomological model:
\[
[\omega]=0
\implies
\text{local gluing repair is possible},
\]
while
\[
[\omega]\neq0
\implies
\text{local correction within the current complex is insufficient}.
\]

A nonzero class does not specify which expansion is best. It proves only that the current local repair family cannot eliminate the incompatibility.

\subsection{Rerouting through an existing structure}

Let
\[
\mathcal P(s,g)
\]
be the set of paths from current state \(s\) to goal set \(g\) in the existing routing graph. A selected path \(P\) may fail because it is costly, unstable, or mismatched to the participant, even when another admissible path exists.

Rerouting solves
\[
P^{*}
\in
\argmin_{P\in\mathcal P(s,g)}
\left[
C(P)
+
\lambda B(P)
+
\gamma O(P)
\right],
\]
where \(C\) is resource cost, \(B\) participant burden, and \(O\) obstruction or incompatibility risk.

If
\[
\mathcal P(s,g)\neq\varnothing,
\]
expansion is not required solely because one path failed. If every existing path violates a protected constraint, the goal is unreachable within the current architecture:
\[
\mathcal P_{\mathrm{adm}}(s,g)
=
\varnothing.
\]
This is evidence for expansion or justified deferral.

\subsection{Reachability as the regime boundary}

Let
\[
\Reach_{\mathcal A}(\Xi_t)
\]
be the set of states reachable under authorized operations within the current architecture \(\mathcal A\). Let
\[
\mathcal G_{\mathrm{goal}}
\]
be the set of states satisfying the current conversational objective.

If
\[
\Reach_{\mathcal A}(\Xi_t)
\cap
\mathcal G_{\mathrm{goal}}
\neq
\varnothing,
\]
the target is reachable without expansion. Repair or rerouting may suffice.

If
\[
\Reach_{\mathcal A}(\Xi_t)
\cap
\mathcal G_{\mathrm{goal}}
=
\varnothing,
\]
no authorized sequence of current operations reaches the goal. Expansion, goal revision, or explicit impossibility reporting is required.

\begin{theorem}[Reachability criterion for necessary expansion]
Assume the current architecture \(\mathcal A\) and authorized operation set are fixed. If the goal set is disjoint from the reachable set, then no finite sequence of repair or rerouting operations internal to \(\mathcal A\) can reach the goal.
\end{theorem}

\begin{proof}
By definition, \(\Reach_{\mathcal A}(\Xi_t)\) contains every state obtainable through finite authorized operation sequences internal to \(\mathcal A\). If this set has empty intersection with the goal set, no such sequence ends in a goal state.
\end{proof}

The theorem gives the cleanest distinction between repair and expansion. Repair changes position inside the current reachable space. Expansion changes the reachable space.

\subsection{Expansion debt and uncontrolled complexity}

Expansion is not free. Every added distinction increases storage, estimation burden, audit surface, and possible coupling.

Let
\[
C_M(M)
\]
be representational complexity,
\[
C_G(G)
\]
graph complexity,
\[
C_{\A}(\A)
\]
assembly complexity, and
\[
C_{\Hh}(\Hh)
\]
transition complexity.

Total architectural complexity is
\[
C_{\mathrm{arch}}
=
\alpha_M C_M
+
\alpha_G C_G
+
\alpha_{\A}C_{\A}
+
\alpha_{\Hh}C_{\Hh}.
\]

An expansion
\[
\mathcal X:
\mathcal A\to\mathcal A^{+}
\]
incurs
\[
\Delta C_{\mathrm{arch}}
=
C_{\mathrm{arch}}(\mathcal A^{+})
-
C_{\mathrm{arch}}(\mathcal A).
\]

If every unusual submission creates a permanent new coordinate, the representation approaches the identity map and loses its simplifying function.

\begin{principle}[Minimal expansion]
Expansion should introduce the smallest reusable distinction that resolves the protected failure class, not a submission-specific exception that merely memorizes one case.
\end{principle}

Reusability can be evaluated by expected future reduction:
\[
U_{\mathrm{reuse}}(\varphi)
=
\E_{\mathrm{future}}
\left[
\Delta\Lambda(\varphi)
+
\Delta O(\varphi)
\right].
\]

\subsection{Expansion validation}

A candidate expansion should be tested against held-out and future contexts.

Let
\[
\mathcal C_{\mathrm{train}},
\quad
\mathcal C_{\mathrm{audit}},
\quad
\mathcal C_{\mathrm{future}}
\]
be context distributions. The expansion is valid only if it reduces the target failure without introducing unacceptable new collapse, obstruction, or disparity elsewhere.

Define net protected change
\[
\Delta_{\mathrm{prot}}
=
\Lambda_{\mathrm{prot}}^{\mathrm{before}}
-
\Lambda_{\mathrm{prot}}^{\mathrm{after}}.
\]
Define collateral change
\[
\Delta_{\mathrm{collateral}}
=
R_{\mathrm{other}}^{\mathrm{after}}
-
R_{\mathrm{other}}^{\mathrm{before}}.
\]

A safe expansion requires
\[
\Delta_{\mathrm{prot}}>0
\]
and
\[
\Delta_{\mathrm{collateral}}
\leq
\eps_{\mathrm{collateral}}.
\]

The expansion should remain reversible during evaluation. Let
\[
\mathcal A^{+}
\to
\mathcal A
\]
be a rollback map preserving all decisions and evidence generated while the expansion was active.

\subsection{Repair and expansion under finite attention}

Diagnosis itself consumes attention. Let
\[
C_{\mathrm{diag}},
\quad
C_{\mathrm{repair}},
\quad
C_{\mathrm{expand}}
\]
be expected costs. Let
\[
p_R,
\quad
p_X
\]
be posterior probabilities that the failure lies in repair or expansion regimes.

The expected cost of immediate repair is
\[
\mathcal C_R
=
C_{\mathrm{repair}}
+
p_X C_{\mathrm{repeat}},
\]
where \(C_{\mathrm{repeat}}\) is the cost of failed repair followed by later diagnosis.

The expected cost of immediate expansion is
\[
\mathcal C_X
=
C_{\mathrm{expand}}
+
p_R C_{\mathrm{overcomplex}},
\]
where \(C_{\mathrm{overcomplex}}\) is unnecessary complexity introduced when repair would have sufficed.

A diagnostic test \(q\) has value when its expected reduction in regime-selection loss exceeds its burden:
\[
\operatorname{VOI}(q)
>
C_{\mathrm{diag}}(q).
\]

Finite attention therefore does not justify always choosing repair because repair is initially cheaper. Repeated repair in an expansion regime can consume more attention than early structural revision.

\subsection{The regime-separation theorem}

Let
\[
\mathfrak R_{\mathcal A}
\]
be the closure of states reachable through repair and rerouting inside architecture \(\mathcal A\). Let
\[
\mathfrak X_{\mathcal A}
\]
be states reachable only after at least one proper architectural expansion.

\begin{theorem}[Repair-expansion separation]
For a fixed initial state \(\Xi_0\) and goal set \(\mathcal G\), exactly one of the following holds:
\[
\mathcal G
\cap
\mathfrak R_{\mathcal A}
\neq
\varnothing,
\]
or
\[
\mathcal G
\cap
\mathfrak R_{\mathcal A}
=
\varnothing.
\]
In the first case, expansion is not logically necessary for reachability. In the second case, every successful trajectory must either expand the architecture, revise the goal, or leave the authorized operation set.
\end{theorem}

\begin{proof}
The two conditions are logical complements. If the intersection is nonempty, at least one repair or rerouting trajectory reaches the goal, so expansion is not necessary. If the intersection is empty, no internal trajectory reaches the goal. Therefore any successful trajectory must alter the architecture, alter the goal set, or employ an operation not contained in the authorized internal closure.
\end{proof}

Determining which case holds may be computationally undecidable or empirically uncertain in a rich architecture. The theorem supplies the conceptual boundary even when exact classification is unavailable.

\subsection{The expansion obligation}

\begin{principle}[Expansion obligation]
When a protected failure persists after demonstrated competency, calibrated rerouting, estimator repair, and locally valid gluing correction, the Attention Compiler must test whether the required continuation lies outside the reachable space of its current representation and compatibility architecture.
\end{principle}

The test must distinguish numerical instability from structural recurrence, random fluctuation from low-period persistence, participant misunderstanding from definition mismatch, projection collapse from gluing obstruction, and local correction from target misspecification.

Repair is appropriate when the architecture already contains the needed distinction and continuation. Rerouting is appropriate when an alternative admissible path already exists. Expansion is appropriate when the current quotient, graph, cover, coefficient structure, or branch space cannot represent what the authorized continuation requires.

The central diagnostic is not whether the system converges. It is whether convergence preserves the futures the representation was built to support.

\section{Recipient-Relative Escalation and Finite-Attention Allocation}

The preceding construction determines whether a submission has survived projection, prerequisite repair, local assembly, and obstruction testing. Survival through those transformations does not yet imply that the submission should be delivered to a particular recipient. Escalation is a separate operation because the value of a residual is relational. A residual can be coherent, novel, and technically admissible while remaining irrelevant to the knowledge state, responsibilities, or current research frontier of the person whose attention is being requested. Conversely, a residual of modest global novelty can possess exceptional value for a particular recipient because it corrects an active error, completes an unfinished argument, supplies a missing observation, or creates a complementary interaction that neither party could produce independently.

Let \(\mathcal{I}\) denote the population of participants, let \(\mathcal{R}\) denote the population of possible recipients, and let \(t\in\mathbb{R}_{\geq 0}\) denote conversational time. For each participant \(i\in\mathcal{I}\), let \(x_i\) denote the original submission, let
\[
T(x_i,\K,\thetaState_i(t))
=
\bigl(
a_i,
p_i,
r_i,
u_i,
\omega_i
\bigr)
\]
denote the output of the full transformative process, and interpret \(a_i\) as the answered component, \(p_i\) as the prerequisite traversal, \(r_i\) as the surviving residual, \(u_i\) as classification uncertainty, and \(\omega_i\) as the collection of detected local and global gluing obstructions. The inclusion of \(\omega_i\) is necessary because two residuals with identical normalized claims may impose radically different human costs. One may already be globally assembled, while the other may require a recipient to resolve incompatible local interpretations before its intellectual content can even be evaluated.

For recipient \(r\in\mathcal{R}\), let
\[
\K_r(t)
\]
denote the recipient's recorded knowledge state at time \(t\), including established claims, unresolved conjectures, current projects, previous answers, known objections, institutional responsibilities, and explicit exclusions. Let
\[
\A_r([t,t+\Delta])
\]
denote the amount of attention that recipient \(r\) is willing and able to allocate during the interval \([t,t+\Delta]\). Attention is measured in units of effective evaluative time rather than in message counts. This distinction is indispensable because a one-line counterexample can require several hours of verification, while a long submission may collapse under normalization into a familiar argument requiring only seconds to classify.

\subsection{Importance, relevance, salience, and recipient-relative value}

Three quantities that are frequently collapsed into a single ranking score must be distinguished. Let
\[
I(x_i)\in\mathbb{R}_{\geq 0}
\]
denote the estimated general importance of the residual, understood as the magnitude of its possible consequences independently of any particular recipient. Let
\[
L(x_i,r,t)\in\mathbb{R}_{\geq 0}
\]
denote its relevance to recipient \(r\) at time \(t\), understood as the degree to which the residual intersects \(\K_r(t)\), the recipient's unresolved questions, or the recipient's practical responsibilities. Let
\[
B(x_i,r,t)\in\mathbb{R}_{\geq 0}
\]
denote its present salience, understood as the urgency or contextual priority of addressing the residual now rather than later.

None of these quantities determines either of the others. There may exist \(x_i,x_j\) and \(r\) such that
\[
I(x_i)>I(x_j),
\qquad
L(x_i,r,t)<L(x_j,r,t),
\]
and there may exist times \(t_1<t_2\) such that
\[
I(x_i)=I(x_i),
\qquad
L(x_i,r,t_1)=L(x_i,r,t_2),
\qquad
B(x_i,r,t_1)\neq B(x_i,r,t_2).
\]
The first relation says that a generally important problem may be less relevant to a particular recipient than a narrower correction to that recipient's present work. The second says that the intellectual content and recipient relation can remain unchanged while a deadline, public event, experimental opportunity, or newly discovered dependency changes the value of immediate attention.

Let \(V_{ir}(t)\) be the true intellectual value produced if residual \(r_i\) is escalated to recipient \(r\) at time \(t\). This quantity is latent. Before escalation, the system has access only to an information state
\[
\mathscr{I}_{ir}(t)
=
\sigma\!\left(
x_i,
a_i,
p_i,
r_i,
u_i,
\omega_i,
\K_r(t),
\thetaState_i(t),
\mathcal{Y}_{1:t}
\right),
\]
where \(\mathcal{Y}_{1:t}\) is the relevant interaction history. The system therefore acts on the conditional estimate
\[
\widehat{v}_{ir}(t)
=
\E\!\left[
V_{ir}(t)
\given
\mathscr{I}_{ir}(t)
\right].
\]
A useful decomposition is
\[
\widehat{v}_{ir}(t)
=
\alpha_I \widehat{I}_i
+
\alpha_L \widehat{L}_{ir}(t)
+
\alpha_B \widehat{B}_{ir}(t)
+
\alpha_N \widehat{N}_{ir}(t)
+
\alpha_C \widehat{C}_{ir}(t)
-
\alpha_H \widehat{H}_{ir}(t),
\]
where \(\widehat{N}_{ir}\) is recipient-relative residual novelty, \(\widehat{C}_{ir}\) is expected consequentiality, and \(\widehat{H}_{ir}\) is the expected harm or opportunity cost induced by interruption. The coefficients are policy parameters rather than natural constants. Their visibility is a governance requirement because changing them changes whose messages reach the recipient.

\begin{definition}[Recipient-relative escalation score]
For a transformed submission \(T(x_i,\K,\thetaState_i(t))\) and recipient \(r\), the recipient-relative escalation score is a measurable function
\[
S_{ir}(t)
=
S\!\left(
x_i
\mid
r,
\K_r(t),
\thetaState_i(t),
r_i,
u_i,
\omega_i,
t
\right).
\]
An escalation policy is a measurable map
\[
\pi_E:
\mathscr{I}_{ir}(t)
\longrightarrow
[0,1]
\]
whose value is the probability that \(x_i\) is escalated to \(r\).
\end{definition}

The probabilistic form includes deterministic policies as the special case in which \(\pi_E\) takes values only in \(\{0,1\}\). Randomization is not an implementation defect. It supplies exploration, permits estimation of false suppression among non-escalated submissions, and prevents infinitesimal score differences from becoming absolute barriers to access.

\subsection{Joint complementarity as a nonadditive value term}

The value of direct interaction cannot always be reduced to the independent quality of the submission and the independent expertise of the recipient. Let
\[
V_i^{\mathrm{solo}}
\]
denote the value obtainable by developing \(r_i\) without recipient \(r\), and let
\[
V_r^{\mathrm{solo}}(r_i)
\]
denote the value recipient \(r\) could obtain from the packet without reciprocal interaction with participant \(i\). Let
\[
V_{ir}^{\mathrm{joint}}
\]
denote the value obtainable through reciprocal interaction. Define the joint complementarity term by
\[
\Delta_{\mathrm{joint}}(x_i,r,t)
=
V_{ir}^{\mathrm{joint}}(t)
-
V_i^{\mathrm{solo}}(t)
-
V_r^{\mathrm{solo}}(r_i,t).
\]
If
\[
\Delta_{\mathrm{joint}}(x_i,r,t)>0,
\]
then direct interaction is superadditive. The conversation produces something that cannot be represented as the sum of what either side could obtain without the other. If
\[
\Delta_{\mathrm{joint}}(x_i,r,t)=0,
\]
then direct reciprocal attention adds no value beyond asynchronous transfer. If
\[
\Delta_{\mathrm{joint}}(x_i,r,t)<0,
\]
then interaction introduces friction, misunderstanding, strategic distortion, or distraction sufficient to make direct contact worse than independent development.

The true complementarity term is also latent. Its posterior estimate is
\[
\widehat{\Delta}_{ir}(t)
=
\E\!\left[
\Delta_{\mathrm{joint}}(x_i,r,t)
\given
\mathscr{I}_{ir}(t)
\right].
\]
The variance
\[
\sigma_{\Delta,ir}^{2}(t)
=
\Var\!\left(
\Delta_{\mathrm{joint}}(x_i,r,t)
\given
\mathscr{I}_{ir}(t)
\right)
\]
is independently important. A residual with moderate expected complementarity and large posterior variance may deserve exploratory review because the upper tail contains outcomes unavailable to a certainty-seeking policy.

\begin{definition}[Interaction necessity]
Direct interaction between participant \(i\) and recipient \(r\) is necessary relative to tolerance \(\eps>0\) when every noninteractive continuation operator \(G\) satisfies
\[
\E\!\left[
V_{ir}^{\mathrm{joint}}
-
V\!\left(G(r_i,\K_r)\right)
\given
\mathscr{I}_{ir}
\right]
>
\eps.
\]
\end{definition}

This condition distinguishes escalation from delivery. A residual may deserve to enter the recipient's archive without requiring a conversation. Escalation in the strong sense is reserved for cases in which reciprocal attention is expected to contribute value that one-way transmission cannot recover.

\begin{proposition}[Escalation under positive net complementarity]
Let \(c_{ir}>0\) be the expected attention cost of direct interaction. Suppose that the utility of escalation is risk neutral and additive in independent intellectual gain and attention cost. Then escalation has positive expected net utility if and only if
\[
\widehat{v}_{ir}
+
\omega\widehat{\Delta}_{ir}
>
c_{ir},
\]
where \(\omega>0\) converts joint intellectual gain into the same utility units as \(\widehat{v}_{ir}\) and \(c_{ir}\).
\end{proposition}

\begin{proof}
Under the assumptions, the expected utility of escalation is
\[
U_E
=
\widehat{v}_{ir}
+
\omega\widehat{\Delta}_{ir}
-
c_{ir}.
\]
The expected utility of non-escalation is normalized to zero after all archive-mediated and noninteractive value has been included in the definitions of \(\widehat{v}_{ir}\) and \(\widehat{\Delta}_{ir}\). Hence escalation is preferred precisely when \(U_E>0\), which is equivalent to the stated inequality.
\end{proof}

The proposition is elementary, but its interpretation is not. The attention compiler does not ask whether an idea is good in isolation. It asks whether assigning this recipient's next unit of attention to this transformed residual is expected to produce more value than the best available noninteractive continuation and the competing uses of that attention.

\subsection{Attention as workload rather than message count}

Let submissions directed toward recipient \(r\) arrive according to a stationary counting process \(N_r(t)\) satisfying
\[
\lim_{t\to\infty}\frac{N_r(t)}{t}
=
\lambda_r
\]
almost surely. For the \(i\)-th arrival, let \(E_i\in\{0,1\}\) indicate escalation and let \(C_{ir}\geq 0\) denote the attention required to evaluate the escalated packet. The cumulative admitted workload by time \(t\) is
\[
W_r^{\mathrm{in}}(t)
=
\sum_{i=1}^{N_r(t)}
E_i C_{ir}.
\]
Let the recipient's cumulative available attention be
\[
W_r^{\mathrm{out}}(t)
=
\int_0^t a_r(s)\,ds,
\]
where \(a_r(s)\geq 0\) is the instantaneous service capacity. The unresolved attention backlog is the reflected workload process
\[
Q_r(t)
=
\sup_{0\leq s\leq t}
\left[
W_r^{\mathrm{in}}(t)-W_r^{\mathrm{in}}(s)
-
\bigl(
W_r^{\mathrm{out}}(t)-W_r^{\mathrm{out}}(s)
\bigr)
\right].
\]
For constant average attention supply
\[
A_r
=
\lim_{t\to\infty}\frac{W_r^{\mathrm{out}}(t)}{t},
\]
the elementary load parameter is
\[
\rho_r
=
\frac{
\lambda_r
\Pr(E_i=1)
\E[C_{ir}\mid E_i=1]
}{
A_r
}.
\]

\begin{theorem}[Necessary workload stability condition]
Assume that the arrival process and marked attention costs are stationary and ergodic, that
\[
0<
\E[C_{ir}\mid E_i=1]
<
\infty,
\]
and that the long-run attention supply equals \(A_r\in(0,\infty)\). If \(Q_r(t)\) is stable in the sense that
\[
\limsup_{t\to\infty}\frac{Q_r(t)}{t}=0
\]
almost surely, then
\[
\lambda_r
\Pr(E_i=1)
\E[C_{ir}\mid E_i=1]
\leq
A_r.
\]
Strict inequality is required for positive recurrent stability under ordinary nondegenerate queueing assumptions.
\end{theorem}

\begin{proof}
By stationarity and ergodicity,
\[
\lim_{t\to\infty}
\frac{W_r^{\mathrm{in}}(t)}{t}
=
\lambda_r
\E[E_i C_{ir}]
\]
almost surely. Since
\[
\E[E_iC_{ir}]
=
\Pr(E_i=1)\E[C_{ir}\mid E_i=1],
\]
we obtain
\[
\lim_{t\to\infty}
\frac{W_r^{\mathrm{in}}(t)}{t}
=
\lambda_r
\Pr(E_i=1)
\E[C_{ir}\mid E_i=1].
\]
Also,
\[
\lim_{t\to\infty}
\frac{W_r^{\mathrm{out}}(t)}{t}
=
A_r.
\]
If the admitted workload rate exceeded \(A_r\), then
\[
W_r^{\mathrm{in}}(t)-W_r^{\mathrm{out}}(t)
\]
would grow linearly with positive asymptotic slope. Since the reflected workload \(Q_r(t)\) is bounded below by this difference up to the initial condition, \(Q_r(t)/t\) would have a positive lower limit, contradicting stability. Therefore the admitted workload rate cannot exceed \(A_r\). Equality leaves no slack for stochastic fluctuations. Under nondegenerate arrival or service variability, such fluctuations produce an unbounded expected queue, so strict inequality is required for positive recurrence.
\end{proof}

The resulting stability condition is
\[
\lambda_r
\Pr(E)
\E[C_{ir}\mid E]
<
A_r.
\]
This is the recipient-relative version of the central congestion inequality. It reveals three distinct effects of transformation. Archive resolution reduces \(\Pr(E)\). Prerequisite repair and residual extraction reduce \(\E[C_{ir}\mid E]\). Better packet construction can also increase \(A_r\) in effective terms by allowing the recipient to evaluate more intellectual content per unit of time.

Let \(C_{ir}^{\mathrm{raw}}\) be the attention cost of evaluating the untransformed submission and let \(C_{ir}^{T}\) be the cost of evaluating its transformed escalation packet. Define the compression gain
\[
G_{ir}^{T}
=
C_{ir}^{\mathrm{raw}}-C_{ir}^{T}.
\]
The transformation is attention-saving when
\[
\E[G_{ir}^{T}\mid E]>0.
\]
It is representationally harmful when compression reduces the visible packet while increasing the recipient's reconstruction burden, giving
\[
G_{ir}^{T}<0.
\]
This occurs when projection removes context needed to interpret the residual, when provenance is omitted, or when unresolved gluing obstructions are concealed rather than represented.

\begin{lemma}[Obstruction-sensitive cost decomposition]
Suppose the transformed evaluation cost admits the decomposition
\[
C_{ir}^{T}
=
C_{ir}^{\mathrm{claim}}
+
C_{ir}^{\mathrm{prov}}
+
C_{ir}^{\mathrm{repair}}
+
C_{ir}^{\mathrm{glue}}
+
C_{ir}^{\mathrm{interaction}}.
\]
If \(\omega_i=0\), then \(C_{ir}^{\mathrm{glue}}=0\). If \(\omega_i\neq 0\) and the obstruction has not been localized in the escalation packet, then
\[
C_{ir}^{\mathrm{glue}}
\geq
C_{ir}^{\mathrm{localize}},
\]
where \(C_{ir}^{\mathrm{localize}}\) is the minimum cost of discovering which compatibility condition fails.
\end{lemma}

\begin{proof}
When \(\omega_i=0\), the local sections admit a compatible global assembly, so no attention is required to resolve a gluing failure. Hence \(C_{ir}^{\mathrm{glue}}=0\). When \(\omega_i\neq0\), any correct evaluation must distinguish a genuine substantive disagreement from a failure of compatibility among local representations. If the packet does not identify the obstruction, the recipient must at least perform the localization operation needed to determine where compatibility fails. Consequently the gluing cost cannot be smaller than the minimum cost of that operation.
\end{proof}

It follows that obstruction detection before escalation is valuable even when the system cannot repair the obstruction. A localized failure is a compressed explanation of why global assembly failed. An unlocalized failure is merely additional work.

\subsection{Selective classification and abstention}

Escalation can be formulated as a selective prediction problem. Let
\[
Y_{ir}\in\{0,1\}
\]
indicate whether escalation to recipient \(r\) would be valuable under a specified ex post criterion. Let
\[
h_{ir}\in\{0,1\}
\]
be the system's prediction and let
\[
g_{ir}\in\{0,1\}
\]
be a selection function, where \(g_{ir}=1\) means that the system commits to an automated decision and \(g_{ir}=0\) means that it abstains and routes the case to audit or human triage. The selective risk is
\[
\mathcal{R}_{\mathrm{sel}}(h,g)
=
\frac{
\E\!\left[
\ell(h_{ir},Y_{ir})g_{ir}
\right]
}{
\E[g_{ir}]
},
\]
provided \(\E[g_{ir}]>0\). The coverage is
\[
\phi(g)=\E[g_{ir}].
\]
Reducing coverage can reduce selective risk because the system may abstain on uncertain cases. However, abstention does not eliminate attention cost. It transfers the cost from the recipient-facing queue to an audit or secondary-review queue.

Let
\[
c_{ir}^{+}
\]
denote the cost of a false escalation and let
\[
c_{ir}^{-}
\]
denote the epistemic cost of false suppression. If
\[
p_{ir}
=
\Pr(Y_{ir}=1\mid\mathscr{I}_{ir}),
\]
then the expected loss of escalation is
\[
L_E
=
c_{ir}^{+}(1-p_{ir}),
\]
while the expected loss of non-escalation is
\[
L_D
=
c_{ir}^{-}p_{ir}.
\]
Escalation minimizes conditional expected loss exactly when
\[
c_{ir}^{+}(1-p_{ir})
<
c_{ir}^{-}p_{ir},
\]
which is equivalent to
\[
p_{ir}
>
\frac{c_{ir}^{+}}{c_{ir}^{+}+c_{ir}^{-}}.
\]

\begin{proposition}[Asymmetric escalation threshold]
Under binary escalation loss, the Bayes-optimal escalation threshold is
\[
\tau_{ir}^{*}
=
\frac{c_{ir}^{+}}{c_{ir}^{+}+c_{ir}^{-}}.
\]
Increasing the estimated cost of false suppression decreases the threshold, while increasing the cost of unnecessary recipient attention increases it.
\end{proposition}

\begin{proof}
The threshold follows from the preceding comparison. Differentiation gives
\[
\frac{\partial\tau_{ir}^{*}}{\partial c_{ir}^{-}}
=
-\frac{c_{ir}^{+}}{(c_{ir}^{+}+c_{ir}^{-})^2}
<0
\]
and
\[
\frac{\partial\tau_{ir}^{*}}{\partial c_{ir}^{+}}
=
\frac{c_{ir}^{-}}{(c_{ir}^{+}+c_{ir}^{-})^2}
>0.
\]
\end{proof}

This makes the governance problem explicit. A system that treats recipient inconvenience as extremely costly and false suppression as negligible will set a high escalation threshold. The resulting selectivity is not a neutral consequence of finite attention. It is the mathematical expression of a policy choice about which error matters more.

With abstention, let \(c_{ir}^{A}\) be the cost of secondary review. Abstention is optimal when
\[
c_{ir}^{A}
<
\min\left\{
c_{ir}^{+}(1-p_{ir}),
c_{ir}^{-}p_{ir}
\right\}.
\]
Thus abstention occupies an interval around the decision boundary whenever secondary review is cheaper than either kind of uncertain commitment.

\subsection{Single-interval allocation as a constrained optimization}

Fix recipient \(r\) and interval \([t,t+\Delta]\). Suppose \(n\) transformed residuals are eligible for escalation. Let
\[
z_i\in\{0,1\}
\]
indicate whether residual \(i\) is selected. Let \(c_{ir}>0\) be its expected attention cost. A baseline allocation problem is
\[
\begin{aligned}
\text{maximize}\quad
&
\sum_{i=1}^{n}
z_i
\left(
\widehat{v}_{ir}
+
\omega\widehat{\Delta}_{ir}
\right)
\\
\text{subject to}\quad
&
\sum_{i=1}^{n}
z_i c_{ir}
\leq
\A_r([t,t+\Delta]).
\end{aligned}
\]
This is a recipient-relative knapsack problem. The objective becomes more faithful when uncertainty, exploration, neglected-domain correction, urgency, and obstruction value are included:
\[
s_{ir}
=
\widehat{v}_{ir}
+
\omega\widehat{\Delta}_{ir}
+
\gamma u_i
+
\delta e_{ir}^{x}
+
\eta n_{ir}
+
\zeta B_{ir}
+
\chi o_{ir}
-
\psi h_{ir},
\]
where \(e_{ir}^{x}\) rewards distance from familiar regions of the archive, \(n_{ir}\) compensates for systematic neglect, \(o_{ir}\) rewards potentially informative obstruction cases, and \(h_{ir}\) estimates interaction risk. The optimization becomes
\[
\begin{aligned}
\text{maximize}\quad
&
\sum_{i=1}^{n}z_i s_{ir}
\\
\text{subject to}\quad
&
\sum_{i=1}^{n}z_i c_{ir}\leq A_r,
\qquad
z_i\in\{0,1\}.
\end{aligned}
\]

A ratio policy that sorts by \(s_{ir}/c_{ir}\) is optimal for the fractional relaxation
\[
z_i\in[0,1],
\]
but not generally for the indivisible problem.

\begin{theorem}[Optimality of the ratio rule for divisible attention]
Assume residual evaluations are perfectly divisible, so that assigning \(z_i\in[0,1]\) yields value \(z_i s_{ir}\) at cost \(z_i c_{ir}\). If all \(s_{ir}\geq0\) and all \(c_{ir}>0\), then an optimal solution is obtained by ordering residuals so that
\[
\frac{s_{1r}}{c_{1r}}
\geq
\frac{s_{2r}}{c_{2r}}
\geq
\cdots
\geq
\frac{s_{nr}}{c_{nr}}
\]
and allocating attention in that order until the budget is exhausted, with at most one residual receiving a fractional allocation.
\end{theorem}

\begin{proof}
Suppose a feasible allocation assigns positive attention to residual \(j\) while residual \(i\) is not fully allocated and
\[
\frac{s_{ir}}{c_{ir}}
>
\frac{s_{jr}}{c_{jr}}.
\]
Choose a sufficiently small attention quantity \(\delta>0\). Remove cost \(\delta\) from residual \(j\) and assign the same cost to residual \(i\). This changes \(z_j\) by \(-\delta/c_{jr}\) and \(z_i\) by \(+\delta/c_{ir}\). The resulting change in value is
\[
\delta
\left(
\frac{s_{ir}}{c_{ir}}
-
\frac{s_{jr}}{c_{jr}}
\right)
>0,
\]
contradicting optimality. Hence every higher-ratio residual must be saturated before a lower-ratio residual receives attention. Only the boundary residual can be fractional.
\end{proof}

Human evaluation is rarely perfectly divisible. Reading half a proof may produce less than half its value, and interrupted evaluation may impose restart costs. The binary formulation is therefore generally more realistic. Nevertheless, the ratio rule supplies a useful shadow price interpretation. If \(\nu_r\geq0\) is the Lagrange multiplier associated with the relaxed attention constraint, then the Lagrangian is
\[
\mathcal{L}(z,\nu_r)
=
\sum_i z_i s_{ir}
-
\nu_r
\left(
\sum_i z_i c_{ir}-A_r
\right),
\]
and residual \(i\) receives attention only when
\[
s_{ir}-\nu_r c_{ir}>0.
\]
The multiplier \(\nu_r\) is the marginal intellectual value of one additional unit of recipient attention. When congestion rises, \(\nu_r\) rises, and the effective escalation threshold becomes more demanding even if the intrinsic score of every residual remains unchanged.

\subsection{Nonadditive portfolios and submodular attention}

Residuals can overlap. Two submissions may contain the same counterexample, repair the same missing prerequisite edge, or jointly support a conclusion that neither supports alone. Additive selection therefore fails to represent redundancy and synergy.

Let
\[
\mathcal{S}\subseteq\{1,\dots,n\}
\]
be a selected set and let
\[
F_r(\mathcal{S})
\]
denote its total recipient-relative value. The marginal value of adding residual \(i\notin\mathcal{S}\) is
\[
\Delta_F(i\mid\mathcal{S})
=
F_r(\mathcal{S}\cup\{i\})-F_r(\mathcal{S}).
\]
The function \(F_r\) is submodular when
\[
\Delta_F(i\mid\mathcal{S})
\geq
\Delta_F(i\mid\mathcal{T})
\]
for every
\[
\mathcal{S}\subseteq\mathcal{T}
\]
and \(i\notin\mathcal{T}\). This is the diminishing-returns condition. It is appropriate when earlier selections already cover some of the claims, evidence, or conceptual territory supplied by later submissions.

For equal evaluation costs and cardinality budget \(k\), define the greedy sequence
\[
\mathcal{S}_0=\varnothing,
\qquad
i_j
\in
\argmax_{i\notin\mathcal{S}_{j-1}}
\Delta_F(i\mid\mathcal{S}_{j-1}),
\qquad
\mathcal{S}_j
=
\mathcal{S}_{j-1}\cup\{i_j\}.
\]

\begin{theorem}[Greedy guarantee under diminishing returns]
Suppose \(F_r\) is nonnegative, monotone, and submodular with \(F_r(\varnothing)=0\). Let \(\mathcal{S}^{*}\) maximize \(F_r(\mathcal{S})\) over all sets satisfying \(|\mathcal{S}|\leq k\). Then the greedy set \(\mathcal{S}_k\) satisfies
\[
F_r(\mathcal{S}_k)
\geq
\left(
1-\left(1-\frac{1}{k}\right)^k
\right)
F_r(\mathcal{S}^{*})
\geq
\left(1-\frac{1}{e}\right)
F_r(\mathcal{S}^{*}).
\]
\end{theorem}

\begin{proof}
Fix \(j\in\{0,\dots,k-1\}\). By monotonicity,
\[
F_r(\mathcal{S}^{*})
\leq
F_r(\mathcal{S}_j\cup\mathcal{S}^{*}).
\]
By submodularity,
\[
F_r(\mathcal{S}_j\cup\mathcal{S}^{*})
-
F_r(\mathcal{S}_j)
\leq
\sum_{i\in\mathcal{S}^{*}\setminus\mathcal{S}_j}
\Delta_F(i\mid\mathcal{S}_j).
\]
Since \(|\mathcal{S}^{*}|\leq k\), at least one admissible item has marginal gain at least
\[
\frac{
F_r(\mathcal{S}^{*})-F_r(\mathcal{S}_j)
}{k}.
\]
The greedy choice has marginal gain no smaller than this value. Hence
\[
F_r(\mathcal{S}_{j+1})-F_r(\mathcal{S}_j)
\geq
\frac{
F_r(\mathcal{S}^{*})-F_r(\mathcal{S}_j)
}{k}.
\]
Let
\[
D_j
=
F_r(\mathcal{S}^{*})-F_r(\mathcal{S}_j).
\]
Then
\[
D_{j+1}
\leq
\left(1-\frac{1}{k}\right)D_j.
\]
Iteration gives
\[
D_k
\leq
\left(1-\frac{1}{k}\right)^k
F_r(\mathcal{S}^{*}).
\]
Therefore
\[
F_r(\mathcal{S}_k)
\geq
\left(
1-\left(1-\frac{1}{k}\right)^k
\right)
F_r(\mathcal{S}^{*}).
\]
Finally,
\[
\left(1-\frac{1}{k}\right)^k
\leq e^{-1},
\]
which yields the second inequality.
\end{proof}

Positive synergy violates submodularity. Two individually weak residuals can become decisive together. If
\[
F_r(\{i,j\})
>
F_r(\{i\})
+
F_r(\{j\}),
\]
then treating submissions independently can suppress both members of a valuable pair. The routing system must therefore search not only for high-value residuals but also for complementary bundles. Define pairwise synergy by
\[
\Gamma_{ijr}
=
F_r(\{i,j\})-F_r(\{i\})-F_r(\{j\}).
\]
A quadratic approximation to portfolio value is
\[
F_r(\mathcal{S})
=
\sum_{i\in\mathcal{S}}s_{ir}
+
\sum_{\substack{i,j\in\mathcal{S}\\i<j}}
\Gamma_{ijr}.
\]
The resulting allocation problem is
\[
\begin{aligned}
\text{maximize}\quad
&
\sum_i z_i s_{ir}
+
\sum_{i<j}z_i z_j\Gamma_{ijr}
\\
\text{subject to}\quad
&
\sum_i z_i c_{ir}\leq A_r.
\end{aligned}
\]
This expression formalizes a central reason for collective recurrence analysis. A collection of individually non-escalated residuals may jointly define a pattern whose value exceeds the sum of its parts.

\subsection{The exploration reserve}

A policy that maximizes only \(\widehat{v}_{ir}\) exploits its present model of value. Because that model was trained on previously recognized contributions, pure exploitation creates a self-confirming attention boundary. Residuals unlike earlier successes receive low predicted value, are not reviewed, and therefore never generate the evidence needed to correct the model.

Partition the recipient's attention budget as
\[
A_r=A_{r,p}+A_{r,x},
\qquad
A_{r,x}=\xi_r A_r,
\qquad
\xi_r\in(0,1),
\]
where \(A_{r,p}\) supports predicted high-value selections and \(A_{r,x}\) supports exploratory review. The exploration policy may be random, uncertainty-weighted, diversity-seeking, obstruction-seeking, or stratified across participant classes and knowledge-graph regions.

Let \(\mathcal{D}_t\) be the eligible residual population during interval \(t\). Assign each residual a strictly positive exploratory inclusion probability
\[
\pi_{ir}^{x}(t)>0.
\]
If \(Y_{ir}\) is an ex post value label obtained after review, then the Horvitz-Thompson estimator
\[
\widehat{T}_{Y}
=
\sum_{i\in\mathcal{S}_x}
\frac{Y_{ir}}{\pi_{ir}^{x}}
\]
is unbiased for the population total
\[
T_Y=\sum_{i\in\mathcal{D}_t}Y_{ir},
\]
provided the inclusion probabilities are known.

\begin{proposition}[Unbiasedness of inverse-probability exploratory estimation]
Let \(Z_i\in\{0,1\}\) indicate exploratory inclusion with
\[
\E[Z_i]=\pi_i>0.
\]
Then
\[
\widehat{T}
=
\sum_{i=1}^{n}\frac{Z_iY_i}{\pi_i}
\]
satisfies
\[
\E[\widehat{T}]=\sum_{i=1}^{n}Y_i.
\]
\end{proposition}

\begin{proof}
By linearity of expectation,
\[
\E[\widehat{T}]
=
\sum_{i=1}^{n}
\frac{Y_i}{\pi_i}\E[Z_i]
=
\sum_{i=1}^{n}
\frac{Y_i}{\pi_i}\pi_i
=
\sum_{i=1}^{n}Y_i.
\]
No independence assumption among the \(Z_i\) is required for unbiasedness.
\end{proof}

Strictly positive inclusion probabilities are an epistemic access condition. If a domain has
\[
\pi_{ir}^{x}=0
\]
for every residual assigned to it, then the system cannot estimate the value it suppresses there. The absence of observed successes from that domain becomes uninterpretable because it is observationally equivalent to a policy that never looked.

A diversity-sensitive exploratory objective can be written
\[
\max_{\mathcal{S}_x}
\left[
\sum_{i\in\mathcal{S}_x}u_i
+
\alpha
\sum_{i\in\mathcal{S}_x}
\dist(r_i,\mathcal{K}_{r}^{\mathrm{familiar}})
+
\beta
\sum_{\substack{i,j\in\mathcal{S}_x\\i<j}}
\dist(r_i,r_j)
\right]
\]
subject to
\[
\sum_{i\in\mathcal{S}_x}c_{ir}
\leq
A_{r,x}.
\]
The first term favors uncertain classifications. The second favors distance from familiar archive regions. The third prevents the exploration set from concentrating all of its attention in one unfamiliar region.

\subsection{Constrained access and neglect correction}

Let participants be partitioned into policy-relevant classes
\[
\mathcal{I}
=
\bigcup_{g\in\mathcal{G}}\mathcal{I}_g.
\]
The classes may represent submission domains, institutional positions, linguistic registers, geographic regions, or other categories relevant to auditing systematic neglect. Let
\[
\alpha_g
\]
denote a minimum review share for class \(g\). A constrained allocation may require
\[
\sum_{i\in\mathcal{I}_g}z_i c_{ir}
\geq
\alpha_g A_r
\]
for selected groups, subject to
\[
\sum_{g\in\mathcal{G}}\alpha_g\leq1.
\]
Such constraints do not assert equal intrinsic value across all groups. They prevent the estimator from converting historical absence of review into evidence of low prospective value.

An alternative is to introduce a neglect multiplier. Let
\[
N_{gr}(t)
=
\frac{
\text{eligible workload from group }g
}{
\text{reviewed workload from group }g+\eps
}
\]
and define
\[
n_{ir}(t)
=
\varphi\!\left(N_{g(i)r}(t)\right)
\]
for a monotone function \(\varphi\). The score
\[
s_{ir}
=
s_{ir}^{0}
+
\eta n_{ir}
\]
then increases as a class remains systematically under-reviewed.

\begin{proposition}[Finite neglect under persistent arrivals]
Suppose group \(g\) produces eligible residuals infinitely often, the neglect bonus \(n_{ir}(t)\) diverges as the time since the last review from \(g\) diverges, and all competing base scores remain uniformly bounded above. If at least one residual from \(g\) eventually has cost below the available interval budget, then a score-maximizing policy must select a residual from \(g\) in finite time.
\end{proposition}

\begin{proof}
Let \(M<\infty\) be a uniform upper bound on all competing base scores and bounded bonuses. By assumption, the neglect-adjusted score of an eligible residual from \(g\) eventually exceeds \(M\). Once this occurs, that residual dominates every feasible competitor in score. Since its cost is eventually feasible, a score-maximizing policy must select it. Therefore indefinite neglect is impossible under the stated assumptions.
\end{proof}

The proposition also identifies its own limitation. If residuals from a class are systematically assigned inflated costs, feasibility can fail before ranking begins. Fair access therefore requires auditing both value estimates and cost estimates. A policy can appear value-neutral while suppressing a participant class through exaggerated predictions of the attention required to understand it.

\subsection{Recipient pacing and the consumption-side queue}

Escalation controls admission to the recipient-facing queue. It does not determine the rate at which the recipient consumes already admitted packets. Let
\[
\lambda_{e,r}
=
\lambda_r\Pr(E)
\]
be the escalation arrival rate, measured in packets per unit time, and let
\[
\mu_r
\]
be the recipient's chosen packet-viewing rate. If packet costs are normalized to a common display unit, the consumption queue is stable only if
\[
\lambda_{e,r}<\mu_r.
\]
With heterogeneous costs, the correct condition is again workload-based:
\[
\lambda_{e,r}
\E[C_{ir}\mid E]
<
\mu_r
\E[S_r],
\]
where \(S_r\) is the amount of evaluative work completed per viewing opportunity.

A pacing operator is a map
\[
\mathcal{P}_{r,t}:
\mathcal{Q}_r(t)
\longrightarrow
\mathcal{V}_r(t),
\]
where \(\mathcal{Q}_r(t)\) is the current escalation queue and \(\mathcal{V}_r(t)\) is the subset displayed during the next viewing interval. The cardinality or total expected cost of \(\mathcal{V}_r(t)\) is bounded by the recipient's chosen pace.

Let the waiting time of item \(i\) at time \(t\) be
\[
w_i(t)=t-t_i^{E},
\]
where \(t_i^{E}\) is its escalation time. A priority index may take the form
\[
J_{ir}(t)
=
\frac{
s_{ir}
+
\alpha w_i(t)
+
\beta B_{ir}(t)
}{
c_{ir}
}.
\]
The waiting-time term prevents starvation. If \(\alpha>0\), then the priority of a continuously waiting feasible residual eventually exceeds any uniformly bounded newcomer score.

\begin{theorem}[No starvation under linear aging]
Assume every packet has finite cost bounded by
\[
0<c_{\min}\leq c_{ir}\leq c_{\max}<\infty,
\]
all non-aging score components are uniformly bounded, and the pacer selects a feasible packet maximizing \(J_{ir}(t)\) with \(\alpha>0\). If the queue is workload-stable and packet \(i\) remains eligible, then packet \(i\) is selected in finite time.
\end{theorem}

\begin{proof}
Suppose packet \(i\) is never selected. Its waiting time tends to infinity, so
\[
J_{ir}(t)
\geq
\frac{\alpha w_i(t)-M}{c_{\max}}
\longrightarrow\infty
\]
for some finite bound \(M\) on the magnitude of all non-aging terms. A newly arriving or previously waiting competitor \(j\) can have greater priority only if its own waiting term or bounded base score compensates for this divergence. Since the queue is stable, it cannot contain an infinite amount of older unresolved workload ahead of \(i\). Consequently there is a finite time after which \(i\) has maximal feasible priority. The pacer must then select it, contradicting the assumption.
\end{proof}

A purely value-ranked pacer and a purely value-ranked escalation policy share the same structural defect. Both can indefinitely defer residuals that do not resemble previously rewarded material. Exploration must therefore occur on both sides of admission. Some packets should be escalated despite model uncertainty, and some already-escalated packets should be shown despite having lower predicted immediate value than their neighbors.

\subsection{Representation-sensitive scheduling}

The state of a residual is not fixed while it waits. New archive material may answer it, new participant responses may repair it, and new submissions may create a recurrent pattern that changes its value. Let
\[
r_i(t)
\]
denote the residual state and let
\[
\mathcal{U}_{t\to t+\Delta}
\]
be the update operator induced by archive growth, participant calibration, and collective recurrence:
\[
r_i(t+\Delta)
=
\mathcal{U}_{t\to t+\Delta}
\bigl(
r_i(t),
\K(t+\Delta),
\thetaState_i(t+\Delta)
\bigr).
\]
The score is consequently dynamic:
\[
s_{ir}(t+\Delta)
\neq
s_{ir}(t)
\]
even in the absence of direct review.

Define the expected decay rate
\[
d_{ir}(t)
=
-
\frac{d}{dt}
\E[V_{ir}(t)\mid\mathscr{I}_{ir}(t)]
\]
when delayed attention destroys opportunity, and define the expected maturation rate
\[
m_{ir}(t)
=
\frac{d}{dt}
\E[V_{ir}(t)\mid\mathscr{I}_{ir}(t)]
\]
when delay permits useful development. A scheduling policy that treats all delay as harmful will escalate immature submissions too early. A policy that treats all delay as harmless will miss time-sensitive opportunities.

For differentiable expected value, the first-order cost of delaying residual \(i\) by \(\Delta t\) is
\[
\mathcal{D}_{ir}(\Delta t)
=
V_{ir}(t)-V_{ir}(t+\Delta t)
=
d_{ir}(t)\Delta t+o(\Delta t)
\]
when \(d_{ir}(t)>0\). A time-sensitive priority index is therefore
\[
J_{ir}^{\mathrm{time}}(t)
=
\frac{
s_{ir}(t)
+
\kappa d_{ir}(t)
-
\ell m_{ir}(t)
}{
c_{ir}(t)
}.
\]
Positive decay increases urgency. Positive maturation can justify continued automated development before escalation.

\subsection{The shadow price of a discarded distinction}

Finite attention creates pressure toward representational simplicity. Every additional distinction preserved in an escalation packet can increase reading cost. Yet discarded distinctions can create false suppression or force expensive reconstruction. Let \(D\) be a candidate distinction that may be retained or removed. Let
\[
\Delta C_r(D)>0
\]
be the additional recipient attention required to preserve and present it. Let
\[
\Delta V_r(D)
\]
be the expected value preserved by retaining it, including the reduction in false suppression and the reduction in downstream reconstruction cost.

If the shadow price of attention is \(\nu_r\), then retaining \(D\) is locally optimal when
\[
\Delta V_r(D)
>
\nu_r\Delta C_r(D).
\]
This inequality is the finite-attention form of representational sufficiency. A distinction is not dispensable merely because the projection operator can remove it. It is dispensable only when the value it preserves is smaller than the recipient-relative opportunity cost of carrying it forward.

Let \(q_0\) be the quotient that discards \(D\) and \(q_1\) the refinement that preserves it. Suppose
\[
q_0=\pi\circ q_1
\]
for a surjective map \(\pi\). Let
\[
\mathcal{L}_r(q)
=
\E\!\left[
\ell_r\bigl(
f_r(X),
g_r(q(X))
\bigr)
\right]
+
\nu_r C_r(q)
\]
be the recipient-relative representational loss. Then \(q_1\) is preferred to \(q_0\) exactly when
\[
\mathcal{L}_r(q_1)<\mathcal{L}_r(q_0),
\]
or equivalently,
\[
\E\!\left[
\ell_r(f_r(X),g_r(q_0(X)))
-
\ell_r(f_r(X),g_r(q_1(X)))
\right]
>
\nu_r
\left(
C_r(q_1)-C_r(q_0)
\right).
\]

\begin{theorem}[Recipient-relative refinement criterion]
Assume \(q_1\) refines \(q_0\), the downstream decoder is optimized separately for each representation, and attention cost is additive. Then preserving the distinction represented by the refinement \(q_1\) is optimal if and only if its reduction in Bayes risk exceeds its shadow-priced attention cost.
\end{theorem}

\begin{proof}
For \(k\in\{0,1\}\), define the optimized downstream risk
\[
R_r(q_k)
=
\inf_{g}
\E\!\left[
\ell_r(f_r(X),g(q_k(X)))
\right].
\]
The total optimized loss is
\[
\mathcal{L}_r^{*}(q_k)
=
R_r(q_k)+\nu_r C_r(q_k).
\]
The refinement is preferred exactly when
\[
R_r(q_1)+\nu_r C_r(q_1)
<
R_r(q_0)+\nu_r C_r(q_0).
\]
Rearranging yields
\[
R_r(q_0)-R_r(q_1)
>
\nu_r\bigl(C_r(q_1)-C_r(q_0)\bigr).
\]
The left side is the reduction in Bayes risk produced by preserving the distinction. The right side is its attention cost evaluated at the recipient's marginal shadow price.
\end{proof}

Because \(q_1\) refines \(q_0\), the optimized prediction risk satisfies
\[
R_r(q_1)\leq R_r(q_0).
\]
More information cannot increase Bayes risk when the decoder is free to ignore it. It can nevertheless increase total loss because human attention is not free. Representational simplicity is therefore not minimization of descriptive size alone. It is minimization of total recipient-relative loss under a finite attention price.

\subsection{Escalation after repair, expansion, and obstruction localization}

Let \(\mathcal{R}_i\) denote the set of available repair operators, let \(\mathcal{X}_i\) denote the set of representation-expansion operators, and let \(\mathcal{G}_i\) denote the set of obstruction-localization operators. For an operator
\[
\Omega\in
\mathcal{R}_i\cup\mathcal{X}_i\cup\mathcal{G}_i,
\]
define its expected escalation gain by
\[
\Delta S_{ir}(\Omega)
=
\E\!\left[
S_{ir}\bigl(\Omega(r_i)\bigr)-S_{ir}(r_i)
\given
\mathscr{I}_{ir}
\right]
\]
and its participant-side cost by \(k_i(\Omega)\). The system should apply \(\Omega\) before escalation when
\[
\Delta S_{ir}(\Omega)
+
\Delta C_{ir}^{T}(\Omega)
>
k_i(\Omega),
\]
where \(\Delta C_{ir}^{T}(\Omega)\) is the expected reduction in recipient evaluation cost.

This produces three distinct pre-escalation regimes. Repair is appropriate when the representation is adequate but the participant has not yet demonstrated the prerequisites needed to use it. Expansion is appropriate when the current quotient has collapsed a task-relevant distinction. Obstruction localization is appropriate when compatible local formulations fail to assemble globally and the failure itself may be intellectually informative.

Let
\[
\Omega_i^{*}
=
\argmax_{\Omega}
\left[
\Delta S_{ir}(\Omega)
+
\Delta C_{ir}^{T}(\Omega)
-
k_i(\Omega)
\right].
\]
Immediate escalation is optimal only when
\[
\sup_{\Omega}
\left[
\Delta S_{ir}(\Omega)
+
\Delta C_{ir}^{T}(\Omega)
-
k_i(\Omega)
\right]
\leq0.
\]
Thus a residual should reach the recipient when additional automated transformation has nonpositive expected net value, not merely when the system has exhausted a fixed sequence of steps.

\begin{definition}[Transformative stopping boundary]
The transformative stopping boundary for participant \(i\) and recipient \(r\) is the set of states \(\sigma\) satisfying
\[
\sup_{\Omega\in\mathfrak{O}(\sigma)}
\E\!\left[
V_{ir}(\Omega(\sigma))
-
V_{ir}(\sigma)
-
k_i(\Omega)
\given
\sigma
\right]
\leq0,
\]
where \(\mathfrak{O}(\sigma)\) is the set of admissible nonhuman continuation operators at state \(\sigma\).
\end{definition}

At this boundary, further automated mediation is not expected to improve the packet enough to justify its cost. The residual has either become ready for recipient attention, become ready for some other recipient, become resolvable without escalation, or become irrecoverably inadmissible under the present task.

\subsection{A unified stochastic control formulation}

The complete allocation problem can be expressed as a partially observed stochastic control process. Let the system state be
\[
\Sigma_t
=
\left(
\K_t,
\{\thetaState_i(t)\}_{i\in\mathcal{I}},
\{r_i(t)\}_{i\in\mathcal{I}},
\{\omega_i(t)\}_{i\in\mathcal{I}},
Q_r(t),
A_r(t)
\right).
\]
The state is only partially observed because participant competencies, true residual values, future complementarities, and false-suppression outcomes are latent. Let the action space include answer retrieval, diagnostic prompting, prerequisite routing, representational expansion, obstruction testing, archival resolution, escalation, exploratory review, deferral, and rejection:
\[
a_t
\in
\mathfrak{A}(\Sigma_t).
\]
Let
\[
R(\Sigma_t,a_t)
\]
be the immediate recipient-relative intellectual reward minus participant burden, attention cost, delay cost, and governance penalties. With discount factor \(\beta\in(0,1)\), the value function is
\[
J^{*}(\Sigma_t)
=
\sup_{\pi}
\E_{\pi}
\left[
\sum_{k=0}^{\infty}
\beta^k
R(\Sigma_{t+k},a_{t+k})
\given
\Sigma_t
\right].
\]
The Bellman equation is
\[
J^{*}(\Sigma)
=
\sup_{a\in\mathfrak{A}(\Sigma)}
\left[
R(\Sigma,a)
+
\beta
\E\!\left[
J^{*}(\Sigma')
\given
\Sigma,a
\right]
\right].
\]
Finite attention appears as a state constraint:
\[
Q_r(t)\leq Q_r^{\max}
\]
or as a long-run average constraint:
\[
\limsup_{T\to\infty}
\frac{1}{T}
\E\!\left[
\sum_{t=1}^{T}
E_tC_{ir,t}
\right]
\leq
A_r.
\]
False suppression, domain neglect, provenance loss, and unreviewable decisions enter as additional constrained costs.

Introducing multipliers
\[
\nu_A,\nu_D,\nu_N,\nu_P\geq0
\]
for attention, false suppression, neglect, and provenance violations yields the constrained Lagrangian reward
\[
\widetilde{R}
=
R
-
\nu_A(C-A_r)
-
\nu_D(D-\eps)
-
\nu_N(N-\delta)
-
\nu_P(P-\pi_0).
\]
The resulting policy does not optimize attention efficiency in isolation. It optimizes intellectual value subject to explicit prices for the system's characteristic failure modes.

\subsection{Conditional stability with bounded suppression}

Let \(V_i\) be the event that submission \(i\) contains a genuinely valuable residual, let \(D_i\) be the event that it is discarded or indefinitely delayed, and let \(E_i\) be the event that it is escalated. Let the estimated false-suppression rate satisfy
\[
\widehat{\Pr}(D_i\mid V_i)\leq\eps.
\]
Let the exploration reserve satisfy
\[
A_{r,x}\geq\xi_r A_r.
\]
Let provenance completeness be represented by
\[
P_i\geq p_0,
\]
and let every automated resolution be contestable through an appeal operator \(\mathfrak{a}\).

\begin{theorem}[Conditional finite-attention viability]
Suppose that the following conditions hold:
\[
\lambda_r\Pr(E_i)\E[C_{ir}\mid E_i]<A_r,
\]
\[
\widehat{\Pr}(D_i\mid V_i)\leq\eps,
\]
\[
A_{r,x}\geq\xi_rA_r>0,
\]
and
\[
P_i\geq p_0
\]
for every escalated packet. Suppose further that cost estimates are integrable, exploratory inclusion probabilities are strictly positive on every eligible stratum, and appeals can return a discarded residual to the transformative process. Then the recipient-facing workload is stable, every eligible stratum remains statistically observable, and the system possesses a corrective path for at least some false suppressions.
\end{theorem}

\begin{proof}
The strict workload inequality and the queueing assumptions imply that the long-run admitted workload is smaller than the recipient's service capacity, yielding stability of the recipient-facing workload process. Since the exploration budget is strictly positive and every eligible stratum has strictly positive inclusion probability, no eligible stratum is completely excluded from observation. Consequently inverse-probability estimators can be constructed for audited outcomes within each stratum. The false-suppression bound constrains the estimated rate at which valuable residuals are discarded or indefinitely delayed. Provenance completeness ensures that escalated packets retain the information required to inspect how their residuals were produced. Finally, the appeal operator supplies a transition from at least some discarded states back into the active transformative state space. Therefore suppression decisions are not universally absorbing. These conclusions are conditional because the true event \(V_i\) remains latent for unaudited residuals and because the estimated bound depends on the adequacy of exploratory sampling.
\end{proof}

The theorem does not establish that the system identifies all valuable ideas. No finite-attention mechanism can supply such a guarantee when value is latent, evaluation is costly, and the arrival process is unbounded. It establishes the narrower and more defensible result that stable attention allocation need not require making every rejection invisible, unauditable, or irreversible.

\subsection{The recipient-relative principle of representational simplicity}

The mathematical structure of this section can be condensed into a single constrained criterion. For a candidate representation \(q\), transformation policy \(T\), escalation policy \(\pi_E\), and pacing policy \(\mathcal{P}_r\), define
\[
\begin{aligned}
\mathfrak{J}_r(q,T,\pi_E,\mathcal{P}_r)
={}&
\E\!\left[
V_{ir}^{\mathrm{joint}}
+
V_{ir}^{\mathrm{archive}}
+
V_{ir}^{\mathrm{learning}}
\right]
\\
&-
\nu_r
\E\!\left[
C_{ir}^{\mathrm{transform}}
+
C_{ir}^{\mathrm{recipient}}
\right]
\\
&-
\alpha_r
\Pr(D_i\mid V_i)
-
\beta_r
\E[\Burden_i]
-
\gamma_r
\E[\mathrm{ProvenanceLoss}_i].
\end{aligned}
\]
The integrated design problem is
\[
\begin{aligned}
\argmax_{q,T,\pi_E,\mathcal{P}_r}
\quad&
\mathfrak{J}_r(q,T,\pi_E,\mathcal{P}_r)
\\
\text{subject to}\quad&
\lambda_r\Pr(E)
\E[C_{ir}\mid E]
<
A_r,
\\
&
A_{r,x}\geq\xi_rA_r,
\\
&
\widehat{\Pr}(D_i\mid V_i)\leq\eps,
\\
&
q\text{ preserves every distinction required by the admitted task class},
\\
&
T\text{ exposes unresolved gluing obstructions},
\\
&
\pi_E\text{ remains appealable and auditable}.
\end{aligned}
\]

Representational simplicity is therefore recipient-relative but not recipient-sovereign. The recipient's limited capacity determines the shadow price of additional distinctions, yet it does not authorize arbitrary projection. A representation is simple in the relevant sense when it removes distinctions whose preservation would cost more attention than the downstream value they protect, while retaining, repairing, or explicitly marking every distinction whose loss would corrupt evaluation, suppress complementarity, or conceal an obstruction.

The central escalation inequality can now be written
\[
\E\!\left[
\Delta_{\mathrm{joint}}(x_i,r)
+
\Delta_{\mathrm{recipient}}(r_i,\K_r)
+
\Delta_{\mathrm{audit}}(u_i,\omega_i)
\given
\mathscr{I}_{ir}(t)
\right]
>
c_{ir}
+
\nu_r^{-1}\mathcal{O}_{ir},
\]
where \(\mathcal{O}_{ir}\) is the opportunity cost imposed on competing residuals. The left side contains the expected value of direct complementarity, recipient-specific contribution, and information gained by resolving uncertainty or obstruction. The right side contains both the immediate attention cost and the value displaced elsewhere in the queue.

This is the point at which CLIO, HYDRA, and the attention compiler become one architecture. CLIO determines which distinctions survive projection into an admissible residual. HYDRA determines whether the surviving local components can participate in a coherent transition and whether their coupling is merely stored, jointly addressed, or operationally constitutive. The attention compiler determines whether the resulting object merits one particular recipient's scarce attention under a workload constraint. Failure at the CLIO level calls for representational expansion. Failure at the HYDRA level calls for coupling analysis or obstruction localization. Failure at the allocation level calls for rerouting, delayed maturation, exploratory sampling, or selection of another recipient. None of these failures is equivalent to intellectual worthlessness.

The governing condition is consequently not
\[
\text{send the highest-scoring message},
\]
but
\[
\text{preserve the distinctions required for evaluation, assemble what can be assembled, expose what cannot be glued, and spend recipient attention only where its marginal joint value exceeds its full opportunity cost}.
\]
Under this condition, finite attention does not disappear. It becomes an explicit constraint on representation, transformation, and access rather than an invisible privilege exercised by an inbox.

\section{Audit, Appeal, and the Estimation of False Suppression}

The allocation theory of the preceding section depends upon quantities that cannot be observed directly at the moment when they matter most. The true recipient-relative value \(V_{ir}\) is latent before escalation. The value of a discarded residual is more severely latent because discarding it prevents the interaction through which its value might have become observable. The event that a residual was correctly suppressed and the event that a valuable residual was never permitted to develop can therefore produce the same administrative record. Both appear as non-escalation. This observational equivalence creates the central epistemic difficulty of auditing an attention compiler.

Let \(V_i\in\{0,1\}\) indicate that submission \(i\) contains a genuinely valuable residual relative to some declared evaluation horizon, recipient class, and task family. Let \(D_i\in\{0,1\}\) indicate that the system discards the residual or delays it beyond the horizon at which its value could be realized. The central safety quantity is
\[
\psi
=
\Prob(D_i=1\mid V_i=1).
\]
This is the false-suppression rate. It is analogous to a false-negative rate, but the analogy is incomplete because the label \(V_i\) is not independently available for all rejected cases. A medical classifier can sometimes be audited against a later diagnostic outcome. An intellectual triage system may alter the outcome by deciding whether the residual receives the development, comparison, and reciprocal attention needed to become legible. Value is therefore partly revealed by intervention and partly produced through intervention.

The complementary quantity
\[
\Prob(E_i=1\mid V_i=0)
\]
measures false escalation, where \(E_i\) denotes delivery to scarce recipient attention. False escalation consumes finite capacity and can destabilize the recipient-facing queue. False suppression is nevertheless the more difficult quantity because an escalated residual creates evidence about itself, while a discarded residual ordinarily does not.

\subsection{Potential value under alternative routing policies}

A rigorous account must distinguish observed outcomes from potential outcomes. Let
\[
A_i\in\mathfrak{A}
\]
denote the routing action applied to submission \(i\), where the action space includes archival resolution, prerequisite repair, representational expansion, obstruction localization, escalation, exploratory review, deferral, and discard. For each action \(a\in\mathfrak{A}\), let
\[
Y_i(a)
\]
denote the potential intellectual outcome that would be observed if action \(a\) were applied. Only
\[
Y_i=Y_i(A_i)
\]
is observed.

For a particular recipient \(r\), define the potential escalation value
\[
Y_{ir}(E)
\]
and the potential non-escalation value
\[
Y_{ir}(D).
\]
The individual treatment effect of escalation is
\[
\tau_{ir}
=
Y_{ir}(E)-Y_{ir}(D).
\]
A residual is escalation-dependent at threshold \(\delta>0\) when
\[
\tau_{ir}>\delta.
\]
This condition is stronger than saying that the residual is valuable. A submission may have substantial archive-mediated value even without escalation. The false-suppression event relevant to scarce recipient attention is therefore
\[
V_{ir}^{E}
=
\ind\!\left\{
Y_{ir}(E)-Y_{ir}(D)>\delta
\right\},
\]
and the corresponding recipient-relative false-suppression rate is
\[
\psi_r^{E}
=
\Prob(D_i=1\mid V_{ir}^{E}=1).
\]

The fundamental problem is that
\[
Y_{ir}(E)
\quad\text{and}\quad
Y_{ir}(D)
\]
cannot both be observed for the same submission. Even if a discarded case is later escalated through an audit, the delayed escalation occurs under a changed archive, a changed recipient state, and a changed participant state. The delayed outcome is not generally identical to the outcome that immediate escalation would have produced.

Let
\[
Y_{ir}(E,t)
\]
denote the potential outcome under escalation at time \(t\). If an audit escalates at \(t+\Delta\), then it observes
\[
Y_{ir}(E,t+\Delta),
\]
not \(Y_{ir}(E,t)\). Define the temporal treatment distortion by
\[
\Lambda_{ir}(t,\Delta)
=
Y_{ir}(E,t)-Y_{ir}(E,t+\Delta).
\]
Positive \(\Lambda_{ir}\) represents value destroyed by delay. Negative \(\Lambda_{ir}\) represents value gained through maturation before delayed review. An audit estimator that ignores \(\Lambda_{ir}\) confounds the quality of the original decision with the consequences of the audit delay.

\begin{definition}[Horizon-relative value]
For evaluation horizon \(H>0\), recipient \(r\), and threshold \(\delta\), define
\[
V_{ir}^{H,\delta}
=
\ind\!\left\{
\sup_{t'\in[t_i,t_i+H]}
Y_{ir}(E,t')
-
Y_{ir}(D,t_i)
>
\delta
\right\}.
\]
The horizon-relative false-suppression rate is
\[
\psi_r(H,\delta)
=
\Prob\!\left(
D_i=1
\given
V_{ir}^{H,\delta}=1
\right).
\]
\end{definition}

There is no context-free false-suppression rate. Its value depends on which recipient is considered, what degree of gain counts as valuable, which intervention family is allowed, and how long a delayed opportunity remains actionable.

\subsection{The nonidentifiability of unaudited suppression}

Let \(X_i\) denote every feature observed before the routing decision, including the transformed residual, provenance, participant-state estimate, obstruction state, estimated evaluation cost, and recipient relation. Let
\[
e(X_i)
=
\Prob(E_i=1\mid X_i)
\]
be the escalation propensity. If the policy is deterministic and there exists a region \(\mathcal{B}\) of feature space such that
\[
e(x)=0
\qquad
\text{for all }x\in\mathcal{B},
\]
then no escalated outcomes are observed from \(\mathcal{B}\).

\begin{theorem}[Nonidentifiability under zero review probability]
Suppose there exists a measurable region \(\mathcal{B}\) with
\[
\Prob(X_i\in\mathcal{B})>0
\]
and
\[
\Prob(E_i=1\mid X_i=x)=0
\]
for every \(x\in\mathcal{B}\). Without additional structural assumptions connecting outcomes inside \(\mathcal{B}\) to outcomes outside \(\mathcal{B}\), the distribution of
\[
Y_i(E)\mid X_i\in\mathcal{B}
\]
is not identifiable from observed data.
\end{theorem}

\begin{proof}
Observed data contain \(X_i\), the action \(E_i\), and the realized outcome
\[
Y_i=E_iY_i(E)+(1-E_i)Y_i(D).
\]
For \(X_i\in\mathcal{B}\), the propensity is zero, so \(E_i=0\) almost surely and
\[
Y_i=Y_i(D).
\]
Consequently the observed-data likelihood over \(\mathcal{B}\) depends on the conditional distribution of \(Y_i(D)\) but does not depend on the conditional distribution of \(Y_i(E)\). Let \(\mathbb{P}_1\) and \(\mathbb{P}_2\) agree on the joint distribution of
\[
(X_i,E_i,Y_i(D))
\]
and on all potential outcomes outside \(\mathcal{B}\), while assigning different conditional distributions to
\[
Y_i(E)\mid X_i\in\mathcal{B}.
\]
The two data-generating processes induce exactly the same distribution over observed data but different escalation values inside \(\mathcal{B}\). Hence the latter distribution is not identifiable.
\end{proof}

The theorem has a direct design consequence. A deterministic rejection region is not merely harsh. It is epistemically sealed. No quantity computed solely from the policy's ordinary operation can reveal whether that region contains valuable residuals. Claims that the region is low value are then supported by the absence of successes from a population that was never allowed to produce observable successes.

Positive exploratory inclusion probabilities restore the possibility of estimation. They do not guarantee accurate estimation because rare inclusion can produce enormous variance, but they prevent complete nonidentifiability.

\subsection{Random audits of non-escalated residuals}

Let
\[
\mathcal{D}_t
=
\{i:D_i=1\}
\]
denote the set of residuals not escalated during audit period \(t\). For each \(i\in\mathcal{D}_t\), let
\[
Z_i\in\{0,1\}
\]
indicate selection for audit and let
\[
\pi_i
=
\Prob(Z_i=1\mid X_i,D_i=1)
\]
be its known inclusion probability. Assume
\[
0<\pi_i\leq1.
\]
After audit, let
\[
\widetilde{V}_i\in\{0,1\}
\]
be the audit determination that the residual was valuable under the declared horizon and criterion.

The total number of valuable suppressed residuals in the period is
\[
T_V
=
\sum_{i\in\mathcal{D}_t}V_i.
\]
The inverse-probability estimator is
\[
\widehat{T}_V
=
\sum_{i\in\mathcal{D}_t}
\frac{Z_i\widetilde{V}_i}{\pi_i}.
\]
If audit labels are correct so that
\[
\widetilde{V}_i=V_i,
\]
then
\[
\E[\widehat{T}_V\mid\mathcal{D}_t]
=
T_V.
\]

To estimate the conditional rate
\[
\Prob(D=1\mid V=1),
\]
one also needs the total number of valuable residuals, including escalated and non-escalated cases. Let
\[
T_V^{E}
=
\sum_{i:E_i=1}V_i
\]
be the number of valuable escalated residuals and define
\[
\widehat{\psi}
=
\frac{
\widehat{T}_V
}{
\widehat{T}_V+T_V^{E}
}.
\]
This ratio is not exactly unbiased, even when \(\widehat{T}_V\) is unbiased, because the expectation of a ratio is not generally the ratio of expectations. It is nevertheless consistent under standard increasing-sample conditions.

\begin{proposition}[Consistency of the audited suppression estimator]
Suppose audit inclusion probabilities satisfy
\[
\inf_i\pi_i\geq\pi_{\min}>0,
\]
audit labels are correct, the relevant empirical totals grow linearly with the audit population size \(n\), and the inverse-probability estimator obeys
\[
\frac{\widehat{T}_V-T_V}{n}
\xrightarrow{p}
0.
\]
If
\[
\frac{T_V+T_V^{E}}{n}
\xrightarrow{p}
\vartheta>0,
\]
then
\[
\widehat{\psi}
-
\frac{T_V}{T_V+T_V^{E}}
\xrightarrow{p}
0.
\]
\end{proposition}

\begin{proof}
Write
\[
\widehat{\psi}
=
g\!\left(
\frac{\widehat{T}_V}{n},
\frac{T_V^{E}}{n}
\right),
\qquad
g(a,b)=\frac{a}{a+b}.
\]
By assumption,
\[
\frac{\widehat{T}_V}{n}
-
\frac{T_V}{n}
\xrightarrow{p}0.
\]
The limiting denominator is bounded away from zero because
\[
\frac{T_V+T_V^{E}}{n}
\xrightarrow{p}\vartheta>0.
\]
The function \(g\) is continuous wherever \(a+b>0\). The continuous mapping theorem therefore gives
\[
g\!\left(
\frac{\widehat{T}_V}{n},
\frac{T_V^{E}}{n}
\right)
-
g\!\left(
\frac{T_V}{n},
\frac{T_V^{E}}{n}
\right)
\xrightarrow{p}0,
\]
which is the desired result.
\end{proof}

When audit inclusion is independent Bernoulli sampling, the conditional variance of the total estimator is
\[
\Var(\widehat{T}_V\mid\mathcal{D}_t)
=
\sum_{i\in\mathcal{D}_t}
\frac{1-\pi_i}{\pi_i}V_i^2.
\]
Since \(V_i\in\{0,1\}\),
\[
\Var(\widehat{T}_V\mid\mathcal{D}_t)
=
\sum_{i:V_i=1}
\frac{1-\pi_i}{\pi_i}.
\]
Small inclusion probabilities make the estimator unstable. A nominal exploration mechanism that samples a class with vanishing probability technically preserves identifiability while providing practically useless precision.

\subsection{Optimal audit allocation under finite audit capacity}

Suppose discarded residuals are partitioned into \(G\) strata. Let \(N_g\) be the number of discarded residuals in stratum \(g\), let
\[
p_g
=
\Prob(V_i=1\mid i\in g,D_i=1)
\]
be the stratum-specific suppressed-value prevalence, let
\[
\sigma_g^2=p_g(1-p_g),
\]
and let \(c_g^{A}\) be the cost of auditing one residual in the stratum. If \(n_g\) cases are audited from stratum \(g\), the approximate variance of the stratified estimator for the suppressed-value total is
\[
\Var(\widehat{T}_V)
\approx
\sum_{g=1}^{G}
N_g^2
\left(
1-\frac{n_g}{N_g}
\right)
\frac{\sigma_g^2}{n_g}.
\]
Ignoring finite-population corrections, the audit-allocation problem is
\[
\begin{aligned}
\text{minimize}\quad&
\sum_{g=1}^{G}
\frac{N_g^2\sigma_g^2}{n_g}
\\
\text{subject to}\quad&
\sum_{g=1}^{G}n_gc_g^{A}
\leq
A_{\Audit},
\qquad
n_g>0.
\end{aligned}
\]

\begin{theorem}[Cost-sensitive Neyman allocation]
Under the continuous relaxation \(n_g>0\), the variance-minimizing audit allocation satisfies
\[
n_g^{*}
=
\frac{
A_{\Audit}
N_g\sigma_g/\sqrt{c_g^{A}}
}{
\sum_{h=1}^{G}
N_h\sigma_h\sqrt{c_h^{A}}
}.
\]
Equivalently,
\[
n_g^{*}
\propto
\frac{N_g\sigma_g}{\sqrt{c_g^{A}}}.
\]
\end{theorem}

\begin{proof}
Introduce a multiplier \(\lambda_A>0\) and form
\[
\mathcal{L}
=
\sum_{g=1}^{G}
\frac{N_g^2\sigma_g^2}{n_g}
+
\lambda_A
\left(
\sum_{g=1}^{G}n_gc_g^{A}-A_{\Audit}
\right).
\]
Differentiation with respect to \(n_g\) gives
\[
-\frac{N_g^2\sigma_g^2}{n_g^2}
+
\lambda_Ac_g^{A}
=
0.
\]
Hence
\[
n_g
=
\frac{N_g\sigma_g}{\sqrt{\lambda_Ac_g^{A}}}.
\]
Applying the binding budget constraint,
\[
\sum_{g=1}^{G}
\frac{N_g\sigma_g\sqrt{c_g^{A}}}{\sqrt{\lambda_A}}
=
A_{\Audit},
\]
so
\[
\sqrt{\lambda_A}
=
\frac{
\sum_{h=1}^{G}
N_h\sigma_h\sqrt{c_h^{A}}
}{
A_{\Audit}
}.
\]
Substitution yields the stated allocation.
\end{proof}

The formula allocates more audits to large strata, uncertain strata, and inexpensive strata. Taken without a minimum-inclusion constraint, it can allocate almost no review to a small, expensive, historically neglected class. Precision optimization must therefore be supplemented by an observability condition
\[
n_g\geq n_g^{\min}>0
\]
or
\[
\pi_i\geq\pi_g^{\min}>0.
\]
Statistical efficiency and epistemic access are related but nonidentical objectives.

\subsection{Auditor error and latent audit truth}

An audit is not an oracle. Let \(V_i\) be the latent value label and let \(\widetilde{V}_i\) be the auditor's observed judgment. Define audit sensitivity and specificity by
\[
\alpha_A
=
\Prob(\widetilde{V}_i=1\mid V_i=1),
\]
\[
\beta_A
=
\Prob(\widetilde{V}_i=0\mid V_i=0).
\]
If the true suppressed-value prevalence is
\[
p=\Prob(V_i=1\mid D_i=1),
\]
then the observed positive-audit rate is
\[
\widetilde{p}
=
\Prob(\widetilde{V}_i=1\mid D_i=1)
=
\alpha_Ap+(1-\beta_A)(1-p).
\]
Solving for \(p\) gives
\[
p
=
\frac{
\widetilde{p}+\beta_A-1
}{
\alpha_A+\beta_A-1
},
\]
provided
\[
\alpha_A+\beta_A>1.
\]

\begin{proposition}[Identifiability under known audit error]
If audit sensitivity \(\alpha_A\) and specificity \(\beta_A\) are known and satisfy
\[
\alpha_A+\beta_A>1,
\]
then the latent suppressed-value prevalence \(p\) is identifiable from the observed positive-audit rate \(\widetilde p\) through
\[
p
=
\frac{\widetilde p+\beta_A-1}{\alpha_A+\beta_A-1}.
\]
If
\[
\alpha_A+\beta_A=1,
\]
then the audit judgment is statistically independent of the truth and \(p\) is not identifiable from \(\widetilde p\).
\end{proposition}

\begin{proof}
The first statement follows by rearranging
\[
\widetilde p
=
(\alpha_A+\beta_A-1)p+1-\beta_A.
\]
When the coefficient of \(p\) is nonzero, the equation has a unique solution. When
\[
\alpha_A+\beta_A-1=0,
\]
the observed rate reduces to
\[
\widetilde p=1-\beta_A=\alpha_A
\]
for every possible value of \(p\). Hence the observation contains no information about the latent prevalence.
\end{proof}

Auditor error is unlikely to be homogeneous. A polished conventional residual and an unusual terminologically immature residual can receive different sensitivities even when their eventual values are comparable. Let
\[
\alpha_A(x)
=
\Prob(\widetilde V=1\mid V=1,X=x)
\]
and
\[
\beta_A(x)
=
\Prob(\widetilde V=0\mid V=0,X=x).
\]
If the system estimates only global averages, it can satisfy an overall audit-accuracy target while remaining nearly blind in a particular representational region. Audit validation must therefore be stratified over the same features that may influence initial routing.

Multiple auditors permit partial estimation of audit uncertainty. Let
\[
\widetilde V_{ij}
\]
be auditor \(j\)'s judgment of residual \(i\). A latent-class model may assume
\[
\Prob(\widetilde V_{ij}\mid V_i,\widetilde V_{ik})
=
\Prob(\widetilde V_{ij}\mid V_i)
\]
for \(j\neq k\). Under sufficient numbers of conditionally independent auditors and suitable nondegeneracy assumptions, latent value prevalence and auditor confusion parameters can be estimated jointly. The conditional-independence assumption is fragile because auditors may share the same archive, vocabulary, institutional priors, and model-generated reconstruction. Apparent agreement can therefore measure common dependence rather than independent validation.

Let the residual correlation between auditors \(j\) and \(k\), conditional on latent value, be
\[
\rho_{jk\mid V}
=
\Corr(
\widetilde V_{ij},
\widetilde V_{ik}
\mid
V_i
).
\]
Positive \(\rho_{jk\mid V}\) causes naive independence models to overstate the evidential force of agreement. Audit diversity should consequently include diversity of methods, disciplinary priors, and representations, not merely multiple instances of the same evaluator.

\subsection{Appeal as a state transition rather than a complaint channel}

An appeal is not merely an opportunity to repeat the original submission. It is an operator acting on the routing state. Let
\[
\sigma_i
=
\left(
x_i,
a_i,
p_i,
r_i,
u_i,
\omega_i,
\thetaState_i,
d_i,
\K
\right)
\]
be the complete state associated with submission \(i\), including its decision \(d_i\). An appeal operator is
\[
\mathfrak{a}:
(\sigma_i,\chi_i)
\longrightarrow
\sigma_i',
\]
where \(\chi_i\) is the appeal content. The output state may alter the charitable reconstruction, prerequisite diagnosis, competency estimate, archive comparison, residual extraction, obstruction state, recipient match, or decision.

A valid appeal must be capable of changing at least one decision-relevant quantity. If
\[
S(\sigma_i')=S(\sigma_i)
\]
for every admissible \(\chi_i\), then the purported appeal mechanism is procedurally decorative. It accepts input without admitting any transition capable of reversing the decision.

\begin{definition}[Substantive appealability]
A routing decision at state \(\sigma_i\) is substantively appealable when there exists an admissible appeal \(\chi_i\) such that
\[
\Prob\!\left(
d_i'\neq d_i
\mid
\sigma_i,\chi_i
\right)>0
\]
and the transition probability depends on the evidential content of \(\chi_i\), not merely on random reconsideration.
\end{definition}

The appeal content may assert that the archive answer does not entail the reconstructed claim, that a prerequisite has already been demonstrated, that the system has conflated understanding with assent, that the quotient removed a decisive distinction, that an alleged inconsistency is a local-coordinate mismatch, that the assigned recipient is wrong, or that new evidence changes the residual. These are mathematically different challenges because they target different components of the state.

Let
\[
\chi_i
=
\left(
\chi_i^{Q},
\chi_i^{K},
\chi_i^{\theta},
\chi_i^{G},
\chi_i^{r},
\chi_i^{c}
\right),
\]
where the components challenge reconstruction, archive subsumption, participant-state inference, gluing analysis, recipient relation, and cost estimation. The appeal update can be expressed as
\[
\sigma_i'
=
\mathcal{B}\!\left(
\sigma_i,
\chi_i
\right),
\]
where \(\mathcal{B}\) is a posterior update operator.

If the original decision was based on posterior
\[
\Prob(V_i=1\mid\mathscr{I}_i),
\]
then appeal evidence \(\chi_i\) produces
\[
\Prob(V_i=1\mid\mathscr{I}_i,\chi_i)
=
\frac{
\Prob(\chi_i\mid V_i=1,\mathscr{I}_i)
\Prob(V_i=1\mid\mathscr{I}_i)
}{
\Prob(\chi_i\mid\mathscr{I}_i)
}.
\]
The appeal changes the decision when the posterior crosses the recipient-relative threshold
\[
\tau_{ir}^{*}
=
\frac{c_{ir}^{+}}{c_{ir}^{+}+c_{ir}^{-}}.
\]

\begin{proposition}[Minimum Bayes factor for reversal]
Let the pre-appeal probability of value be \(p_i\), let the escalation threshold be \(\tau\in(0,1)\), and suppose
\[
p_i<\tau.
\]
Define the appeal Bayes factor
\[
B_i
=
\frac{
\Prob(\chi_i\mid V_i=1,\mathscr{I}_i)
}{
\Prob(\chi_i\mid V_i=0,\mathscr{I}_i)
}.
\]
The appeal reverses the non-escalation decision exactly when
\[
B_i
>
\frac{\tau(1-p_i)}{p_i(1-\tau)}.
\]
\end{proposition}

\begin{proof}
The posterior odds equal the prior odds multiplied by the Bayes factor:
\[
\frac{p_i'}{1-p_i'}
=
B_i\frac{p_i}{1-p_i}.
\]
Reversal occurs when \(p_i'>\tau\), equivalently when
\[
\frac{p_i'}{1-p_i'}
>
\frac{\tau}{1-\tau}.
\]
Substituting the posterior-odds identity gives
\[
B_i\frac{p_i}{1-p_i}
>
\frac{\tau}{1-\tau}.
\]
Rearrangement yields
\[
B_i
>
\frac{\tau(1-p_i)}{p_i(1-\tau)}.
\]
\end{proof}

The proposition shows why vague requests for reconsideration often fail even under a fair appeal mechanism. An appeal must supply evidence whose likelihood under a valuable-residual hypothesis is sufficiently larger than its likelihood under a nonvaluable-residual hypothesis. The system should therefore expose which component of the original decision requires evidence rather than forcing the participant to guess how to change an opaque score.

\subsection{Appeal accessibility and burden}

An appeal mechanism can exist formally while remaining inaccessible in practice. Let
\[
B_i^{A}
\]
denote the participant burden of constructing an appeal. Let
\[
G_i^{A}
\]
denote the expected gain from successful reconsideration, and let
\[
p_i^{R}
\]
denote the participant's subjective probability that an appeal will reverse or materially improve the decision. A rational participant appeals when
\[
p_i^{R}G_i^{A}
>
B_i^{A}.
\]
If the system requires specialized legal, mathematical, or institutional language to challenge an erroneous reconstruction, then \(B_i^{A}\) can be correlated with the same surface competencies that produced the original false suppression.

Let \(Z_i\) represent the participant's access resources and suppose
\[
B_i^{A}=b(Z_i,X_i).
\]
If
\[
\frac{\partial b}{\partial Z_i}<0,
\]
then better-resourced participants face lower appeal burden. Observed appeal rates will be selection-biased:
\[
\Prob(\text{appeal}\mid D,V,Z)
\neq
\Prob(\text{appeal}\mid D,V).
\]
Successful appeals alone cannot estimate false suppression because many suppressed participants will never appeal.

\begin{theorem}[Appeal-only underestimation]
Suppose a discarded valuable residual is appealed with probability
\[
a_i\in[0,1],
\]
and an appealed false suppression is correctly reversed with probability
\[
r_i\in[0,1].
\]
Let
\[
\psi=\Prob(D=1\mid V=1)
\]
be the true false-suppression rate. The observed rate of successfully reversed false suppressions among valuable residuals is
\[
\psi_{\mathrm{rev}}
=
\E[a_ir_i\mid D=1,V=1]\psi.
\]
Consequently
\[
\psi_{\mathrm{rev}}\leq\psi,
\]
with equality only when
\[
a_ir_i=1
\]
almost surely among falsely suppressed valuable residuals.
\end{theorem}

\begin{proof}
A successful reversal requires the joint occurrence of false suppression, appeal, and correct reconsideration. Conditional on \(V=1\),
\[
\Prob(\text{successful reversal})
=
\Prob(D=1\mid V=1)
\E[a_ir_i\mid D=1,V=1].
\]
This gives the stated identity. Since \(a_ir_i\in[0,1]\), its conditional expectation is at most one, proving the inequality. Equality requires \(a_ir_i=1\) almost surely on the conditioning event.
\end{proof}

Appeals are therefore corrective mechanisms but not sufficient audit samples. Random review captures cases whose participants lack the time, vocabulary, confidence, or continued interest required to appeal.

\subsection{Decision provenance and reconstructible audit trails}

An audit requires more than the final decision. It requires the state from which the decision followed. Let the decision record be
\[
\Pi_i
=
\left(
x_i^{0},
\widehat{x}_i,
Q(x_i),
K_i^{\mathrm{ret}},
p_i,
\thetaState_i,
q_i,
r_i,
\omega_i,
s_{ir},
c_{ir},
\tau_{ir},
d_i,
\mathcal{M}_i,
t_i
\right),
\]
where \(x_i^{0}\) is the original expression, \(\widehat{x}_i\) is the charitable reconstruction, \(Q(x_i)\) is the normalized claim set, \(K_i^{\mathrm{ret}}\) is the retrieved archive material, \(p_i\) is the prerequisite route, \(q_i\) is the applied quotient or projection, \(r_i\) is the residual, \(\omega_i\) is the obstruction state, \(s_{ir}\) is the escalation score, \(c_{ir}\) is the predicted attention cost, \(\tau_{ir}\) is the operative threshold, \(d_i\) is the decision, \(\mathcal{M}_i\) identifies the model and policy versions, and \(t_i\) records relevant times.

The record must separate observed, retrieved, inferred, and generated content. Let
\[
\Pi_i
=
\Pi_i^{\mathrm{obs}}
\sqcup
\Pi_i^{\mathrm{ret}}
\sqcup
\Pi_i^{\mathrm{inf}}
\sqcup
\Pi_i^{\mathrm{gen}}.
\]
If these components are collapsed, an audit cannot determine whether a disputed premise originated with the participant, the archive, the reconstruction model, or the routing policy.

Define provenance completeness by
\[
\mathrm{PC}(\Pi_i)
=
\frac{
\sum_{j=1}^{m}w_j
\ind\{\text{required provenance field }j\text{ is present}\}
}{
\sum_{j=1}^{m}w_j
}.
\]
A packet satisfies the provenance constraint when
\[
\mathrm{PC}(\Pi_i)\geq p_0.
\]
Presence alone is insufficient if the record cannot reproduce the decision.

\begin{definition}[Decision reconstructibility]
A decision record \(\Pi_i\) is reconstructible under policy version \(\mathcal{M}_i\) when there exists a deterministic replay operator \(\mathcal{R}_{\mathcal{M}_i}\) such that
\[
\mathcal{R}_{\mathcal{M}_i}(\Pi_i)=d_i.
\]
For stochastic policies, reconstructibility additionally requires the random seed or sampled action probability needed to reproduce the realized decision.
\end{definition}

\begin{proposition}[Necessity of versioned provenance]
Suppose the routing map changes from \(\mathcal{M}_t\) to \(\mathcal{M}_{t'}\). If the decision record stores neither the earlier map nor a sufficient version identifier, then a disagreement between replayed and recorded decisions cannot generally be attributed uniquely to record corruption, policy drift, model drift, or nondeterminism.
\end{proposition}

\begin{proof}
Let the recorded state be \(\sigma_i\) and decision be
\[
d_i=\mathcal{M}_t(\sigma_i).
\]
A later replay produces
\[
d_i'=\mathcal{M}_{t'}(\sigma_i).
\]
If \(d_i'\neq d_i\) and \(\mathcal{M}_t\) is unavailable, then the evidence is compatible with at least four hypotheses: the state record differs from the original state, the policy changed, the underlying model changed while retaining the same nominal policy, or the original policy was stochastic. Since the missing version information prevents these hypotheses from being distinguished, the cause is not uniquely identifiable.
\end{proof}

Audit logs must therefore preserve not only what was decided but the representational and computational conditions under which the decision became admissible.

\subsection{Projection audits and distinction-loss tests}

A routing audit that examines only the final score can miss the earliest and most consequential error. If CLIO projection collapses two states that require different routing actions, every downstream score is computed on an already corrupted representation.

Let
\[
q:\mathcal{X}\to\mathcal{Z}
\]
be the operative projection and let
\[
f_r:\mathcal{X}\to\mathcal{Y}_r
\]
be the ideal recipient-relative routing disposition. Exact sufficiency requires
\[
q(x)=q(y)
\Longrightarrow
f_r(x)=f_r(y).
\]
Define the projection-collision risk by
\[
\mathcal{C}_r(q)
=
\Prob\!\left(
f_r(X)\neq f_r(X')
\given
q(X)=q(X')
\right).
\]
A practical audit samples pairs or neighborhoods inside quotient fibers and tests whether distinctions erased by \(q\) predict divergent downstream outcomes.

Let
\[
\mathcal{F}_z=q^{-1}(z)
\]
be the fiber above \(z\). For a routing outcome \(Y\), define within-fiber variance
\[
\Var(Y\mid q(X)=z).
\]
The expected quotient insufficiency is
\[
\mathcal{I}_r(q)
=
\E_{Z}
\left[
\Var(Y\mid Z)
\right].
\]
If \(Y\) is deterministic given \(X\), then
\[
\mathcal{I}_r(q)=0
\]
is equivalent to \(Y\) being almost surely measurable with respect to \(\sigma(q(X))\).

\begin{theorem}[Conditional-variance criterion for functional sufficiency]
Let \(Y=f_r(X)\in L^2\), and let \(Z=q(X)\). Then the following conditions are equivalent:
\[
Y=g(Z)
\quad\text{almost surely for some measurable }g,
\]
\[
\Var(Y\mid Z)=0
\quad\text{almost surely},
\]
and
\[
\E[\Var(Y\mid Z)]=0.
\]
\end{theorem}

\begin{proof}
If \(Y=g(Z)\) almost surely, then conditioning on \(Z\) fixes \(Y\), so
\[
\Var(Y\mid Z)=0
\]
almost surely. This immediately implies
\[
\E[\Var(Y\mid Z)]=0.
\]
Conversely, conditional variance is nonnegative. If its expectation is zero, then
\[
\Var(Y\mid Z)=0
\]
almost surely. Since
\[
\Var(Y\mid Z)
=
\E\!\left[
\left(
Y-\E[Y\mid Z]
\right)^2
\given Z
\right],
\]
zero conditional variance implies
\[
Y=\E[Y\mid Z]
\]
almost surely. The conditional expectation \(\E[Y\mid Z]\) is measurable with respect to \(\sigma(Z)\), so by the Doob-Dynkin lemma there exists a measurable \(g\) such that
\[
\E[Y\mid Z]=g(Z).
\]
Therefore \(Y=g(Z)\) almost surely.
\end{proof}

An audit may estimate \(\mathcal{I}_r(q)\) by deliberately soliciting or reconstructing counterfactual variants that preserve \(q(x)\) while changing discarded coordinates. For \(L_2\)-normalization,
\[
q(x)=\frac{x}{\lVert x\rVert_2},
\]
the fiber contains positive scalar multiples
\[
\mathcal{F}_{q(x)}
=
\{\alpha x:\alpha>0\}.
\]
A magnitude-sensitive routing outcome violates sufficiency whenever there exist \(\alpha,\beta>0\) such that
\[
f_r(\alpha x)\neq f_r(\beta x).
\]
Projection audit therefore requires interventions along the discarded radial direction, not merely evaluation on naturally occurring normalized states.

\subsection{Coupling audits and hidden shared parameters}

HYDRA introduces a different audit problem. An apparently separable transition may conceal coupling through shared parameters. Let
\[
\mathcal{H}(P,R,M;\vartheta)
\]
be the HYDRA transition operator, where
\[
\vartheta
=
(\mu_c,\Sigma_c,\FieldPhi,\FieldV,\FieldS,\ldots)
\]
collects shared state and field parameters. A superficial factorization may appear as
\[
\mathcal{H}(P,R,M;\vartheta)
=
h\!\left(
h_P(P;\vartheta),
h_R(R;\vartheta),
h_M(M;\vartheta)
\right).
\]
This does not establish operational separability because each component depends on the same \(\vartheta\).

Define the mixed intervention derivative
\[
\mathcal{J}_{PR}
=
\frac{\partial^2\mathcal{H}}{\partial P\,\partial R},
\]
with analogous terms \(\mathcal{J}_{PM}\) and \(\mathcal{J}_{RM}\). Nonzero mixed derivatives provide local evidence of coupling. Yet even vanishing direct mixed derivatives can miss coupling transmitted through shared parameters. If
\[
\vartheta=\vartheta(P,R,M),
\]
then the total derivative satisfies
\[
\frac{d\mathcal{H}}{dP}
=
\frac{\partial\mathcal{H}}{\partial P}
+
\frac{\partial\mathcal{H}}{\partial\vartheta}
\frac{\partial\vartheta}{\partial P}.
\]
The total mixed derivative contains terms such as
\[
\frac{\partial^2\mathcal{H}}
{\partial\vartheta^2}
\frac{\partial\vartheta}{\partial P}
\frac{\partial\vartheta}{\partial R}
+
\frac{\partial\mathcal{H}}{\partial\vartheta}
\frac{\partial^2\vartheta}{\partial P\,\partial R}.
\]
An audit that holds \(\vartheta\) numerically fixed tests only explicit coupling, while an audit that allows \(\vartheta\) to update tests effective coupling under the actual transition dynamics.

Let the separable baseline be
\[
\mathcal{H}_{0}(P,R,M)
=
h\!\left(
h_P(P),h_R(R),h_M(M)
\right),
\]
and define the coupling residual
\[
\mathcal{E}_{\mathcal{H}}(P,R,M)
=
\mathcal{H}(P,R,M)-\mathcal{H}_{0}(P,R,M).
\]
A useful coupling audit estimates
\[
\Gamma_{\mathcal{H}}
=
\E\!\left[
\lVert
\mathcal{E}_{\mathcal{H}}(P,R,M)
\rVert^2
\right].
\]
A nonzero value rejects the proposed separable baseline. It does not prove that the coupling improves performance, preserves admissibility, or supports the intended cognitive interpretation.

To test usefulness, let \(\mathcal{L}_{\mathrm{task}}\) be downstream task loss and compare
\[
\Delta_{\mathrm{coupling}}
=
\E[
\mathcal{L}_{\mathrm{task}}(\mathcal{H}_0)
]
-
\E[
\mathcal{L}_{\mathrm{task}}(\mathcal{H})
].
\]
Operational coupling is beneficial relative to the tested baseline only when
\[
\Delta_{\mathrm{coupling}}>0
\]
and the improvement is not purchased by unacceptable instability, opacity, or false suppression elsewhere.

\subsection{Obstruction audits}

A gluing obstruction can be mishandled in two symmetric ways. The system can suppress a valuable residual because its local formulations do not assemble under the current cover, or it can escalate a merely inconsistent residual because local sophistication is mistaken for global coherence.

Let \(\mathfrak{U}=\{U_\alpha\}\) be a cover of the representational domain and let
\[
s_\alpha\in\mathcal{F}(U_\alpha)
\]
be local sections. On overlaps \(U_{\alpha\beta}\), define discrepancy
\[
\delta_{\alpha\beta}
=
s_\alpha|_{U_{\alpha\beta}}
-
s_\beta|_{U_{\alpha\beta}}.
\]
When the local data form a Čech cocycle, let
\[
[\delta]\in\check{H}^{1}(\mathfrak{U},\mathcal{F})
\]
be its obstruction class. An audit asks whether the detected class is intrinsic or an artifact of the chosen representation, cover, reconstruction, or transition maps.

Let \(q\) be the current representation and let \(q'\) be a refinement. There is an induced comparison map
\[
\varphi_{q'q}:
\check{H}^{1}(\mathfrak{U}',\mathcal{F}')
\longrightarrow
\check{H}^{1}(\mathfrak{U},\mathcal{F}).
\]
If an obstruction disappears under refinement, so that
\[
[\delta'] = 0
\qquad
\text{while}
\qquad
[\delta]\neq0,
\]
the original obstruction may have been created by projection collapse. If it persists across justified refinements and compatible covers, confidence increases that the failure is structural.

\begin{definition}[Representation-stable obstruction]
An obstruction class \([\delta_q]\) detected under representation \(q\) is stable over refinement family \(\mathfrak{Q}\) when every admissible refinement \(q'\succeq q\) induces a nonzero compatible class
\[
[\delta_{q'}]\neq0
\]
whose image under the comparison map equals \([\delta_q]\).
\end{definition}

An obstruction audit should therefore record
\[
\left(
[\delta_q],
\mathfrak{U},
\mathcal{F},
q,
\mathfrak{Q}_{\mathrm{tested}}
\right)
\]
rather than the bare statement that the submission is inconsistent. The same visible incompatibility can indicate incoherence, a missing transition map, an inadequate cover, or a quotient that erased the coordinate needed to glue local sections.

\subsection{Sequential audit and posterior monitoring}

False-suppression auditing is continuous rather than terminal. Let
\[
\psi_t
\]
be the false-suppression probability during period \(t\). A stationary model assumes
\[
\psi_t=\psi,
\]
but model updates, archive growth, shifting submission populations, and changing recipient priorities make this assumption fragile.

A Bayesian monitoring model may assign
\[
\psi_t\sim\operatorname{Beta}(\alpha_t,\beta_t).
\]
If audits in period \(t\) reveal \(k_t\) valuable suppressions among \(n_t\) effectively sampled valuable cases under a simplified binomial observation model, the posterior update is
\[
\alpha_{t+1}
=
\delta\alpha_t+k_t,
\]
\[
\beta_{t+1}
=
\delta\beta_t+n_t-k_t,
\]
where \(\delta\in(0,1]\) is a forgetting factor. When \(\delta<1\), old evidence loses weight, permitting adaptation to policy drift.

The posterior mean is
\[
\E[\psi_t\mid\mathcal{D}_{1:t}]
=
\frac{\alpha_t}{\alpha_t+\beta_t},
\]
and the posterior variance is
\[
\Var(\psi_t\mid\mathcal{D}_{1:t})
=
\frac{
\alpha_t\beta_t
}{
(\alpha_t+\beta_t)^2(\alpha_t+\beta_t+1)
}.
\]
A safety rule may require
\[
\Prob(\psi_t>\eps\mid\mathcal{D}_{1:t})
<
\delta_{\mathrm{risk}}.
\]
If the posterior violation probability exceeds \(\delta_{\mathrm{risk}}\), the system must increase exploration, lower the escalation threshold, revise a projection, expand a prerequisite route, or suspend an affected policy region.

For importance-weighted audits with unequal inclusion probabilities, simple Beta-binomial conjugacy does not apply exactly. One may instead use weighted likelihood
\[
L(\psi)
\propto
\prod_{i:Z_i=1}
\left[
\psi^{V_i}
(1-\psi)^{1-V_i}
\right]^{1/\pi_i},
\]
or a design-based estimator with asymptotic confidence intervals. The weighted likelihood is a pseudo-likelihood because inverse-probability exponents need not correspond to independent replicated observations. Its uncertainty must be calibrated to the sampling design rather than read directly from an ordinary conjugate posterior.

\subsection{Change detection in routing behavior}

Let
\[
\widehat{\psi}_{g,t}
\]
be the estimated false-suppression rate for stratum \(g\) during period \(t\). Define a cumulative deviation process
\[
C_{g,t}^{+}
=
\max\left\{
0,
C_{g,t-1}^{+}
+
\widehat{\psi}_{g,t}
-
\eps_g
-
k_g
\right\},
\]
where \(k_g>0\) is a tolerance parameter. An alarm occurs when
\[
C_{g,t}^{+}>h_g.
\]
This detects persistent upward shifts rather than isolated noisy estimates.

Alternatively, let \(P_t\) be the distribution of routing features among escalated cases and \(Q_t\) the distribution among eligible submissions. Distributional exclusion can be measured by
\[
\KL(P_t\Vert Q_t)
=
\int
\log\!\left(
\frac{dP_t}{dQ_t}
\right)dP_t
\]
when \(P_t\) is absolutely continuous with respect to \(Q_t\). A growing divergence can indicate that escalation is narrowing toward a small representational region even before enough audited outcomes exist to estimate false suppression precisely.

Support failure is more severe than large divergence. If
\[
Q_t(B)>0
\qquad\text{and}\qquad
P_t(B)=0
\]
for some region \(B\), then no eligible submissions from \(B\) are being escalated. The exploration policy must determine whether this is justified or whether \(B\) has become an epistemically sealed rejection region.

\subsection{Calibration of escalation probabilities}

Suppose the routing system assigns each residual a predicted escalation value probability
\[
\widehat p_i
=
\Prob(V_i=1\mid X_i).
\]
Calibration requires
\[
\Prob(V_i=1\mid \widehat p_i=p)=p.
\]
Define the calibration function
\[
m(p)
=
\E[V_i\mid\widehat p_i=p].
\]
Perfect calibration means
\[
m(p)=p
\]
almost everywhere.

Observed calibration among escalated submissions estimates
\[
\E[V_i\mid\widehat p_i=p,E_i=1],
\]
which need not equal
\[
\E[V_i\mid\widehat p_i=p].
\]
If selection depends only on \(\widehat p_i\), then conditioning on \(\widehat p_i\) can make the two coincide. If selection also depends on cost, participant class, recipient relation, obstruction state, or hidden policy variables, escalated-only calibration is biased.

Under positivity and conditional exchangeability,
\[
V_i(E)\ind E_i\mid X_i,
\]
one may estimate population calibration with weights
\[
w_i
=
\frac{1}{e(X_i)}
\]
for escalated cases. The weighted calibration moment is
\[
\E\!\left[
\frac{E_i}{e(X_i)}
\bigl(
V_i-\widehat p_i
\bigr)
\phi(\widehat p_i)
\right]
=
0
\]
for every suitable test function \(\phi\) under correct calibration.

\begin{proof}
Conditioning on \(X_i\),
\[
\E\!\left[
\frac{E_i}{e(X_i)}
(V_i-\widehat p_i)
\phi(\widehat p_i)
\given X_i
\right]
=
\frac{\E[E_i\mid X_i]}{e(X_i)}
\E[
V_i-\widehat p_i
\mid X_i
]
\phi(\widehat p_i),
\]
where conditional exchangeability permits separation of the selection indicator and potential value. Since
\[
\E[E_i\mid X_i]=e(X_i),
\]
this becomes
\[
\E[
V_i-\widehat p_i
\mid X_i
]
\phi(\widehat p_i).
\]
Taking expectations yields the population calibration moment, which vanishes when
\[
\widehat p_i=\E[V_i\mid X_i].
\]
\end{proof}

Calibration is not sufficiency. A perfectly calibrated system can still suppress nearly every valuable case if it assigns low probabilities accurately to a population in which value is rare but consequential. Nor does calibration establish fairness across representation classes. A global calibration identity can conceal opposite errors that cancel between groups.

\subsection{Doubly robust estimation}

Let
\[
m(x)
=
\E[V_i\mid X_i=x,D_i=1]
\]
be an outcome model for value among suppressed residuals, and let \(\pi(x)\) be the audit inclusion probability. A doubly robust estimator for the mean suppressed value is
\[
\widehat{\mu}_{\mathrm{DR}}
=
\frac{1}{N_D}
\sum_{i:D_i=1}
\left[
\widehat m(X_i)
+
\frac{Z_i}{\widehat\pi(X_i)}
\left(
\widetilde V_i-\widehat m(X_i)
\right)
\right],
\]
where \(N_D\) is the number of suppressed residuals.

\begin{theorem}[Double robustness]
Assume audit labels equal latent values and regularity conditions permit interchange of expectations. The estimator's population analogue
\[
\mu_{\mathrm{DR}}
=
\E\left[
\widehat m(X)
+
\frac{Z}{\widehat\pi(X)}
(V-\widehat m(X))
\given D=1
\right]
\]
equals
\[
\mu
=
\E[V\mid D=1]
\]
if either
\[
\widehat m(X)=\E[V\mid X,D=1]
\]
almost surely or
\[
\widehat\pi(X)=\Prob(Z=1\mid X,D=1)
\]
almost surely.
\end{theorem}

\begin{proof}
Condition on \(X\) and \(D=1\). Let
\[
m_0(X)=\E[V\mid X,D=1]
\]
and
\[
\pi_0(X)=\Prob(Z=1\mid X,D=1).
\]
The conditional expectation of the augmentation term is
\[
\E\left[
\frac{Z}{\widehat\pi(X)}
(V-\widehat m(X))
\given X,D=1
\right]
=
\frac{\pi_0(X)}{\widehat\pi(X)}
\left(
m_0(X)-\widehat m(X)
\right).
\]
Hence the conditional expectation of the full expression is
\[
\widehat m(X)
+
\frac{\pi_0(X)}{\widehat\pi(X)}
\left(
m_0(X)-\widehat m(X)
\right).
\]
If \(\widehat\pi=\pi_0\), this reduces to \(m_0(X)\). If \(\widehat m=m_0\), the correction term vanishes and the expression again equals \(m_0(X)\). Taking expectations conditional on \(D=1\) yields
\[
\E[m_0(X)\mid D=1]
=
\E[V\mid D=1].
\]
\end{proof}

Double robustness does not mean robustness to arbitrary failure. If both the audit propensity model and the value model are wrong, the estimator may be biased. If some region has true audit probability zero, inverse weighting cannot recover it. If audit labels systematically inherit the initial model's representational blindness, neither component observes the missing distinction.

\subsection{Delayed discovery and survival analysis}

Some suppressed residuals are later independently rediscovered, cited, formalized, or recognized as valuable. Let
\[
T_i^{R}
\]
be the time from the original decision to later recognition, with
\[
T_i^{R}=\infty
\]
when recognition never occurs during observation. The recognition hazard is
\[
h_i(t)
=
\lim_{\Delta t\downarrow0}
\frac{
\Prob(
t\leq T_i^{R}<t+\Delta t
\mid
T_i^{R}\geq t,
X_i
)
}{
\Delta t
}.
\]
Observed late recognition supplies lower-bound evidence about false suppression. If a residual discarded at \(t_i\) is later recognized as valuable, then the earlier decision was at least temporally inefficient and may have been false under the original horizon.

Censoring complicates inference. A residual not recognized by the end of observation may be worthless, permanently invisible, independently rediscovered under another name, or valuable but still unrecognized. Let
\[
C_i
\]
be the censoring time and observe
\[
\widetilde T_i
=
\min(T_i^{R},C_i),
\qquad
\Delta_i
=
\ind\{T_i^{R}\leq C_i\}.
\]
Under independent censoring conditional on \(X_i\), survival methods estimate
\[
S_R(t\mid X_i)
=
\Prob(T_i^{R}>t\mid X_i).
\]
Independent censoring is doubtful when observation length depends on the same prestige, vocabulary, or institutional access that affects recognition. A residual from a well-connected participant may be easier to track and more likely to be independently surfaced. Delayed-recognition data must therefore be combined with randomized audit rather than treated as a self-sufficient ground truth.

\subsection{Adversarial adaptation to the audit mechanism}

Once audit criteria become known, participants and routing components can adapt to them. Let the participant choose expression strategy
\[
s_i\in\mathcal{S}
\]
to maximize
\[
U_i(s_i)
=
\Prob(E_i=1\mid s_i)G_i
-
B_i(s_i).
\]
If the escalation classifier rewards conventional vocabulary, participants who can imitate that vocabulary will do so whether or not it improves conceptual precision. The audit distribution then shifts from the natural submission distribution to a strategically selected distribution.

The routing system itself can also game aggregate constraints. Suppose it must satisfy
\[
\widehat{\Prob}(D\mid V)\leq\eps
\]
under an audit sampler concentrated in easily assessed strata. It can improve the measured rate by performing well on those strata while leaving unobserved regions unchanged.

Let \(P_A\) be the audit distribution and \(P_T\) the target submission distribution. The measured audit risk is
\[
R_A
=
\E_{P_A}[\ell(D,V)],
\]
while the target risk is
\[
R_T
=
\E_{P_T}[\ell(D,V)].
\]
If
\[
P_A\neq P_T,
\]
then low \(R_A\) does not imply low \(R_T\). Under absolute continuity,
\[
R_T
=
\E_{P_A}
\left[
\frac{dP_T}{dP_A}\ell(D,V)
\right].
\]
Large density ratios produce unstable estimates and expose regions poorly covered by audit.

An audit design must therefore mix predictable stratification with concealed or randomized sampling. Predictability supports transparent governance. Randomness prevents the audited sample from collapsing into the subset on which the system has learned to look good.

\subsection{An impossibility result for complete validation}

The strongest conceivable requirement would be to guarantee simultaneously that every valuable residual is escalated, no valueless residual consumes recipient attention, and the recipient-facing queue remains stable under unbounded arrivals. These conditions are generally incompatible.

\begin{theorem}[Finite-attention validation impossibility]
Suppose submissions arrive at unbounded rate \(\lambda\), determining whether \(V_i=1\) requires positive expected attention
\[
\E[C_i^{V}]\geq c_0>0,
\]
and no pre-evaluation observable feature perfectly separates valuable from valueless residuals. A system with finite attention capacity \(A<\infty\) cannot simultaneously guarantee
\[
\Prob(D_i=1\mid V_i=1)=0,
\]
\[
\Prob(E_i=1\mid V_i=0)=0,
\]
and stable total evaluation workload for every admissible arrival distribution.
\end{theorem}

\begin{proof}
Because pre-evaluation features do not perfectly separate the classes, zero false suppression and zero false escalation require resolving the value label of every arrival through an evaluation consuming expected attention at least \(c_0\). The expected incoming evaluation workload is therefore at least
\[
\lambda c_0.
\]
For an unbounded arrival rate, choose an admissible \(\lambda>A/c_0\). Then
\[
\lambda c_0>A,
\]
so the incoming workload exceeds finite capacity and the evaluation queue is unstable. Avoiding instability requires declining to evaluate some arrivals or deciding them from imperfect features. The former permits valuable residuals to be discarded without evaluation, while the latter produces either false suppression or false escalation for some admissible distribution. Hence the three guarantees cannot hold simultaneously.
\end{proof}

The purpose of audit is not to defeat this impossibility. It is to make the unavoidable error distribution visible, estimable, contestable, and revisable. Finite attention forces selective evaluation. It does not force the system to conceal where selectivity fails.

\subsection{Audit as representational repair}

Audit is commonly imagined as an external check applied after the system has acted. In the integrated architecture, audit is also a learning operator that modifies the representations on which future routing depends.

Let an audited false suppression produce diagnostic record
\[
\zeta_i
=
\left(
q_i,
\delta_i^{\mathrm{lost}},
p_i^{\mathrm{mis}},
\omega_i^{\mathrm{mis}},
r_i^{\mathrm{mis}},
c_{ir}^{\mathrm{mis}},
\thetaState_i^{\mathrm{mis}}
\right),
\]
where the components identify a lost distinction, misassigned prerequisite, misclassified obstruction, incorrect residual, inflated cost, or mistaken participant-state estimate. Let the system parameters be
\[
\Theta_t
=
\left(
q_t,
G_t,
\K_t,
\mathcal{H}_t,
S_t,
\pi_t
\right).
\]
The audit update is
\[
\Theta_{t+1}
=
\mathcal{U}_{\Audit}(\Theta_t,\zeta_i).
\]

A local correction changes only the audited case. A structural correction changes the map that generated the case. Define the generalization set
\[
\mathcal{N}(\zeta_i)
=
\left\{
j:
\dist_{\Theta_t}(x_j,x_i)\leq\varepsilon_i
\right\}.
\]
If the diagnosed failure arises from a shared quotient, prerequisite edge, or scoring rule, then every \(j\in\mathcal{N}(\zeta_i)\) may require reconsideration.

Let
\[
\Delta_{\Audit}(\zeta_i)
=
\sum_{j\in\mathcal{N}(\zeta_i)}
\left[
\Prob_{\Theta_{t+1}}(D_j=1\mid V_j=1)
-
\Prob_{\Theta_t}(D_j=1\mid V_j=1)
\right].
\]
A successful structural repair has
\[
\Delta_{\Audit}(\zeta_i)<0.
\]
The value of auditing case \(i\) therefore includes not only the correction of \(i\) but the expected reduction of future suppression across its failure neighborhood.

Define the expected audit value
\[
\mathcal{V}_{\Audit}(i)
=
\Prob(\text{decision error}\mid\mathscr{I}_i)
\left[
G_i^{\mathrm{case}}
+
G_i^{\mathrm{struct}}
\right]
-
c_i^{\Audit},
\]
where \(G_i^{\mathrm{case}}\) is the value of correcting the individual case, \(G_i^{\mathrm{struct}}\) is the value of repairing the shared mechanism, and \(c_i^{\Audit}\) is the audit cost. This resembles active learning, but the target is broader than predictive accuracy. The audit may revise the quotient space, the prerequisite topology, the coupling model, the recipient relation, or the meaning of an admissible residual.

\begin{proposition}[Structural audit dominance]
Suppose cases \(i\) and \(j\) have equal probabilities of individual decision error and equal individual correction values. If auditing \(i\) is expected to repair a shared failure affecting \(m>1\) future cases while auditing \(j\) affects only itself, and audit costs are equal, then
\[
\mathcal{V}_{\Audit}(i)
>
\mathcal{V}_{\Audit}(j)
\]
whenever the expected per-case structural improvement is positive.
\end{proposition}

\begin{proof}
Let the common individual expected correction value be \(g\), the audit cost be \(c\), and the error probability be \(p\). Let the per-case expected structural gain from auditing \(i\) be \(s>0\) over \(m\) future cases. Then
\[
\mathcal{V}_{\Audit}(i)
=
p(g+ms)-c,
\]
while
\[
\mathcal{V}_{\Audit}(j)
=
pg-c.
\]
Their difference is
\[
pm s>0.
\]
\end{proof}

This creates a principled reason to audit recurrent edge failures, unusual projection collisions, and gluing obstructions even when their immediate recipient-relative values are uncertain. They may reveal a defect in the transformation architecture rather than merely an isolated classification error.

\subsection{A constrained audit and appeal objective}

Let \(A_{\Audit}\) be the audit attention budget and let \(A_{\Appeal}\) be the appeal-processing budget. Let \(z_i^{A}\in\{0,1\}\) indicate audit selection and \(z_i^{P}\in\{0,1\}\) indicate appeal review. Define
\[
v_i^{A}
=
G_i^{\mathrm{case}}
+
G_i^{\mathrm{struct}}
+
G_i^{\mathrm{est}}
+
G_i^{\mathrm{coverage}},
\]
where \(G_i^{\mathrm{est}}\) is the value of reducing uncertainty in the false-suppression estimate and \(G_i^{\mathrm{coverage}}\) is the value of preserving observability in an under-sampled region. Define appeal value
\[
v_i^{P}
=
G_i^{\mathrm{reversal}}
+
G_i^{\mathrm{state\ correction}}
+
G_i^{\mathrm{procedural}}.
\]
The joint allocation problem is
\[
\begin{aligned}
\text{maximize}\quad&
\sum_i z_i^{A}v_i^{A}
+
\sum_i z_i^{P}v_i^{P}
\\
\text{subject to}\quad&
\sum_i z_i^{A}c_i^{A}
\leq
A_{\Audit},
\\
&
\sum_i z_i^{P}c_i^{P}
\leq
A_{\Appeal},
\\
&
\pi_i^{A}\geq\pi_{\min,g(i)},
\\
&
\operatorname{delay}_i^{P}\leq H_i^{P},
\\
&
\mathrm{PC}(\Pi_i)\geq p_0.
\end{aligned}
\]
The minimum audit probability preserves identifiability. The appeal-delay constraint prevents a formally reversible decision from becoming practically irreversible through queueing delay. The provenance constraint ensures that reviewers can reconstruct what they are being asked to reconsider.

\subsection{The auditability condition for representational simplicity}

Representational simplicity is safe only when the consequences of discarded distinctions remain discoverable. Let \(q\) be the operative quotient and let
\[
\mathcal{L}(q)
\]
be the set of distinctions erased by \(q\). For a discarded distinction \(\delta\in\mathcal{L}(q)\), define its downstream relevance event
\[
R_\delta
=
\ind\!\left\{
\exists f\in\mathcal{F}_{\mathrm{future}}
:
f(x)\neq f(y)
\text{ for some }q(x)=q(y)
\text{ distinguished by }\delta
\right\}.
\]
Let
\[
O_\delta
\]
be the event that the audit architecture can observe evidence of the resulting divergence. Define auditability of \(q\) over future task family \(\mathcal{F}_{\mathrm{future}}\) by
\[
\mathfrak{A}(q)
=
\Prob(O_\delta=1\mid R_\delta=1).
\]
A projection can be task-sufficient under the present task family while having low auditability under future task expansion. Such a representation is locally efficient but brittle.

\begin{definition}[Repairable simplicity]
A quotient \(q\) is repairably simple at tolerance \((\eps,\delta)\) when its present task-relative loss is at most \(\eps\),
\[
\inf_g
\E[
\ell(f(X),g(q(X)))
]
\leq\eps,
\]
and every discarded distinction whose future relevance exceeds \(\delta\) has positive detection probability,
\[
\Prob(O_\delta=1\mid R_\delta=1)>0.
\]
\end{definition}

Exact sufficiency eliminates the need for repair within the declared domain because no task-relevant distinction is lost. Task-relative sufficiency accepts bounded loss. Repairable simplicity adds a temporal condition: when the task family changes and a previously discarded distinction becomes relevant, the architecture retains some path by which the failure can be detected and the quotient refined.

\begin{theorem}[No repair without an observable residual]
Let \(q\) collapse states \(x\) and \(y\) with
\[
q(x)=q(y)
\]
but
\[
f'(x)\neq f'(y)
\]
for a later task \(f'\). If every observable available to the audit and appeal system is measurable with respect to \(q(X)\), then no audit rule can distinguish the failure at \(x\) from the failure at \(y\).
\end{theorem}

\begin{proof}
Let every audit observable \(O\) have the form
\[
O=h(q(X))
\]
for some measurable \(h\). Since
\[
q(x)=q(y),
\]
we have
\[
O(x)=h(q(x))=h(q(y))=O(y)
\]
for every available audit observable. Therefore the complete observed audit state is identical at \(x\) and \(y\). Any audit rule measurable with respect to that state must assign the same output distribution to both cases. Since \(f'(x)\neq f'(y)\), the rule cannot correctly distinguish their task outcomes. Repair requires an observable that is not measurable solely through the failed quotient.
\end{proof}

This theorem explains why provenance must preserve more than the final compressed residual. If the original expression, discarded coordinates, alternative reconstructions, or local sections are destroyed, later audit remains trapped inside the same quotient that generated the error. A system cannot repair a distinction it has made impossible to observe.

\subsection{The central epistemic accounting identity}

Let \(N_V\) be the total number of valuable residuals arriving during a period. Partition them into valuable escalations \(N_{VE}\), valuable archive-mediated resolutions \(N_{VA}\), valuable delayed developments \(N_{VM}\), detected false suppressions \(N_{VD}^{\mathrm{det}}\), and undetected false suppressions \(N_{VD}^{\mathrm{undet}}\). Then
\[
N_V
=
N_{VE}
+
N_{VA}
+
N_{VM}
+
N_{VD}^{\mathrm{det}}
+
N_{VD}^{\mathrm{undet}}.
\]
Ordinary operational records observe the first four terms imperfectly and do not directly observe the fifth. Hence
\[
N_{VD}^{\mathrm{undet}}
=
N_V
-
N_{VE}
-
N_{VA}
-
N_{VM}
-
N_{VD}^{\mathrm{det}}
\]
cannot be computed because \(N_V\) is unknown.

Audit supplies an estimator
\[
\widehat N_{VD}^{\mathrm{undet}}
\]
through randomized review, delayed discovery, appeals, and model-based correction. The false-suppression estimate becomes
\[
\widehat{\psi}
=
\frac{
N_{VD}^{\mathrm{det}}
+
\widehat N_{VD}^{\mathrm{undet}}
}{
N_{VE}
+
N_{VA}
+
N_{VM}
+
N_{VD}^{\mathrm{det}}
+
\widehat N_{VD}^{\mathrm{undet}}
}.
\]
Every term depends upon a declared value criterion and time horizon. The estimate should therefore be reported with its audit design, inclusion probabilities, label model, uncertainty interval, and known regions of nonidentifiability.

The governing constraint is not
\[
\Prob(D\mid V)\leq\eps,
\]
because the left side is unavailable. The operational constraint is
\[
\Prob\!\left(
\psi>\eps
\given
\mathcal{D}_{\Audit}
\right)
\leq
\delta_{\mathrm{risk}},
\]
together with explicit lower bounds on audit coverage:
\[
\pi_i^{A}\geq\pi_{\min,g(i)}>0.
\]
The first condition controls posterior risk under the audit model. The second prevents the audit model from manufacturing confidence by excluding the regions in which errors are hardest to observe.

\subsection{Conclusion}

Audit and appeal complete the temporal logic of representational simplicity. Projection discards distinctions. Routing acts on the quotient. HYDRA assembles local components into an operational state. Escalation spends finite recipient attention on selected residuals. Audit asks whether the distinctions discarded earlier were in fact irrelevant to the futures the system later attempted to support. Appeal allows a participant to introduce evidence that the original representation, prerequisite diagnosis, coupling analysis, recipient match, or attention-cost estimate was wrong.

The false-suppression rate is not directly measurable because suppression removes the ordinary path by which value becomes visible. It can only be estimated through interventions that reopen a sample of closed paths. Random audit supplies identifiability. Stratification supplies coverage. Inverse-probability estimation corrects unequal sampling. Multiple evaluators expose judgment uncertainty. Delayed recognition supplies additional lower-bound evidence. Appeals correct individual cases and reveal failure modes that a sampling design may not have anticipated. Provenance preserves access to distinctions outside the quotient that made the original decision.

The resulting safety condition is
\[
\begin{aligned}
\text{maintain}\quad&
\lambda_r\Prob(E)\E[C_{ir}\mid E]<A_r,
\\
\text{while requiring}\quad&
\Prob\!\left(
\psi_r(H,\delta)>\eps
\given
\mathcal{D}_{\Audit}
\right)
\leq
\delta_{\mathrm{risk}},
\\
&
\pi_i^{A}\geq\pi_{\min,g(i)}>0,
\\
&
\mathrm{PC}(\Pi_i)\geq p_0,
\\
&
\Prob(d_i'\neq d_i\mid\chi_i)>0
\text{ for some evidentially admissible appeal}.
\end{aligned}
\]
Stability controls how much attention enters the recipient-facing system. Audit controls what can be learned about what remained outside it. Appeal controls whether an individual routing decision is an absorbing state. Provenance controls whether later repair can reach beyond the representation that produced the failure.

A finite-attention architecture cannot promise never to suppress value. It can promise something structurally meaningful: no rejection region is treated as self-validating merely because it produces no observed successes; no compressed representation is treated as permanently sufficient merely because its discarded distinctions have become invisible; and no decision is treated as epistemically complete while the system lacks a mechanism for discovering how that decision could have been wrong.
\section{Collective Recurrence, Distributional Novelty, and Lossy Observation}

The preceding sections treated the residual \(r_i\) as the principal unit of analysis. That perspective is necessary but incomplete. A submission can possess negligible residual novelty when considered individually while participating in a collectively novel pattern when considered as one event in a population. Repeated questions, repeated prerequisite failures, independently rediscovered counterexamples, correlated representational collapses, and geographically or temporally concentrated appeals can reveal changes in the domain that no single message establishes.

This distinction separates semantic novelty from distributional novelty. Semantic novelty concerns what remains after a particular submission has been reconstructed and compared with the archive. Distributional novelty concerns a change in the arrival process, dependency structure, failure topology, or joint configuration of many submissions. A familiar claim can therefore be semantically nonnovel and distributionally significant. Conversely, an individually unusual claim can be semantically novel while having no detectable population-level recurrence.

Let
\[
\mathcal{X}
\]
be the submission space and let
\[
N(B\times[0,t])
\]
count submissions arriving by time \(t\) whose transformed representations lie in measurable region \(B\subseteq\mathcal{X}\). The marked submission process is
\[
\mathcal{N}
=
\sum_{i\geq1}
\delta_{(t_i,z_i,m_i)},
\]
where \(t_i\) is the arrival time, \(z_i\) is the normalized representational state, and \(m_i\) contains marks such as participant state, prerequisite route, residual class, obstruction type, recipient relation, appeal outcome, and archive disposition.

For measurable \(B\), define the counting process
\[
N_B(t)
=
\mathcal{N}(B\times[0,t]).
\]
Under a reference model \(\mathbb{P}_0\), let its predictable intensity be
\[
\lambda_0(B,t).
\]
Distributional novelty occurs when the observed marked process becomes insufficiently explained by the reference intensity, mark distribution, dependency structure, or higher-order interaction model.

\subsection{Semantic novelty and distributional novelty}

Let
\[
Q(x_i)
\]
be the normalized claim set extracted from submission \(x_i\), and let
\[
R(x_i\mid\K_t)
=
\left\{
q\in Q(x_i):\K_t\nvDash q
\right\}
\]
be its semantic residual relative to archive state \(\K_t\). Define semantic novelty by a functional
\[
\nu_s(x_i,t)
=
\mathcal{N}_s\!\left(
R(x_i\mid\K_t)
\right).
\]
The value may depend on formal, empirical, combinatorial, applicative, or obstruction-sensitive novelty, but it remains a property of the transformed individual submission relative to the archive.

Let
\[
Z_t
\]
denote a statistic of the submission population during interval \(t\), and let
\[
P_t
\]
be its current distribution. Let
\[
P_0
\]
be a reference distribution inferred from an earlier period or a declared generative model. Distributional novelty can be represented by
\[
\nu_d(t)
=
D(P_t,P_0),
\]
where \(D\) is a divergence, discrepancy, change statistic, or model-relative surprise measure.

\begin{definition}[Distributionally novel recurrence]
A collection of submissions arriving during interval \(I_t\) is distributionally novel at threshold \(\tau_d\) when
\[
D(P_t,P_0)>\tau_d,
\]
even if
\[
\nu_s(x_i,t)\leq\tau_s
\]
for every individual submission \(i\) arriving during \(I_t\).
\end{definition}

The possibility described by the definition is not exceptional. Suppose every submission asks a question already answered by node \(v\in\K_t\), giving
\[
R(x_i\mid\K_t)=\varnothing
\]
and hence
\[
\nu_s(x_i,t)=0.
\]
If the arrival rate at \(v\) changes from \(\lambda_0(v)\) to \(\lambda_1(v)\gg\lambda_0(v)\), the individual questions remain semantically resolved while their recurrence becomes evidence of a changed environment, failed explanation, external event, or newly salient application.

\begin{proposition}[Independence of semantic and distributional novelty]
Neither semantic novelty nor distributional novelty implies the other.
\end{proposition}

\begin{proof}
For semantic novelty without distributional novelty, let submissions be independent draws from a stationary distribution \(P_0\) assigning positive probability to a large family of individually novel residuals. Then some \(x_i\) satisfy
\[
\nu_s(x_i,t)>0,
\]
while the population distribution remains \(P_t=P_0\), giving
\[
\nu_d(t)=0.
\]

For distributional novelty without semantic novelty, let every arriving submission map to an archive-resolved node \(v\), so
\[
\nu_s(x_i,t)=0
\]
for all \(i\). Change only the arrival intensity from \(\lambda_0(v)\) to \(\lambda_1(v)\neq\lambda_0(v)\). Then the population law differs from its reference law and
\[
\nu_d(t)>0
\]
for every divergence that separates the two count distributions. Thus neither form implies the other.
\end{proof}

\subsection{Recurrence as evidence under a point-process model}

Fix a node \(v\) in the prerequisite or answer graph. Suppose submissions associated with \(v\) follow a Poisson process under the reference model:
\[
N_v(t+\Delta)-N_v(t)
\sim
\operatorname{Poisson}(\lambda_{0,v}\Delta).
\]
If \(n_v\) arrivals are observed during an interval of length \(\Delta\), the reference likelihood is
\[
L_0(n_v)
=
e^{-\lambda_{0,v}\Delta}
\frac{(\lambda_{0,v}\Delta)^{n_v}}{n_v!}.
\]
Under an alternative rate \(\lambda_{1,v}\), the likelihood ratio is
\[
\Lambda_v
=
\frac{
e^{-\lambda_{1,v}\Delta}
(\lambda_{1,v}\Delta)^{n_v}/n_v!
}{
e^{-\lambda_{0,v}\Delta}
(\lambda_{0,v}\Delta)^{n_v}/n_v!
}
=
\exp\!\left(
-(\lambda_{1,v}-\lambda_{0,v})\Delta
\right)
\left(
\frac{\lambda_{1,v}}{\lambda_{0,v}}
\right)^{n_v}.
\]
The log-likelihood ratio is
\[
\log\Lambda_v
=
n_v
\log\!\left(
\frac{\lambda_{1,v}}{\lambda_{0,v}}
\right)
-
(\lambda_{1,v}-\lambda_{0,v})\Delta.
\]
The alternative is favored when
\[
n_v
>
\frac{
(\lambda_{1,v}-\lambda_{0,v})\Delta
}{
\log(\lambda_{1,v}/\lambda_{0,v})
}.
\]

When \(\lambda_{1,v}\) is unknown, its maximum-likelihood estimate is
\[
\widehat{\lambda}_{1,v}
=
\frac{n_v}{\Delta}.
\]
Substitution gives the generalized log-likelihood ratio
\[
G_v
=
n_v
\log\!\left(
\frac{n_v}{\lambda_{0,v}\Delta}
\right)
-
n_v
+
\lambda_{0,v}\Delta,
\]
with the convention that the first term is zero when \(n_v=0\). This statistic measures rate surprise without requiring the alternative intensity to be specified in advance.

The same construction extends to a nonhomogeneous reference intensity \(\lambda_{0,v}(t)\). If arrival times in interval \(I=[a,b]\) are
\[
t_1,\ldots,t_n,
\]
then the point-process log likelihood is
\[
\ell(\lambda)
=
\sum_{j=1}^{n}\log\lambda(t_j)
-
\int_a^b\lambda(u)\,du.
\]
A recurrence detector compares a candidate intensity \(\lambda_1\) with \(\lambda_0\):
\[
\log\Lambda
=
\sum_{j=1}^{n}
\log
\frac{\lambda_1(t_j)}{\lambda_0(t_j)}
-
\int_a^b
\left[
\lambda_1(u)-\lambda_0(u)
\right]du.
\]
This prevents predictable daily, seasonal, publication-driven, or event-driven variation from being mistaken for novelty merely because raw counts are high.

\subsection{Bayesian recurrence inference}

Let the unknown arrival rate at node \(v\) have Gamma prior
\[
\lambda_v
\sim
\operatorname{Gamma}(\alpha_0,\beta_0),
\]
with density
\[
p(\lambda_v)
=
\frac{\beta_0^{\alpha_0}}{\Gamma(\alpha_0)}
\lambda_v^{\alpha_0-1}
e^{-\beta_0\lambda_v}.
\]
After observing \(n_v\) arrivals during exposure \(\Delta\), conjugacy yields
\[
\lambda_v\mid n_v
\sim
\operatorname{Gamma}
\left(
\alpha_0+n_v,
\beta_0+\Delta
\right).
\]
The posterior mean is
\[
\E[\lambda_v\mid n_v]
=
\frac{\alpha_0+n_v}{\beta_0+\Delta},
\]
and the posterior variance is
\[
\Var(\lambda_v\mid n_v)
=
\frac{\alpha_0+n_v}{(\beta_0+\Delta)^2}.
\]
The probability of an increase above a consequential threshold \(\lambda_v^{*}\) is
\[
\Prob(\lambda_v>\lambda_v^{*}\mid n_v)
=
1-
F_{\operatorname{Gamma}}
\left(
\lambda_v^{*};
\alpha_0+n_v,
\beta_0+\Delta
\right).
\]
A distributional escalation can be triggered when
\[
\Prob(\lambda_v>\lambda_v^{*}\mid n_v)
>
1-\delta.
\]

The threshold \(\lambda_v^{*}\) should be tied to an operational consequence. If each arrival at node \(v\) consumes expected automated cost \(c_v^{A}\), then the automated workload is
\[
W_v
=
\lambda_v c_v^{A}.
\]
If the available automated capacity for that region is \(A_v\), the critical rate is
\[
\lambda_v^{*}
=
\frac{A_v}{c_v^{A}}.
\]
The posterior instability probability becomes
\[
\Prob(
\lambda_v c_v^{A}>A_v
\mid n_v
).
\]
Recurrence is then escalated not because it is merely surprising but because the inferred rate threatens the stability or adequacy of the current answer node.

\subsection{Recurrence can indict the archive rather than the participant}

Let \(v\) be an answer node with canonical explanation \(e_v\). Suppose \(n_v\) participants are routed through \(v\), and let
\[
F_v
=
\frac{
\text{participants who remain unresolved after }e_v
}{
\text{participants exposed to }e_v
}
\]
be the observed post-explanation failure rate. A high \(F_v\) admits several hypotheses. The participants may lack a prerequisite. The explanation may be poor. The graph may omit an intermediate distinction. The retrieval rule may be routing heterogeneous questions into one node. The archive may incorrectly claim that \(e_v\) resolves them.

Let
\[
H_P,\quad H_E,\quad H_G,\quad H_R,\quad H_K
\]
denote participant-deficiency, explanation-deficiency, graph-deficiency, routing-deficiency, and archive-error hypotheses. Given recurrence evidence \(\mathcal{D}_v\),
\[
\Prob(H_j\mid\mathcal{D}_v)
\propto
\Prob(\mathcal{D}_v\mid H_j)\Prob(H_j).
\]
The system should not assign all repeated failures to \(H_P\) by default. Such a policy makes the prerequisite graph unfalsifiable because every failure of the graph is reclassified as evidence about the participant.

Suppose the node has participant-specific competency covariates \(C_i\), route variant \(P_i\), and explanation version \(E_i\). A logistic diagnostic model may take the form
\[
\log
\frac{
\Prob(Y_i=1)
}{
1-\Prob(Y_i=1)
}
=
\beta_0
+
\beta_C^{\top}C_i
+
\beta_P^{\top}P_i
+
\beta_E^{\top}E_i
+
\beta_{PE}^{\top}(P_i\otimes E_i),
\]
where \(Y_i=1\) indicates successful resolution. If explanation effects remain large after competency adjustment, recurrence supplies evidence about the archive rather than merely the participant population.

Randomized route variants provide stronger evidence. Let participants judged admissible for either explanation be assigned
\[
A_i\in\{0,1\},
\]
where \(A_i=0\) receives the current explanation and \(A_i=1\) receives a candidate repair. The average explanation effect is
\[
\tau_E
=
\E[Y_i(1)-Y_i(0)].
\]
Under random assignment,
\[
\widehat{\tau}_E
=
\overline{Y}_{A=1}
-
\overline{Y}_{A=0}
\]
is unbiased for \(\tau_E\). A positive value shows that some recurrence previously attributed to participant failure was reducible by changing the archive.

\subsection{Prerequisite-edge failure as a graph-valued signal}

Let
\[
G_t=(V_t,E_t)
\]
be the prerequisite graph. For edge
\[
e=(u\to v),
\]
let
\[
A_e(t)
\]
be the number of traversal attempts during period \(t\), and let
\[
F_e(t)
\]
be the number of failures. The empirical edge-failure rate is
\[
\widehat p_e(t)
=
\frac{F_e(t)}{A_e(t)}
\]
when \(A_e(t)>0\).

To avoid treating low-exposure edges as highly certain, assign
\[
p_e
\sim
\operatorname{Beta}(\alpha_e,\beta_e).
\]
The posterior is
\[
p_e\mid F_e,A_e
\sim
\operatorname{Beta}
\left(
\alpha_e+F_e,
\beta_e+A_e-F_e
\right).
\]
For an acceptable failure threshold \(\tau_e\), define the edge-alarm probability
\[
\zeta_e
=
\Prob(p_e>\tau_e\mid F_e,A_e).
\]
The edge enters structural audit when
\[
\zeta_e>1-\delta_e.
\]

A high failure rate does not uniquely identify the repair. Let \(w\) be a candidate intermediate node and consider replacing
\[
u\to v
\]
with
\[
u\to w\to v.
\]
Let the original expected traversal burden be
\[
C_{\mathrm{old}}
=
c_{uv}
+
p_{uv}c_{\mathrm{fail}},
\]
and the repaired burden be
\[
C_{\mathrm{new}}
=
c_{uw}
+
c_{wv}
+
p_{uw}c_{\mathrm{fail},1}
+
(1-p_{uw})p_{wv}c_{\mathrm{fail},2}.
\]
The insertion is burden-reducing when
\[
C_{\mathrm{new}}<C_{\mathrm{old}}.
\]
A longer route can therefore be optimal when it sufficiently decreases the probability or cost of failure.

\begin{proposition}[Recurrence-driven edge repair]
Suppose inserting \(w\) changes successful traversal probability from
\[
1-p_{uv}
\]
to
\[
(1-p_{uw})(1-p_{wv}).
\]
The insertion improves success probability exactly when
\[
p_{uv}
>
p_{uw}+p_{wv}-p_{uw}p_{wv}.
\]
\end{proposition}

\begin{proof}
The repaired route has higher success probability when
\[
(1-p_{uw})(1-p_{wv})
>
1-p_{uv}.
\]
Expanding the left side gives
\[
1-p_{uw}-p_{wv}+p_{uw}p_{wv}
>
1-p_{uv}.
\]
Rearranging yields
\[
p_{uv}
>
p_{uw}+p_{wv}-p_{uw}p_{wv}.
\]
\end{proof}

The proposition does not say that every high-failure edge should be subdivided. The intermediate node must reduce combined failure enough to compensate for added burden and dropout risk. Recurrence identifies where the topology may be wrong. It does not determine the replacement topology by itself.

\subsection{Collective residual formation}

Individual residual extraction produces sets
\[
R_i=R(x_i\mid\K_t).
\]
Let
\[
\mathcal{R}_t
=
\{R_i:t_i\in I_t\}
\]
be the residual family during interval \(I_t\). A collective residual is not generally the union
\[
\bigcup_iR_i,
\]
because the union retains claims without representing their frequency, dependence, complementarity, or mutual obstruction.

Let \(\mathcal{C}_t\) be a simplicial complex whose vertices are normalized claims and whose simplex
\[
\{q_{i_0},\ldots,q_{i_k}\}
\]
is included when those claims co-occur, mutually support one another, or participate in a jointly admissible assembly. Let \(w(\sigma)\) be the number or weighted frequency of independent submissions supporting simplex \(\sigma\).

A collective residual may be represented as
\[
R_{\mathrm{coll}}(I_t)
=
\left(
\mathcal{C}_t,
w_t,
\Delta w_t,
\Omega_t
\right),
\]
where \(\Delta w_t\) records changes from the reference period and \(\Omega_t\) records collective gluing obstructions. Two periods can contain the same claim vertices while differing in their higher-order topology. An increase in co-occurrence among previously independent claims can therefore be distributionally novel even when no new claim appears.

Let \(A_t\) be a claim co-occurrence matrix:
\[
(A_t)_{jk}
=
\sum_{i:t_i\in I_t}
\ind\{q_j,q_k\in Q(x_i)\}.
\]
A structural change can be measured by
\[
D_A(t)
=
\lVert
\widetilde A_t-\widetilde A_0
\rVert_F,
\]
where \(\widetilde A\) is normalized for total submission volume. Normalization is necessary because raw co-occurrence counts increase automatically when the arrival rate rises.

For spectral comparison, let
\[
L_t=D_t-\widetilde A_t
\]
be the graph Laplacian. Changes in its eigenvalues can reveal altered cluster structure:
\[
D_{\mathrm{spec}}(t)
=
\sum_{j=1}^{m}
w_j
\left(
\lambda_j(L_t)-\lambda_j(L_0)
\right)^2.
\]
A new low-conductance cluster of mutually recurring claims can appear before any one claim becomes individually prominent.

\subsection{Independence and evidential multiplicity}

Ten repetitions from one copied source do not provide the same evidence as ten independent rediscoveries. Let submissions \(x_1,\ldots,x_n\) support normalized claim \(q\). Let
\[
D_i\in\{1,\ldots,K\}
\]
identify dependence clusters derived from shared sources, near-duplicate language, common model outputs, citation ancestry, coordinated campaigns, or direct copying.

If cluster \(k\) contributes \(n_k\) submissions, a crude effective sample size is
\[
n_{\mathrm{eff}}
=
\frac{
\left(\sum_{k=1}^{K}n_k\right)^2
}{
\sum_{k=1}^{K}n_k^2
}.
\]
If every submission is independent so that \(n_k=1\), then
\[
n_{\mathrm{eff}}=n.
\]
If all submissions belong to one dependence cluster, then
\[
n_{\mathrm{eff}}=1.
\]

For correlated binary support indicators \(Y_i\) with common pairwise correlation \(\rho\), the variance of the sample mean is
\[
\Var(\overline Y)
=
\frac{\sigma^2}{n}
\left[
1+(n-1)\rho
\right].
\]
The corresponding effective sample size is
\[
n_{\mathrm{eff}}
=
\frac{n}{1+(n-1)\rho}.
\]
As \(n\to\infty\),
\[
n_{\mathrm{eff}}
\to
\frac{1}{\rho}
\]
for fixed \(\rho>0\). Arbitrarily many correlated repetitions therefore provide only bounded independent evidence.

\begin{definition}[Independence-weighted recurrence]
For recurrence cluster sizes \(n_1,\ldots,n_K\), define
\[
\mathcal{R}_{\ind}
=
\sum_{k=1}^{K}
\varphi(n_k),
\]
where \(\varphi\) is nondecreasing, concave, and satisfies \(\varphi(0)=0\). The choice
\[
\varphi(n)=1-e^{-\alpha n}
\]
gives rapidly diminishing credit to repetitions within one dependence cluster while preserving approximately additive credit across independent clusters.
\end{definition}

Independence cannot be inferred solely from linguistic dissimilarity. Two differently worded submissions may descend from the same source, while two nearly identical formulations may have been independently derived from a highly constrained problem. Dependence estimation must therefore remain probabilistic and provenance-sensitive.

Let \(G_S\) be a source-dependency graph over submissions. A path
\[
x_i\leadsto x_j
\]
indicates possible informational ancestry. The recurrence packet should preserve both raw frequency
\[
n(q)
\]
and effective independent frequency
\[
n_{\mathrm{eff}}(q).
\]
The difference
\[
n(q)-n_{\mathrm{eff}}(q)
\]
is itself informative. It distinguishes broad independent recurrence from concentrated replication or amplification.

\subsection{Burst detection and temporal concentration}

A fixed-window count can miss bursts whose boundaries do not align with the chosen intervals. Let event times associated with topic \(v\) be
\[
t_1<t_2<\cdots<t_n.
\]
Define interarrival times
\[
\Delta_i=t_i-t_{i-1}.
\]
Under a homogeneous Poisson reference with rate \(\lambda_0\),
\[
\Delta_i\sim\operatorname{Exponential}(\lambda_0).
\]
A run of unusually short interarrival times supplies burst evidence.

A simple exponentially weighted count satisfies
\[
Z_v(t)
=
\int_{(-\infty,t]}
e^{-\kappa(t-s)}\,dN_v(s),
\]
with differential form
\[
dZ_v(t)
=
-\kappa Z_v(t)\,dt+dN_v(t).
\]
Under constant reference intensity \(\lambda_0\),
\[
\E[Z_v(t)]
\to
\frac{\lambda_0}{\kappa}
\]
in stationarity. Define standardized burst score
\[
B_v(t)
=
\frac{
Z_v(t)-\lambda_0/\kappa
}{
\sqrt{\lambda_0/(2\kappa)}
},
\]
using the stationary variance of exponentially filtered Poisson shot noise. A burst alarm occurs when
\[
B_v(t)>\tau_B.
\]

Burst significance must be corrected for simultaneous monitoring of many nodes. If \(m\) independent nodes are tested at marginal level \(\alpha\), the probability of at least one false alarm is
\[
1-(1-\alpha)^m.
\]
The Bonferroni condition
\[
\alpha_{\mathrm{node}}
=
\frac{\alpha_{\mathrm{family}}}{m}
\]
controls familywise error conservatively. False-discovery-rate procedures can retain greater power when many nodes are monitored and some true changes are expected.

\subsection{Distributional novelty as divergence}

Suppose transformed submissions during reference and current periods have densities \(p_0(z)\) and \(p_t(z)\). The Kullback-Leibler divergence is
\[
D_{\KL}(P_t\Vert P_0)
=
\int
p_t(z)
\log
\frac{p_t(z)}{p_0(z)}
\,dz.
\]
It is infinite if \(P_t\) assigns positive mass to a region absent from the support of \(P_0\). Such support expansion is particularly important because it represents a genuinely new representational region rather than a reweighting of familiar regions.

Because \(D_{\KL}\) is asymmetric, one may use Jensen-Shannon divergence:
\[
D_{\mathrm{JS}}(P_t,P_0)
=
\frac{1}{2}
D_{\KL}(P_t\Vert M)
+
\frac{1}{2}
D_{\KL}(P_0\Vert M),
\]
where
\[
M=\frac{1}{2}(P_t+P_0).
\]
This quantity is finite and symmetric.

For kernel \(k\), the squared maximum mean discrepancy is
\[
\operatorname{MMD}^2(P_t,P_0)
=
\E[k(X,X')]
+
\E[k(Y,Y')]
-
2\E[k(X,Y)],
\]
where
\[
X,X'\sim P_t,
\qquad
Y,Y'\sim P_0.
\]
With a characteristic kernel,
\[
\operatorname{MMD}(P_t,P_0)=0
\]
if and only if
\[
P_t=P_0.
\]
MMD permits distributional change detection without explicit density estimation, which is useful in high-dimensional residual spaces.

The representation \(q\) constrains which distributional changes can be detected. If
\[
q_{\#}P_t=q_{\#}P_0,
\]
where \(q_{\#}\) denotes pushforward measure, then every detector operating solely on \(q(X)\) sees no change even if
\[
P_t\neq P_0.
\]

\begin{theorem}[Distributional invisibility under projection]
Let \(q:\mathcal{X}\to\mathcal{Z}\) be measurable. If
\[
q_{\#}P_1=q_{\#}P_0,
\]
then every measurable statistic
\[
T=T(q(X_1),\ldots,q(X_n))
\]
has the same distribution under \(P_1\) as under \(P_0\). Consequently no test based only on the projected sample can distinguish \(P_1\) from \(P_0\) with power exceeding its size.
\end{theorem}

\begin{proof}
Under \(P_j\), each projected observation \(q(X_i)\) has distribution \(q_{\#}P_j\). Equality of the pushforward measures implies equality of the joint projected-sample distributions under independent sampling. Every measurable function \(T\) of the projected sample therefore has the same distribution under both hypotheses. For any rejection region \(A\),
\[
\Prob_{P_1}(T\in A)
=
\Prob_{P_0}(T\in A).
\]
The power equals the test size.
\end{proof}

This is the population-level counterpart of CLIO insufficiency. A quotient can erase not only an individual task-relevant distinction but an entire distributional shift. Audit must therefore retain selected pre-projection statistics or randomized access to original submissions.

\subsection{Quota-limited observation and apparent recurrence}

Distributional novelty can be estimated only through an observation channel whose sampling geometry is understood. Let an external message stream emit relevant events according to a counting process \(N(t)\), while an ingestion interface permits at most \(Q_{\max}\) quota units during an interval of duration \(T_Q\). Suppose further that the interface provides no push channel and imposes a minimum polling interval
\[
\Delta_{\min}>0.
\]
If a basic request costs one quota unit, the maximum floor-limited polling frequency is
\[
\nu_{\mathrm{floor}}
=
\frac{1}{\Delta_{\min}},
\]
while the quota-limited sustainable frequency is
\[
\nu_Q
=
\frac{Q_{\max}}{T_Q}.
\]
The admissible long-run polling rate satisfies
\[
\nu_{\mathrm{adm}}
\leq
\min\left\{
\frac{1}{\Delta_{\min}},
\frac{Q_{\max}}{T_Q}
\right\}.
\]

Polling continuously at the interface floor exhausts a unit-cost quota after
\[
T_{\mathrm{exhaust}}
=
Q_{\max}\Delta_{\min}.
\]
Whenever
\[
Q_{\max}\Delta_{\min}<T_Q,
\]
continuous minimum-interval polling is impossible even though each individual request is locally admissible. The observation policy is locally valid but globally unstable.

For a representative quota
\[
Q_{\max}=10^4
\]
and polling floor
\[
\Delta_{\min}\approx5\ \mathrm{s},
\]
unit-cost exhaustion occurs after
\[
T_{\mathrm{exhaust}}
\approx
5\times10^4\ \mathrm{s}
\approx
13.9\ \mathrm{h}.
\]
If one effective poll consumes approximately six quota units, exhaustion instead occurs after
\[
T_{\mathrm{exhaust}}
\approx
\frac{10^4}{6}\cdot5\ \mathrm{s}
\approx
8.33\times10^3\ \mathrm{s}
\approx
2.31\ \mathrm{h}.
\]
The structural result does not depend upon the identity of the service. A request-count quota, a nonzero polling floor, and the absence of a push channel jointly constrain which recurrence patterns can be observed.

A polling-frequency constraint and a bounded response batch are distinct restrictions. The first determines which emitted events remain recoverable when an observation occurs, while the second limits how many recoverable events can be returned by that observation.

Let event times during polling interval
\[
I_k=[t_k,t_{k+1})
\]
be
\[
T_{k,1},\ldots,T_{k,X_k},
\]
where
\[
X_k=N(t_{k+1})-N(t_k).
\]
Let
\[
\sigma(a)\in[0,1]
\]
be the probability that an event remains recoverable at age \(a\geq0\). Assume
\[
\sigma(0)=1
\]
and that \(\sigma\) is nonincreasing. Define
\[
I_{k,j}
=
\ind\!\left\{
\text{event }j\text{ remains recoverable at }t_{k+1}
\right\}.
\]
Conditional on its emission time,
\[
\Prob(I_{k,j}=1\mid T_{k,j})
=
\sigma(t_{k+1}-T_{k,j}).
\]
The recoverable count is
\[
R_k
=
\sum_{j=1}^{X_k}I_{k,j},
\]
and temporal-sampling loss is
\[
S_k
=
X_k-R_k.
\]
Its conditional expectation is
\[
\E[S_k\mid T_{k,1},\ldots,T_{k,X_k}]
=
\sum_{j=1}^{X_k}
\left[
1-\sigma(t_{k+1}-T_{k,j})
\right].
\]

A hard retention window \(\tau>0\) is represented by
\[
\sigma_{\tau}(a)
=
\ind\{a\leq\tau\}.
\]
Then
\[
R_k
=
\sum_{j=1}^{X_k}
\ind\{t_{k+1}-T_{k,j}\leq\tau\}.
\]
If
\[
\Delta_k=t_{k+1}-t_k\leq\tau,
\]
then
\[
R_k=X_k,
\qquad
S_k=0.
\]
If
\[
\Delta_k>\tau,
\]
events emitted before \(t_{k+1}-\tau\) are lost:
\[
S_k
=
N(t_{k+1}-\tau)-N(t_k).
\]

Stochastic expiration with constant hazard \(\gamma>0\) is represented by
\[
\sigma_{\gamma}(a)=e^{-\gamma a}.
\]
If the arrival process is conditionally Poisson with intensity \(\lambda(u)\), then
\[
\E[R_k]
=
\int_{t_k}^{t_{k+1}}
\lambda(u)
\sigma(t_{k+1}-u)\,du,
\]
and
\[
\E[S_k]
=
\int_{t_k}^{t_{k+1}}
\lambda(u)
\left[
1-\sigma(t_{k+1}-u)
\right]du.
\]
For constant intensity \(\lambda\),
\[
\E[R_k]
=
\frac{\lambda}{\gamma}
\left(
1-e^{-\gamma\Delta_k}
\right),
\]
while
\[
\E[S_k]
=
\lambda\Delta_k
-
\frac{\lambda}{\gamma}
\left(
1-e^{-\gamma\Delta_k}
\right).
\]
Since
\[
1-e^{-x}<x
\]
for every \(x>0\),
\[
\E[S_k]>0
\]
for every \(\gamma>0\) and \(\Delta_k>0\). At the degenerate boundary \(\Delta_k=0\),
\[
\E[S_k]=0.
\]
The expansion
\[
1-e^{-\gamma\Delta_k}
=
\gamma\Delta_k
-
\frac{\gamma^2\Delta_k^2}{2}
+
\frac{\gamma^3\Delta_k^3}{6}
+
O(\Delta_k^4)
\]
gives
\[
\E[S_k]
=
\frac{\lambda\gamma}{2}\Delta_k^2
-
\frac{\lambda\gamma^2}{6}\Delta_k^3
+
O(\Delta_k^4).
\]
Hence
\[
\E[S_k]
=
\frac{\lambda\gamma}{2}\Delta_k^2
+
O(\Delta_k^3)
\qquad
\text{as }\Delta_k\downarrow0.
\]

If the interface returns at most \(B_k\) recoverable events, then
\[
Y_k=\min\{R_k,B_k\}.
\]
Batch-truncation loss is
\[
H_k
=
R_k-Y_k
=
(R_k-B_k)_{+},
\]
and total interval loss is
\[
\Lambda_k
=
X_k-Y_k
=
S_k+H_k.
\]
Over \(m\) intervals,
\[
R_{\mathrm{true}}(m)
=
\sum_{k=0}^{m-1}X_k,
\]
\[
R_{\mathrm{recoverable}}(m)
=
\sum_{k=0}^{m-1}R_k,
\]
and
\[
R_{\mathrm{obs}}(m)
=
\sum_{k=0}^{m-1}Y_k.
\]
Therefore
\[
L_{\mathrm{time}}(m)
=
\sum_{k=0}^{m-1}S_k,
\]
\[
L_{\mathrm{batch}}(m)
=
\sum_{k=0}^{m-1}(R_k-B_k)_{+},
\]
and
\[
L_{\mathrm{total}}(m)
=
R_{\mathrm{true}}(m)
-
R_{\mathrm{obs}}(m)
=
L_{\mathrm{time}}(m)
+
L_{\mathrm{batch}}(m).
\]

Let
\[
c:\mathbb{N}_{>0}\to\mathbb{R}_{>0}
\]
be the quota cost of a request with batch capacity \(B\). Assume \(c\) is nondecreasing:
\[
B_1\leq B_2
\Longrightarrow
c(B_1)\leq c(B_2).
\]
One may additionally assume convexity on a continuous relaxation:
\[
c\!\left(
\alpha B_1+(1-\alpha)B_2
\right)
\leq
\alpha c(B_1)+(1-\alpha)c(B_2)
\]
for every \(\alpha\in[0,1]\). This is distinct from diminishing average cost, which requires
\[
\frac{c(B_2)}{B_2}
\leq
\frac{c(B_1)}{B_1}
\qquad
\text{whenever }B_2\geq B_1.
\]
Monotonicity alone is sufficient for the observation-loss argument.

Define the admissible schedule class
\[
\mathfrak{S}(Q_{\max},\Delta_{\min},T_Q)
=
\left\{
\mathcal{S}
=
\{(t_k,B_k)\}_{k=0}^{m-1}
:
\sum_{k=0}^{m-1}c(B_k)\leq Q_{\max},
\quad
\Delta_k\geq\Delta_{\min},
\quad
t_m\leq T_Q
\right\}.
\]
The loss-minimizing schedule satisfies
\[
\mathcal{S}^{*}
\in
\argmin_{
\mathcal{S}\in
\mathfrak{S}(Q_{\max},\Delta_{\min},T_Q)
}
\E\!\left[
L_{\mathrm{time}}(\mathcal{S})
+
L_{\mathrm{batch}}(\mathcal{S})
\right].
\]

Under hard retention, zero temporal loss requires
\[
\Delta_{\min}\leq\tau.
\]
Covering the full quota interval requires at least
\[
m_{\tau}
=
\left\lceil
\frac{T_Q}{\tau}
\right\rceil
\]
polls. If
\[
c_{\min}
=
\inf_{B\geq1}c(B),
\]
another necessary condition is
\[
m_{\tau}c_{\min}\leq Q_{\max}.
\]
When either
\[
\Delta_{\min}>\tau
\]
or
\[
\left\lceil
\frac{T_Q}{\tau}
\right\rceil c_{\min}>Q_{\max},
\]
some temporal loss is structurally unavoidable under continuous arrivals.

Under exponential recoverability and constant polling interval
\[
\Delta_k=\Delta_{\min},
\]
expected temporal loss per interval is exactly
\[
\E[S_k]
=
\lambda\Delta_{\min}
-
\frac{\lambda}{\gamma}
\left(
1-e^{-\gamma\Delta_{\min}}
\right)
>
0.
\]
For small \(\Delta_{\min}\),
\[
\E[S_k]
=
\frac{\lambda\gamma}{2}
\Delta_{\min}^2
+
O(\Delta_{\min}^3).
\]
Thus the positive polling floor converts a loss approaching zero in the continuous-observation limit into a strictly positive per-interval loss. This is the exponential-recoverability counterpart of the hard-window condition
\[
\Delta_{\min}>\tau.
\]

Absence of observed recurrence therefore does not entail absence of underlying recurrence unless retention, polling, batching, and quota constraints jointly establish
\[
L_{\mathrm{total}}=0
\]
over the relevant arrival regime.

\subsection{Debiasing recurrence estimates under observation loss}

Suppose each emitted event of type \(z\) is observed independently with probability
\[
\pi_{\mathrm{obs}}(z,t)>0.
\]
Let \(O_i\in\{0,1\}\) indicate observation. The true count in region \(B\) is
\[
N_B
=
\sum_i\ind\{z_i\in B\},
\]
and the inverse-probability estimator is
\[
\widehat N_B
=
\sum_i
\frac{
O_i\ind\{z_i\in B\}
}{
\pi_{\mathrm{obs}}(z_i,t_i)
}.
\]
Then
\[
\E[\widehat N_B]=N_B.
\]
The result follows from
\[
\E\!\left[
\frac{O_i}{\pi_{\mathrm{obs}}(z_i,t_i)}
\given z_i,t_i
\right]
=1.
\]

Batch truncation generally violates independent observation because inclusion depends on queue position, recency, ranking, or interface ordering. Let
\[
\pi_i
=
\Prob(
i\text{ is returned}
\mid
\mathcal{F}_{t_{k+1}^{-}}
)
\]
be the conditional inclusion probability given the pre-poll event history. If \(\pi_i\) is known and positive, inverse-probability estimation remains possible. If the interface deterministically returns only the most recent \(B\) events, older excess events have
\[
\pi_i=0.
\]
Their contribution is not identifiable from returned events alone.

A lower bound follows without full identifiability. Since
\[
Y_k\leq R_k\leq X_k,
\]
one always has
\[
R_{\mathrm{obs}}(m)
\leq
R_{\mathrm{recoverable}}(m)
\leq
R_{\mathrm{true}}(m).
\]
If external metadata provide an upper bound
\[
X_k\leq U_k,
\]
then
\[
Y_k
\leq X_k\leq U_k.
\]
Interval-censored recurrence can be propagated into novelty estimates:
\[
D_{\min}
=
\inf_{
X_k\in[Y_k,U_k]
}
D(P_X,P_0),
\]
\[
D_{\max}
=
\sup_{
X_k\in[Y_k,U_k]
}
D(P_X,P_0).
\]
If
\[
D_{\min}>\tau_d,
\]
distributional novelty is established despite observation loss. If
\[
D_{\max}\leq\tau_d,
\]
novelty is ruled out within the uncertainty set. Otherwise the observation channel is insufficient to decide.

\subsection{Collective escalation packets}

A recurrent pattern should not be escalated as a pile of individually resolved messages. Let
\[
\mathcal{I}_q(t)
=
\left\{
i:t_i\in I_t,\ q\in Q(x_i)
\right\}
\]
be the submissions associated with recurring claim or question \(q\). Define the collective packet
\[
\Pi_q^{\mathrm{coll}}
=
\left(
q,
n_q,
n_{\mathrm{eff},q},
\widehat{\lambda}_{q,t},
\Delta\widehat{\lambda}_{q,t},
\mathcal{C}_{q,t},
\Omega_{q,t},
L_{\mathrm{obs},q},
\mathcal{H}_{q}
\right),
\]
where \(n_q\) is raw recurrence, \(n_{\mathrm{eff},q}\) is independence-adjusted recurrence, \(\widehat{\lambda}_{q,t}\) is current intensity, \(\Delta\widehat{\lambda}_{q,t}\) is the inferred change, \(\mathcal{C}_{q,t}\) records co-occurrence structure, \(\Omega_{q,t}\) records collective obstructions, \(L_{\mathrm{obs},q}\) records observation-loss uncertainty, and \(\mathcal{H}_q\) contains competing causal hypotheses for the recurrence.

Let
\[
V_q^{\mathrm{coll}}
\]
be the value of reviewing the collective pattern and let
\[
C_q^{\mathrm{coll}}
\]
be its attention cost. Reviewing the compressed collective packet is preferable to reviewing individual messages when
\[
C_q^{\mathrm{coll}}
<
\sum_{i\in\mathcal{I}_q(t)}C_{ir}
\]
and the compression preserves every distinction required to diagnose the recurrence.

The collective packet must retain representative variation. Let \(d\) be a distance over reconstructed submissions. Select representatives
\[
J_q^{*}
\in
\argmax_{
J\subseteq\mathcal{I}_q,\ |J|\leq k
}
\left[
\sum_{i\in J}w_i
+
\alpha
\sum_{\substack{i,j\in J\\i<j}}d(x_i,x_j)
\right].
\]
This prevents the canonical example from erasing minority formulations. A recurrence cluster can have a dominant familiar form and a small variant containing the actual residual novelty.

\subsection{From recurrence to structural repair}

Let the system architecture at time \(t\) be
\[
\Theta_t
=
\left(
q_t,
\K_t,
G_t,
\mathcal{H}_t,
S_t,
\pi_t
\right).
\]
A collective recurrence statistic
\[
Z_t
\]
induces posterior over failure hypotheses
\[
\Prob(H_j\mid Z_t,\Theta_t).
\]
A structural repair action
\[
a\in\mathfrak{A}_{\mathrm{repair}}
\]
changes the architecture:
\[
\Theta_{t+1}
=
\mathcal{U}_{\mathrm{coll}}(\Theta_t,Z_t,a).
\]
The expected net value of repair is
\[
U(a\mid Z_t)
=
\E\!\left[
\sum_{s=t+1}^{t+H}
\Delta V_s(a)
-
\Delta C_s(a)
\given
Z_t
\right]
-
C(a),
\]
where \(C(a)\) is immediate repair cost and the summation captures future reductions in burden, false suppression, and repeated labor.

A recurrence alarm should trigger repair only when
\[
\sup_{a\in\mathfrak{A}_{\mathrm{repair}}}
U(a\mid Z_t)>0.
\]
Otherwise the correct response may be continued observation. Statistical surprise alone does not establish that intervention is worthwhile.

\begin{proposition}[Amortized value of recurrence repair]
Suppose a repair costs \(C_0>0\), reduces expected per-arrival processing cost by \(\delta c>0\), and applies to a future arrival process with expected count \(\Lambda_H\) over horizon \(H\). Ignoring discounting and secondary effects, the repair has positive expected net value exactly when
\[
\Lambda_H>\frac{C_0}{\delta c}.
\]
\end{proposition}

\begin{proof}
Expected total savings are
\[
\Lambda_H\delta c.
\]
Net value is
\[
\Lambda_H\delta c-C_0.
\]
This is positive exactly when
\[
\Lambda_H>C_0/\delta c.
\]
\end{proof}

The same calculation applies to explanation repair, prerequisite insertion, quotient refinement, observation-channel improvement, and creation of a new answer node. Recurrence converts a one-time repair cost into an amortized investment.

\subsection{Distributional novelty and HYDRA coupling}

Collective recurrence can alter HYDRA state even when no individual residual escalates. Let aggregate perception, relevance, and memory states be
\[
P_t^{\mathrm{agg}},
\qquad
R_t^{\mathrm{agg}},
\qquad
M_t^{\mathrm{agg}}.
\]
Define
\[
P_t^{\mathrm{agg}}
=
\int
\phi_P(z)\,d\mathcal{N}_t(z),
\]
\[
R_t^{\mathrm{agg}}
=
\mathcal{R}\!\left(
P_t^{\mathrm{agg}},
\K_t,
\FieldPhi_t,
\FieldV_t,
\FieldS_t
\right),
\]
and
\[
M_{t+1}^{\mathrm{agg}}
=
\mathcal{M}\!\left(
M_t^{\mathrm{agg}},
P_t^{\mathrm{agg}},
R_t^{\mathrm{agg}}
\right).
\]
The global transition is
\[
Z_{t+1}^{\mathrm{agg}}
=
\mathcal{H}\!\left(
P_t^{\mathrm{agg}},
R_t^{\mathrm{agg}},
M_t^{\mathrm{agg}};
\FieldPhi_t,
\FieldV_t,
\FieldS_t
\right).
\]

If recurrence merely increments a counter stored beside the archive, it is co-located metadata. If the counter changes which nodes are retrieved or jointly addressed, it participates in joint addressing. If recurrence changes the admissible transition itself by modifying relevance fields, memory activation, route topology, or escalation policy, it is operationally coupled.

Let \(\mathcal{H}_0\) be the transition with recurrence signal suppressed and \(\mathcal{H}_1\) the transition with recurrence active. Define
\[
\Delta_{\mathrm{rec}}
=
\E\!\left[
\mathcal{L}_{\mathrm{task}}(\mathcal{H}_0)
-
\mathcal{L}_{\mathrm{task}}(\mathcal{H}_1)
\right].
\]
A nonzero difference establishes operational influence. Positive
\[
\Delta_{\mathrm{rec}}>0
\]
indicates task improvement under the tested loss. It does not establish universal benefit because recurrence-sensitive routing can amplify coordinated noise, fashionable misconceptions, or platform-induced sampling artifacts.

\subsection{Adversarial recurrence and endogenous popularity}

Let the system increase visibility or structural priority as a function of observed recurrence:
\[
a_{t+1}
=
g(N_{\mathrm{obs}}(t)),
\qquad
g'>0.
\]
Suppose visibility increases future arrival intensity:
\[
\lambda_{t+1}
=
\lambda_{\mathrm{base}}
+
\beta a_{t+1}.
\]
Then recurrence and visibility form a feedback loop:
\[
N_{\mathrm{obs}}(t)
\longrightarrow
a_{t+1}
\longrightarrow
\lambda_{t+1}
\longrightarrow
N_{\mathrm{obs}}(t+1).
\]
Observed popularity is partly endogenous.

Linearizing around equilibrium with
\[
a_{t+1}=\alpha N_t,
\qquad
\E[N_{t+1}\mid N_t]
=
\lambda_{\mathrm{base}}\Delta
+
\beta\alpha\Delta N_t,
\]
the expected dynamics are stable only if
\[
\beta\alpha\Delta<1.
\]
If
\[
\beta\alpha\Delta>1,
\]
small recurrence fluctuations can be amplified rather than damped.

\begin{proposition}[Self-amplifying recurrence]
Under the linear expected recurrence model
\[
m_{t+1}
=
b+\rho m_t,
\qquad
b>0,
\]
the expected recurrence converges to
\[
m^{*}=\frac{b}{1-\rho}
\]
when \(0\leq\rho<1\). It grows without a finite stable equilibrium when \(\rho\geq1\).
\end{proposition}

\begin{proof}
Iteration gives
\[
m_t
=
\rho^tm_0
+
b\sum_{j=0}^{t-1}\rho^j.
\]
For \(0\leq\rho<1\),
\[
\rho^t\to0
\]
and
\[
\sum_{j=0}^{t-1}\rho^j
\to
\frac{1}{1-\rho},
\]
so
\[
m_t\to b/(1-\rho).
\]
If \(\rho=1\), then
\[
m_t=m_0+bt
\]
diverges linearly. If \(\rho>1\), the geometric term diverges exponentially.
\end{proof}

Distributional novelty must therefore be estimated against the system's own intervention history. Let \(A_t\) denote exposure allocated by the system. The counterfactual baseline is not
\[
\lambda_t=\lambda_0,
\]
but
\[
\lambda_t
=
\lambda_0(t)
+
h(A_{1:t},X_t).
\]
A recurrence increase unexplained after conditioning on exposure is stronger evidence of exogenous change than a raw increase produced by the routing system itself.

\subsection{False collective novelty under shared generation}

Suppose a generative system emits \(n\) superficially independent submissions derived from one latent template \(Z\). Conditional on \(Z\),
\[
X_i\sim P(\cdot\mid Z),
\]
but marginally the submissions are correlated. A detector assuming independence computes log evidence
\[
\ell_{\mathrm{naive}}
=
\sum_{i=1}^{n}
\log
\frac{p_1(X_i)}{p_0(X_i)}.
\]
The correct joint evidence is
\[
\ell_{\mathrm{joint}}
=
\log
\frac{
\int
\prod_{i=1}^{n}p_1(X_i\mid z)\,dP_1(z)
}{
\int
\prod_{i=1}^{n}p_0(X_i\mid z)\,dP_0(z)
}.
\]
These quantities can differ arbitrarily. Repeated model-generated paraphrases can make \(\ell_{\mathrm{naive}}\) grow approximately linearly in \(n\) even though the effective independent evidence remains close to one latent draw.

Let
\[
H(X_{1:n})
\]
be joint entropy and
\[
\operatorname{TC}(X_{1:n})
=
\sum_{i=1}^{n}H(X_i)-H(X_{1:n})
\]
be total correlation. Large total correlation indicates that raw multiplicity overstates informational multiplicity. An independence-adjusted recurrence score may take the form
\[
R_{\mathrm{adj}}
=
R_{\mathrm{raw}}
-
\kappa\operatorname{TC}(X_{1:n}),
\]
or may use cluster-level aggregation before recurrence testing.

The correction must not erase genuine collective discovery merely because participants use the same technical vocabulary. Dependence evidence should be grounded in provenance, timing, source overlap, model fingerprints, or causal ancestry rather than superficial similarity alone.

\subsection{Collective recurrence as an exploration signal}

Let \(u_i\) be classification uncertainty for residual \(i\), and let
\[
r_i^{\mathrm{freq}}
\]
be its recurrence contribution. An exploration policy can assign priority
\[
a_i^{x}
=
\gamma u_i
+
\eta r_i^{\mathrm{freq}}
+
\delta d_i^{\mathrm{ind}}
+
\zeta o_i,
\]
where \(d_i^{\mathrm{ind}}\) rewards independence from existing recurrence clusters and \(o_i\) rewards obstruction value.

A familiar question with high recurrence but low uncertainty may deserve structural review rather than individual escalation. An unusual residual with low recurrence and high uncertainty may deserve random review. A repeated unusual residual supported by independent arrivals can deserve both. The channels should remain distinct:
\[
\text{individual escalation}
\neq
\text{collective structural escalation}
\neq
\text{audit selection}.
\]
Collapsing them would force every recurrent pattern into the recipient's personal queue or cause every individually weak residual to disappear from aggregate analysis.

Let
\[
z_i^{E},
\quad
z_q^{C},
\quad
z_i^{A}
\]
indicate individual escalation, collective escalation, and audit selection. The joint allocation problem is
\[
\begin{aligned}
\text{maximize}\quad&
\sum_i z_i^{E}v_{ir}^{E}
+
\sum_q z_q^{C}v_{qr}^{C}
+
\sum_i z_i^{A}v_i^{A}
\\
\text{subject to}\quad&
\sum_i z_i^{E}c_{ir}^{E}
+
\sum_q z_q^{C}c_{qr}^{C}
\leq A_r,
\\
&
\sum_i z_i^{A}c_i^{A}
\leq A_{\Audit},
\\
&
z_i^{E},z_q^{C},z_i^{A}\in\{0,1\}.
\end{aligned}
\]
A single submission may contribute to a collective packet without being individually escalated. Its provenance remains attached to the aggregate evidence, but its full conversational burden does not enter the recipient-facing queue.

\subsection{A safety bound under lossy recurrence observation}

Let \(C_t\) be the event that a consequential collective pattern exists during period \(t\), and let \(M_t\) be the event that the observation and recurrence system misses it. The collective false-suppression rate is
\[
\psi_{\mathrm{coll}}
=
\Prob(M_t=1\mid C_t=1).
\]
Decompose miss events into observation failure \(O_t\), projection failure \(P_t\), dependence-model failure \(D_t\), detector failure \(G_t\), and allocation failure \(A_t\). Then
\[
M_t
\subseteq
O_t\cup P_t\cup D_t\cup G_t\cup A_t.
\]
By the union bound,
\[
\psi_{\mathrm{coll}}
\leq
\Prob(O_t\mid C_t)
+
\Prob(P_t\mid C_t)
+
\Prob(D_t\mid C_t)
+
\Prob(G_t\mid C_t)
+
\Prob(A_t\mid C_t).
\]

\begin{theorem}[Sufficient collective-detection bound]
Suppose that, conditional on a consequential collective pattern,
\[
\Prob(O_t)\leq\eps_O,
\qquad
\Prob(P_t)\leq\eps_P,
\qquad
\Prob(D_t)\leq\eps_D,
\]
\[
\Prob(G_t)\leq\eps_G,
\qquad
\Prob(A_t)\leq\eps_A.
\]
Then
\[
\psi_{\mathrm{coll}}
\leq
\eps_O+\eps_P+\eps_D+\eps_G+\eps_A.
\]
\end{theorem}

\begin{proof}
The miss event is contained in the union of the five failure events. Conditional application of Boole's inequality yields
\[
\Prob(M_t\mid C_t)
\leq
\sum_{J\in\{O,P,D,G,A\}}
\Prob(J_t\mid C_t),
\]
and substitution of the assumed bounds gives the result.
\end{proof}

The bound can be loose because the failure events interact. Observation loss can cause projection estimates to fail, and projection collapse can corrupt dependence estimation. The decomposition remains useful because it prevents all missed recurrence from being assigned to the statistical detector. A perfect detector cannot recover events the ingestion channel never exposed or distinctions the quotient erased.

\subsection{The collective update equation}

The archive should evolve under both individually validated contributions and aggregate recurrence. Let
\[
\K_t
\]
be the archive state and let
\[
U_t^{\mathrm{ind}}
\]
be accepted individual updates. Let
\[
U_t^{\mathrm{coll}}
\]
be structural updates justified by collective evidence. Let
\[
U_t^{\mathrm{audit}}
\]
be repairs arising from audited suppressions. The archive transition is
\[
\K_{t+1}
=
\mathcal{K}\!\left(
\K_t,
U_t^{\mathrm{ind}},
U_t^{\mathrm{coll}},
U_t^{\mathrm{audit}}
\right).
\]
To preserve provenance, these update classes are not merged before storage.

A collective update may alter answer nodes, edge thresholds, route alternatives, observation schedules, recurrence baselines, or escalation policies. It should not silently convert frequency into truth. Let
\[
q
\]
be a recurrent claim. Frequency changes the priority of investigating \(q\), not its truth value:
\[
\Prob(q\mid n_q)
\neq
1
\]
merely because \(n_q\) is large. Under independent evidence items \(e_1,\ldots,e_n\),
\[
\frac{
\Prob(q\mid e_{1:n})
}{
\Prob(\neg q\mid e_{1:n})
}
=
\frac{\Prob(q)}{\Prob(\neg q)}
\prod_{i=1}^{n}
\frac{\Prob(e_i\mid q)}{\Prob(e_i\mid\neg q)}.
\]
If repeated messages are likely under both \(q\) and \(\neg q\), their likelihood ratios are near one and multiplicity provides little truth evidence. Recurrence can nevertheless provide strong evidence that the archive must address the claim.

This yields the separation
\[
\text{recurrence of }q
\Longrightarrow
\text{increased value of investigating or explaining }q,
\]
but not
\[
\text{recurrence of }q
\Longrightarrow
q.
\]

\subsection{Conclusion}

The individual residual and the collective pattern occupy different levels of the architecture. Individual semantic novelty asks what remains after one submission is compared with existing knowledge. Distributional novelty asks whether the population of submissions has changed in rate, composition, dependence, topology, or obstruction structure. The first can justify direct engagement with a particular participant. The second can justify repairing an answer, inserting a prerequisite, refining a quotient, changing an observation schedule, or escalating a compact aggregate packet.

The complete recurrence transformation is
\[
\left\{
x_i
\right\}_{i:t_i\in I_t}
\xrightarrow{T}
\left\{
r_i,p_i,u_i,\omega_i
\right\}
\xrightarrow{\mathcal{D}}
\left(
P_t,
\lambda_t,
\mathcal{C}_t,
n_{\mathrm{eff},t},
\Omega_t
\right)
\xrightarrow{\mathcal{G}}
R_{\mathrm{coll}}(I_t),
\]
where \(\mathcal{D}\) performs dependence-sensitive aggregation and \(\mathcal{G}\) extracts the collective residual. The observation channel sits before this transformation and can irreversibly alter its input:
\[
\mathcal{N}_{\mathrm{true}}
\xrightarrow{\mathcal{O}_{\Delta,B,Q,\sigma}}
\mathcal{N}_{\mathrm{obs}}
\xrightarrow{q}
q_{\#}\mathcal{N}_{\mathrm{obs}}.
\]
Loss can occur through temporal expiration, batch truncation, quota exhaustion, projection collapse, dependence misclassification, statistical detection failure, or allocation failure.

Representational simplicity at the collective level therefore requires more than compact aggregation. It requires preservation of rate, independence, minority variation, observation uncertainty, and structural obstruction. A recurrence packet that reports only the most common wording is not sufficient. It may erase the independent multiplicity that made the pattern evidentially significant, or erase the minority variant that contained the unresolved remainder.

The governing collective condition is
\[
\text{retain enough population structure to distinguish repeated noise, copied amplification, independent rediscovery, curricular failure, archive failure, and genuine domain change}.
\]
Under finite attention, individual repetition should usually be compressed. Its distributional structure should not be discarded with it. Repetition is often semantically empty, but the shape of repetition can be new.


\section{Governance Conditions and Impossibility Results}

The architecture developed thus far can reduce repeated labor, detect prerequisite deficiencies, preserve residual novelty, assemble distributed conversational state, expose gluing obstructions, allocate scarce attention, audit suppression, and extract collective recurrence. None of these capacities determines who is entitled to define the archive, choose the quotient, set the prerequisite graph, specify the admissible task family, assign recipients, price attention, interpret an obstruction, or decide when an appeal has supplied enough evidence. Those are governance decisions.

Governance enters the mathematics wherever a representational or allocative choice is treated as fixed. The projection
\[
q:\mathcal{X}\to\mathcal{Z}
\]
selects which distinctions remain visible. The prerequisite graph
\[
G=(V,E)
\]
selects which competencies are treated as dependencies. The archive
\[
\K
\]
selects which claims count as established, answered, subsumed, rejected, or unresolved. The escalation score
\[
S(x_i\mid r,\K_r,\thetaState_i,t)
\]
selects which forms of value justify access. The audit policy determines which rejected regions remain observable. Governance is not an external ethical supplement added after the technical system is complete. It is the process by which these mathematical objects are authorized, revised, contested, and constrained.

Let the governed attention compiler be represented by
\[
\mathfrak{C}
=
\left(
\mathcal{X},
q,
\K,
G,
\ThetaSpace,
\mathcal{H},
S,
\pi_E,
\pi_A,
\mathfrak{a},
\mathcal{P},
\mathcal{L}
\right),
\]
where \(q\) is the representational projection, \(\K\) is the archive, \(G\) is the prerequisite graph, \(\ThetaSpace\) is the participant-state space, \(\mathcal{H}\) is the HYDRA transition operator, \(S\) is the escalation valuation, \(\pi_E\) is the escalation policy, \(\pi_A\) is the audit policy, \(\mathfrak{a}\) is the appeal operator, \(\mathcal{P}\) is the provenance system, and \(\mathcal{L}\) is the policy ledger. A governance regime is a family of update rules
\[
\Gamma
=
\left(
\Gamma_q,
\Gamma_{\K},
\Gamma_G,
\Gamma_{\mathcal{H}},
\Gamma_S,
\Gamma_E,
\Gamma_A,
\Gamma_{\mathfrak{a}}
\right)
\]
specifying who may alter each component, under what evidence, through which procedure, with which record, and subject to which reversibility constraints.

\subsection{Concentrated transmission and concentrated convergence}

A one-to-many broadcast concentrates control over what a population receives. A many-to-one attention compiler concentrates control over what a recipient is required to encounter. Reversing the direction of message flow does not by itself decentralize authority.

Let \(\mathcal{I}\) be the set of participants and let \(r\) be a protected recipient. In a broadcast system, an editorial operator
\[
B_r:
\mathcal{M}_r
\to
\mathcal{Y}^{\mathcal{I}}
\]
maps messages available to speaker \(r\) into outputs delivered to many listeners. In an attention compiler, a convergence operator
\[
C_r:
\mathcal{X}^{\mathcal{I}}
\to
\mathcal{Q}_r
\]
maps many incoming submissions into a recipient-facing queue.

The broadcast operator controls inclusion on the output side. The convergence operator controls inclusion on the input side. Both can be represented as selective maps:
\[
B_r(m)=
\begin{cases}
\operatorname{transmit}(m),&b_r(m)=1,\\
\varnothing,&b_r(m)=0,
\end{cases}
\]
and
\[
C_r(x_i)=
\begin{cases}
\operatorname{escalate}(x_i),&e_r(x_i)=1,\\
\operatorname{mediate}(x_i),&e_r(x_i)=0.
\end{cases}
\]
The asymmetry lies in which side of the relation bears the selection cost. In broadcast, listeners cannot compel transmission. In convergence, participants cannot compel reception.

\begin{definition}[Convergence authority]
The convergence authority of recipient \(r\) is the set
\[
\mathfrak{G}_r
=
\left\{
q,\K_r,G_r,S_r,\pi_{E,r},\pi_{A,r},\mathfrak{a}_r
\right\}
\]
of representational, archival, curricular, allocative, audit, and appeal components whose values determine whether and in what form an incoming submission can reach \(r\).
\end{definition}

If the recipient directly or indirectly controls every element of \(\mathfrak{G}_r\), then the architecture is procedurally centralized even if participation at the outer boundary is universal. This does not establish that the system is illegitimate. Scarce personal attention supplies a valid reason for recipient control. It establishes that access to the outer conversation and influence over the inner routing criteria are different rights.

\subsection{Standing, authority, and the separation of governance roles}

Let
\[
\mathcal{R}_{\mathrm{gov}}
=
\{
\mathrm{author},
\mathrm{recipient},
\mathrm{participant},
\mathrm{curator},
\mathrm{auditor},
\mathrm{reviewer},
\mathrm{operator}
\}
\]
be governance roles. A single person may occupy several roles, but the permissions associated with them should remain logically distinct.

Let
\[
\mathrm{Perm}(\rho,o)
\in
\{
\mathrm{read},
\mathrm{propose},
\mathrm{challenge},
\mathrm{modify},
\mathrm{approve},
\mathrm{reverse}
\}
\]
denote the permission held by role \(\rho\) over governed object \(o\). A permission system is a relation
\[
\mathfrak{P}
\subseteq
\mathcal{R}_{\mathrm{gov}}
\times
\mathcal{O}_{\mathrm{gov}}
\times
\mathcal{A}_{\mathrm{gov}},
\]
where \(\mathcal{O}_{\mathrm{gov}}\) contains the quotient, archive, graph, models, thresholds, audit schedules, and appeal records.

A minimal separation condition requires that no actor be the sole authority to perform all of the following operations on one disputed decision:
\[
\operatorname{construct},
\qquad
\operatorname{decide},
\qquad
\operatorname{audit},
\qquad
\operatorname{deny\ appeal}.
\]
Formally, if \(a(o)\) denotes the set of actors authorized for operation \(a\) on object \(o\), require
\[
\left|
\operatorname{construct}(o)
\cup
\operatorname{decide}(o)
\cup
\operatorname{audit}(o)
\cup
\operatorname{deny}(o)
\right|
\geq2,
\]
together with the stronger nonidentity condition
\[
\operatorname{decide}(o)
\not\supseteq
\operatorname{audit}(o)
\]
for policy-critical decisions.

This need not require an external institution for every personal attention compiler. It requires that an audit not be defined as the same inference repeated by the same model under the same representation. An independently parameterized review path, a randomized human sample, an alternative quotient, or a participant-supplied reconstruction can instantiate functional separation.

\begin{definition}[Functional independence]
Two review procedures
\[
\mathcal{R}_1,\mathcal{R}_2
\]
are functionally independent relative to failure family \(\mathfrak{F}\) when there exists no \(F\in\mathfrak{F}\) such that
\[
F
\Longrightarrow
\left(
\mathcal{R}_1\text{ fails}
\iff
\mathcal{R}_2\text{ fails}
\right)
\]
with probability one.
\end{definition}

Different model instances with identical training, archive access, prompts, and quotient structure may be numerically distinct while remaining functionally dependent. Governance requires independence with respect to plausible failure mechanisms, not merely multiplicity of components.

\subsection{Visible reasons and the right to locate disagreement}

Let a routing decision be
\[
d_i
=
D\!\left(
x_i,
q_i,
\K_i,
G_i,
\thetaState_i,
\omega_i,
r,
t
\right).
\]
A participant cannot challenge \(d_i\) effectively unless the system reveals which premises and transformations were decision-relevant.

Define a routing explanation
\[
\mathcal{E}_i
=
\left(
\widehat x_i,
K_i^{\mathrm{ret}},
M_i,
q_i,
P_i,
R_i,
\omega_i,
s_{ir},
\tau_{ir},
d_i
\right),
\]
where \(\widehat x_i\) is the charitable reconstruction, \(K_i^{\mathrm{ret}}\) is retrieved archive support, \(M_i\) is the missing-prerequisite set, \(q_i\) identifies the applied quotient, \(P_i\) is the selected route, \(R_i\) is the residual, \(\omega_i\) is the obstruction state, \(s_{ir}\) is the score, and \(\tau_{ir}\) is the threshold.

A routing explanation is actionably complete when it permits the participant to identify at least one decision-relevant proposition whose revision could alter the disposition. Let
\[
\mathfrak{C}_i
\]
be the set of contestable claims in \(\mathcal{E}_i\), and let
\[
\partial d_i
\]
denote the set of state variables to which the decision is locally sensitive. Actionable completeness requires
\[
\mathfrak{C}_i\cap\partial d_i\neq\varnothing.
\]
An explanation that reveals only variables to which the decision is insensitive is formally descriptive but procedurally useless.

\begin{definition}[Contestable routing]
A routing policy is contestable at state \(\sigma_i\) when there exists a finite evidential intervention \(\chi_i\) such that
\[
D(\sigma_i)\neq
D\!\left(
\mathcal{B}(\sigma_i,\chi_i)
\right)
\]
for at least one admissible correction to a disclosed decision premise.
\end{definition}

Contestability does not require that every disagreement reverse the outcome. It requires that the policy distinguish disagreement from correction. A participant can understand a prerequisite and reject it. They can accept the archive's cited answer while denying that it answers their actual question. They can agree that two local sections fail to glue while arguing that the failure is the intended result rather than a defect.

\subsection{Comprehension, assent, and ideological fixed points}

Let \(v\in V\) be a prerequisite node. Let
\[
C_i(v)\in\{0,1\}
\]
indicate demonstrated comprehension, and let
\[
A_i(v)\in\{0,1\}
\]
indicate assent. A legitimate competency gate uses \(C_i(v)\), not \(A_i(v)\).

\begin{definition}[Assent independence]
A prerequisite policy is assent-independent when
\[
\Prob(
\operatorname{pass}_i(v)=1
\mid
C_i(v),A_i(v)
)
=
\Prob(
\operatorname{pass}_i(v)=1
\mid
C_i(v)
)
\]
for every admissible participant state.
\end{definition}

Assent independence can be tested through counterfactual response pairs. Let \(y_i^{+}\) and \(y_i^{-}\) demonstrate equal command of a distinction while reaching opposite conclusions. A gate satisfies local assent symmetry when
\[
\Prob(\operatorname{pass}\mid y_i^{+})
=
\Prob(\operatorname{pass}\mid y_i^{-}).
\]
Systematic asymmetry supplies evidence that the prerequisite graph is functioning as ideological screening.

A stronger failure occurs when the archive defines understanding in terms of agreement with its conclusions:
\[
C_i(v)
=
\ind\{A_i(v)=1\}.
\]
Then disagreement implies incompetence by construction. The system becomes self-sealing because no critic can both reject the node and be recognized as understanding it.

\begin{theorem}[Self-sealing prerequisite impossibility]
Suppose access to recipient \(r\) requires passing node \(v\), and the system defines
\[
\operatorname{pass}_i(v)=A_i(v).
\]
Suppose a valid criticism of \(v\) necessarily entails
\[
A_i(v)=0.
\]
Then no valid criticism of \(v\) can reach \(r\) through the ordinary routing path.
\end{theorem}

\begin{proof}
Every valid criticism entails \(A_i(v)=0\). By the gate definition,
\[
\operatorname{pass}_i(v)=0.
\]
Passing \(v\) is necessary for access to \(r\), so the submission is blocked before escalation. Therefore no valid criticism of \(v\) reaches \(r\) through that path.
\end{proof}

Audit and appeal can break the fixed point only if they bypass the same assent condition. An appeal that again requires assent to the disputed node reproduces the impossibility at a second layer.

\subsection{Provenance authority and archive nonmonotonicity}

The archive distinguishes participant claims, recipient-authored claims, retrieved sources, model inferences, and generated reconstructions. Let provenance type be
\[
\pi(q)
\in
\{
\mathrm{participant},
\mathrm{recipient},
\mathrm{source},
\mathrm{inference},
\mathrm{generation},
\mathrm{audit}
\}.
\]
Let confidence be
\[
c(q)\in[0,1].
\]
Let status be
\[
s(q)
\in
\{
\mathrm{asserted},
\mathrm{supported},
\mathrm{contested},
\mathrm{defeated},
\mathrm{unresolved}
\}.
\]
An archive item is therefore represented by
\[
\kappa
=
(q,\pi(q),c(q),s(q),t,\mathcal{D}_q),
\]
where \(\mathcal{D}_q\) records dependencies and evidence.

Archive growth is nonmonotonic. A later correction can defeat an earlier answer:
\[
\K_t\models q,
\qquad
\K_{t+1}\nvDash q.
\]
The system must therefore preserve the distinction between append-only provenance and revisable epistemic status. Let
\[
\mathcal{P}_{t+1}
=
\mathcal{P}_t\cup\{\text{new records}\}
\]
be monotone provenance storage, while
\[
\K_{t+1}
=
\operatorname{Revise}(\K_t,\Delta_t)
\]
may be nonmonotone.

\begin{definition}[Epistemic reversibility]
An archive is epistemically reversible when every conclusion \(q\) has a dependency record sufficient to recompute its status after one or more supporting premises are withdrawn, defeated, or reclassified.
\end{definition}

Let
\[
\operatorname{Supp}(q)
\subseteq\K
\]
be the support set for \(q\). If support item \(p\in\operatorname{Supp}(q)\) is defeated, the affected closure is
\[
\operatorname{Dep}^{+}(p)
=
\left\{
q:p\leadsto q
\text{ in the dependency graph}
\right\}.
\]
Every item in \(\operatorname{Dep}^{+}(p)\) requires reconsideration. Without dependency provenance, a correction can be recorded while its downstream consequences remain silently unchanged.

\subsection{Representation pluralism}

A single quotient \(q\) creates a single family of invisible distinctions. Governance can reduce this risk by maintaining a family
\[
\mathfrak{Q}
=
\{q_{\alpha}:\mathcal{X}\to\mathcal{Z}_{\alpha}\}_{\alpha\in A}
\]
of admissible representations.

For task \(f\), define the insufficiency of \(q_\alpha\) by
\[
\mathcal{I}_f(q_\alpha)
=
\inf_g
\E\!\left[
\ell(
f(X),
g(q_\alpha(X))
)
\right].
\]
No representation is uniformly preferred unless
\[
\mathcal{I}_f(q_{\alpha^{*}})
\leq
\mathcal{I}_f(q_\alpha)
\]
for every \(f\) in the admitted task class and every \(\alpha\). Such dominance is rare under heterogeneous future tasks.

Define the plural representation
\[
Q_{\mathfrak{Q}}(x)
=
\left(
q_\alpha(x)
\right)_{\alpha\in A}.
\]
It refines every component representation, but its storage and attention costs can be large. Governance must choose a subset
\[
A^{*}\subseteq A
\]
that balances risk reduction against burden:
\[
A^{*}
\in
\argmin_{A'\subseteq A}
\left[
\sum_{f\in\mathcal{F}}
w_f
\inf_g
\E[
\ell(f(X),g(Q_{A'}(X)))
]
+
\nu C(Q_{A'})
\right].
\]

A participant-supplied alternative representation \(q_i^{\mathrm{alt}}\) should be admissible when it exposes a decision-relevant distinction hidden by the default quotient. Let
\[
q_0(x)=q_0(y)
\]
but
\[
q_i^{\mathrm{alt}}(x)\neq q_i^{\mathrm{alt}}(y).
\]
If the routing disposition differs,
\[
f_r(x)\neq f_r(y),
\]
then \(q_i^{\mathrm{alt}}\) witnesses a failure of default sufficiency.

\begin{proposition}[Witnessed quotient insufficiency]
If there exist \(x,y\in\mathcal{X}\) such that
\[
q_0(x)=q_0(y)
\]
and
\[
f_r(x)\neq f_r(y),
\]
then no decoder
\[
g:\mathcal{Z}_0\to\mathcal{Y}_r
\]
can satisfy
\[
g(q_0(z))=f_r(z)
\]
for every \(z\in\mathcal{X}\).
\end{proposition}

\begin{proof}
Suppose such a decoder exists. Since \(q_0(x)=q_0(y)\),
\[
g(q_0(x))=g(q_0(y)).
\]
Exact decoding would imply
\[
f_r(x)=g(q_0(x))=g(q_0(y))=f_r(y),
\]
contradicting the hypothesis.
\end{proof}

Governance must therefore treat a valid collision witness as evidence for representational expansion, not as a request to restate the same case more clearly inside the failed quotient.

\subsection{Threshold governance and policy sensitivity}

Let the escalation decision be
\[
E_i(\tau)
=
\ind\{S_i\geq\tau\}.
\]
The selected set is
\[
\mathcal{E}(\tau)
=
\{i:S_i\geq\tau\}.
\]
If the score distribution has density \(f_S\), then expected coverage is
\[
\kappa(\tau)
=
\Prob(S\geq\tau)
=
1-F_S(\tau),
\]
with derivative
\[
\kappa'(\tau)
=
-f_S(\tau).
\]
When \(f_S(\tau)\) is large, small threshold changes alter many decisions.

Let \(G\) be a participant class and define group-specific coverage
\[
\kappa_G(\tau)
=
\Prob(S\geq\tau\mid G).
\]
The disparity between groups \(G\) and \(H\) is
\[
\Delta_{GH}(\tau)
=
\kappa_G(\tau)-\kappa_H(\tau).
\]
Its derivative is
\[
\Delta_{GH}'(\tau)
=
-f_{S\mid G}(\tau)+f_{S\mid H}(\tau).
\]
A globally small policy change can produce a large group-specific effect when score densities differ near the threshold.

Threshold governance should therefore report a sensitivity interval. For perturbation radius \(\delta>0\), define
\[
\operatorname{Flip}_i(\delta)
=
\ind\{
|S_i-\tau|\leq\delta
\}.
\]
The expected decision instability is
\[
\mathcal{J}(\delta)
=
\Prob(|S-\tau|\leq\delta).
\]
A policy with high \(\mathcal{J}(\delta)\) is brittle. Many dispositions depend upon score differences smaller than the uncertainty in the score itself.

Let the score estimate have posterior distribution
\[
S_i\mid\mathscr{I}_i
\sim P_i.
\]
A confidence-aware decision escalates when
\[
\Prob(S_i\geq\tau\mid\mathscr{I}_i)
>
1-\delta_E,
\]
rejects when
\[
\Prob(S_i<\tau\mid\mathscr{I}_i)
>
1-\delta_D,
\]
and abstains otherwise. This exposes uncertainty that a point estimate would conceal.

\subsection{Privacy and the minimum sufficient participant model}

The participant-state model
\[
\thetaState_i(t)
\]
can improve routing, but every additional latent variable creates surveillance capacity. Let
\[
\Theta_i^{\mathrm{full}}
\]
be a broad behavioral model and let
\[
\thetaState_i^{\mathrm{task}}
=
m_r(\Theta_i^{\mathrm{full}})
\]
be the information required for the current routing task.

A governance constraint should require task sufficiency with minimal excess information. Let \(Y_r\) be the routing outcome and \(Z\) a candidate participant representation. Define task risk
\[
R_r(Z)
=
\inf_g
\E[
\ell(Y_r,g(Z))
].
\]
Let \(I(Z;W)\) be the mutual information between \(Z\) and sensitive participant attribute \(W\). The privacy-constrained representation problem is
\[
\begin{aligned}
\text{minimize}\quad&
I(Z;W)
\\
\text{subject to}\quad&
R_r(Z)\leq R_r^{*}+\eps,
\end{aligned}
\]
where \(R_r^{*}\) is the minimum achievable routing risk under the available observations.

More generally, an information-bottleneck objective is
\[
\min_{p(z\mid x)}
\left[
I(X;Z)-\beta I(Z;Y_r)
\right].
\]
The first term penalizes retention of participant information. The second rewards task-relevant information. The solution is not automatically ethically sufficient because \(Y_r\) may itself encode a biased policy. Minimal information relative to an unjust target remains unjust.

\begin{definition}[Epistemic locality]
A participant-state variable \(\theta_j\) is epistemically local to routing task \(r\) when changing \(\theta_j\) while holding every task-relevant variable fixed does not change the routing disposition:
\[
D_r(\theta_j,\theta_{-j})
=
D_r(\theta_j',\theta_{-j})
\]
for all admissible \(\theta_j,\theta_j'\).
\end{definition}

Variables that are epistemically local in this sense need not be stored for the decision. If they are retained for research or aggregate improvement, that secondary use requires separate authorization and retention limits.

\subsection{Gaming, strategic expression, and mechanism dependence}

Participants adapt to routing rules. Let type
\[
\vartheta_i
\]
encode the participant's actual residual, competence, and intended contribution. Let
\[
m_i
\]
be the submitted representation. The participant chooses \(m_i\) to maximize
\[
U_i(m_i;\vartheta_i)
=
p_E(m_i)G_i(\vartheta_i)
-
c_i(m_i,\vartheta_i)
-
r_i(m_i),
\]
where \(p_E\) is escalation probability, \(c_i\) is expression cost, and \(r_i\) is the expected penalty for detected manipulation.

A routing mechanism is truthfully expressive when the participant's best response is to provide the most decision-relevant faithful representation of \(\vartheta_i\). Formally, if
\[
m^{*}(\vartheta_i)
\]
is the faithful report, incentive compatibility requires
\[
U_i(m^{*}(\vartheta_i);\vartheta_i)
\geq
U_i(m;\vartheta_i)
\]
for every admissible \(m\).

This condition is unlikely to hold exactly when linguistic polish, credentials, citations, or conventional formatting increase the score independently of residual value. Let
\[
S(m_i)
=
V(\vartheta_i)
+
\alpha P(m_i)
+
\varepsilon_i,
\]
where \(P(m_i)\) is polish. If \(\alpha>0\), participants can increase scores through polish even when \(V(\vartheta_i)\) is fixed. Developmental transformation reduces this incentive by reconstructing low-polish submissions and adversarially testing high-polish ones.

A cost gate changes the strategic problem:
\[
U_i(m_i)
=
p_E(m_i)G_i-c_i(m_i)-\gamma.
\]
A fixed monetary cost \(\gamma\) suppresses submissions whose expected private gain is below \(\gamma\), not submissions whose social value is below \(\gamma\). If wealth changes the marginal burden of payment, the gate selects jointly on confidence, private benefit, and ability to pay.

Let \(w_i\) be wealth and let the subjective cost be
\[
\gamma_i=\frac{\gamma}{u'(w_i)}
\]
under a local utility conversion. Even if two participants have equal expected contribution, their submission probabilities can differ because
\[
\gamma_i\neq\gamma_j.
\]
An economic gate can reduce arrival rate without being value-neutral.

\subsection{Fairness criteria and incompatibility}

Let \(G\in\{0,1\}\) denote a participant group, let \(V\in\{0,1\}\) denote true value, and let \(E\in\{0,1\}\) denote escalation. Several fairness conditions are plausible.

Demographic parity requires
\[
\Prob(E=1\mid G=0)
=
\Prob(E=1\mid G=1).
\]
Equal opportunity requires
\[
\Prob(E=1\mid V=1,G=0)
=
\Prob(E=1\mid V=1,G=1).
\]
Equalized odds additionally requires
\[
\Prob(E=1\mid V=0,G=0)
=
\Prob(E=1\mid V=0,G=1).
\]
Predictive parity requires
\[
\Prob(V=1\mid E=1,G=0)
=
\Prob(V=1\mid E=1,G=1).
\]
Calibration requires
\[
\Prob(V=1\mid S=s,G=g)=s
\]
for all relevant \(s,g\).

When base rates
\[
\Prob(V=1\mid G=g)
\]
differ, these conditions can be mutually incompatible.

\begin{theorem}[Predictive parity and equalized-odds incompatibility]
Suppose
\[
0<\Prob(E=1\mid V=v,G=g)<1
\]
for \(v,g\in\{0,1\}\), and suppose the value base rates differ:
\[
\pi_0
=
\Prob(V=1\mid G=0)
\neq
\Prob(V=1\mid G=1)
=
\pi_1.
\]
If equalized odds holds with common true-positive rate \(t\) and common false-positive rate \(f\), where \(t,f\in(0,1)\), then predictive parity cannot hold unless
\[
t=f.
\]
When \(t=f\), the decision contains no information about \(V\).
\end{theorem}

\begin{proof}
Under equalized odds,
\[
\Prob(E=1\mid V=1,G=g)=t
\]
and
\[
\Prob(E=1\mid V=0,G=g)=f.
\]
By Bayes' rule,
\[
\Prob(V=1\mid E=1,G=g)
=
\frac{
t\pi_g
}{
t\pi_g+f(1-\pi_g)
}.
\]
Define
\[
h(\pi)
=
\frac{t\pi}{t\pi+f(1-\pi)}.
\]
For \(t,f>0\),
\[
h'(\pi)
=
\frac{tf}{[t\pi+f(1-\pi)]^2}
>0.
\]
Hence \(h\) is strictly increasing. Since \(\pi_0\neq\pi_1\),
\[
h(\pi_0)\neq h(\pi_1),
\]
so predictive parity fails. Equality can be recovered only in degenerate cases, including \(t=f\), where
\[
E\ind V\mid G
\]
and the decision has no discriminatory value.
\end{proof}

Governance cannot select every fairness criterion simultaneously. It must state which error distribution is being constrained and why. In an attention compiler, equal opportunity and bounded false suppression may often be more closely related to the central epistemic objective than demographic parity, but that choice remains context-dependent.

\subsection{The impossibility of recipient-neutral value}

The escalation value
\[
V_{ir}
\]
depends on recipient \(r\). There is generally no scalar
\[
V_i
\]
that preserves every recipient's ordering.

\begin{theorem}[No universal recipient-independent ranking]
Suppose there exist two recipients \(r_1,r_2\) and submissions \(i,j\) such that
\[
V_{ir_1}>V_{jr_1}
\]
while
\[
V_{ir_2}<V_{jr_2}.
\]
Then no scalar recipient-independent score \(s_i\) can reproduce both recipients' strict orderings.
\end{theorem}

\begin{proof}
To reproduce \(r_1\)'s ordering, the score must satisfy
\[
s_i>s_j.
\]
To reproduce \(r_2\)'s ordering, it must satisfy
\[
s_i<s_j.
\]
No scalar pair satisfies both inequalities.
\end{proof}

A universal novelty ranking therefore cannot replace recipient-relative routing. The system may compute general importance, but direct access remains a matching problem involving the residual, participant, recipient, and present context.

\subsection{The impossibility of lossless universal compression}

Let \(\mathcal{F}\) be a downstream task family. A quotient \(q\) is sufficient for \(\mathcal{F}\) when every \(f\in\mathcal{F}\) factors through \(q\):
\[
f=g_f\circ q.
\]
If \(\mathcal{F}\) contains every function
\[
f:\mathcal{X}\to\{0,1\},
\]
then any sufficient quotient must be injective.

\begin{theorem}[Universal-task sufficiency requires injectivity]
Let
\[
q:\mathcal{X}\to\mathcal{Z}.
\]
If every binary function
\[
f:\mathcal{X}\to\{0,1\}
\]
factors through \(q\), then \(q\) is injective.
\end{theorem}

\begin{proof}
Assume \(q\) is not injective. Then there exist distinct \(x,y\in\mathcal{X}\) such that
\[
q(x)=q(y).
\]
Define
\[
f(z)
=
\begin{cases}
1,&z=x,\\
0,&z\neq x.
\end{cases}
\]
If \(f=g\circ q\), then
\[
f(x)=g(q(x))=g(q(y))=f(y),
\]
but
\[
f(x)=1
\qquad\text{and}\qquad
f(y)=0,
\]
a contradiction. Therefore \(q\) is injective.
\end{proof}

Universal future-proof compression is impossible. Any nontrivial quotient sacrifices sufficiency for some possible downstream task. Governance must therefore declare the task family relative to which compression is justified and preserve a repair path for task expansion.

\subsection{The impossibility of simultaneous privacy and unrestricted future repair}

Suppose the system deletes every discarded coordinate after computing \(q(X)\). Let \(W\) be a sensitive attribute contained entirely in the discarded information:
\[
I(W;q(X))=0.
\]
This offers strong privacy relative to \(W\). Suppose a future task \(f'\) depends on \(W\). Then no decoder using only \(q(X)\) can be sufficient for \(f'\).

\begin{theorem}[Privacy-repair tension]
If
\[
I(W;Z)=0
\]
for \(Z=q(X)\), and a binary future task satisfies
\[
Y=f'(X)=W
\]
with
\[
0<\Prob(W=1)<1,
\]
then every predictor \(\widehat Y=g(Z)\) has Bayes error at least
\[
\min\{\Prob(W=1),\Prob(W=0)\}.
\]
\end{theorem}

\begin{proof}
Independence gives
\[
\Prob(W=1\mid Z)=\Prob(W=1)
\]
almost surely. Hence observing \(Z\) does not change the posterior distribution of \(W\). Under zero-one loss, the Bayes-optimal predictor selects the more probable constant label. Its error is the probability of the less common label:
\[
\min\{\Prob(W=1),\Prob(W=0)\}.
\]
No function of \(Z\) can do better.
\end{proof}

Preserving all discarded coordinates supports future repair but weakens deletion-based privacy. Destroying them supports privacy but forecloses some later tasks. Governance must specify retention horizons, protected provenance stores, access conditions, and which future repairs are deliberately made impossible.

\subsection{The impossibility of perfect explanation under bounded attention}

A complete explanation of a routing decision may be as complex as the decision process itself. Let
\[
K(D,\sigma)
\]
be the description length of the decision computation at state \(\sigma\), and let
\[
K(\mathcal{E})
\]
be the description length of an explanation. Suppose an explanation is strongly faithful when it permits exact reconstruction of the decision for every admissible perturbation in neighborhood \(U(\sigma)\).

\begin{theorem}[Explanation compression limit]
For a decision family whose restriction to \(U(\sigma)\) has Kolmogorov complexity \(K_U(D)\), every strongly faithful explanation \(\mathcal{E}\) from which that restriction is computable satisfies
\[
K(\mathcal{E})
\geq
K_U(D)-c,
\]
where \(c\) is a constant depending only on the fixed reconstruction procedure.
\end{theorem}

\begin{proof}
By assumption, a fixed reconstruction program \(R\) computes the decision restriction \(D|_{U(\sigma)}\) from \(\mathcal{E}\). Hence
\[
K_U(D)
\leq
K(\mathcal{E})+K(R)+O(1).
\]
Absorb \(K(R)+O(1)\) into constant \(c\) and rearrange:
\[
K(\mathcal{E})
\geq
K_U(D)-c.
\]
\end{proof}

A short explanation cannot always be both complete and strongly faithful. Governance should require actionable explanations rather than pretending that every complex decision admits a simple exhaustive rationale. The explanation must reveal the contestable boundary, provenance, and decisive factors, even when it cannot reproduce the entire computation in ordinary language.

\subsection{The impossibility of simultaneous complete access and bounded delay}

Let submissions arrive with rate \(\lambda\), let every submission be entitled to positive recipient attention \(c_0>0\), and let the recipient have capacity \(A\). If
\[
\lambda c_0>A,
\]
the queue is unstable.

\begin{theorem}[Access-delay impossibility]
Suppose every submission must receive at least \(c_0>0\) units of recipient attention and submissions arrive at long-run rate \(\lambda\). If the recipient supplies attention at rate \(A<\lambda c_0\), then no scheduling policy can guarantee finite waiting time for every submission.
\end{theorem}

\begin{proof}
By time \(t\), incoming required workload approaches
\[
\lambda c_0t,
\]
while supplied attention approaches
\[
At.
\]
The unserved workload grows at asymptotic rate
\[
\lambda c_0-A>0.
\]
If every submission had finite waiting time under a work-conserving policy, the persistent linearly growing backlog would have to be cleared in finite time, contradicting the negative capacity slack. Therefore some waiting times diverge.
\end{proof}

Universal outer access is compatible with finite recipient attention only because transformation permits many submissions to be answered, developed, aggregated, rerouted, or resolved without direct recipient contact. Universal direct access is not.

\subsection{The impossibility of exact false-suppression knowledge}

Section 10 established that false suppression must be estimated. The impossibility can be stated more strongly.

\begin{theorem}[No exact suppression census without exhaustive counterfactual evaluation]
Suppose value \(V_i\) is observable only after applying an evaluation intervention of positive cost to submission \(i\). If any discarded submission remains unevaluated, then the exact finite-population false-suppression count
\[
T_{VD}
=
\sum_{i:D_i=1}V_i
\]
is not identifiable without additional assumptions.
\end{theorem}

\begin{proof}
Let \(j\) be a discarded unevaluated submission. Construct two latent value assignments agreeing on every evaluated submission and every observed variable, but set
\[
V_j=0
\]
under the first assignment and
\[
V_j=1
\]
under the second. Both assignments generate identical observed records because \(j\) was not evaluated. Their false-suppression totals differ by one. Therefore the exact total is not identifiable.
\end{proof}

Random audit estimates the total but does not convert it into an observed census. Governance must prohibit presentation of an estimated false-suppression rate as exact.

\subsection{The impossibility of neutral governance}

Suppose a governance regime selects weights
\[
w
=
(w_\rho,w_h,w_g,w_\tau,w_e,w_u,w_c)
\]
for novelty, coherence, consequence, tractability, evidence, uncertainty, and cost. Different weights produce different escalation orderings.

\begin{theorem}[No policy-neutral scalarization under conflicting objectives]
Let \(x\) and \(y\) have objective vectors
\[
v(x),v(y)\in\mathbb{R}^{d}
\]
such that neither Pareto-dominates the other. Then there exist positive weight vectors \(w,w'\) satisfying
\[
w^{\top}v(x)>w^{\top}v(y)
\]
and
\[
w'^{\top}v(x)<w'^{\top}v(y).
\]
\end{theorem}

\begin{proof}
Because neither vector dominates the other, there exist coordinates \(j,k\) such that
\[
v_j(x)>v_j(y)
\]
and
\[
v_k(x)<v_k(y).
\]
Choose a positive weight vector \(w\) assigning sufficiently large weight to coordinate \(j\) and sufficiently small positive weights to the others. Then
\[
w^{\top}(v(x)-v(y))>0.
\]
Choose \(w'\) analogously with sufficiently large weight on coordinate \(k\). Then
\[
w'^{\top}(v(x)-v(y))<0.
\]
\end{proof}

A scalar score necessarily embeds priority judgments when objectives conflict. Making the weights visible does not make them neutral. It makes their consequences inspectable.

\subsection{Moniform markets and appropriation constraints}

A residual classified as a tradable partial can be finished by someone other than its originator. Let \(m_i\) be a moniform residual, let \(o_i\) be its originator, let \(f_j\) be a finisher, and let \(r\) be an eventual recipient. Let finished value be
\[
v(m_i,f_j,r),
\]
finishing cost be
\[
F(m_i,f_j),
\]
and transfer price be
\[
p(m_i).
\]
A positive surplus exists when
\[
v(m_i,f_j,r)-F(m_i,f_j)-p(m_i)>0.
\]
This economic condition does not determine attribution, consent, or control over derivative use.

Let
\[
a_i\in\{0,1\}
\]
indicate whether the originator authorized market transfer, and let
\[
\pi_i^{\mathrm{attr}}\in[0,1]
\]
measure attribution completeness. A governed trade requires
\[
a_i=1
\]
and
\[
\pi_i^{\mathrm{attr}}\geq\pi_0.
\]
If the system can reclassify an unsuccessful escalation attempt as a tradable residual without participant authorization, then developmental triage becomes an extraction mechanism.

Define originator surplus
\[
U_o=p(m_i)-C_o,
\]
finisher surplus
\[
U_f=v_f-F-p_f,
\]
and recipient surplus
\[
U_r=v_r-p_r.
\]
Individual rationality requires
\[
U_o\geq0,
\qquad
U_f\geq0,
\qquad
U_r\geq0.
\]
A market may clear while allocating most surplus to one party. Governance must distinguish market feasibility from equitable distribution.

\subsection{Randomization and accountable arbitrariness}

Random review protects against model closure, but randomization also produces decisions that cannot be justified by comparative merit alone. Let
\[
\pi_i^{x}
\]
be exploratory inclusion probability. Two identical residuals may receive different outcomes because one is sampled and the other is not.

This arbitrariness is accountable when the sampling law is declared, inclusion probabilities are positive, and randomization serves an auditable epistemic purpose. Let
\[
Z_i\sim\operatorname{Bernoulli}(\pi_i^{x}).
\]
The realized difference
\[
Z_i-Z_j
\]
for equal cases is not a value judgment. It is a sampling event.

A governance record must therefore distinguish
\[
\mathrm{not\ selected\ by\ value\ policy}
\]
from
\[
\mathrm{not\ sampled\ under\ exploration}.
\]
Without this distinction, participants may interpret random exclusion as substantive rejection, and later auditors may incorrectly treat the absence of review as evidence of low value.

\subsection{Governed silence}

Finite attention sometimes requires deliberate noncontinuation. Let
\[
\mathfrak{A}(\sigma)
\]
be the set of reachable actions at state \(\sigma\). Define silence as
\[
\operatorname{Silence}(\sigma)
=
\text{the deliberate selection of no outward continuation}
\]
despite
\[
\mathfrak{A}(\sigma)\neq\varnothing.
\]
Silence differs from failure because an action is available. It differs from discard because the state may remain preserved. It differs from delay because no resumption time is promised.

A governed silence record is
\[
\mathcal{S}_i
=
\left(
\sigma_i,
\mathfrak{A}(\sigma_i),
\operatorname{reason}_i,
\operatorname{scope}_i,
\operatorname{retention}_i,
\operatorname{reviewability}_i
\right).
\]
Silence may be justified by privacy, safety, legal constraint, ethical sensitivity, recipient incapacity, or the risk that response itself amplifies harm. It becomes unaccountable when it is indistinguishable from invisible rejection.

Let \(C_{\mathrm{respond}}\) be expected cost of response and \(H_{\mathrm{respond}}\) expected harm. Let \(V_{\mathrm{respond}}\) be expected value. Silence has positive expected utility when
\[
V_{\mathrm{respond}}
<
C_{\mathrm{respond}}
+
H_{\mathrm{respond}}.
\]
This inequality does not erase provenance or appeal rights unless disclosure itself creates the relevant harm.

\subsection{Governance invariants}

The governed system should preserve a set of invariants under every authorized update. Let
\[
\Theta_t
=
(q_t,\K_t,G_t,\mathcal{H}_t,S_t,\pi_t)
\]
be the architecture state and let
\[
\Gamma_a
\]
be an authorized governance transition.

The provenance invariant requires
\[
\mathcal{P}_t
\subseteq
\mathcal{P}_{t+1}.
\]
The contestability invariant requires that every adverse decision retain an admissible challenge path unless a documented exception applies:
\[
d_i\in\mathcal{D}_{\mathrm{adverse}}
\Longrightarrow
\exists\chi_i\in\mathfrak{A}_i^{\mathrm{appeal}}.
\]
The exploration invariant requires
\[
A_{x,t}\geq\xi A_t.
\]
The observability invariant requires
\[
\pi_{g,t}^{A}>0
\]
for every eligible audit stratum \(g\). The assent-independence invariant requires
\[
\operatorname{pass}\ind A_i(v)\mid C_i(v).
\]
The recipient-capacity invariant requires
\[
\lambda_r\Prob(E)\E[C\mid E]<A_r.
\]
The representation-repair invariant requires either original-state retention under protected access or an explicit declaration that the discarded distinction is irrecoverable:
\[
q_t(x)=q_t(y)
\Longrightarrow
\left[
\operatorname{Recoverable}_t(x,y)
\lor
\operatorname{DeclaredIrrecoverable}_t(x,y)
\right].
\]

\begin{definition}[Governance-admissible transition]
An architecture update
\[
\Theta_{t+1}=\Gamma_a(\Theta_t)
\]
is governance-admissible when it preserves every applicable invariant, records its authorization and evidence, and supplies a reversal or migration procedure for every nontrivial change whose downstream effects are not yet established.
\end{definition}

\subsection{Versioned governance and retroactive review}

A policy change at time \(t^{*}\) creates two decision regimes:
\[
D_t
=
\begin{cases}
D^{(0)},&t<t^{*},\\
D^{(1)},&t\geq t^{*}.
\end{cases}
\]
If \(D^{(1)}\) corrects a defect in \(D^{(0)}\), earlier decisions affected by the same defect may require replay.

Let
\[
\mathcal{B}_{t^{*}}
=
\left\{
i:
D^{(0)}(\sigma_i)\neq D^{(1)}(\sigma_i)
\right\}
\]
be the retroactive discrepancy set. Recomputing every historical decision may be infeasible. Define estimated correction value
\[
v_i^{R}
=
\Prob(
D^{(0)}(\sigma_i)\text{ was harmful}
\mid
\mathscr{I}_i,D^{(1)}
)
G_i
\]
and replay cost \(c_i^{R}\). Retroactive review becomes
\[
\begin{aligned}
\text{maximize}\quad&
\sum_{i\in\mathcal{B}_{t^{*}}}
z_i^{R}v_i^{R}
\\
\text{subject to}\quad&
\sum_{i\in\mathcal{B}_{t^{*}}}
z_i^{R}c_i^{R}
\leq
A_R.
\end{aligned}
\]
Random sampling from the nonreplayed discrepancy set remains necessary to estimate residual historical harm.

A policy ledger records
\[
\mathcal{L}
=
\left\{
(t_k,\Theta_{t_k},\Delta\Theta_{t_k},
\operatorname{authority}_k,
\operatorname{evidence}_k,
\operatorname{scope}_k)
\right\}_{k\geq0}.
\]
Without versioning, a later replay cannot distinguish changed judgment from inconsistent execution.

\subsection{A governance trilemma}

Three desirable properties recur throughout the architecture. Let \(U\) denote universal direct access to recipient attention, let \(B\) denote bounded recipient workload, and let \(Z\) denote zero false suppression of valuable residuals.

\begin{theorem}[Attention governance trilemma]
Under positive evaluation costs, unbounded admissible arrivals, and imperfect pre-evaluation value information, no system can guarantee \(U\), \(B\), and \(Z\) simultaneously.
\end{theorem}

\begin{proof}
Universal direct access requires that every arrival receive recipient attention of positive expected cost. Under unbounded arrival rates, there exists an arrival process whose required workload exceeds finite recipient capacity, violating \(B\). To preserve \(B\), the system must mediate, delay, or suppress some arrivals before direct recipient evaluation. Since pre-evaluation information is imperfect, there exists an admissible distribution under which some mediated or suppressed arrivals are valuable, violating \(Z\). Hence no system guarantees all three properties.
\end{proof}

The architecture selects bounded workload, broad outer access, and an audited upper bound on estimated false suppression. It does not promise universal direct access or zero epistemic loss.

\subsection{A representation governance trilemma}

Let \(C\) denote nontrivial compression, \(F\) denote sufficiency for every possible future task, and \(P\) denote irreversible deletion of discarded information for privacy.

\begin{theorem}[Representation governance trilemma]
For a domain containing at least two distinguishable states, no representation can guarantee \(C\), \(F\), and \(P\) simultaneously.
\end{theorem}

\begin{proof}
Nontrivial compression implies that the quotient \(q\) is noninjective, so there exist \(x\neq y\) with
\[
q(x)=q(y).
\]
Define a future task \(f\) that distinguishes \(x\) and \(y\). Then \(f\) cannot factor through \(q\), violating \(F\). Preserving auxiliary information sufficient to distinguish \(x\) and \(y\) would permit repair but violate irreversible deletion \(P\). Thus all three properties cannot hold simultaneously.
\end{proof}

The system must choose task-relative compression, protected recoverability, or irreversible privacy according to declared contexts. There is no universal technical resolution.

\subsection{A governance objective under explicit constraints}

Let
\[
\mathcal{V}(\Gamma)
\]
be expected intellectual value under governance regime \(\Gamma\). Let
\[
\mathcal{C}_{\mathrm{att}}(\Gamma)
\]
be recipient attention cost,
\[
\mathcal{C}_{\mathrm{bur}}(\Gamma)
\]
participant burden,
\[
\mathcal{R}_{\mathrm{sup}}(\Gamma)
\]
false-suppression risk,
\[
\mathcal{R}_{\mathrm{priv}}(\Gamma)
\]
privacy risk,
\[
\mathcal{R}_{\mathrm{cap}}(\Gamma)
\]
capture risk, and
\[
\mathcal{R}_{\mathrm{irr}}(\Gamma)
\]
irreversibility risk.

A scalarized governance objective is
\[
\begin{aligned}
J(\Gamma)
={}&
\mathcal{V}(\Gamma)
-
\alpha\mathcal{C}_{\mathrm{att}}(\Gamma)
-
\beta\mathcal{C}_{\mathrm{bur}}(\Gamma)
\\
&
-
\gamma\mathcal{R}_{\mathrm{sup}}(\Gamma)
-
\delta\mathcal{R}_{\mathrm{priv}}(\Gamma)
-
\eta\mathcal{R}_{\mathrm{cap}}(\Gamma)
-
\zeta\mathcal{R}_{\mathrm{irr}}(\Gamma).
\end{aligned}
\]
Because the weights are governance judgments, a constrained formulation is often more transparent:
\[
\begin{aligned}
\text{maximize}\quad&
\mathcal{V}(\Gamma)
\\
\text{subject to}\quad&
\lambda_r\Prob(E)\E[C\mid E]<A_r,
\\
&
\Prob(
\psi_r>\eps
\mid
\mathcal{D}_{\Audit}
)
\leq\delta_{\mathrm{risk}},
\\
&
A_x\geq\xi A,
\\
&
\pi_g^{A}\geq\pi_{\min,g}>0,
\\
&
I(Z;W)\leq\kappa_{\mathrm{priv}},
\\
&
\mathrm{PC}(\Pi_i)\geq p_0,
\\
&
\mathcal{J}(\delta)\leq\kappa_{\mathrm{brittle}},
\\
&
\operatorname{pass}\ind A_i(v)\mid C_i(v),
\\
&
\Gamma\text{ preserves governance invariants}.
\end{aligned}
\]
The constraints do not remove normative judgment. They expose where it enters and make violations legible.

\subsection{Conditional legitimacy}

Let \(\mathfrak{C}\) be a governed attention compiler. Define legitimacy relative to declared population \(\mathcal{I}\), recipient set \(\mathcal{R}\), task family \(\mathcal{F}\), and evaluation horizon \(H\).

\begin{definition}[Procedural legitimacy]
The system is procedurally legitimate at tolerance vector
\[
\varepsilon
=
(\eps_{\mathrm{sup}},
\eps_{\mathrm{priv}},
\eps_{\mathrm{bur}},
\eps_{\mathrm{delay}},
\eps_{\mathrm{capture}})
\]
when its policies are versioned and visible, adverse decisions are actionably explained, comprehension is distinguished from assent, provenance is reconstructible, audits have positive coverage over every eligible stratum, appeals are substantively capable of revision, representation changes admit collision challenges, and all declared risk estimates remain within their respective tolerances.
\end{definition}

This legitimacy is conditional rather than absolute. It depends on the declared task family, the adequacy of audit labels, the representativeness of sampled strata, the correctness of cost estimates, and the availability of viable alternatives. A system can satisfy its own procedures while the procedures encode unjust priorities. Governance must therefore permit challenge not only to individual applications of the rules but to the rules themselves.

Let
\[
\mathfrak{a}_{\mathrm{case}}
\]
challenge a decision and let
\[
\mathfrak{a}_{\mathrm{rule}}
\]
challenge the governing policy. Rule-level contestability requires
\[
\exists\chi:
\Gamma'=
\mathfrak{a}_{\mathrm{rule}}(\Gamma,\chi)
\neq
\Gamma
\]
for sufficiently strong evidence. If every policy challenge is reclassified as a case-level misunderstanding, governance becomes self-sealing at the architectural level.

\subsection{The governance boundary}

A recipient possesses a legitimate claim over their own finite attention. A participant possesses a legitimate claim not to have their contribution misrepresented, silently appropriated, ideologically screened, or declared answered by material that does not address it. An archive possesses no independent authority. Its authority is inherited from provenance, evidence, continuing review, and the procedures through which its claims remain revisable.

These claims do not combine into a frictionless right of access. They define a boundary. The recipient can refuse direct attention without acquiring the right to falsify why the submission was refused. The participant can contest reconstruction without acquiring the right to compel unlimited reciprocal engagement. The system can compress repeated material without acquiring the right to erase the distributional pattern created by repetition. The archive can route through prerequisites without acquiring the right to equate disagreement with incompetence. The audit can estimate false suppression without claiming to observe it exactly.

The governed transformation can be written
\[
x_i
\xrightarrow[\Gamma_q]{q}
z_i
\xrightarrow[\Gamma_{\K},\Gamma_G]{T}
(a_i,p_i,r_i,u_i,\omega_i)
\xrightarrow[\Gamma_{\mathcal{H}}]{\mathcal{H}}
\sigma_i
\xrightarrow[\Gamma_S,\Gamma_E]{S,\pi_E}
d_i
\xrightarrow[\Gamma_A,\Gamma_{\mathfrak{a}}]{\Audit,\mathfrak{a}}
d_i'.
\]
Every arrow is both a computation and an authorized intervention. Every output inherits uncertainty from earlier arrows. A later appeal can repair an earlier state only if provenance retains a path back to the distinction that was lost or misclassified.

The principal governance requirement is therefore not that every decision be correct. That requirement is impossible under finite attention, latent value, nontrivial compression, and changing future tasks. The requirement is that error not become authority merely by becoming difficult to observe.

\subsection{Conclusion}

The mathematical architecture cannot govern itself. It can expose the locations at which governance decisions occur, quantify some consequences, preserve evidence, and make certain failures impossible to conceal cheaply. It cannot derive the correct weights, task family, privacy boundary, fairness criterion, or distribution of institutional authority from the fact that attention is finite.

The impossibility results establish the boundaries of any credible promise. Universal direct access, bounded workload, and zero false suppression cannot coexist under unbounded arrivals. Nontrivial compression cannot remain sufficient for every future task. Irreversible privacy can block later representational repair. Exact false-suppression knowledge requires exhaustive evaluation. Complete explanations can be as complex as the decisions they explain. Recipient-relative values need not admit a universal ranking. Multiple fairness criteria can conflict when base rates differ. Scalarization cannot avoid policy judgment when objectives are noncomparable.

Within these limits, a defensible system preserves several conditional guarantees. It keeps provenance monotone while allowing epistemic status to change. It separates comprehension from assent. It reserves positive audit probability for every eligible region. It makes adverse routing decisions contestable at both case and rule levels. It exposes threshold sensitivity and uncertainty. It treats participant-state inference as task-local rather than as permission for unrestricted psychological classification. It distinguishes random nonselection from substantive rejection. It treats moniform transfer as consensual exchange rather than automatic appropriation. It records silence as an intentional governance operation when silence is chosen.

Representational simplicity is legitimate only when simplicity remains answerable to what it excludes. A quotient must admit collision witnesses. A prerequisite must admit competent dissent. A score must disclose its policy dependence. An audit must reach beyond the system's preferred examples. An appeal must be capable of changing something. An archive must preserve the distinction between what was written, what was inferred, what was generated, and what was later corrected.

The governing condition is
\[
\text{finite attention may justify selective continuation, but it does not justify untraceable transformation}.
\]
A governed attention compiler does not eliminate concentrated authority. It converts authority from an invisible fact of inbox congestion into a versioned, constrained, auditable, and partially reversible set of operations. Whether that conversion is sufficient remains a political and institutional question, but without it the mathematical sophistication of the routing system would merely make exclusion harder to see.

\section{The Complete Architecture and Its Principal Conditional Guarantees}

The preceding sections introduced the components of representational simplicity separately because each solves a different failure problem. CLIO addresses the admissibility of projection. HYDRA addresses the operational assembly of perception, relevance, and memory. The attention compiler addresses developmental transformation, recipient-relative escalation, and finite-capacity allocation. Obstruction analysis addresses failures of global assembly. Audit and appeal address the invisibility and reversibility of suppression. Collective recurrence addresses information that exists only at the population level. Governance constrains the authority exercised through every one of these operations.

The integrated system is not a sequence of independent filters. It is a partially observed, dynamically governed network of transformations with feedback, repair, expansion, and audit paths. An error introduced at an early representational stage can propagate through every later decision. A correction discovered at a later stage can require revision of the earlier quotient, archive, prerequisite graph, participant-state model, coupling operator, or attention policy. The full architecture must therefore be described as a stateful dynamical system rather than as a static pipeline.

Let
\[
\mathfrak{R}
=
\left(
\mathcal{X},
\mathcal{Z},
\mathcal{K},
\mathcal{G},
\ThetaSpace,
\mathcal{Y},
\mathcal{Q},
\mathcal{U},
\mathcal{L}
\right)
\]
be the ambient representational system. Here \(\mathcal{X}\) is the raw submission space, \(\mathcal{Z}\) is the family of admissible projected representations, \(\mathcal{K}\) is the space of archive states, \(\mathcal{G}\) is the space of prerequisite graphs, \(\ThetaSpace\) is the participant-state space, \(\mathcal{Y}\) is the space of developmental and evaluative outputs, \(\mathcal{Q}\) is the collection of recipient-facing attention queues, \(\mathcal{U}\) is the space of governance-authorized updates, and \(\mathcal{L}\) is the policy and provenance ledger.

At time \(t\), the complete architecture state is
\[
\Xi_t
=
\left(
\K_t,
G_t,
\mathfrak{Q}_t,
\{\thetaState_i(t)\}_{i\in\mathcal{I}},
\mathcal{H}_t,
\mathcal{Q}_t,
\mathcal{A}_t,
\mathcal{D}_t,
\Gamma_t,
\mathcal{P}_t
\right),
\]
where \(\mathfrak{Q}_t\) is the active family of quotients, \(\mathcal{H}_t\) is the HYDRA transition operator, \(\mathcal{Q}_t\) is the collection of queues, \(\mathcal{A}_t\) is the audit state, \(\mathcal{D}_t\) is the collective recurrence state, \(\Gamma_t\) is the governance regime, and \(\mathcal{P}_t\) is provenance.

A new submission arrives as
\[
\xi_i^{0}
=
\left(
x_i,
i,
r_i^{\mathrm{intended}},
t_i,
\pi_i^{0}
\right),
\]
where \(x_i\) is the native expression, \(i\) is the participant identity or pseudonymous identifier, \(r_i^{\mathrm{intended}}\) is the intended recipient, \(t_i\) is the arrival time, and \(\pi_i^{0}\) records origin provenance.

\subsection{The composed transformation}

The native submission first enters a charitable reconstruction operator
\[
\mathcal{C}_{\mathrm{char}}:
\mathcal{X}\times\ThetaSpace\times\K
\to
\widehat{\mathcal{X}}\times\mathcal{U}_{\mathrm{char}},
\]
giving
\[
\left(
\widehat x_i,
u_i^{\mathrm{char}}
\right)
=
\mathcal{C}_{\mathrm{char}}
\left(
x_i,
\thetaState_i(t_i),
\K_{t_i}
\right).
\]
The reconstruction is not substituted for the original. Provenance retains
\[
x_i
\quad\text{and}\quad
\widehat x_i
\]
as distinct objects.

Claim extraction produces
\[
Q_i
=
\mathcal{C}_{\mathrm{claim}}(\widehat x_i),
\]
where
\[
Q_i
=
\{q_{i1},\ldots,q_{im_i}\}
\]
contains normalized claims, questions, definitions, commitments, requests, and uncertainty markers.

A context-relative CLIO operator selects quotient
\[
q_{\alpha_i}
\in
\mathfrak{Q}_{t_i}
\]
according to the current task family:
\[
z_i
=
q_{\alpha_i}
\left(
\widehat x_i,Q_i
\right).
\]
The quotient creates equivalence relation
\[
x\sim_{\alpha_i}y
\iff
q_{\alpha_i}(x)=q_{\alpha_i}(y).
\]
Its admissibility is not inferred from its existence. It is represented by a posterior sufficiency estimate
\[
\chi_i
=
\Prob\!\left(
f(x)=f(y)
\text{ for all admitted }f
\given
x\sim_{\alpha_i}y,\mathcal{D}_{t_i}
\right).
\]

Archive comparison computes
\[
\mathcal{K}_{i}^{\mathrm{ret}}
=
\operatorname{Retrieve}
\left(
z_i,
\K_{t_i}
\right)
\]
and decomposes the claim set into
\[
Q_i
=
Q_i^{\mathrm{answered}}
\sqcup
Q_i^{\mathrm{related}}
\sqcup
Q_i^{\mathrm{unresolved}}
\sqcup
Q_i^{\mathrm{uncertain}}.
\]
The answered component is
\[
a_i
=
Q_i^{\mathrm{answered}},
\]
while the preliminary residual is
\[
r_i^{0}
=
Q_i^{\mathrm{unresolved}}
\cup
Q_i^{\mathrm{uncertain}}.
\]

The required concept set is
\[
R_i
=
R(x_i)
=
\left\{
v\in V_t:
\operatorname{sim}(x_i,v)\geq\tau_v
\right\}.
\]
Given demonstrated competencies \(C_i(t)\), the missing prerequisite set is
\[
M_i
=
R_i
\setminus
\closure(C_i(t)).
\]
The route planner selects
\[
P_i^{*}
\in
\argmin_{P}
\left[
\operatorname{cost}(P)
+
\alpha\operatorname{redundancy}(P)
+
\beta\operatorname{dropout\text{-}risk}(P)
+
\gamma\operatorname{projection\text{-}risk}(P)
\right]
\]
subject to
\[
M_i
\subseteq
\closure(C_i(t)\cup P).
\]

Participant calibration is sequential. At conversational time \(s\), the state posterior is
\[
p_i^{s}(\theta)
=
\Prob(
\thetaState_i(s)=\theta
\mid
y_{1:s}
).
\]
The next prompt is selected by
\[
q_{s+1}^{*}
\in
\argmax_q
\left[
\IG(q)
+
\alpha\Prog(q)
-
\beta\Burden(q)
+
\gamma\operatorname{RepairValue}(q)
\right].
\]
The updated state satisfies
\[
p_i^{s+1}(\theta)
\propto
\Prob(
y_{s+1}\mid\theta,q_{s+1}^{*}
)
p_i^{s}(\theta).
\]

Developmental transformation yields
\[
T_i
=
T(
x_i,
\K_t,
G_t,
\thetaState_i(t),
q_{\alpha_i}
)
=
\left(
a_i,
p_i,
r_i,
u_i,
\omega_i,
\Pi_i
\right),
\]
where \(r_i\) is the developed residual, \(u_i\) is classification uncertainty, \(\omega_i\) is the gluing-obstruction state, and \(\Pi_i\) is the provenance-bearing transformation record.

\subsection{The HYDRA transition}

The developed submission is assembled into HYDRA state
\[
Z_i(t)
=
\left(
P_i(t),
R_i(t),
M_i(t)
\right),
\]
where \(P_i\) is the PERSCEN perceptual state, \(R_i\) is the RAT relevance state, and \(M_i\) is the CoM memory state. The transition is
\[
Z_i(t+1)
=
\mathcal{H}_t
\left(
P_i(t),
R_i(t),
M_i(t);
\FieldPhi_t,
\FieldV_t,
\FieldS_t
\right).
\]

A more explicit sequential form is
\[
Z_i(t+1)
=
\operatorname{GLU}_{\RSVP}
\circ
\mathcal{M}
\circ
\mathcal{T}
\circ
\mathcal{F}_{a}
\circ
\mathcal{G}_{a}
\circ
\rho_c
\left(
Z_i(t)
\right).
\]
Each component changes the admissibility conditions of the next. Consequently the composition cannot generally be replaced by an unordered collection of independent lookups.

Operational coupling is measured relative to a separable baseline
\[
\mathcal{H}_0(P,R,M)
=
h\!\left(
h_P(P),
h_R(R),
h_M(M)
\right).
\]
The coupling residual is
\[
\mathcal{E}_{\mathcal{H}}(P,R,M)
=
\mathcal{H}(P,R,M)-\mathcal{H}_0(P,R,M),
\]
and its expected magnitude is
\[
\Gamma_{\mathcal{H}}
=
\E
\left[
\left\|
\mathcal{E}_{\mathcal{H}}(P,R,M)
\right\|^2
\right].
\]
Operational influence requires
\[
\Gamma_{\mathcal{H}}>0.
\]
Operational usefulness requires the stronger condition
\[
\Delta_{\mathrm{coupling}}
=
\E[
\mathcal{L}_{\mathrm{task}}(\mathcal{H}_0)
]
-
\E[
\mathcal{L}_{\mathrm{task}}(\mathcal{H})
]
>0.
\]

The combined state must also satisfy admissibility:
\[
\mathcal{H}(P,R,M;\FieldPhi,\FieldV,\FieldS)
\in
\Adm_{t+1}.
\]
A transition can be nonseparable yet harmful. Coupling is a structural property. Beneficial coupling is a task-relative empirical claim.

\subsection{The obstruction layer}

Let
\[
\mathfrak{U}_i
=
\{U_{\alpha}\}_{\alpha\in A_i}
\]
be a cover of the submission's representational domain. Local reconstructions are sections
\[
s_{\alpha}\in\mathcal{F}(U_{\alpha}).
\]
On overlaps,
\[
\delta_{\alpha\beta}
=
s_{\alpha}|_{U_{\alpha\beta}}
-
s_{\beta}|_{U_{\alpha\beta}}.
\]
If
\[
\delta_{\alpha\beta}=0
\]
for every nonempty overlap and higher compatibility conditions hold, a global section
\[
s\in\mathcal{F}\!\left(
\bigcup_{\alpha}U_{\alpha}
\right)
\]
exists with
\[
s|_{U_\alpha}=s_\alpha.
\]

If local compatibility fails, the system computes obstruction state
\[
\omega_i
=
\left(
\{\delta_{\alpha\beta}\},
[\delta_i],
\operatorname{loc}_i,
\operatorname{stability}_i
\right),
\]
where \([\delta_i]\) is the relevant cohomological class when defined, \(\operatorname{loc}_i\) localizes the failure, and \(\operatorname{stability}_i\) records whether the obstruction persists under admissible refinement.

The obstruction induces a branching decision:
\[
\Omega_i
=
\begin{cases}
\operatorname{repair},&
[\delta_i]=0
\text{ after missing transition data are supplied},
\\
\operatorname{expand},&
[\delta_i]\neq0
\text{ under }q
\text{ but vanishes under refinement }q',
\\
\operatorname{escalate\ obstruction},&
[\delta_i]\neq0
\text{ and remains stable under justified refinements},
\\
\operatorname{reject\ incoherence},&
\text{local claims fail their own admissibility conditions}.
\end{cases}
\]
The final branch is not inferred merely from nonzero obstruction. It requires failure that cannot be interpreted as a meaningful global impossibility, missing coordinate, or representation-induced artifact.

\subsection{Repair and expansion as distinct control actions}

Let
\[
\mathfrak{O}_i
=
\mathfrak{R}_i
\cup
\mathfrak{X}_i
\cup
\mathfrak{G}_i
\]
be the set of repair, representation-expansion, and obstruction-localization operators.

Repair acts within a representation:
\[
\mathcal{R}:
(z_i,\thetaState_i,G)
\to
(z_i,\thetaState_i',G).
\]
Expansion changes the representation:
\[
\mathcal{X}:
(z_i,q)
\to
(z_i',q'),
\qquad
q'\succeq q.
\]
Obstruction localization changes the information available about failed assembly:
\[
\mathcal{G}:
(\omega_i,\mathfrak{U}_i)
\to
(\omega_i',\operatorname{loc}_i).
\]

The optimal pre-escalation operator is
\[
\Omega_i^{*}
\in
\argmax_{\Omega\in\mathfrak{O}_i}
\left[
\Delta S_{ir}(\Omega)
+
\Delta C_{ir}^{T}(\Omega)
+
\Delta V_i^{\mathrm{future}}(\Omega)
-
k_i(\Omega)
\right].
\]
The system reaches the transformative stopping boundary when
\[
\sup_{\Omega\in\mathfrak{O}_i}
\E\!\left[
V_{ir}(\Omega(\sigma_i))
-
V_{ir}(\sigma_i)
-
k_i(\Omega)
\given
\sigma_i
\right]
\leq0.
\]

At this boundary, one of four dispositions is appropriate:
\[
d_i
\in
\{
\operatorname{resolve},
\operatorname{reroute},
\operatorname{market},
\operatorname{escalate}
\}.
\]
Resolution returns an archive-mediated answer. Rerouting identifies another recipient, curriculum, or domain. Market routing offers an authorized moniform residual to a competent finisher. Escalation assigns scarce reciprocal attention.

\subsection{Recipient matching and escalation}

For candidate recipient \(r\), define the available information
\[
\mathscr{I}_{ir}(t)
=
\sigma\!\left(
x_i,
\widehat x_i,
a_i,
p_i,
r_i,
u_i,
\omega_i,
\Pi_i,
\K_r(t),
\thetaState_i(t)
\right).
\]
Recipient-relative value is
\[
\widehat v_{ir}(t)
=
\E[
V_{ir}(t)
\mid
\mathscr{I}_{ir}(t)
].
\]
Joint complementarity is
\[
\widehat{\Delta}_{ir}(t)
=
\E[
\Delta_{\mathrm{joint}}(x_i,r,t)
\mid
\mathscr{I}_{ir}(t)
].
\]
The recipient match is
\[
r_i^{*}
\in
\argmax_{r\in\mathcal{R}_i}
\left[
\widehat v_{ir}
+
\omega\widehat{\Delta}_{ir}
-
c_{ir}
-
h_{ir}
\right],
\]
where \(h_{ir}\) is expected interaction harm or mismatch cost.

For fixed recipient \(r\), define score
\[
s_{ir}
=
\widehat v_{ir}
+
\omega\widehat{\Delta}_{ir}
+
\gamma u_i
+
\delta e_{ir}^{x}
+
\eta n_{ir}
+
\zeta B_{ir}
+
\chi o_{ir}
-
\psi h_{ir}.
\]
The interval allocation problem is
\[
\begin{aligned}
\text{maximize}\quad&
F_r(\mathcal{S})
\\
\text{subject to}\quad&
\sum_{i\in\mathcal{S}}c_{ir}
\leq A_r,
\end{aligned}
\]
where \(F_r\) may be additive, submodular, or synergy-sensitive.

The recipient-facing workload is stable only when
\[
\lambda_r
\Prob(E_i=1)
\E[C_{ir}\mid E_i=1]
<
A_r.
\]
The attention budget is partitioned:
\[
A_r=A_{r,p}+A_{r,x},
\qquad
A_{r,x}\geq\xi_rA_r,
\]
where \(A_{r,p}\) supports predicted value and \(A_{r,x}\) supports exploration.

The recipient's consumption rate is governed separately by pacer
\[
\mathcal{P}_{r,t}:
\mathcal{Q}_r(t)
\to
\mathcal{V}_r(t).
\]
If the escalated workload arrival rate is \(\lambda_{e,r}\), stable consumption requires
\[
\lambda_{e,r}
\E[C_{ir}\mid E]
<
A_r^{\mathrm{view}}.
\]
Admission stability does not imply consumption stability unless both inequalities hold.

\subsection{Audit, appeal, and retroactive repair}

For non-escalated case \(i\), let audit indicator be
\[
Z_i^{A}
\sim
\operatorname{Bernoulli}(\pi_i^{A}),
\qquad
\pi_i^{A}>0.
\]
Let \(\widetilde V_i\) be the audit label. The inverse-probability estimator of valuable suppressions is
\[
\widehat T_{VD}
=
\sum_{i:D_i=1}
\frac{
Z_i^{A}\widetilde V_i
}{
\pi_i^{A}
}.
\]
The estimated false-suppression rate is
\[
\widehat\psi
=
\frac{
\widehat T_{VD}
}{
\widehat T_{VD}+T_{VE}+T_{VA}+T_{VM}
},
\]
where \(T_{VE}\) counts valuable escalations, \(T_{VA}\) valuable archive resolutions, and \(T_{VM}\) valuable mediated developments.

The operative safety statement is posterior:
\[
\Prob(
\psi>\eps
\mid
\mathcal{D}_{\Audit}
)
\leq
\delta_{\mathrm{risk}}.
\]
It is not a claim that the population rate is directly observed.

An appeal acts on full decision state:
\[
\sigma_i'
=
\mathfrak{a}(\sigma_i,\chi_i).
\]
Substantive appealability requires
\[
\exists\chi_i:
\Prob(
d_i'\neq d_i
\mid
\sigma_i,\chi_i
)>0.
\]
An audit-discovered structural failure generates architecture update
\[
\Theta_{t+1}
=
\mathcal{U}_{\Audit}
\left(
\Theta_t,
\zeta_i
\right),
\]
where \(\zeta_i\) may revise the quotient, graph, archive, participant-state model, coupling operator, score, or audit sampler.

If the update changes historical decisions, define discrepancy set
\[
\mathcal{B}_{t+1}
=
\left\{
j:
D_{\Theta_t}(\sigma_j)
\neq
D_{\Theta_{t+1}}(\sigma_j)
\right\}.
\]
Retroactive review is allocated according to expected correction value under finite replay budget.

\subsection{Collective recurrence and archive evolution}

Individual transformed records enter marked process
\[
\mathcal{N}_t
=
\sum_{i:t_i\leq t}
\delta_{(t_i,z_i,m_i)}.
\]
The collective detector computes
\[
\mathcal{D}_{\mathrm{coll}}
\left(
\mathcal{N}_{[t-\Delta,t]},
\mathcal{N}_{\mathrm{ref}}
\right)
=
\left(
\nu_d(t),
\widehat\lambda_t,
n_{\mathrm{eff},t},
\mathcal{C}_t,
\Omega_t,
L_{\mathrm{obs},t}
\right),
\]
where \(\nu_d\) is distributional novelty, \(\widehat\lambda_t\) is the estimated arrival intensity, \(n_{\mathrm{eff},t}\) is independence-adjusted recurrence, \(\mathcal{C}_t\) is co-occurrence topology, \(\Omega_t\) is collective obstruction structure, and \(L_{\mathrm{obs},t}\) records observation loss.

The collective residual is
\[
R_{\mathrm{coll}}(t)
=
\mathcal{G}_{\mathrm{coll}}
\left(
\nu_d(t),
\widehat\lambda_t,
n_{\mathrm{eff},t},
\mathcal{C}_t,
\Omega_t
\right).
\]
It can trigger an answer-node repair, prerequisite-edge insertion, quotient refinement, observation-policy change, audit expansion, or collective escalation packet.

Archive evolution is
\[
\K_{t+1}
=
\mathcal{K}
\left(
\K_t,
U_t^{\mathrm{ind}},
U_t^{\mathrm{coll}},
U_t^{\mathrm{audit}}
\right).
\]
Provenance evolves monotonically:
\[
\mathcal{P}_{t}
\subseteq
\mathcal{P}_{t+1},
\]
while epistemic status remains nonmonotonic:
\[
\K_t\models q
\quad\not\Longrightarrow\quad
\K_{t+1}\models q.
\]

\subsection{The global transition system}

The architecture can now be represented as controlled stochastic transition
\[
\Xi_{t+1}
=
\mathfrak{F}
\left(
\Xi_t,
X_{t:t+1},
A_t,
W_t
\right),
\]
where \(X_{t:t+1}\) is the new submission batch, \(A_t\) is the authorized action, and \(W_t\) is exogenous uncertainty.

The action space includes
\[
A_t
\in
\left\{
\Ans,
\Req,
\operatorname{diagnose},
\operatorname{repair},
\operatorname{expand},
\operatorname{localize},
\operatorname{test},
\operatorname{reroute},
\operatorname{market},
\Esc,
\Audit,
\operatorname{appeal},
\operatorname{aggregate},
\operatorname{revise},
\operatorname{silence}
\right\}.
\]
These are not equally available at every state. Let
\[
\mathfrak{A}(\Xi_t)
\]
be the admissible action set. Governance requires
\[
A_t\in\mathfrak{A}(\Xi_t).
\]

Define immediate reward
\[
\begin{aligned}
\mathfrak{r}(\Xi_t,A_t)
={}&
V_t^{\mathrm{resolved}}
+
V_t^{\mathrm{learning}}
+
V_t^{\mathrm{joint}}
+
V_t^{\mathrm{repair}}
+
V_t^{\mathrm{audit}}
\\
&
-
C_t^{\mathrm{participant}}
-
C_t^{\mathrm{recipient}}
-
C_t^{\mathrm{computation}}
-
H_t^{\mathrm{suppression}}
-
H_t^{\mathrm{privacy}}
-
H_t^{\mathrm{governance}}.
\end{aligned}
\]
The constrained value function is
\[
J^{*}(\Xi)
=
\sup_{\pi}
\E_{\pi}
\left[
\sum_{t=0}^{\infty}
\beta^t
\mathfrak{r}(\Xi_t,A_t)
\given
\Xi_0=\Xi
\right]
\]
subject to workload, audit, privacy, provenance, exploration, and governance constraints.

The Bellman equation is
\[
J^{*}(\Xi)
=
\sup_{a\in\mathfrak{A}(\Xi)}
\left[
\mathfrak{r}(\Xi,a)
+
\beta
\E[
J^{*}(\Xi')
\mid
\Xi,a
]
\right].
\]
An exact solution is generally infeasible because \(\Xi\) contains latent participant states, changing archives, graph-valued structures, variable recipient capacities, and population-level processes. The equation serves as a declaration of what a locally optimized component omits. A routing rule that minimizes current evaluation cost can increase future archive defects. A diagnostic prompt that maximizes information gain can increase participant dropout. A quotient that minimizes present description length can increase future repair cost. A greedy escalation decision can destroy portfolio complementarity.

\subsection{The principal invariants}

Let a trajectory be
\[
\Xi_0
\to
\Xi_1
\to
\cdots.
\]
The provenance invariant is
\[
\mathcal{P}_t
\subseteq
\mathcal{P}_{t+1}.
\]
The attribution invariant requires that every stored claim preserve origin type:
\[
\pi(q,t+1)
\supseteq
\pi(q,t).
\]
The workload invariant is
\[
\lambda_r(t)
\Prob(E_t)
\E[C_t\mid E_t]
<
A_r(t).
\]
The exploration invariant is
\[
A_{r,x}(t)\geq\xi_rA_r(t).
\]
The audit-positivity invariant is
\[
\pi_{g,t}^{A}\geq\pi_{\min,g}>0
\]
for every eligible stratum.

The assent-independence invariant is
\[
\operatorname{pass}_i(v)
\ind
A_i(v)
\mid
C_i(v).
\]
The appeal invariant is
\[
d_i\in\mathcal{D}_{\mathrm{adverse}}
\Longrightarrow
\mathfrak{A}_i^{\mathrm{appeal}}\neq\varnothing.
\]
The representation-accounting invariant is
\[
q_t(x)=q_t(y)
\Longrightarrow
\left[
\operatorname{TaskEquivalent}_t(x,y)
\lor
\operatorname{CollisionRiskRecorded}_t(x,y)
\right].
\]
The obstruction-accounting invariant is
\[
\omega_i\neq0
\Longrightarrow
\left[
\operatorname{localized}(\omega_i)
\lor
\operatorname{uncertaintyRecorded}(\omega_i)
\right].
\]
The collective-preservation invariant requires
\[
\text{individual resolution}
\not\Longrightarrow
\text{deletion from aggregate recurrence statistics}.
\]

\subsection{A conditional stability theorem}

\begin{theorem}[Recipient-facing workload stability]
Fix recipient \(r\). Suppose transformed submissions arrive according to a stationary ergodic marked point process with rate \(\lambda_r\). Let \(E_i\) indicate escalation and let \(C_{ir}\) be the required recipient attention, with
\[
0<
\E[C_{ir}\mid E_i=1]
<
\infty.
\]
Suppose recipient service capacity has long-run rate \(A_r>0\), and suppose
\[
\lambda_r
\Prob(E_i=1)
\E[C_{ir}\mid E_i=1]
<
A_r.
\]
Under a work-conserving service policy and ordinary regenerative assumptions, the recipient-facing workload queue is positive recurrent.
\end{theorem}

\begin{proof}
The long-run admitted workload rate is
\[
\lambda_r
\E[E_iC_{ir}]
=
\lambda_r
\Prob(E_i=1)
\E[C_{ir}\mid E_i=1].
\]
By assumption this is strictly smaller than service capacity \(A_r\). Let the capacity slack be
\[
\delta_r
=
A_r
-
\lambda_r
\Prob(E_i=1)
\E[C_{ir}\mid E_i=1]
>0.
\]
Consider the workload sampled at regeneration times. Outside a bounded set, expected service during a regeneration cycle exceeds expected admitted workload by an amount proportional to \(\delta_r\). The workload therefore has negative drift outside that bounded set. A Foster-Lyapunov criterion with workload as Lyapunov function gives positive recurrence, subject to the stated integrability and regeneration assumptions.
\end{proof}

This theorem guarantees queue stability, not intellectual adequacy. A policy that escalates nothing satisfies the inequality. The remaining guarantees prevent that degenerate solution from being presented as success.

\subsection{A conditional nonclosure theorem}

\begin{theorem}[Statistical nonclosure under positive exploration]
Suppose every eligible audit stratum \(g\) has inclusion probability
\[
\pi_g^{A}\geq\pi_{\min,g}>0,
\]
and every eligible escalation region has exploratory review probability
\[
\pi_g^{x}\geq\pi_{\min,g}^{x}>0.
\]
If stratum \(g\) produces infinitely many eligible cases, then with probability one infinitely many cases from \(g\) are eventually audited and infinitely many are eventually explored, assuming conditionally independent sampling or an ergodic sampling design with the same positive lower bounds.
\end{theorem}

\begin{proof}
Let \(Z_{g,n}^{A}\) indicate audit inclusion for the \(n\)-th eligible case from \(g\). Under conditional independence,
\[
\sum_{n=1}^{\infty}
\Prob(Z_{g,n}^{A}=1)
\geq
\sum_{n=1}^{\infty}
\pi_{\min,g}
=
\infty.
\]
The second Borel-Cantelli lemma implies that
\[
Z_{g,n}^{A}=1
\]
infinitely often with probability one. The same argument applies to exploratory review. Under an ergodic design, the positive long-run inclusion frequency yields the same conclusion by the ergodic theorem.
\end{proof}

The theorem establishes observational nonclosure, not fair outcomes. A region can be sampled and still systematically misjudged. It nevertheless prevents the strongest form of self-sealing exclusion in which a region is never reviewed and its absence of successes is then treated as evidence of worthlessness.

\subsection{A conditional quotient-repair theorem}

\begin{theorem}[Repairability under retained collision witnesses]
Let
\[
q:\mathcal{X}\to\mathcal{Z}
\]
be a current quotient, let \(f'\) be a later task, and suppose there exist \(x,y\) such that
\[
q(x)=q(y)
\]
but
\[
f'(x)\neq f'(y).
\]
Suppose provenance retains a witness variable
\[
w:\mathcal{X}\to\mathcal{W}
\]
with
\[
w(x)\neq w(y),
\]
and suppose the refined representation
\[
q'(z)
=
(q(z),w(z))
\]
is admissible. Then the specific collision between \(x\) and \(y\) is removed under \(q'\).
\end{theorem}

\begin{proof}
Since
\[
w(x)\neq w(y),
\]
one has
\[
q'(x)
=
(q(x),w(x))
\neq
(q(y),w(y))
=
q'(y).
\]
Thus \(x\) and \(y\) lie in distinct fibers of \(q'\), so the collision created by \(q\) is removed.
\end{proof}

The theorem is deliberately local. The refinement may still collapse other task-relevant distinctions. Repair proceeds by witnessed expansion rather than by claiming immediate universal sufficiency.

\subsection{A conditional routing-sufficiency theorem}

\begin{theorem}[Exact downstream preservation under functional factorization]
Let
\[
q:\mathcal{X}\to\mathcal{Z}
\]
be a projection and let
\[
\mathcal{F}
=
\{f_j:\mathcal{X}\to\mathcal{Y}_j\}_{j\in J}
\]
be the declared downstream task family. Suppose that for every \(j\in J\) there exists
\[
g_j:\mathcal{Z}\to\mathcal{Y}_j
\]
such that
\[
f_j=g_j\circ q.
\]
Then replacing \(x\) by \(q(x)\) preserves every downstream task in \(\mathcal{F}\) exactly.
\end{theorem}

\begin{proof}
For every \(x\in\mathcal{X}\) and \(j\in J\),
\[
g_j(q(x))
=
(g_j\circ q)(x)
=
f_j(x).
\]
Therefore every task output is unchanged by projection.
\end{proof}

This is the exact JPEG-like case generalized to an arbitrary declared function family. Outside finite exhaustively verifiable domains, establishing the premise is difficult. The theorem clarifies what must be proved rather than allowing the projection to certify itself.

\subsection{A conditional obstruction theorem}

\begin{theorem}[Global assembly under vanishing compatibility obstruction]
Let \(\mathcal{F}\) be a sheaf over domain \(X\), let
\[
\mathfrak{U}=\{U_\alpha\}
\]
be an open cover, and let
\[
s_\alpha\in\mathcal{F}(U_\alpha)
\]
be local sections satisfying
\[
s_\alpha|_{U_\alpha\cap U_\beta}
=
s_\beta|_{U_\alpha\cap U_\beta}
\]
for every \(\alpha,\beta\). Then there exists a unique global section
\[
s\in\mathcal{F}(X)
\]
such that
\[
s|_{U_\alpha}=s_\alpha
\]
for every \(\alpha\).
\end{theorem}

\begin{proof}
Existence follows from the gluing axiom of the sheaf. Uniqueness follows from locality: if \(s\) and \(s'\) have identical restrictions to every member of a cover, then
\[
s=s'.
\]
\end{proof}

The theorem identifies the exact condition under which local conversational reconstructions can be treated as one global claim. When the compatibility condition fails, global assembly is not licensed. The failure may become the residual rather than being smoothed away.

\subsection{A conditional audit-consistency theorem}

\begin{theorem}[Consistency of estimated false suppression]
Suppose non-escalated cases are sampled for audit with known probabilities satisfying
\[
\pi_i^{A}\geq\pi_{\min}>0.
\]
Suppose audit labels are conditionally unbiased for horizon-relative value, sampling moments are finite, and the number of valuable cases grows proportionally with the eligible population. Then the inverse-probability ratio estimator
\[
\widehat\psi_n
=
\frac{
\widehat T_{VD,n}
}{
\widehat T_{VD,n}+T_{V\bar D,n}
}
\]
converges in probability to the finite-population false-suppression ratio as \(n\to\infty\).
\end{theorem}

\begin{proof}
Inverse-probability estimation gives
\[
\frac{
\widehat T_{VD,n}-T_{VD,n}
}{n}
\xrightarrow{p}0.
\]
By proportional growth,
\[
\frac{
T_{VD,n}+T_{V\bar D,n}
}{n}
\xrightarrow{p}\vartheta>0.
\]
The ratio map
\[
g(a,b)=\frac{a}{a+b}
\]
is continuous where the denominator is positive. The continuous mapping theorem yields
\[
\widehat\psi_n
-
\frac{
T_{VD,n}
}{
T_{VD,n}+T_{V\bar D,n}
}
\xrightarrow{p}0.
\]
\end{proof}

The result depends on audit-label adequacy and positive coverage. It does not convert counterfactual value into a directly observed property.

\subsection{A conditional recurrence-detection theorem}

\begin{theorem}[Consistent rate-shift detection under loss-corrected observation]
Suppose true events of type \(v\) follow a Poisson process with rate \(\lambda_v\). Suppose every event is observed independently with known probability
\[
\pi_v\in(0,1],
\]
giving observed rate
\[
\lambda_v^{\mathrm{obs}}=\pi_v\lambda_v.
\]
Let
\[
\widehat\lambda_v
=
\frac{N_v^{\mathrm{obs}}(T)}{\pi_vT}.
\]
Then
\[
\widehat\lambda_v
\xrightarrow{\mathrm{a.s.}}
\lambda_v
\]
as \(T\to\infty\). Consequently every fixed rate shift
\[
|\lambda_v-\lambda_{0,v}|>\delta
\]
can be detected consistently by a threshold test whose tolerance converges to zero more slowly than \(T^{-1/2}\).
\end{theorem}

\begin{proof}
Independent thinning of a Poisson process produces a Poisson process with rate \(\pi_v\lambda_v\). By the strong law for Poisson processes,
\[
\frac{
N_v^{\mathrm{obs}}(T)
}{T}
\xrightarrow{\mathrm{a.s.}}
\pi_v\lambda_v.
\]
Division by \(\pi_v>0\) gives
\[
\widehat\lambda_v
\xrightarrow{\mathrm{a.s.}}
\lambda_v.
\]
For a fixed nonzero shift, consistency of the estimator implies that an acceptance band shrinking to zero will eventually exclude the reference rate with probability approaching one, provided it shrinks slowly enough to accommodate sampling fluctuations.
\end{proof}

If the observation probability is zero in a region or unknown because of deterministic batch truncation, the conclusion fails. Observation-channel governance is therefore part of recurrence validity.

\subsection{A conditional governance-preservation theorem}

\begin{theorem}[Invariant preservation under governance-admissible updates]
Let
\[
\mathfrak{I}
=
\{
I_1,\ldots,I_m
\}
\]
be governance invariants. Suppose
\[
I_j(\Xi_0)=1
\]
for every \(j\), and every authorized transition
\[
\Xi_{t+1}
=
\Gamma_{a_t}(\Xi_t)
\]
satisfies
\[
\bigwedge_{j=1}^{m}I_j(\Xi_t)=1
\Longrightarrow
\bigwedge_{j=1}^{m}I_j(\Xi_{t+1})=1.
\]
Then every invariant holds along every authorized finite trajectory.
\end{theorem}

\begin{proof}
The result follows by induction. The invariants hold at \(t=0\). Assume they hold at time \(t\). Governance admissibility implies they hold at \(t+1\). Hence they hold for every finite \(t\).
\end{proof}

The theorem is formally simple because the substantive work lies in specifying invariants and verifying each transition. It nevertheless establishes that governance constraints must be transition properties, not periodic aspirations.

\subsection{The principal integrated guarantee}

The earlier results can be composed only conditionally. Queue stability does not imply audit validity. Audit validity does not imply quotient sufficiency. Quotient sufficiency does not imply beneficial HYDRA coupling. Beneficial coupling does not imply legitimate governance. The full guarantee must state every dependency.

\begin{theorem}[Conditional viability of the integrated architecture]
Consider an attention compiler with CLIO projection, HYDRA assembly, prerequisite-sensitive development, obstruction localization, recipient-relative escalation, positive exploratory review, audited suppression estimation, collective recurrence detection, and governance-constrained updates. Suppose the following conditions hold:
\[
f=g_f\circ q
\]
for every presently admitted downstream task \(f\);
\[
\mathcal{H}(P,R,M)\in\Adm
\]
for every admitted HYDRA transition;
\[
\Delta_{\mathrm{coupling}}\geq0;
\]
\[
\lambda_r
\Prob(E)
\E[C\mid E]
<
A_r
\]
for every protected recipient;
\[
A_{r,x}\geq\xi_rA_r>0;
\]
\[
\pi_{g}^{A}\geq\pi_{\min,g}>0
\]
for every eligible audit stratum;
\[
\Prob(
\psi_r>\eps
\mid
\mathcal{D}_{\Audit}
)
\leq
\delta_{\mathrm{risk}};
\]
\[
\mathrm{PC}(\Pi_i)\geq p_0
\]
for every decision-bearing record;
\[
\operatorname{pass}_i(v)
\ind
A_i(v)
\mid
C_i(v);
\]
every adverse decision has a substantively capable appeal path;
every representation change preserves a collision-witness or declares the relevant information irrecoverable;
every stable gluing obstruction is preserved rather than silently normalized away;
and every architecture update is governance-admissible.

Then the architecture has the following conditional properties. Presently admitted task outputs are preserved under projection. Admitted global states arise only from compatible local assemblies or explicitly represented obstructions. Recipient-facing queues are stable. No eligible audit stratum is observationally closed. Estimated false suppression remains within the declared posterior risk tolerance. Adverse decisions remain contestable. Provenance remains reconstructible. Competent dissent is not excluded merely for being dissent. Distributional recurrence survives individual archival resolution. Architecture updates preserve the declared governance invariants.
\end{theorem}

\begin{proof}
Task preservation follows from functional factorization through \(q\). Admissible global assembly follows from the HYDRA admissibility condition and the sheaf gluing criterion, with nonvanishing obstructions retained explicitly. Queue stability follows from the strict workload inequalities. Positive exploration and positive audit inclusion prevent observational closure by the nonclosure theorem. The posterior suppression inequality gives the declared risk control relative to the audit model. Provenance completeness and versioning permit decision reconstruction. Assent independence ensures that demonstrated comprehension, rather than agreement, controls prerequisite passage. Substantive appealability preserves a state transition capable of revising adverse decisions. Collision-witness retention permits local quotient repair when a later task exposes insufficiency. Collective-preservation rules prevent resolved individual records from disappearing from population statistics. Governance-invariant preservation follows inductively from governance-admissible transitions. Each conclusion is conditional on its corresponding premise, and no conclusion supplies a premise omitted elsewhere.
\end{proof}

\subsection{What the theorem does not establish}

The integrated theorem does not prove that every valuable residual is recognized. It does not prove that the archive is true, that the participant-state model is psychologically complete, that the prerequisite graph is culturally neutral, that HYDRA coupling corresponds to a biological or physical mechanism, that RSVP field semantics are uniquely correct, or that the recipient's declared priorities are just.

It does not prove sufficiency for undeclared future tasks. The factorization condition applies to the present admitted task family. It does not prove that audit labels reveal intrinsic intellectual value. They reveal value relative to an operational criterion, recipient, horizon, and review process. It does not prove that positive sampling probability gives adequate finite-sample power. It prevents complete exclusion while possibly leaving large uncertainty. It does not prove that every obstruction is meaningful. It proves only that a genuine failure of compatibility should not be replaced by fictional global coherence.

It does not make the architecture neutral. The quotient, task family, archive, graph, score weights, cost model, risk tolerances, privacy limits, and governance invariants remain authored choices. The theorem shows what follows if those choices and empirical premises hold. It does not derive them from first principles.

\subsection{Representational simplicity as a variational principle}

The complete theory can be compressed into a variational problem. Let \(\rho\) denote a representation and let \(\pi\) denote a transformation and routing policy. Define representational complexity
\[
C_{\mathrm{rep}}(\rho),
\]
recipient attention cost
\[
C_{\mathrm{att}}(\rho,\pi),
\]
participant burden
\[
C_{\mathrm{part}}(\rho,\pi),
\]
task loss
\[
L_{\mathrm{task}}(\rho,\pi),
\]
suppression risk
\[
R_{\mathrm{sup}}(\rho,\pi),
\]
obstruction-concealment risk
\[
R_{\mathrm{obs}}(\rho,\pi),
\]
privacy risk
\[
R_{\mathrm{priv}}(\rho,\pi),
\]
and governance irreversibility
\[
R_{\mathrm{irr}}(\rho,\pi).
\]

The unconstrained functional is
\[
\begin{aligned}
\mathcal{J}(\rho,\pi)
={}&
C_{\mathrm{rep}}(\rho)
+
\alpha C_{\mathrm{att}}(\rho,\pi)
+
\beta C_{\mathrm{part}}(\rho,\pi)
+
\gamma L_{\mathrm{task}}(\rho,\pi)
\\
&
+
\delta R_{\mathrm{sup}}(\rho,\pi)
+
\eta R_{\mathrm{obs}}(\rho,\pi)
+
\zeta R_{\mathrm{priv}}(\rho,\pi)
+
\kappa R_{\mathrm{irr}}(\rho,\pi).
\end{aligned}
\]
Representational simplicity is
\[
(\rho^{*},\pi^{*})
\in
\argmin_{\rho,\pi}
\mathcal{J}(\rho,\pi)
\]
subject to the architecture's explicit stability, audit, provenance, appeal, exploration, and governance constraints.

The representation is not simplest when it has the fewest coordinates. It is simplest when no retained distinction costs more than the downstream loss its removal would create, and no discarded distinction creates an unpriced obligation elsewhere.

For candidate distinction \(d\), let
\[
\Delta C(d)
\]
be the cost of retaining it and let
\[
\Delta L(d)
\]
be the expected increase in total loss if it is removed. Retention is locally optimal when
\[
\Delta L(d)>\Delta C(d).
\]
Removal is locally optimal when
\[
\Delta L(d)<\Delta C(d).
\]
At an interior optimum,
\[
\Delta L(d)=\Delta C(d).
\]

The cost terms include more than storage. Retaining a distinction can increase explanation burden, route complexity, privacy exposure, model variance, and recipient attention. Removing it can increase false suppression, repair cost, gluing failure, distributional blindness, and future task insufficiency.

\subsection{CLIO and HYDRA as dual constraints}

CLIO and HYDRA occupy opposite sides of the variational problem. CLIO asks which distinctions can be removed before downstream continuation. HYDRA asks which surviving components must remain jointly operative during continuation.

Let
\[
q_{\CLIO}:\mathcal{X}\to\mathcal{Z}
\]
be the quotient and let
\[
\mathcal{H}_{\HYDRA}:
\mathcal{Z}_P
\times
\mathcal{Z}_R
\times
\mathcal{Z}_M
\to
\mathcal{Z}'
\]
be the assembled transition.

Their composition is
\[
\mathcal{X}
\xrightarrow{q_{\CLIO}}
\mathcal{Z}
\xrightarrow{\Delta}
\mathcal{Z}_P
\times
\mathcal{Z}_R
\times
\mathcal{Z}_M
\xrightarrow{\mathcal{H}_{\HYDRA}}
\mathcal{Z}',
\]
where \(\Delta\) distributes the projected state into its operational components.

CLIO failure occurs when
\[
q_{\CLIO}(x)=q_{\CLIO}(y)
\]
while the required continuation differs:
\[
\mathcal{H}_{\HYDRA}(\Delta q_{\CLIO}(x))
\neq
\mathcal{H}_{\HYDRA}(\Delta q_{\CLIO}(y)).
\]
Exact deterministic composition would prevent this inequality when the projected inputs are identical. In practice, the difference arises because discarded information re-enters through context, hidden state, participant memory, or later task expansion. The apparent equality at the CLIO level was therefore not equality in the full operational state.

HYDRA failure occurs when the representation preserves the necessary distinctions but the assembly operator cannot combine them admissibly:
\[
q_{\CLIO}(x)\text{ is task-sufficient},
\]
yet
\[
\mathcal{H}_{\HYDRA}(P,R,M)\notin\Adm.
\]
The first failure requires expansion. The second requires repair of coupling, transition maps, or global compatibility.

This yields the duality:
\[
\CLIO:
\quad
\text{remove every distinction that continuation does not require},
\]
\[
\HYDRA:
\quad
\text{preserve every coupling that continuation cannot decompose}.
\]
Representational simplicity lies between them. Excess distinction creates burden. Insufficient distinction creates collapse. Excess separation destroys operational coupling. Excess fusion destroys independent variation and auditability.

\subsection{The hierarchy of representational integration}

The hierarchy developed in the companion sufficiency analysis can now be stated inside the complete system:
\[
\text{co-location}
\subsetneq
\text{joint addressing}
\subsetneq
\text{operational coupling}
\subsetneq
\text{globally admissible integration}.
\]

Co-location means components share a storage or display environment. Joint addressing means a single operation retrieves or manipulates them together. Operational coupling means the transition cannot be reproduced by independently evaluating the components and combining their outputs under the proposed factorization. Globally admissible integration additionally requires that the coupled transition satisfy local constraints, overlap compatibility, governance invariants, and downstream task conditions.

The inclusions are strict. Two independently recoverable fields can be co-located without joint addressing. Two fields can share an address while remaining independently decoded. Two components can be operationally coupled while producing an inadmissible transition. A globally admissible integration therefore demands more than compact packing or nonseparability.

Let predicates
\[
C_0(x),C_1(x),C_2(x),C_3(x)
\]
represent these four levels. Then
\[
C_3(x)\Longrightarrow C_2(x)
\Longrightarrow C_1(x)
\Longrightarrow C_0(x),
\]
while the converses fail in general.

\subsection{Failure as a diagnostic output}

A failed transition is not automatically noise. Let
\[
\pi_n^{*}
\]
be the best continuation under representation level \(n\), and let
\[
\pi_{\Adm}
\]
be an admissible target. Convergence
\[
\pi_n^{*}\to\pi_{\Adm}
\]
is desirable only when the target representation preserves the distinctions required by the task.

Suppose TARTAN-like adaptation alternates between two projected regimes:
\[
z_{2k}\to z_A,
\qquad
z_{2k+1}\to z_B.
\]
If
\[
q(z_A)=q(z_B)
\]
but their unprojected continuations differ, the oscillation witnesses quotient insufficiency. Forcing convergence inside the quotient can erase a reachability-relevant distinction.

Let
\[
\mathcal{O}_{\mathrm{osc}}
=
\limsup_{n\to\infty}
\dist(\pi_{n+1}^{*},\pi_n^{*}).
\]
A nonzero value can indicate unstable optimization, but if the oscillation aligns with a discarded coordinate \(w\), so that
\[
\Cov(
\pi_{n+1}^{*}-\pi_n^{*},
w_{n+1}-w_n
)\neq0,
\]
the failure contains diagnostic information about the representation.

The repair decision is
\[
\begin{cases}
\operatorname{repair\ dynamics},&
\text{if oscillation persists within a sufficient representation},
\\
\operatorname{expand\ representation},&
\text{if oscillation tracks a discarded task-relevant distinction}.
\end{cases}
\]
Convergence is not an unconditional success criterion. It is success only after sufficiency has been established relative to the continuation being optimized.

\subsection{The complete preservation, constraint, and continuation law}

The architecture can be summarized by three operators. Let
\[
\operatorname{Preserve}_t
\]
retain distinctions, provenance, and unresolved obstructions needed by the admitted future. Let
\[
\operatorname{Constrain}_t
\]
apply archive knowledge, prerequisite conditions, task admissibility, finite attention, privacy, and governance. Let
\[
\operatorname{Continue}_t
\]
select and execute an admissible next transition.

The state evolves as
\[
\Xi_{t+1}
=
\operatorname{Continue}_t
\circ
\operatorname{Constrain}_t
\circ
\operatorname{Preserve}_t
(\Xi_t,X_t).
\]
The order matters. Constraint applied before preservation can erase the evidence needed to challenge the constraint. Continuation applied before constraint can spend attention on inadmissible states. Preservation without later constraint accumulates unbounded representational burden.

For distinction \(d\), define survival indicator
\[
s_t(d)\in\{0,1\}.
\]
Define future relevance
\[
r_t(d)
=
\Prob(
d\text{ changes an admitted continuation}
\mid
\mathscr{I}_t
).
\]
Define retention cost \(c_t(d)\) and loss from deletion \(\ell_t(d)\). The preservation rule is
\[
s_t(d)
=
\ind\{
\ell_t(d)\geq c_t(d)
\},
\]
subject to privacy and governance constraints. Since \(\ell_t(d)\) is estimated, the rule must reserve uncertainty:
\[
s_t(d)
=
1
\quad
\text{when}
\quad
\Prob(
\ell_t(d)\geq c_t(d)
\mid
\mathscr{I}_t
)
>
1-\delta_d.
\]
A distinction can be stored in protected provenance without being exposed in every recipient-facing packet. Preservation and presentation are separate decisions.

\subsection{Final synthesis}

The original compression question asked which distinctions can be removed, combined, or made indirectly computable without preventing a downstream system from performing its intended operation. The integrated attention question asks the same thing at institutional scale.

The raw many-to-one communication space is too large to deliver directly. CLIO projects it into task-relative admissible coordinates. The archive removes claims already disposed of by prior reasoning. The prerequisite graph identifies distinctions the participant must demonstrate rather than merely encounter. Sequential calibration estimates which developmental intervention will both advance the idea and reduce uncertainty. HYDRA binds perception, relevance, and memory into an operational continuation. Obstruction analysis prevents incompatible local interpretations from being disguised as a coherent global state. Residual extraction preserves what remains unresolved. Recipient-relative allocation spends finite attention where direct complementarity is expected to exceed its opportunity cost. Audit reopens rejected paths so that suppression does not become observationally self-validating. Collective recurrence preserves information carried by the shape of repetition. Governance makes every transformation attributable, contestable, and conditionally reversible.

The integrated transformation is
\[
\begin{aligned}
\mathcal{X}
&\xrightarrow{\mathcal{C}_{\mathrm{char}}}
\widehat{\mathcal{X}}
\xrightarrow{q_{\CLIO}}
\mathcal{Z}
\xrightarrow{\operatorname{Retrieve}_{\K}}
(a,r^{0})
\xrightarrow{\operatorname{Route}_{G,\theta}}
(a,p,r^{1})
\\
&\xrightarrow{\mathcal{H}_{\HYDRA}}
(a,p,r^{2},\omega)
\xrightarrow{\operatorname{Repair/Expand}}
(a,p,r,u,\omega)
\xrightarrow{S(\cdot\mid r,\K_r,\theta,t)}
\mathcal{Q}_r
\\
&\xrightarrow{\mathcal{P}_r}
\operatorname{RecipientAttention}
\xrightarrow{\Audit,\mathfrak{a},\mathcal{D}_{\mathrm{coll}}}
(\K',G',q',\mathcal{H}',S').
\end{aligned}
\]
The final arrow returns to the beginning. Audit, appeal, and recurrence modify the structures through which later submissions will be interpreted. The architecture is therefore reflexive:
\[
\Xi_{t+1}
=
\mathfrak{F}(\Xi_t,X_t),
\]
and \(\mathfrak{F}\) itself is among the objects that later evidence can change.

Representational simplicity does not mean retaining as little as possible. It means retaining exactly those distinctions whose removal would cost more than their preservation, while keeping a path by which mistaken removals can be detected. It does not mean fusing every related subsystem. It means coupling precisely those components whose independent treatment fails to reproduce the admissible transition. It does not mean forwarding only polished novelty. It means transforming submissions until known work, missing prerequisites, representational artifacts, and genuine residuals become distinguishable.

The architecture's strongest defensible promise is conditional:
\[
\begin{aligned}
&
\text{if the declared quotient is sufficient for the admitted task,}
\\
&
\text{if HYDRA coupling produces admissible transitions,}
\\
&
\text{if workload remains below recipient capacity,}
\\
&
\text{if exploration and audit preserve positive coverage,}
\\
&
\text{if provenance permits reconstruction and repair,}
\\
&
\text{if competent dissent is not equated with incompetence,}
\\
&
\text{if obstruction is exposed rather than normalized away,}
\\
&
\text{and if governance invariants survive every update,}
\end{aligned}
\]
then the system can reduce repeated human labor while preserving a bounded, auditable path by which unfamiliar value may still reach the frontier of scarce attention.

The conditions cannot be removed without weakening the conclusion. That dependence is not a defect of the theory. It is the theory's principal refusal. The system does not inherit exact sufficiency from the elegance of a projection, useful coupling from nonseparability, intellectual value from novelty, truth from recurrence, fairness from randomization, or legitimacy from efficiency. Each claim requires its own evidence.

The final variational statement is
\[
\boxed{
\begin{aligned}
\text{minimize}\quad&
\text{representational burden}
+
\text{repeated cognitive labor}
+
\text{recipient attention cost}
\\
&
+
\text{participant burden}
+
\text{repair debt}
+
\text{suppression risk}
+
\text{governance irreversibility}
\\
\text{subject to}\quad&
\text{task-relative sufficiency},
\\
&
\text{admissible operational coupling},
\\
&
\text{localized gluing obstruction},
\\
&
\text{stable finite-attention queues},
\\
&
\text{positive exploratory and audit coverage},
\\
&
\text{provenance completeness},
\\
&
\text{substantive appealability},
\\
&
\text{collective-pattern preservation},
\\
&
\text{privacy constraints},
\\
&
\text{governance invariants}.
\end{aligned}
}
\]

A representation is simple only relative to the continuations it still supports. An archive is thin-walled only if neighboring relevance can cross its boundaries without losing provenance. A conversation is scalable only if repetition can be compressed without erasing recurrence. A gate is legitimate only if it tests understanding without requiring assent. A coupling is useful only if it changes the transition for the better. An escalation is justified only if the particular recipient's attention contributes something the surrounding machinery cannot supply. A rejection is defensible only if the system preserves some method of discovering that it was wrong.

The governing principle of the completed architecture is therefore
\[
\boxed{
\text{compress what can be reconstructed, preserve what can change continuation, expose what cannot be glued, and audit what finite attention leaves unseen}.
}
\]

\newpage

\begin{thebibliography}{99}

\bibitem{anderson2009dialog}
Anderson, M. (2009).
System and Method for Providing an Intelligent Multi-Step Dialog with a User.
U.S. Patent No. 7,539,656 B2.

\bibitem{anderson2011thinwalls}
Anderson, M. (2011).
Graphical User Interface for Chat Room with Thin Walls.
U.S. Patent No. 8,015,246 B1.

\bibitem{anderson2022bubblecity}
Anderson, M. (2022).
\textit{Bubble City: Design Proposal, A Twitter Alternative Which Is Not a Social Medium, It Is a Real Time Idea Router}.
Version 2.0.
Syntience Inc.

\bibitem{anderson2022um1}
Anderson, M. (2022).
Understanding Machine One.
\textit{Experimental Epistemology}, post 9, 25 July 2022.
Syntience Inc.

\bibitem{bang2005}
Bang, H. and Robins, J.~M. (2005).
Doubly robust estimation in missing data and causal inference models.
\textit{Biometrics}, 61(4), 962--973.

\bibitem{bertsekas2017}
Bertsekas, D.~P. (2017).
\textit{Dynamic Programming and Optimal Control}.
Fourth edition.
Athena Scientific, Belmont, Massachusetts.

\bibitem{billingsley1995}
Billingsley, P. (1995).
\textit{Probability and Measure}.
Third edition.
John Wiley \& Sons, New York.

\bibitem{brown2005worldcafe}
Brown, J., Isaacs, D., and the World Caf\'e Community (2005).
\textit{The World Caf\'e: Shaping Our Futures Through Conversations That Matter}.
Berrett-Koehler Publishers, San Francisco.

\bibitem{chouldechova2017}
Chouldechova, A. (2017).
Fair prediction with disparate impact: A study of bias in recidivism prediction instruments.
\textit{Big Data}, 5(2), 153--163.

\bibitem{cochran1977}
Cochran, W.~G. (1977).
\textit{Sampling Techniques}.
Third edition.
John Wiley \& Sons, New York.

\bibitem{corbett1994}
Corbett, A.~T. and Anderson, J.~R. (1994).
Knowledge tracing: Modeling the acquisition of procedural knowledge.
\textit{User Modeling and User-Adapted Interaction}, 4(4), 253--278.

\bibitem{cover2006}
Cover, T.~M. and Thomas, J.~A. (2006).
\textit{Elements of Information Theory}.
Second edition.
John Wiley \& Sons, Hoboken, New Jersey.

\bibitem{curry2014}
Curry, J.~M. (2014).
\textit{Sheaves, Cosheaves and Applications}.
Doctoral dissertation, University of Pennsylvania.

\bibitem{daley2003}
Daley, D.~J. and Vere-Jones, D. (2003).
\textit{An Introduction to the Theory of Point Processes, Volume I: Elementary Theory and Methods}.
Second edition.
Springer, New York.

\bibitem{degyurky2014}
de Gyurky, S.~M. and Tarbell, M.~A. (2014).
\textit{The Autonomous System: A Foundational Synthesis of the Sciences of the Mind}.
John Wiley \& Sons, Hoboken, New Jersey.

\bibitem{durrett2019}
Durrett, R. (2019).
\textit{Probability: Theory and Examples}.
Fifth edition.
Cambridge University Press, Cambridge.

\bibitem{fisher1922}
Fisher, R.~A. (1922).
On the mathematical foundations of theoretical statistics.
\textit{Philosophical Transactions of the Royal Society of London, Series A}, 222, 309--368.

\bibitem{geifman2017}
Geifman, Y. and El-Yaniv, R. (2017).
Selective classification for deep neural networks.
In \textit{Advances in Neural Information Processing Systems 30}, 4878--4887.

\bibitem{godement1958}
Godement, R. (1958).
\textit{Topologie alg\'ebrique et th\'eorie des faisceaux}.
Hermann, Paris.

\bibitem{halmos1949}
Halmos, P.~R. and Savage, L.~J. (1949).
Application of the Radon-Nikodym theorem to the theory of sufficient statistics.
\textit{The Annals of Mathematical Statistics}, 20(2), 225--241.

\bibitem{hernan2020}
Hern\'an, M.~A. and Robins, J.~M. (2020).
\textit{Causal Inference: What If}.
Chapman \& Hall/CRC, Boca Raton, Florida.

\bibitem{horvitz1952}
Horvitz, D.~G. and Thompson, D.~J. (1952).
A generalization of sampling without replacement from a finite universe.
\textit{Journal of the American Statistical Association}, 47(260), 663--685.

\bibitem{kleinberg2017fairness}
Kleinberg, J., Mullainathan, S., and Raghavan, M. (2017).
Inherent trade-offs in the fair determination of risk scores.
In \textit{Proceedings of the Eighth Innovations in Theoretical Computer Science Conference}, Article 43, 1--23.

\bibitem{kleinrock1975}
Kleinrock, L. (1975).
\textit{Queueing Systems, Volume 1: Theory}.
John Wiley \& Sons, New York.

\bibitem{lehmann1998}
Lehmann, E.~L. and Casella, G. (1998).
\textit{Theory of Point Estimation}.
Second edition.
Springer, New York.

\bibitem{li2019}
Li, M. and Vit\'anyi, P. (2019).
\textit{An Introduction to Kolmogorov Complexity and Its Applications}.
Fourth edition.
Springer, Cham.

\bibitem{maclane1992}
Mac Lane, S. and Moerdijk, I. (1992).
\textit{Sheaves in Geometry and Logic: A First Introduction to Topos Theory}.
Springer-Verlag, New York.

\bibitem{markou2003a}
Markou, M. and Singh, S. (2003).
Novelty detection: A review, part 1, statistical approaches.
\textit{Signal Processing}, 83(12), 2481--2497.

\bibitem{markou2003b}
Markou, M. and Singh, S. (2003).
Novelty detection: A review, part 2, neural-network-based approaches.
\textit{Signal Processing}, 83(12), 2499--2521.

\bibitem{mitchell2002}
Mitchell, J.~L., Cazes, A.~N., and Leeder, N.~M. (2002).
Fast JPEG Huffman Encoding and Decoding.
U.S. Patent No. 6,373,412 B1.

\bibitem{nemhauser1978}
Nemhauser, G.~L., Wolsey, L.~A., and Fisher, M.~L. (1978).
An analysis of approximations for maximizing submodular set functions, I.
\textit{Mathematical Programming}, 14, 265--294.

\bibitem{neyman1935}
Neyman, J. (1935).
Sur un teorema concernente le cosiddette statistiche sufficienti.
\textit{Giornale dell'Istituto Italiano degli Attuari}, 6, 320--334.

\bibitem{pearl2009}
Pearl, J. (2009).
\textit{Causality: Models, Reasoning, and Inference}.
Second edition.
Cambridge University Press, Cambridge.

\bibitem{puterman1994}
Puterman, M.~L. (1994).
\textit{Markov Decision Processes: Discrete Stochastic Dynamic Programming}.
John Wiley \& Sons, New York.

\bibitem{robins1994}
Robins, J.~M., Rotnitzky, A., and Zhao, L.~P. (1994).
Estimation of regression coefficients when some regressors are not always observed.
\textit{Journal of the American Statistical Association}, 89(427), 846--866.

\bibitem{settles2009}
Settles, B. (2009).
\textit{Active Learning Literature Survey}.
Computer Sciences Technical Report 1648.
University of Wisconsin--Madison.

\bibitem{tishby1999}
Tishby, N., Pereira, F.~C., and Bialek, W. (1999).
The information bottleneck method.
In \textit{Proceedings of the Thirty-Seventh Annual Allerton Conference on Communication, Control, and Computing}, 368--377.

\bibitem{vapnik1998}
Vapnik, V.~N. (1998).
\textit{Statistical Learning Theory}.
John Wiley \& Sons, New York.

\end{thebibliography}

\end{document}